% =========================================================================
% modenanalyse_2fe2s -- Supplementary Technical Documentation
% Version 1.1.0 (June 2026)
% =========================================================================
%
% This document is the technical supplement to the user manual.
% It is intended to be self-contained: a reviewer should be able to
% read it without consulting the main manual. Every formula in the
% code is derived, every algorithm is specified by pseudocode, every
% configuration field is enumerated, every Excel column is described.
% Where the manual summarises, this supplement goes deeper.
%
% Compilation: pdflatex Supplement.tex (x3 for refs / TOC)
% =========================================================================

\documentclass[10pt,a4paper,twoside]{article}

\usepackage[T1]{fontenc}
\usepackage[utf8]{inputenc}
\usepackage[main=english,provide=*]{babel}
\usepackage{lmodern}
\usepackage{microtype}

\usepackage{amsmath, amssymb, amsthm}
\usepackage{mathtools}
\usepackage{bm}
\usepackage{siunitx}
\sisetup{range-phrase = --, range-units = single}

\usepackage[a4paper, left=2.3cm, right=2.3cm, top=2.5cm, bottom=2.8cm]{geometry}
\usepackage{setspace}
\onehalfspacing
\setlength{\parindent}{0pt}
\setlength{\parskip}{0.5em}

\usepackage{xcolor}
\definecolor{accent}{RGB}{0,80,140}
\definecolor{lightgrey}{RGB}{230,230,230}
\definecolor{darkgrey}{RGB}{60,60,60}
\definecolor{boxshade}{RGB}{245,245,250}

\usepackage{booktabs}
\usepackage{array}
\usepackage{tabularx}
\usepackage{longtable}
\usepackage{multirow}

\usepackage{enumitem}
\setlist{itemsep=2pt, topsep=4pt}

\usepackage{listings}
\lstdefinestyle{pycode}{
  language=Python,
  basicstyle=\ttfamily\footnotesize,
  keywordstyle=\color{accent},
  commentstyle=\color{darkgrey}\itshape,
  stringstyle=\color{darkgrey},
  showstringspaces=false,
  breaklines=true,
  frame=single,
  framesep=4pt,
  rulecolor=\color{lightgrey},
  backgroundcolor=\color{boxshade},
}
\lstdefinestyle{pseudo}{
  basicstyle=\ttfamily\footnotesize,
  keywordstyle=\color{accent}\bfseries,
  morekeywords={for,while,if,else,then,return,function,procedure,foreach,do,end,where,otherwise,break,continue,yield,let,output,input,require,assert},
  commentstyle=\color{darkgrey}\itshape,
  showstringspaces=false,
  breaklines=true,
  frame=single,
  framesep=4pt,
  rulecolor=\color{lightgrey},
  backgroundcolor=\color{boxshade},
}

\usepackage{tcolorbox}
\tcbuselibrary{breakable, skins}
\newtcolorbox{infobox}[1][]{
  enhanced, breakable, boxrule=0.4pt,
  colback=boxshade, colframe=accent!50!black,
  arc=2pt, left=8pt, right=8pt, top=6pt, bottom=6pt,
  #1
}
\newtcolorbox{derivbox}[1][]{
  enhanced, breakable, boxrule=0.4pt,
  colback=white, colframe=darkgrey,
  arc=2pt, left=8pt, right=8pt, top=6pt, bottom=6pt,
  title={#1},
  fonttitle=\bfseries\color{accent},
}
\newtcolorbox{workedexample}[1][]{
  enhanced, breakable, boxrule=0.6pt,
  colback=yellow!5!white, colframe=orange!70!black,
  arc=2pt, left=8pt, right=8pt, top=6pt, bottom=6pt,
  title={\textbf{Worked numerical example:} #1},
  fonttitle=\color{orange!50!black},
}
\newtcolorbox{bytebox}[1][]{
  enhanced, breakable, boxrule=0.4pt,
  colback=gray!5!white, colframe=gray!50!black,
  arc=1pt, left=6pt, right=6pt, top=4pt, bottom=4pt,
  fontupper=\ttfamily\footnotesize,
  #1
}

\usepackage{hyperref}
\hypersetup{
  colorlinks=true,
  linkcolor=accent,
  citecolor=accent,
  urlcolor=accent,
  pdftitle={modenanalyse\_2fe2s -- Supplementary Technical Documentation},
  pdfauthor={Lukas Knauer, AG Schuenemann, RPTU Kaiserslautern-Landau},
}
\usepackage[noabbrev]{cleveref}

\usepackage{fancyhdr}
\pagestyle{fancy}
\fancyhf{}
\fancyhead[LE,RO]{\textsf{\small\thepage}}
\fancyhead[RE]{\textsf{\small\nouppercase{\leftmark}}}
\fancyhead[LO]{\textsf{\small\nouppercase{\rightmark}}}
\renewcommand{\headrulewidth}{0.2pt}

% ---------- Theorem environments ------------------------------------------
\newtheorem{definition}{Definition}[section]
\newtheorem{satz}{Statement}[section]
\newtheorem{lemma}[satz]{Lemma}

% ---------- Math macros (same names as in Manual.tex for portability) -----
\newcommand{\R}{\mathbb{R}}
\newcommand{\rvec}{\mathbf{r}}
\newcommand{\Rvec}{\mathbf{R}}
\newcommand{\Rmat}{\mathbf{R}}
\newcommand{\evec}{\mathbf{e}}
\newcommand{\qvec}{\mathbf{q}}
\newcommand{\xvec}{\mathbf{x}}
\newcommand{\xhat}{\hat{\mathbf{x}}}
\newcommand{\yhat}{\hat{\mathbf{y}}}
\newcommand{\zhat}{\hat{\mathbf{z}}}
\newcommand{\nhat}{\hat{\mathbf{n}}}
\newcommand{\bhat}{\hat{\mathbf{b}}}
\newcommand{\Mmat}{\mathbf{M}}
\newcommand{\Hmat}{\mathbf{H}}
\newcommand{\Umat}{\mathbf{U}}
\newcommand{\Vmat}{\mathbf{V}}
\newcommand{\Smat}{\boldsymbol{\Sigma}}
\newcommand{\Lmat}{\boldsymbol{\Lambda}}
\newcommand{\Pmat}{\mathbf{P}}
\newcommand{\Qmat}{\mathbf{Q}}
\newcommand{\kB}{k_{\mathrm{B}}}
\newcommand{\urms}{u_{\mathrm{rms}}}
\newcommand{\Fei}{\mathrm{Fe}}
\newcommand{\OOP}{\mathrm{OOP}}
\newcommand{\INP}{\mathrm{INP}}
\newcommand{\Ag}{A_{\mathrm{g}}}
\newcommand{\Bgi}[1]{B_{\mathrm{#1g}}}
\newcommand{\Bui}[1]{B_{\mathrm{#1u}}}
\newcommand{\Aui}{A_{\mathrm{u}}}
\newcommand{\PCET}{\mathrm{PCET}}
\newcommand{\ET}{\mathrm{ET}}
\newcommand{\PT}{\mathrm{PT}}
\newcommand{\ket}[1]{\left|#1\right\rangle}
\newcommand{\bra}[1]{\left\langle#1\right|}

% ---------- Formula-to-code reference macro -------------------------------
% Used identically to the Manual: \codref{module.function} points to the
% function in the source tree that implements a given formula. The companion
% test tests/test_manual_code_refs.py is being extended to also verify the
% Supplement's codref entries.
\newcommand{\codref}[1]{%
  \quad\text{\textcolor{gray!70!black}{\footnotesize $\rightsquigarrow$~\texttt{#1}}}}

% ---------- Source-location macro for prose -------------------------------
% \src{core.py}{131-191} prints as a compact "core.py:131-191" reference.
\newcommand{\src}[2]{\textcolor{gray!70!black}{\footnotesize\texttt{[#1:#2]}}}

% ---------- File / config / function typesetters --------------------------
\newcommand{\fn}[1]{\texttt{#1}}
\newcommand{\file}[1]{\texttt{#1}}
\newcommand{\cfgkey}[1]{\texttt{#1}}

% =========================================================================
\title{%
  \textcolor{accent}{\Huge\textbf{modenanalyse\_2fe2s}}\\[0.3cm]
  \LARGE Supplementary Technical Documentation\\[0.2cm]
  \large for reviewers, methodology auditors, and developers
}
\author{Lukas Knauer\\
\small AG Sch\"unemann, RPTU Kaiserslautern-Landau\\
\small \texttt{lknauer@rptu.de}}
\date{Software version 1.1.0 \quad\textbar\quad Document compiled \today}

\begin{document}

\maketitle
\thispagestyle{empty}

\vfill
\begin{center}
\fbox{\parbox{0.86\textwidth}{
\small
\textbf{Scope of this document.}\\[2pt]
This supplement is the technical companion to the user manual
(\file{Manual.pdf}). The manual is the recommended entry point and
covers installation, configuration, workflow, output reference, and a
physically motivated theory overview with literature citations. This
supplement is \emph{independent of the manual} -- a reviewer can read
it standalone -- and goes \emph{deeper}: every formula in the code is
derived from first principles, every algorithm is given in pseudocode,
every data structure is enumerated, every configuration field is
described with its effect on the code path, every Excel column is
specified with its source function and uncertainty-propagation rule,
and the audit trail of v1.0.4 is presented in full. Where the manual
summarises, the supplement spells out.
\\[6pt]
\textbf{Code citations.} Every formula in this document carries
$\rightsquigarrow$~\texttt{module.function} markers pointing to the
implementing function in the source tree. These are verified by the
test suite (\file{tests/test\_manual\_code\_refs.py}) and cannot
silently drift out of sync with the code.\\[6pt]
\textbf{Repository.}
\url{https://github.com/lknauer/modenanalyse_2fe2s}
}}
\end{center}
\vfill

\clearpage
\tableofcontents
\clearpage

% =========================================================================
% PART I -- ARCHITECTURE AND CONVENTIONS
% =========================================================================
\part{Architecture and Conventions}

\section{System overview}
\label{sec:overview}

The package \fn{modenanalyse\_2fe2s} processes the output of a
quantum-chemical normal-mode calculation on a \mbox{[2Fe-2S]} cluster
embedded in a protein (or, for benchmark systems, in a small molecular
ligand environment), and produces a suite of mode-by-mode geometric,
energetic, and spectroscopic descriptors. The primary input is a
Gaussian frequency log (\file{.log}) with \fn{Freq=hpmodes} or, with
slightly reduced precision, a regular frequency log; an ORCA
\file{.hess} file is also accepted via an independent parser. The
optional second input is a Protein Data Bank file (\file{.pdb}) of
the underlying protein, used to map Gaussian atom indices to residue
labels and to detect secondary-structure elements when residue-level
analysis is requested.

The principal outputs are a hierarchy of Excel workbooks plus a
matching \file{run-log.txt} that records every decision the
pipeline took. The Excel workbooks group the per-mode descriptors by
chemical and physical theme:

\begin{itemize}
  \item \emph{Cluster-core} descriptors: OOP/INP fractions, cluster
        centre-of-mass displacement, cluster expansion, cluster
        rotation, $D_{2\mathrm{h}}$-symmetry-coordinate amplitudes
        (SCSD), kernel-classification scores.
  \item \emph{Fe-ligand bond} descriptors: stretching, bending
        in-plane, bending out-of-plane, all in \AA{}, with uncertainty.
  \item \emph{His-NH} descriptors: N-H stretching contribution per
        protonated histidine.
  \item \emph{Group} descriptors: pre-defined or auto-generated atom
        groups (residues, cofactor fragments) with OOP/INP/angle/
        torsion contributions.
  \item \emph{Marcus-Hush reorganization}: per-mode
        $\lambda_X^{\text{mode}}$ and per-channel
        $\Lambda_X^{\text{total}}$ along the
        $X \in \{\text{FeFe}, \text{FeN}, \text{FeS}, \text{NH},
                \text{HA}\}$ bond channels, plus frequency-resolved
        modulation spectra $M_X(\omega)$, co-modulation spectra
        $C_{\PCET}(\omega)$, $C_{\PT\text{-FeN}}(\omega)$,
        $C_{\ET\text{-FeS}}(\omega)$, and cumulative
        $\Lambda_X(\omega)$.
  \item \emph{Secondary-structure} descriptors: per-SSE-element
        amplitudes, centre-of-mass amplitudes, axial amplitudes,
        tilting angles.
  \item \emph{Embeddings}: three parallel UMAP-HDBSCAN clusterings on
        complementary feature spaces (Marcus-Hush reorganization
        features, SSE-element fingerprint, C$_\alpha$-amplitude
        fingerprint).
  \item \emph{Diagnostic} sheets: the \texttt{Coordination} sheet
        documents the Gaussian-to-PDB atom mapping used by all
        downstream calculations.
\end{itemize}

What the tool is \emph{not} designed to compute (and what is left
to specialised partner tools) includes: the NRVS spectrum itself
(the tool computes Fe-projected amplitudes and modulations, not the
NRVS cross-section), the Lamb-M\"o\ss bauer factor, recoilless
fractions, Fe-projected thermodynamic functions
($\Delta U$, $\Delta S$, $\Delta F$), spin-orbit-coupled
g-tensors, M\"o\ss bauer parameters, and absolute reorganization
energies in the original Marcus-classical sense (the tool computes
spectroscopic \emph{bond} reorganization energies, which are related
but distinct -- the distinction is spelled out in
\cref{sec:lambda-derivation}).

\subsection{Versions and supported file formats}

\begin{longtable}{p{0.18\textwidth} p{0.20\textwidth} p{0.55\textwidth}}
\toprule
\textbf{Property} & \textbf{Value} & \textbf{Notes} \\
\midrule
Package version           & 1.1.0              & SemVer; the audit trail
                                                  is \cref{sec:audit-trail}. \\
Python                    & $\geq$ 3.11        & Type-hint syntax used. \\
Primary input             & Gaussian \file{.log} \texttt{Freq=hpmodes}
                                              & 5-decimal eigenvector
                                                precision. Standard
                                                3-decimal eigenvectors
                                                are accepted as fallback. \\
Alternative input         & ORCA \file{.hess}  & Independent parser
                                                in \file{orca\_io.py}. \\
Secondary input           & PDB                & Optional. Enables
                                                residue labels,
                                                SSE analysis, embedding-
                                                feature C$_\alpha$
                                                amplitudes. \\
Output format             & \file{.xlsx} (openpyxl) & One main
                                                workbook plus optional
                                                interp / SSE /
                                                embedding workbooks. \\
Test count                & 147 unit + 4 E2E   & \cref{sec:test-suite} \\
Manual length             & 104 pages          & \file{Manual.pdf}\\
Anleitung length          & 107 pages          & \file{Anleitung.pdf}
                                                (German translation) \\
This supplement           & this document      & target \mbox{$\sim$}~180 pages \\
\bottomrule
\end{longtable}

% -------------------------------------------------------------------------
\section{Data flow architecture}
\label{sec:dataflow}

The pipeline is a sequence of well-defined phases that pass data
between each other by means of explicit Python data structures (no
hidden global state, except for the per-cluster deduplication sets
\fn{\_WARNED\_EMPTY\_GROUP}, \fn{\_WARNED\_EMPTY\_LIG},
\fn{\_WARNED\_HN\_SKIP}, which are reset at the start of every cluster
run by \fn{core.reset\_warning\_state}). This section walks through
the seven phases. The entry point is \fn{runner.run\_analysis(cfg)}
which dispatches to \fn{runner.\_run\_analysis\_single(cfg)} once per
cluster.

\subsection{Phase 1 -- byte-offset scan of the log file}
\label{sec:phase1-scan}

\codref{logio.scan\_log}

\fn{logio.scan\_log} performs a streaming, byte-level scan of the
Gaussian \file{.log} file in fixed-size chunks (default 4 MiB).
Three byte-pattern markers are matched against each chunk:

\begin{itemize}
  \item \texttt{Standard orientation:} -- the heading of every
        Cartesian-coordinate dump. The very first such dump (after
        SCF convergence and before any geometry optimisation step)
        is the equilibrium geometry used by all downstream
        geometric analyses.
  \item \texttt{and normal coordinates:} -- the heading of every
        section that lists displacement eigenvectors. Gaussian writes
        these in groups of three or five modes per row, depending on
        the \texttt{hpmodes} option, and writes one such section per
        frequency analysis. For an unscaled job, there is exactly one
        section; for a job with thermochemistry at multiple
        temperatures, there can be multiple sections, but they all
        contain the same eigenvectors.
  \item \texttt{Frequencies --} -- the headers that list the
        per-mode frequency, reduced mass, force constant, and IR
        intensity. One header per 3-mode (standard) or 5-mode
        (\texttt{hpmodes}) block.
\end{itemize}

For each match, the byte offset at the \emph{start} of the matched
pattern is recorded. This is a \texttt{O}($N_\text{bytes}$) scan with
constant memory; for a 4-GB Gaussian log it completes in a few seconds
on a modern SSD. A 200-byte rolling tail buffer ensures that pattern
matches spanning chunk boundaries are not missed.

The scan results are returned as three sorted lists of integers:
\fn{so\_off}, \fn{nc\_off}, \fn{fr\_off}. The cache layer
(\cref{sec:cache}) stores these lists in a pickle keyed by the log
file's full path, size, mtime, and Python major version; on cache
hit, the scan is skipped entirely, which makes subsequent runs
(e.g.\ for different frequency windows) order-of-magnitude faster
than the first run.

\begin{infobox}
\textbf{Clarification of a common misconception.}
\fn{scan\_log} \emph{reads the entire log file from beginning to end}
to find these offsets. It does not honour any frequency filter --
the frequency filter is applied much later in the pipeline. The
metadata of every mode in the file therefore becomes available in
RAM, regardless of \cfgkey{freq\_min}/\cfgkey{freq\_max}. What the
frequency filter \emph{does} guard is the expensive eigenvector
analysis in Phase 4: only modes whose frequency falls inside
$[\cfgkey{freq\_min}, \cfgkey{freq\_max}]$ trigger a full
\fn{analyze\_mode\_with\_fallback} call. This distinction matters for
the context-mode loading in
\cref{sec:phase5-context}: the candidate set for context modes is
the full set of modes in the log, not a re-scan.
\end{infobox}

\subsection{Phase 2 -- equilibrium geometry}
\label{sec:phase2-geom}

\codref{logio.read\_std\_orient}

The first Standard-Orientation block (selected by smallest byte
offset in \fn{so\_off}) is parsed by \fn{logio.read\_std\_orient}.
This block lists, for every atom in the calculation:

\begin{itemize}
  \item the \emph{Center number} (1-based, called the \emph{Gaussian
        centre} throughout this document),
  \item the \emph{Atomic number} (used to look up the element symbol
        via the dictionary \fn{logio.\_ELEM}),
  \item the equilibrium Cartesian coordinates $(x, y, z)$ in \AA{}
        (Gaussian writes them in the standard orientation, with the
        principal-axis frame chosen by Gaussian's symmetry detector).
\end{itemize}

The output is a list of \fn{Atom} dictionaries
$\{\texttt{symbol}, x, y, z, \texttt{center}\}$ and an
\fn{idx\_map} from Gaussian centre to position in the atoms list. By
convention, \emph{Gaussian centre} = position in the list $+ 1$, so
the mapping is trivial -- but it is computed explicitly because
later filtering (heavy-atom-only mode, hydrogen-included mode) breaks
the trivial correspondence.

The choice of equilibrium frame is significant: every subsequent
geometric operation -- the SVD-based cluster-normal computation in
\cref{sec:cluster-normal}, the Kabsch alignment for the PDB-to-Gaussian
mapping in \cref{sec:kabsch}, the SCSD reference geometry in
\cref{sec:scsd-derivation} -- is expressed in this same frame.

\subsection{Phase 3 -- block metadata}
\label{sec:phase3-meta}

\codref{logio.read\_all\_meta}

\fn{logio.read\_all\_meta} iterates over the \fn{fr\_off} list from
Phase 1, opens the log file at each frequency-header offset, and reads
the per-mode metadata of the corresponding mode block:

\begin{itemize}
  \item the mode numbers in the block (3 or 5 modes per block,
        depending on \texttt{hpmodes}),
  \item the frequency of each mode in cm$^{-1}$,
  \item the reduced mass in amu,
  \item the force constant in mDyne/\AA{},
  \item the IR intensity (parsed but currently unused downstream),
  \item the irreducible-representation label (\texttt{A}, \texttt{A'},
        $A_g$, etc.; used for kernel scoring),
  \item a flag indicating whether the block is from a \texttt{hpmodes}
        section (precision $\sim 10^{-5}$~\AA) or a standard section
        (precision $\sim 10^{-3}$~\AA).
\end{itemize}

The eigenvector matrix itself is \emph{not} loaded here -- only the
metadata. The reason is RAM: a single $N_\text{atoms} \times 3$
eigenvector matrix for a $10^4$-atom system is 240 kB; for $10^4$
modes that would be 2.4 GB. The metadata structure
(\fn{BlockInfo}, \cref{sec:dataclass-blockinfo}) carries the byte
offset and the column index within the block, so a subsequent
\fn{logio.get\_eigvec} call can stream the eigenvector data in
3$N$-row chunks on demand.

The function returns three structures:

\begin{itemize}
  \item \fn{all\_blocks}: list of \fn{BlockInfo}, one per
        \texttt{Frequencies --} header in the log;
  \item \fn{best\_block}: mapping from mode number to its preferred
        \fn{BlockInfo} (HP block preferred over standard for the same
        mode number, used to reconcile mixed \texttt{hpmodes} +
        standard logs);
  \item \fn{cand\_map}: mapping from mode number to all
        \fn{BlockInfo}s that contain that mode (used as a fallback
        ordering by \fn{analyze\_mode\_with\_fallback} if the
        preferred block fails to parse).
\end{itemize}

\subsection{Phase 4 -- main mode-analysis loop}
\label{sec:phase4-loop}

\codref{runner.\_run\_analysis\_single}

The main loop iterates over the modes selected by the frequency
filter:

\begin{lstlisting}[style=pseudo,caption={Main mode-analysis loop, schematic.}]
selected <- []
for bi in all_blocks:
    for col, (mn, freq) in enumerate(zip(bi.mode_nums, bi.freqs)):
        if mn not in best_block or best_block[mn] is not bi:
            continue                # only count preferred block once
        if freq <= 0:
            log "imaginary or zero-frequency mode"; continue
        if cfg.freq_min and freq < cfg.freq_min: continue
        if cfg.freq_max and freq > cfg.freq_max: continue
        if cfg.freq_windows and freq not in any window:
            continue
        selected.append((bi, col, mn, freq))
selected.sort(by=freq)
for (bi, col, mn, freq) in selected:
    r, _ = analyze_mode_with_fallback(
              candidates=cand_map[mn], col=col,
              log_file=cfg.log_file, atoms=atoms, idx_map=idx_map,
              normal=cluster_normal, coord_info=coord_info,
              fe_c=fe_c, s_c=s_c, cfg=cfg, mode_num=mn)
    accumulate_b_factor(b_accum, r)
    results.append(r)
\end{lstlisting}

For every selected mode, the pipeline calls \fn{core.analyze\_mode\_with\_fallback},
which in turn calls \fn{logio.get\_eigvec} to stream
the eigenvector from the log, scales it by $\urms$ (\cref{sec:urms-derivation}),
projects it onto the cluster-normal (\cref{sec:oop-derivation}), and
runs all the per-mode analyses summarised at the start of
\cref{sec:overview}. The return value is a result dictionary with
the keys listed in \cref{sec:dataclass-modedict}.

Phase 4 is the dominant computation cost. For a 1000-mode system,
phase 4 typically takes \mbox{$\sim$}~30 s; phases 1--3 together take
\mbox{$\sim$}~2--3 s; phases 5--7 take \mbox{$\sim$}~1 s. The
\fn{cfg.n\_jobs} option can parallelise phase 4 over multiple cores
(default: serial, deterministic).

\subsection{Phase 5 -- context-mode loading}
\label{sec:phase5-context}

\codref{runner.\_run\_analysis\_single}

After the main loop has produced \fn{results} (a list of per-mode
dicts for the modes inside the frequency window), two additional
small mode sets are loaded for the sole purpose of anchoring the
linear interpolation used in the \file{*\_analysis\_interp\#\#.xlsx}
workbook. These are the \emph{context modes}.

The right-edge (upper-frequency) context set is selected from the
window $(\cfgkey{freq\_max},\, \cfgkey{freq\_max}+
\cfgkey{interp\_context\_cm1}]$. If that window contains no mode --
because the next mode above \cfgkey{freq\_max} is farther than
\cfgkey{interp\_context\_cm1} away -- a single fallback mode is
loaded: the closest mode above \cfgkey{freq\_max} (this is
\emph{FUND 12}, \cref{sec:fund-12}). If even that fails, because the
DFT spectrum genuinely terminates inside the analysis range, a
\emph{synthetic zero anchor} is inserted at
$\cfgkey{freq\_max} + \cfgkey{interp\_context\_cm1}$ with all
observables equal to zero (this is \emph{FUND 13},
\cref{sec:fund-13}, produced by \fn{runner.\_make\_synthetic\_zero\_mode}). The synthetic mode carries sentinel values
$\texttt{number} = -1$, $\texttt{mode\_type} = \texttt{"synthetic\_zero"}$,
$\texttt{precision} = \texttt{"synthetic"}$.

The left-edge (lower-frequency) context set is selected symmetrically
in $[\cfgkey{freq\_min} - \cfgkey{interp\_context\_cm1},\,
\cfgkey{freq\_min})$, again with the FUND 12 single-anchor fallback.
At the lower edge no synthetic zero is inserted: if no real mode
exists below \cfgkey{freq\_min}, the interpolation's
\texttt{left}~$= 0.0$ boundary is already the physically correct
behaviour (the spectrum starts at higher frequency).

\subsection{Phase 6 -- system-wide aggregates}
\label{sec:phase6-aggregate}

After every mode has been analysed, the per-mode results are
aggregated into system-wide quantities. These aggregations are
exposed in the \file{*\_analysis\_main.xlsx} as separate sheets:

\begin{itemize}
  \item $B$-factors (Debye-Waller): \fn{b\_factors} = $8\pi^2/3
        \cdot \fn{b\_accum}$, with \fn{b\_accum} accumulated mode by
        mode during Phase 4 as $\sum_l \urms^2 \cdot
        \lVert\evec_{l,i}\rVert^2$ (\cref{sec:b-factor-derivation}).
  \item Total reorganization $\Lambda_X^{\text{total}}$ per channel
        ($X \in \{\text{FeFe, FeN, FeS, NH, HA}\}$), summed over all
        modes in the window (\cref{sec:lambda-total}).
  \item Modulation spectra $M_X(\omega)$, co-modulation spectra,
        cumulative-reorganization curves $\Lambda_X(\omega)$,
        gaussian-broadened on a uniform frequency grid
        (\cref{sec:modulation-spectra}).
\end{itemize}

\subsection{Phase 7 -- export}

\codref{export.export\_all}

All per-mode results plus the system-wide aggregates are written to
the configured Excel workbooks. The export layer uses openpyxl with
a custom sheet-formatting helper that applies:

\begin{itemize}
  \item per-cell colour-coding by value/sigma ratio (the
        \emph{significance colouring} discussed in
        \cref{sec:uncertainty-overview}),
  \item per-column number-format strings, chosen by quantity type,
  \item header rows with bold typography and merged-cell groupings
        for related columns.
\end{itemize}

When multi-cluster mode is active (\cfgkey{analyze\_all\_clusters} = true),
phases 1--7 are repeated for each detected cluster, with the cluster
index appearing in the output-directory name. Inter-cluster results
are not currently aggregated -- each cluster's outputs are independent.

\subsection{Cache}
\label{sec:cache}

\codref{logio.load\_scan\_cache, logio.save\_scan\_cache}

The byte-offset scan from Phase 1 dominates the runtime for large
log files. To avoid re-scanning when the user runs a second analysis
on the same log (e.g.\ a different frequency window), the scan
results are cached in a pickle file in the current working directory
named \file{.scan\_cache\_\#\#.pkl}, where \verb|##| is a hash of:

\begin{itemize}
  \item the absolute path of the log file,
  \item the log file's size in bytes,
  \item the log file's modification time (\fn{os.stat.st\_mtime}),
  \item the Python major version (3.11, 3.12, \ldots; pickle format
        compatibility between Python major versions is not
        guaranteed),
  \item a hard-coded cache-format version string
        (\fn{logio.\_CACHE\_VERSION}) bumped whenever the format
        changes.
\end{itemize}

A cache hit skips the entire scan in Phase 1 and reduces
end-to-end runtime by 40--80\% depending on log file size. The cache
is invalidated automatically by any of the keying changes; manual
invalidation is to delete the \file{.scan\_cache\_\#\#.pkl} file or
set \cfgkey{use\_cache} = false in the TOML.

% -------------------------------------------------------------------------
\section{Conventions used throughout the code}
\label{sec:conventions}

Quantum-chemistry codes have a number of subtle conventions that
differ between programs, and ad-hoc choices have to be made about how
those interact with the geometric/spectroscopic analysis we
perform. This section enumerates every such choice explicitly. A
reviewer should be able to verify each one against the original
source code via the \codref{} markers.

\subsection{Eigenvector normalisation}
\label{sec:conv-eigvec}

\codref{logio.get\_eigvec}

Gaussian writes Cartesian \emph{mass-unweighted} displacement
eigenvectors $\evec_{l,i} \in \R^3$ for atom $i$ in mode $l$ (with
$l = 1, \ldots, 3N - 6$ for a non-linear molecule of $N$ atoms),
normalised so that
\begin{equation}
  \sum_{i=1}^{N} \lVert\evec_{l,i}\rVert^2 \;\approx\; 1.
  \label{eq:eigvec-norm}
\end{equation}
The approximate sign reflects rounding to 5 (\texttt{hpmodes}) or 3
(\texttt{standard}) decimals. This is the \emph{geometric-Cartesian}
normalisation: the dimensionless eigenvector component
$\evec_{l,i}^{(\alpha)}$ describes the cartesian displacement of atom
$i$ along axis $\alpha$ in some abstract dimensionless ``mode-unit''
system. The mass-weighted eigenvector, which is the one for which
$\sum_i m_i \lVert \bm{u}_{l,i}\rVert^2 = 1$ holds, is obtained from
the displacement form as $\bm{u}_{l,i} = \sqrt{m_i}\,\evec_{l,i}$
(after re-normalisation).

\begin{infobox}
\textbf{Why the distinction matters for the OOP/INP fraction.}
The OOP fraction defined in \cref{sec:oop-derivation},
\begin{equation*}
  \alpha_{\OOP}(G; l) = \frac{\sum_{i \in G}\lVert\evec_{l,i}^{\OOP}\rVert^2}
                              {\sum_{i \in G}\lVert\evec_{l,i}\rVert^2},
\end{equation*}
is the \emph{cartesian-displacement} OOP fraction, because Gaussian's
\texttt{hpmodes} are in the displacement form. If one wanted the
\emph{kinetic-energy} fraction
$T_{\OOP}/T_{\text{total}}$, every $\lVert\evec\rVert^2$ would have to
be replaced by $m_i \lVert\evec\rVert^2$. For a [2Fe-2S] cluster the
two quantities are numerically similar -- both Fe and both S have
broadly comparable masses (56 vs 32 amu) -- but the distinction is
critical when comparing $\alpha_{\OOP}$ values between chemically
different systems. The manual flags this caveat;
\cref{sec:oop-derivation} below derives the formula explicitly so
that the convention is unambiguous.
\end{infobox}

\subsection{Heavy-atom-only vs.\ hydrogen-included filtering}
\label{sec:conv-heavy}

\codref{orca\_io.parseresult\_to\_atoms, orca\_io.get\_eigvec\_orca}

By default, the cluster-core analysis operates on the heavy-atom
subset of the molecule -- explicit hydrogens are filtered out before
the eigenvector matrix is reshaped. This reduces RAM and
matches the chemical intuition that the Fe-S core dynamics are
heavy-atom-dominated. Hydrogen-inclusion is opt-in via
\cfgkey{include\_hn\_vibration} = true, which switches the relevant
analysis (the His-NH stretch contribution, \cref{sec:his-nh}) to a
hydrogen-inclusive eigenvector loader.

Both loaders preserve the \emph{Gaussian centre} (1-based index from
the equilibrium-geometry block) as the unique atom identifier --
the position within the heavy-only or all-H atoms list changes
between loaders, but the centre does not. All cluster-atom centres
(\fn{fe\_c}, \fn{s\_c}, \fn{c2l}, \fn{c2l\_h}) are stored in
Gaussian-centre coordinates, never in positional indices. The
positional index is recovered when needed via
$\text{atoms}[\fn{idx\_map}[\text{centre}]]$.

\subsection{Cluster atom ordering}
\label{sec:conv-cluster-order}

\codref{geometry.find\_cluster, geometry.find\_all\_clusters}

The cluster-core atoms are stored in the canonical order
$[\text{Fe}_1, \text{Fe}_2, \text{S}_1, \text{S}_2]$ throughout the
code. This is essential for SCSD: the hard-coded model coordinates
\fn{core.\_SCSD\_MODEL\_COORDS} are also given in this order, and
the SCSD projection assumes positional correspondence. The two Fe
atoms and the two S atoms are sorted by Gaussian centre (smaller
centre first), making the assignment deterministic regardless of how
Gaussian numbers them. The chirality of the cluster (which Fe is
``up'') is then fixed by the SVD-based cluster-normal computation
(\cref{sec:cluster-normal}), which orients the normal in the
$+\zhat$ half-space of the Gaussian frame.

\subsection{Sign conventions for bond modulations}
\label{sec:conv-dr-signs}

\codref{reorganization.signed\_dr\_along\_axis}

For a generic bond modulation $\Delta r_{ab}$ along atoms $a$ and
$b$, the sign convention used everywhere is
\begin{equation}
  \Delta r_{ab}^{\text{(class.)}}
    \;=\; (\evec_a - \evec_b) \cdot \nhat_{ab},
  \qquad
  \nhat_{ab} \;=\; \frac{\rvec_a - \rvec_b}{\lVert\rvec_a - \rvec_b\rVert}.
  \label{eq:conv-dr}
\end{equation}
That is, the bond unit-vector $\nhat_{ab}$ points from $b$ to $a$,
so $\Delta r_{ab} > 0$ means the bond is stretching (atoms $a, b$
move apart) and $\Delta r_{ab} < 0$ means the bond is contracting.
This convention is used uniformly for all five Marcus-Hush channels
(FeFe, FeN, FeS, NH, HA), with the qualification that for HA the
\emph{reaction coordinate} variant (\cref{sec:dr-ha}) is used instead
of the classical $\evec_H - \evec_A$ projection.

\subsection{Units}
\label{sec:conv-units}

\begin{longtable}{p{0.21\textwidth} p{0.18\textwidth} p{0.55\textwidth}}
\toprule
\textbf{Quantity} & \textbf{Unit} & \textbf{Remark} \\
\midrule
Atomic coordinates                 & \AA{} (\unit{\angstrom})
   & Gaussian standard-orientation. \\
Eigenvector components             & dimensionless
   & Normalised to $\sum_i \lVert\evec_{l,i}\rVert^2 \approx 1$. \\
Thermally scaled eigenvectors      & \AA{}
   & $\evec_{l,i} \cdot \urms(\omega_l)$ in \AA{}.\\
Frequencies                        & cm$^{-1}$
   & Standard spectroscopic convention. \\
Reduced mass                       & amu
   & As Gaussian writes it. \\
Force constants (Gaussian)         & mDyne/\AA{}
   & Parsed but not currently used in \fn{core}. \\
Temperatures                       & K        & SI. \\
Bond modulations $\Delta r_X$      & \AA{}
   & \cref{sec:conv-dr-signs}.\\
Reorganization energies            & cm$^{-1}$
   & Spectroscopic convention; conversion to other units in
     \cref{sec:lambda-derivation}. \\
$B$-factors                        & \AA$^2$  & \cref{sec:b-factor-derivation}. \\
SCSD amplitudes                    & \AA{}    & \cref{sec:scsd-derivation}. \\
SSE-element amplitudes              & \AA{}    & \\
\bottomrule
\end{longtable}

\subsection{Sign of the cluster normal}
\label{sec:conv-normal-sign}

\codref{geometry.cluster\_normal}

The right-singular vector $\mathbf{V}_{:,3}$ produced by the SVD in
\cref{sec:cluster-normal} is determined up to a sign by the SVD
itself. The convention adopted here is to choose the sign so that
$\nhat \cdot \zhat > 0$, i.e.\ the cluster normal points into the
$+\zhat$ half-space of the Gaussian standard orientation. This fixes
the sign of $\Delta r_\text{bend\_oop}$ in the Fe-ligand decomposition
of \cref{sec:fe-ligand-decomp}. The bend\_inp/bend\_oop \emph{absolute
values} are unaffected by the choice of sign for $\nhat$, only the
sign of bend\_oop is.

% =========================================================================
% PART II -- EVERY FORMULA IN THE CODE
% =========================================================================
\part{Every Formula in the Code}

This part derives every non-trivial formula used by the code from
first principles. Each derivation ends with a \codref{} marker that
points to the implementing function in the source tree. The order
follows the data-flow logic of \cref{sec:dataflow}: from the lowest
primitives ($\urms$, OOP/INP) to the highest aggregates
($\Lambda_X^\text{total}$, modulation spectra).

% -------------------------------------------------------------------------
\section{Thermal RMS amplitude $\urms$}
\label{sec:urms-derivation}

\subsection{Quantum-mechanical derivation}

Consider a one-dimensional quantum harmonic oscillator of mass $m$
and angular frequency $\omega$. The Hamiltonian
$\hat{H} = \hat{p}^2/(2m) + \tfrac{1}{2}m\omega^2 \hat{q}^2$ has
energy eigenvalues $E_n = \hbar\omega(n + \tfrac{1}{2})$ and
position-variance expectation values in eigenstate $\ket{n}$
\cite{WilsonDeciusCross1955}:
\begin{equation}
  \langle n | \hat{q}^2 | n \rangle
    \;=\; \frac{\hbar}{m\omega}\left(n + \tfrac{1}{2}\right).
  \label{eq:qho-q2-eigen}
\end{equation}

In thermal equilibrium with reservoir temperature $T$, the
expectation value of $\hat{q}^2$ is the canonical-ensemble average
of \cref{eq:qho-q2-eigen} over the Boltzmann-weighted eigenstates:
\begin{equation}
  \langle \hat{q}^2 \rangle_T
    \;=\;
  \frac{\sum_{n=0}^{\infty}
        e^{-\beta\hbar\omega(n + 1/2)}
        \langle n | \hat{q}^2 | n \rangle}
       {\sum_{n=0}^{\infty}
        e^{-\beta\hbar\omega(n + 1/2)}},
  \qquad
  \beta = \frac{1}{\kB T}.
  \label{eq:qho-q2-thermal}
\end{equation}

Both sums are geometric series in $e^{-\beta\hbar\omega}$. Working
through (writing $x = \beta\hbar\omega/2$):
\begin{align}
  \sum_{n=0}^{\infty} e^{-\beta\hbar\omega(n + 1/2)}
    &= \frac{e^{-x}}{1 - e^{-2x}}
     = \frac{1}{e^{x} - e^{-x}}
     = \frac{1}{2\sinh x}, \\[4pt]
  \sum_{n=0}^{\infty} (n + \tfrac{1}{2})
        e^{-\beta\hbar\omega(n + 1/2)}
    &= -\frac{1}{\beta\hbar\omega}
       \frac{d}{dx}\!\left(\frac{1}{2\sinh x}\right)
     = \frac{\cosh x}{2\sinh^2 x}.
\end{align}
Substituting back,
\begin{equation}
  \langle \hat{q}^2 \rangle_T
    \;=\;
  \frac{\hbar}{m\omega}
  \cdot \frac{\cosh x}{2\sinh^2 x} \cdot 2\sinh x
    \;=\;
  \frac{\hbar}{m\omega}\,\frac{\cosh x}{2\sinh x}
    \;=\;
  \frac{\hbar}{2m\omega}\,\coth\!\left(\frac{\hbar\omega}{2\kB T}\right).
  \label{eq:qho-q2-thermal-final}
\end{equation}

Taking the square root and identifying $m$ with the mode's reduced
mass $m_\text{red}$, we obtain the formula used throughout the code:
\begin{equation}
  \boxed{\;\urms(\omega, m_\text{red}, T)
    \;=\;
   \sqrt{\dfrac{\hbar}{2\,m_\text{red}\,\omega}\,
         \coth\!\left(\dfrac{\hbar\omega}{2\kB T}\right)}\;}
  \codref{core.compute\_thermal\_amplitude}
  \label{eq:urms-final}
\end{equation}
in SI units. The code returns the result in \AA{} by multiplying by
$10^{10}$ before returning.

\subsection{Numerical implementation details}

\codref{core.compute\_thermal\_amplitude}

The function \fn{core.compute\_thermal\_amplitude(freq\_cm1,
red\_mass\_amu, temp\_k, amplitude)} converts inputs into SI
($\omega = 2\pi c \cdot \tilde\nu$, $c$ in cm/s), evaluates
\cref{eq:urms-final}, and returns the result in \AA{}. Several edge
cases are handled defensively:

\begin{itemize}
  \item \emph{Imaginary or zero frequency}: $\omega \leq 0$ returns
        $0$. The mode-loop already filters such modes
        (\cref{sec:phase4-loop}), so this branch is a safeguard.
  \item \emph{Negative or zero reduced mass}: $m_\text{red} \leq 0$
        returns $0$. Should never happen for a valid Gaussian log.
  \item \emph{Classical mode}: $\texttt{temp\_k}$ = \texttt{None}
        triggers the classical-amplitude convention -- $\urms$ is
        replaced by the constant \cfgkey{amplitude} (default 1.0),
        ignoring frequency and mass. This is for benchmarking against
        legacy tools that use a fixed amplitude.
  \item \emph{Numerical $\coth$}: for $\beta\hbar\omega/(2) \ll 1$
        (i.e.\ very low $\omega$ or very high $T$), the
        $\coth(x) \approx 1/x$ regime is approached. The code uses
        \fn{numpy.cosh}/\fn{numpy.sinh} directly; for IEEE-754 doubles
        this is stable down to $\omega \approx 10^{-7}$ cm$^{-1}$,
        well below any physically meaningful regime.
  \item \emph{Zero-point limit}: $T \to 0$ gives
        $\coth(\beta\hbar\omega/2) \to 1$ and
        $\urms \to \sqrt{\hbar/(2 m_\text{red}\omega)}$. This is the
        zero-point amplitude.
\end{itemize}

\begin{derivbox}[Sanity checks (test\_sanity\_physics.py)]
The following property tests verify \cref{eq:urms-final}
analytically:
\begin{itemize}
  \item \fn{test\_urms\_zero\_T\_limit}: with
        $T = 10^{-3}$~K, the result matches
        $\sqrt{\hbar/(2 m_\text{red}\omega)}$ to $10^{-6}$
        relative error.
  \item \fn{test\_urms\_classical\_limit}: with $\omega = 50$~cm$^{-1}$
        and $T = 2000$~K, the result matches the classical
        equipartition limit $\sqrt{\kB T/(m_\text{red}\omega^2)}$ to
        $< 1\%$ relative error.
  \item \fn{test\_urms\_monotonic\_in\_T}: $\urms$ is monotonic
        increasing in $T$ for fixed $\omega, m_\text{red}$.
\end{itemize}
\end{derivbox}

\subsection{Limiting behaviours}

The two important limits are:

\paragraph{Zero-point limit.}
For $\hbar\omega \gg \kB T$, $\coth(x) \to 1$, so
\begin{equation}
  \urms \;\xrightarrow{T \to 0}\; \sqrt{\frac{\hbar}{2 m_\text{red}\omega}}.
  \label{eq:urms-zeropoint}
\end{equation}
For a 200-cm$^{-1}$ mode with $m_\text{red} = 30$~amu, this is
about 0.04 \AA{}.

\paragraph{Classical / equipartition limit.}
For $\hbar\omega \ll \kB T$, $\coth(x) \to 1/x = 2\kB T/(\hbar\omega)$,
so
\begin{equation}
  \urms \;\xrightarrow{T \to \infty}\;
    \sqrt{\frac{\kB T}{m_\text{red}\omega^2}}.
  \label{eq:urms-classical}
\end{equation}
This is the equipartition theorem
$\tfrac{1}{2}m\omega^2 \langle q^2 \rangle = \tfrac{1}{2}\kB T$.
For a 50-cm$^{-1}$ mode at room temperature, $\hbar\omega/\kB =
72$~K, so the classical limit is approximately satisfied.

\begin{workedexample}[$\urms$ for an Fe-S stretch at 80~K]
We compute $\urms$ for a typical Fe-S symmetric-stretch mode of a
\mbox{[2Fe-2S]} cluster, using representative values from a benchmark
log (such values are typical for Cys$_2$His$_2$- or Cys$_4$-coordinated
[2Fe-2S]$^{1+/2+}$ systems):
\begin{itemize}
  \item $\tilde\nu = 320~\text{cm}^{-1}$ (frequency),
  \item $m_\text{red} = 31.97~\text{amu}$ (reduced mass; printed by
        Gaussian as ``Red. masses --'' for this mode),
  \item $T = 80~\text{K}$ (cryogenic NRVS condition).
\end{itemize}

\textbf{Step 1 -- convert to SI:}
\begin{align*}
  \omega &= 2\pi c \cdot \tilde\nu
       = 2\pi \cdot 2.998 \times 10^{10}~\text{cm/s} \cdot 320~\text{cm}^{-1} \\
       &= 6.026 \times 10^{13}~\text{rad/s}. \\[2pt]
  m_\text{red} &= 31.97 \cdot 1.6605 \times 10^{-27}~\text{kg}
       = 5.309 \times 10^{-26}~\text{kg}. \\[2pt]
  \beta\hbar\omega/2 &= \frac{\hbar\omega}{2\kB T}
       = \frac{1.0546 \times 10^{-34} \cdot 6.026 \times 10^{13}}
              {2 \cdot 1.3807 \times 10^{-23} \cdot 80} \\
       &= \frac{6.354 \times 10^{-21}}{2.209 \times 10^{-21}}
       = 2.876.
\end{align*}

\textbf{Step 2 -- evaluate $\coth$:}
\[
  \coth(2.876)
    = \frac{e^{2.876} + e^{-2.876}}{e^{2.876} - e^{-2.876}}
    = \frac{17.74 + 0.0564}{17.74 - 0.0564}
    = 1.0064.
\]
Since $\hbar\omega/\kB T = 5.75 \gg 1$, we are well inside the
zero-point-dominated regime; $\coth$ is close to 1.

\textbf{Step 3 -- assemble:}
\[
  \urms^2 = \frac{\hbar}{2\,m_\text{red}\,\omega}\,\coth(\beta\hbar\omega/2)
   = \frac{1.0546 \times 10^{-34}}
          {2 \cdot 5.309 \times 10^{-26} \cdot 6.026 \times 10^{13}}
   \cdot 1.0064.
\]
\[
  = \frac{1.0546}{6.398} \times 10^{-22} \cdot 1.0064~\text{m}^2
   = 1.659 \times 10^{-23}~\text{m}^2.
\]
\[
  \urms = \sqrt{1.659 \times 10^{-23}}~\text{m}
   = 4.073 \times 10^{-12}~\text{m}
   = \boxed{0.04073~\text{\AA}}.
\]

\textbf{Cross-check (zero-point limit):}
With $\coth \to 1$ exactly, $\urms^\text{ZP} = \sqrt{\hbar/(2 m\omega)}
= 0.04062$~\AA{}. The thermal correction at 80~K is therefore only
$\sim 0.3\%$ -- as expected for $\hbar\omega/\kB T \gg 1$.

\textbf{Cross-check (classical limit):}
$\sqrt{\kB T/(m\omega^2)} = \sqrt{1.105 \times 10^{-21} /
(5.309 \times 10^{-26} \cdot 3.631 \times 10^{27})}
\approx 0.0240$~\AA{}. The true value $0.041$~\AA{} is far above
this -- the classical limit gives \emph{less} amplitude than the
quantum result at 80~K, confirming we are in the quantum regime.

\textbf{In the code:}
\fn{core.compute\_thermal\_amplitude(320.0, 31.97, 80.0, 1.0)} returns
\texttt{0.04073} (the third argument is \cfgkey{temp\_k};
the fourth, \cfgkey{amplitude}, is unused when \cfgkey{temp\_k}
is set).
\end{workedexample}

\subsection{Thermally scaled eigenvector}

The thermally scaled eigenvector $\widetilde{\evec}_{l,i}
\equiv \evec_{l,i} \cdot \urms(\omega_l)$ has units of \AA{} and
\emph{is the displacement of atom $i$ in mode $l$ averaged over the
thermal density matrix}. Every downstream geometric observable
(bond modulation, bend, group-amplitude, SCSD amplitude, SSE-element
amplitude, C$_\alpha$ amplitude) is computed from
$\widetilde{\evec}$, not from the raw $\evec$. This means that:

\begin{itemize}
  \item All observables are in \AA{} and have a direct physical
        interpretation as ``how far does the atom actually move at
        temperature $T$''.
  \item The accompanying $\sigma$ for each observable must scale the
        \emph{raw} per-component eigenvector noise
        $\sigma_e = \cfgkey{sigma\_eigvec}$ by the same $\urms$
        factor, $\sigma_e^\text{therm} = \sigma_e \cdot \urms$
        (\cref{sec:sigma-propagation}). Failure to do this was the
        \emph{FUND 6/7/8} bug class fixed in v1.0.4
        (\cref{sec:fund-6-8}).
\end{itemize}

% -------------------------------------------------------------------------
\section{OOP/INP decomposition}
\label{sec:oop-derivation}

\subsection{Definition}

\codref{core.\_oop, core.classify\_oop\_inp}

Given a unit cluster-normal $\nhat$ (\cref{sec:cluster-normal}) and a
set $G$ of atom indices, define the per-atom out-of-plane component
of the eigenvector $\evec_{l,i}$ (or its thermally scaled form
$\widetilde{\evec}_{l,i}$ -- the ratio is the same):
\begin{equation}
  \evec_{l,i}^{\OOP} \;\equiv\; (\evec_{l,i} \cdot \nhat)\,\nhat.
  \label{eq:oop-component}
\end{equation}
This is the projection of $\evec_{l,i}$ onto the cluster-normal
direction. The orthogonal complement,
$\evec_{l,i}^{\INP} = \evec_{l,i} - \evec_{l,i}^{\OOP}$, lies in
the cluster plane.

The OOP fraction of mode $l$ on group $G$ is the ratio of squared
norms:
\begin{equation}
  \boxed{\;\alpha_{\OOP}(G; l)
    \;=\;
  \frac{\sum_{i \in G}\lVert \evec_{l,i}^{\OOP}\rVert^2}
       {\sum_{i \in G}\lVert \evec_{l,i}\rVert^2}\;\cdot\;100\,\%\;}
  \codref{core.classify\_oop\_inp}
  \label{eq:oop-derivation}
\end{equation}
The INP fraction is $\alpha_{\INP}(G; l) = 1 - \alpha_{\OOP}(G; l)$
(after the percentage conversion).

\begin{derivbox}[Pythagorean orthogonality (test\_sanity\_physics.py)]
The decomposition
\cref{eq:oop-component} satisfies the Pythagorean identity
\begin{equation*}
  \lVert\evec_{l,i}\rVert^2
    = \lVert\evec_{l,i}^{\OOP}\rVert^2
    + \lVert\evec_{l,i}^{\INP}\rVert^2
\end{equation*}
because $\evec^{\OOP}$ and $\evec^{\INP}$ are orthogonal by
construction. Consequently $\alpha_{\OOP} + \alpha_{\INP} = 1$
exactly, with no rounding error beyond machine precision. The unit
test \fn{test\_mixed\_oop\_inp\_satisfies\_pythagoras} verifies this
on a random eigenvector to $< 10^{-12}$ absolute error.
\end{derivbox}

\begin{workedexample}[OOP/INP for a 4-atom mixed mode]
We compute $\alpha_{\OOP}$ for a synthetic cluster mode that
mixes umbrella ($z$-direction) and breathing ($xy$-plane) motion.
The cluster has $\nhat = (0, 0, 1)$ (the $+\zhat$-axis), and the
4 cluster-atom eigenvector components are
\begin{center}
\begin{tabular}{lccc}
\toprule
Atom    & $e_x$    & $e_y$    & $e_z$ \\
\midrule
$\text{Fe}_1$ & $+0.30$  & $+0.20$  & $+0.40$ \\
$\text{Fe}_2$ & $-0.30$  & $-0.20$  & $+0.40$ \\
$\text{S}_1$  & $+0.20$  & $+0.30$  & $-0.30$ \\
$\text{S}_2$  & $-0.20$  & $-0.30$  & $-0.30$ \\
\bottomrule
\end{tabular}
\end{center}

\textbf{Step 1 -- decompose each atom:}
The OOP component of atom $i$ is $(\evec_i \cdot \nhat)\,\nhat
= (e_z)_i \cdot (0,0,1)$. So
$\lVert\evec_i^{\OOP}\rVert^2 = (e_z)_i^2$, and
$\lVert\evec_i\rVert^2 = e_x^2 + e_y^2 + e_z^2$.

\begin{center}
\begin{tabular}{lccccc}
\toprule
Atom & $e_x^2$ & $e_y^2$ & $e_z^2$ & $\lVert\evec\rVert^2$
                                    & $\lVert\evec^{\OOP}\rVert^2$ \\
\midrule
$\text{Fe}_1$ & 0.09 & 0.04 & 0.16 & 0.29 & 0.16 \\
$\text{Fe}_2$ & 0.09 & 0.04 & 0.16 & 0.29 & 0.16 \\
$\text{S}_1$  & 0.04 & 0.09 & 0.09 & 0.22 & 0.09 \\
$\text{S}_2$  & 0.04 & 0.09 & 0.09 & 0.22 & 0.09 \\
\midrule
\textbf{Sum}  &      &      &      & 1.02 & 0.50 \\
\bottomrule
\end{tabular}
\end{center}

\textbf{Step 2 -- ratio:}
\[
  \alpha_{\OOP} = \frac{0.50}{1.02} \cdot 100\%
    = \boxed{49.02\,\%},
  \qquad
  \alpha_{\INP} = 100\% - 49.02\% = 50.98\%.
\]
With $\theta_{\OOP} = 0.7$ (the default \cfgkey{oop\_threshold}),
this mode is classified as \texttt{mixed} (\cref{eq:oop-binary}),
and the 8-bin detailed classifier returns
\texttt{mixed\_inp} (50--70\% INP range).

\textbf{Pythagorean cross-check:}
$\alpha_{\OOP} + \alpha_{\INP} = 49.02 + 50.98 = 100.0\%$ exact.

\textbf{Note on normalisation:}
The sum $\lVert\evec\rVert^2_\text{cluster} = 1.02$ is slightly
above 1, which is normal because the cluster atoms carry only
part of the mode's total amplitude (other atoms -- ligands,
backbone -- carry the rest, summing to $\approx 1$ over the
whole molecule).

\textbf{In the code:}
\fn{core.classify\_oop\_inp(evec\_4x3, normal=[0,0,1], fe\_c, s\_c,
cfg)} returns a dict containing
\texttt{a\_oop\_pct = 49.02},
\texttt{a\_inp\_pct = 50.98},
\texttt{mode\_type = "mixed"},
\texttt{mode\_type\_detail = "mixed\_inp"}.
\end{workedexample}

\subsection{Choice of normalisation: geometric Cartesian fraction}
\label{sec:oop-cartesian-cavet}

The numerator and denominator of \cref{eq:oop-derivation} are both
unweighted sums of squared cartesian-displacement components. They
\emph{do not} carry atomic masses. In the literature one sometimes
sees an alternative definition
\begin{equation}
  \alpha_{\OOP}^{\text{KE}}(G; l)
    \;=\;
  \frac{\sum_{i \in G} m_i \lVert\evec_{l,i}^{\OOP}\rVert^2}
       {\sum_{i \in G} m_i \lVert\evec_{l,i}\rVert^2}
  \label{eq:oop-kinetic}
\end{equation}
where each squared norm is weighted by atomic mass. For a
mass-weighted eigenvector, $\alpha_{\OOP}^{\text{KE}}$ is the
fraction of \emph{kinetic energy} flowing out of plane:
\begin{equation*}
  T_{\OOP}/T_{\text{total}}
    \;=\;
  \frac{\sum_i m_i \lVert\dot\evec_{l,i}^{\OOP}\rVert^2}
       {\sum_i m_i \lVert\dot\evec_{l,i}\rVert^2}
    \;=\;
  \frac{\sum_i m_i \lVert\evec_{l,i}^{\OOP}\rVert^2}
       {\sum_i m_i \lVert\evec_{l,i}\rVert^2}
\end{equation*}
(the time-derivative factors cancel because every mode oscillates
at the same $\omega_l$).

The code uses \cref{eq:oop-derivation} -- the geometric Cartesian
fraction -- not \cref{eq:oop-kinetic}. The reason is operational:
Gaussian's \texttt{hpmodes} writes cartesian-displacement
eigenvectors (\cref{sec:conv-eigvec}), and we want $\alpha_{\OOP}$
to answer the geometric question ``what fraction of the
atom-by-atom displacement points out of plane'', not the energetic
question. For a [2Fe-2S] cluster the two are numerically close
(Fe mass 55.93 amu, S mass 31.97 amu; the mass-weighted version
shifts the answer by no more than a few percent), but for a
chemically dissimilar system (e.g.\ a heavy-metal-light-ligand
complex) the two can differ substantially.

\subsection{Multi-atom aggregation}

\codref{core.classify\_oop\_inp}

For the cluster-core group $G = \{\text{Fe}_1, \text{Fe}_2,
\text{S}_1, \text{S}_2\}$, the OOP fraction in
\cref{eq:oop-derivation} aggregates over all four atoms. For
multi-residue groups defined in \fn{coord\_info.group\_map}, the
aggregation extends over every atom in the group's centre-list.
The numerator and denominator are computed in a single pass,
which is both faster than two passes and exact (no catastrophic
cancellation possible since all summands are non-negative).

\subsection{Threshold-based binary classification}

\codref{core.classify\_oop\_inp}

For diagnostic purposes the cluster-core $\alpha_{\OOP}$ value is
binned into a binary class:
\begin{equation}
  \text{mode\_type}(l)
    \;=\;
  \begin{cases}
    \OOP, & \alpha_{\OOP} \geq \theta_{\OOP},\\
    \INP, & \alpha_{\OOP} \leq 1 - \theta_{\OOP},\\
    \text{mixed}, & \text{otherwise},
  \end{cases}
  \label{eq:oop-binary}
\end{equation}
with the threshold $\theta_{\OOP} = \cfgkey{oop\_threshold}$
(default $0.7$, i.e.\ 70\%). A finer eight-bin classification is
also computed and recorded under the key \fn{mode\_type\_detail}:
\texttt{pure\_oop} (>~95\%), \texttt{strong\_oop} (85--95\%),
\texttt{majority\_oop} (70--85\%), \texttt{mixed\_oop} (50--70\%),
\texttt{mixed\_inp} (30--50\%), \texttt{majority\_inp} (15--30\%),
\texttt{strong\_inp} (5--15\%), \texttt{pure\_inp} (<~5\%).

\subsection{Sanity tests}

\begin{derivbox}[Sanity checks (test\_sanity\_physics.py)]
\begin{itemize}
  \item \fn{test\_pure\_oop\_mode\_yields\_alpha\_oop\_100}: a
        synthetic eigenvector with displacement purely along $\nhat$
        (umbrella mode) yields $\alpha_{\OOP} = 100\%$ exactly.
  \item \fn{test\_pure\_inp\_mode\_yields\_alpha\_oop\_0}: a
        synthetic eigenvector with displacement entirely in the
        cluster plane (breathing mode) yields $\alpha_{\OOP} = 0\%$
        exactly.
  \item \fn{test\_mixed\_oop\_inp\_satisfies\_pythagoras}: for a
        random eigenvector, $\alpha_{\OOP} + \alpha_{\INP} = 1$ to
        $< 10^{-12}$.
\end{itemize}
\end{derivbox}

% -------------------------------------------------------------------------
\section{Cluster-normal via SVD}
\label{sec:cluster-normal}

\subsection{Best-fit plane derivation}

\codref{geometry.cluster\_normal}

For the four cluster-core atom positions
$\{\rvec_i\}_{i=1}^{4} = \{\rvec_{\text{Fe}_1},
\rvec_{\text{Fe}_2}, \rvec_{\text{S}_1}, \rvec_{\text{S}_2}\}$,
we wish to find the plane that best approximates the four points in
the least-squares sense. Equivalently, we seek a unit vector $\nhat$
such that the sum of squared signed distances from the four points
to the plane $\nhat \cdot (\rvec - \bar\rvec) = 0$ is minimised:
\begin{equation}
  \nhat^\star \;=\; \arg\min_{\lVert\nhat\rVert = 1}
    \;\sum_{i=1}^{4} (\nhat \cdot \tilde\rvec_i)^2,
  \qquad
  \tilde\rvec_i \;\equiv\; \rvec_i - \bar\rvec,
  \qquad
  \bar\rvec \;\equiv\; \tfrac{1}{4}\sum_{i=1}^{4}\rvec_i.
  \label{eq:plane-lsq}
\end{equation}
Writing the centred positions as the rows of a $4\times 3$ matrix
$\tilde{\Rmat}$, the objective becomes
$\lVert\tilde{\Rmat}\nhat\rVert^2 = \nhat^\top \tilde{\Rmat}^\top
\tilde{\Rmat}\, \nhat$. By the Courant-Fischer-Weyl min-max theorem,
the minimum (over unit vectors) is the smallest eigenvalue of the
$3\times 3$ matrix $\tilde{\Rmat}^\top \tilde{\Rmat}$, and the
minimiser is the corresponding eigenvector.

The eigendecomposition of $\tilde{\Rmat}^\top \tilde{\Rmat}$ is
read off the SVD of $\tilde{\Rmat}$. Writing
$\tilde{\Rmat} = \Umat\,\Smat\,\Vmat^\top$ with singular values
$\sigma_1 \geq \sigma_2 \geq \sigma_3 \geq 0$, we have
$\tilde{\Rmat}^\top \tilde{\Rmat}
   = \Vmat\,\Smat^2\,\Vmat^\top$, so the smallest-eigenvalue
eigenvector is $\Vmat_{:,3}$:
\begin{equation}
  \boxed{\;\tilde{\Rmat}
    \;=\;
   \Umat\,\Smat\,\Vmat^\top,
   \qquad
   \nhat
    \;=\;
   \pm\,\Vmat_{:,3}\;}
  \codref{geometry.cluster\_normal}
  \label{eq:cluster-normal-svd}
\end{equation}

\subsection{Sign convention}
\label{sec:cluster-normal-sign}

The right-singular vector $\Vmat_{:,3}$ is determined up to a sign by
the SVD. The convention adopted is
\begin{equation}
  \nhat \cdot \zhat \;\geq\; 0,
  \label{eq:nhat-sign-conv}
\end{equation}
i.e.\ the cluster normal points into the upper half-space of the
Gaussian standard orientation. If $\Vmat_{:,3} \cdot \zhat < 0$,
$\nhat$ is flipped. This makes the $\Delta r_\text{bend\_oop}$
sign in the Fe-ligand decomposition (\cref{sec:fe-ligand-decomp})
reproducible across runs.

\subsection{Singular-value spread as planarity diagnostic}

The three singular values $\sigma_1, \sigma_2, \sigma_3$ have a
geometric interpretation: $\sigma_1^2$ is the variance of the
centred points along the first principal axis, $\sigma_2^2$ along
the second, $\sigma_3^2$ along the third (= the cluster-normal
direction). For a perfectly planar four-atom set,
$\sigma_3 = 0$. For real [2Fe-2S] clusters, $\sigma_3$ is typically
$< 0.05$ \AA{} (the rhombic Fe$_2$S$_2$ core is genuinely planar
within DFT precision), and the ratio $\sigma_3/\sigma_1$ provides
a planarity diagnostic that is recorded in the run-log.

\begin{workedexample}[Cluster-normal SVD for a near-planar Fe$_2$S$_2$ core]
We compute $\nhat$ for a typical rhombic Fe$_2$S$_2$ core, using
representative coordinates from a DFT-optimised [2Fe-2S] cluster
(the slight $z$-perturbation reflects the typical out-of-plane
spread of ${\sim}0.02$~\AA{} observed in any real cluster and
makes the SVD non-degenerate):
\begin{center}
\begin{tabular}{lccc}
\toprule
Atom    & $x$~[\AA{}]  & $y$~[\AA{}]   & $z$~[\AA{}]    \\
\midrule
$\text{Fe}_1$ & $+1.350$ & $0.000$  & $+0.012$ \\
$\text{Fe}_2$ & $-1.350$ & $0.000$  & $+0.018$ \\
$\text{S}_1$  & $0.000$  & $+1.075$ & $-0.020$ \\
$\text{S}_2$  & $0.000$  & $-1.075$ & $-0.010$ \\
\bottomrule
\end{tabular}
\end{center}

\textbf{Step 1 -- centroid and centred positions:}
\[
  \bar\rvec = \tfrac{1}{4}(1.35 - 1.35 + 0 + 0,\,0 + 0 + 1.075 - 1.075,\,
                 0.012 + 0.018 - 0.020 - 0.010)
   = (0,\,0,\,0).
\]
The centroid is at the origin (by construction in standard
orientation). The centred-position matrix
$\tilde{\Rmat} \in \R^{4 \times 3}$ is then identical to the
position matrix above.

\textbf{Step 2 -- SVD:}
The Gram matrix $\Hmat = \tilde{\Rmat}^\top \tilde{\Rmat}$ is
\[
\Hmat = \begin{pmatrix}
2 \cdot 1.350^2 & 0 & 0 \\
0 & 2 \cdot 1.075^2 & 0 \\
0 & 0 & 0.000888
\end{pmatrix}
\;=\;
\begin{pmatrix}
3.645 & 0 & 0 \\
0 & 2.311 & 0 \\
0 & 0 & 0.000888
\end{pmatrix}
\]
(the off-diagonals are zero because the atom positions decouple
in $x$, $y$, $z$ for this geometry). Eigenvalues:
$\lambda_1 = 3.645, \lambda_2 = 2.311, \lambda_3 = 8.88 \times 10^{-4}$.
Singular values:
$\sigma_1 = \sqrt{3.645} = 1.909$~\AA{},
$\sigma_2 = 1.520$~\AA{},
$\sigma_3 = 0.0298$~\AA{}.

\textbf{Step 3 -- read off $\nhat$:}
The right-singular vector for $\sigma_3$ is $\Vmat_{:,3} = (0, 0, \pm 1)$.

\textbf{Step 4 -- enforce sign convention.}
Check $\Vmat_{:,3} \cdot \zhat$: if SVD returns $(0,0,+1)$, keep
it; if $(0,0,-1)$, flip. Result: $\boxed{\nhat = (0,\,0,\,+1)}$.

\textbf{Planarity diagnostic:}
$\sigma_3/\sigma_1 = 0.0298/1.909 = 1.56\%$. The cluster is highly
planar; the run-log records this ratio. Values above $5\%$ would
indicate a strongly tetrahedrally-distorted cluster (e.g.\ a
super-reduced [4Fe-4S]-like distortion accidentally tripped by
\cref{sec:algo-cluster-discovery}).

\textbf{In the code:}
\fn{geometry.cluster\_normal([Fe1, Fe2, S1, S2])} returns
\texttt{normal = [0.0, 0.0, 1.0]} and additionally exposes the
three singular values via the diagnostic key
\texttt{cluster\_planarity\_ratio = 0.0156}.
\end{workedexample}

\subsection{Multi-cluster considerations}

\codref{geometry.find\_all\_clusters}

For a system with $K$ \mbox{[2Fe-2S]} clusters, each cluster runs
\fn{cluster\_normal} on its own four atoms. The normals are
\emph{not} forced to align across clusters; each cluster's normal
is fixed independently by the
$\nhat \cdot \zhat \geq 0$ rule. The pipeline's per-cluster
output directories carry the cluster index so the user can identify
which cluster a given $\nhat$ belongs to.

% -------------------------------------------------------------------------
\section{Classical bond modulation $\Delta r_{ab}$}
\label{sec:dr-classical-derivation}

\subsection{Linearisation of a bond length under small displacements}

\codref{reorganization.signed\_dr\_along\_axis}

Consider two atoms $a, b$ with equilibrium positions $\rvec_a, \rvec_b$
and small displacements $\evec_a, \evec_b$. The bond length under
displacement is
\begin{equation}
  r_{ab}(\evec) = \lVert (\rvec_a + \evec_a) - (\rvec_b + \evec_b)\rVert,
  \qquad
  r_{ab}^{(0)} = \lVert\rvec_a - \rvec_b\rVert.
\end{equation}
Expand around $\evec = 0$:
\begin{equation}
  r_{ab}(\evec)
    \;=\;
  \lVert \rvec_a - \rvec_b + (\evec_a - \evec_b)\rVert
    \;=\;
  r_{ab}^{(0)}\sqrt{1 + 2(\evec_a - \evec_b)\cdot\nhat_{ab}/r_{ab}^{(0)}
       + \lVert\evec_a - \evec_b\rVert^2/(r_{ab}^{(0)})^2}.
\end{equation}
With $\nhat_{ab} = (\rvec_a - \rvec_b)/r_{ab}^{(0)}$ and to first
order in the displacement,
\begin{equation}
  r_{ab}(\evec) - r_{ab}^{(0)}
    \;\approx\;
  (\evec_a - \evec_b) \cdot \nhat_{ab}.
  \label{eq:dr-firstorder}
\end{equation}
This is the linearised bond-modulation formula. Defining
\begin{equation}
  \boxed{\;\Delta r_{ab}^\text{(class.)}
    \;\equiv\;
   (\evec_a - \evec_b)\cdot\nhat_{ab}\;}
  \codref{reorganization.signed\_dr\_along\_axis}
  \label{eq:dr-classical-final}
\end{equation}
recovers the formula used throughout the code.

\subsection{Sign and physical interpretation}

The convention $\nhat_{ab} = (\rvec_a - \rvec_b)/r_{ab}^{(0)}$ makes
$\nhat_{ab}$ point from $b$ to $a$. Therefore:
\begin{itemize}
  \item $\Delta r_{ab} > 0$: atom $a$ moves in $+\nhat_{ab}$ direction
        relative to atom $b$, i.e.\ they move apart -- the bond
        \emph{stretches}.
  \item $\Delta r_{ab} < 0$: the bond \emph{contracts}.
  \item $\Delta r_{ab} = 0$: motion is either zero or purely
        perpendicular to the bond.
\end{itemize}

\begin{derivbox}[Sanity checks (test\_sanity\_physics.py)]
\begin{itemize}
  \item \fn{test\_pure\_translation\_mode\_yields\_zero\_bond\_modulation}:
        for any bond, a uniform translation $\evec_a = \evec_b$
        gives $\Delta r_{ab} = 0$ exactly.
  \item \fn{test\_pure\_rotation\_mode\_yields\_zero\_bond\_stretch}:
        a rigid rotation gives $\Delta r_{ab} = 0$ to first order
        (displacements perpendicular to the bond axis).
  \item \fn{test\_pure\_bond\_stretching\_mode\_yields\_full\_modulation}:
        $\evec_a = (+\Delta, 0, 0)$, $\evec_b = (-\Delta, 0, 0)$
        along $\xhat = \nhat_{ab}$ gives $\Delta r_{ab} = 2\Delta$
        exactly.
\end{itemize}
\end{derivbox}

\subsection{Validity of the linear approximation}

The full expression \cref{eq:dr-firstorder} retains a second-order
correction
$\Delta r_{ab}^{(2)} = \tfrac{1}{2}\lVert\evec_a - \evec_b\rVert^2_\perp
  / r_{ab}^{(0)}$
where $\lVert\cdot\rVert_\perp$ is the component perpendicular to
$\nhat_{ab}$. For typical thermally scaled displacements
(\mbox{$\sim$}~0.05 \AA{}) on a typical Fe-S bond
($r_{ab}^{(0)} \approx 2.2$ \AA{}), the second-order term is
$\sim 5 \times 10^{-4}$ \AA{}, which is smaller than the
hpmodes-precision noise floor of $\sim 5 \times 10^{-4}$ \AA{}
times $\urms$. The linearisation is therefore well within the
precision budget of the input data.

\subsection{Application to the five Marcus-Hush channels}

The classical formula \cref{eq:dr-classical-final} is used directly
for four of the five channels:
\begin{itemize}
  \item $X = \text{FeFe}$: $a = \text{Fe}_1$, $b = \text{Fe}_2$, one
        pair per cluster.
  \item $X = \text{FeN}$: $a = \text{Fe}$, $b = \text{N}$, one pair
        per Fe-N(His) coordination.
  \item $X = \text{FeS}$: $a = \text{Fe}$, $b = \text{S}$, one pair
        per Fe-S coordination (both bridging Fe-S$_{\mu_2}$ and
        terminal Fe-S$_{\text{Cys}}$).
  \item $X = \text{NH}$: $a = \text{N}_\text{His}$, $b = \text{H}$,
        one pair per protonated His-N-H bond.
\end{itemize}
For the fifth channel $X = \text{HA}$ (hydrogen-acceptor), the
classical formula would conflate genuine proton-transfer motion
with skeletal motion of the acceptor. The reaction-coordinate
variant \cref{sec:dr-ha} is used instead.

\begin{workedexample}[Fe-S bond modulation for a 320 cm$^{-1}$ mode]
We compute $\Delta r_{\text{Fe}_1\text{-}\text{S}_1}$ for the
synthetic ``symmetric Fe-S stretch'' mode using the equilibrium
geometry from the cluster-normal example
(\cref{sec:cluster-normal}) and a stretch-like eigenvector.

\textbf{Setup:}
\begin{center}
\begin{tabular}{lcc}
\toprule
Atom    & $\rvec$~[\AA{}]                  & $\evec$ (mass-unweighted) \\
\midrule
$\text{Fe}_1$ & $(+1.350,\,0,\,+0.012)$ & $(+0.42,\,0,\,0)$ \\
$\text{S}_1$  & $(0,\,+1.075,\,-0.020)$ & $(0,\,+0.33,\,0)$ \\
\bottomrule
\end{tabular}
\end{center}

\textbf{Step 1 -- Bond geometry:}
\[
  \rvec_{\text{Fe}_1} - \rvec_{\text{S}_1}
   = (+1.350,\,-1.075,\,+0.032),
   \qquad
  \lVert\cdot\rVert
   = \sqrt{1.823 + 1.156 + 0.001}
   = 1.726~\text{\AA{}}.
\]
Bond unit vector:
$\bhat = (1.350, -1.075, +0.032)/1.726 = (+0.782, -0.623, +0.019)$.

\textbf{Step 2 -- Displacement difference:}
\[
  \Delta\evec
   = \evec_{\text{Fe}_1} - \evec_{\text{S}_1}
   = (+0.42 - 0,\,0 - 0.33,\,0 - 0)
   = (+0.42,\,-0.33,\,0).
\]

\textbf{Step 3 -- Project on bond axis (raw, before $\urms$):}
\[
  \Delta r_\text{stretch}^\text{(raw)}
   = \Delta\evec \cdot \bhat
   = (0.42)(0.782) + (-0.33)(-0.623) + (0)(0.019)
\]
\[
   = 0.328 + 0.206 + 0
   = +0.534
   ~~\text{(dimensionless eigenvector units)}.
\]

\textbf{Step 4 -- Apply thermal scaling (80 K, $\urms = 0.0407$~\AA{}):}
\[
  \Delta r_\text{stretch}
   = 0.534 \cdot 0.0407~\text{\AA}
   = \boxed{+0.0217~\text{\AA{}}}.
\]
Positive sign: Fe$_1$ moves \emph{away} from S$_1$ along $\bhat$ -- the
bond is stretching in the upstroke of this oscillation. (At a phase
half-period later the sign would flip; the reported magnitude is the
amplitude of the oscillation.)

\textbf{Step 5 -- Verify orthogonality of bend.}
The perpendicular component $\Delta\evec - 0.534\,\bhat
= (0.42 - 0.417, -0.33 + 0.333, -0.010)
= (+0.003, +0.003, -0.010)$.
Its norm is $0.0108$, contributing the bend\_inp/bend\_oop. The
Pythagorean check
$0.534^2 + 0.0108^2 = 0.285 + 0.000117 \approx 0.285
= \lVert\Delta\evec\rVert^2$ confirms decomposition.

\textbf{Sigma:}
$\sigma_\text{stretch} = \sqrt{2} \cdot 5\!\times\!10^{-4} \cdot
0.0407~\text{\AA} = 2.88 \times 10^{-5}$~\AA{}, so the stretch is
$0.0217/2.88\!\times\!10^{-5} \approx 750\sigma$ significant --
deep above the noise floor.

\textbf{In the code:}
\fn{reorganization.signed\_dr\_along\_axis(evg\_Fe1, evg\_S1,
r\_Fe1, r\_S1)} returns \texttt{0.0217} (the function takes
already-scaled eigenvectors, hence the result is in \AA{} directly).
\end{workedexample}

% -------------------------------------------------------------------------
\section{Reaction-coordinate $\Delta r_{HA}$}
\label{sec:dr-ha}

\subsection{Why a different formula for HA}

The proton-transfer (PT) reaction coordinate is, in the spirit of
Marcus-Hush theory, the displacement of the proton along the
N--A axis (where $A$ is the H-bond acceptor, typically a backbone or
side-chain carbonyl-O or amine-N). The classical formula
\cref{eq:dr-classical-final} applied to the pair (H, A) would read
$(\evec_H - \evec_A)\cdot\nhat_{HA}$, which captures the change in
H-to-A distance. But that distance changes both when the proton moves
(genuine PT character) and when the acceptor moves while the
proton is stationary (skeletal mode, no PT character).

For PCET analysis we want a coordinate that:
\begin{itemize}
  \item is nonzero only when the proton moves relative to its
        \emph{donor} (the N$_\text{His}$);
  \item points along the N-to-A axis (so that the sign distinguishes
        PT toward A from PT away from A).
\end{itemize}

The unique linear functional of $(\evec_N, \evec_H, \evec_A)$
satisfying both is
\begin{equation}
  \boxed{\;\Delta r_{HA}^\text{(react.\ coord.)}
    \;\equiv\;
  (\evec_H - \evec_N)\cdot\nhat_{NA},
  \qquad
  \nhat_{NA} = \frac{\rvec_A - \rvec_N}{\lVert\rvec_A - \rvec_N\rVert}\;}
  \codref{reorganization.compute\_mode\_modulations}
  \label{eq:dr-ha-final}
\end{equation}

The displacement subtracted is $\evec_N$ (the donor), not $\evec_A$.
The projection axis is $\nhat_{NA}$ (donor-to-acceptor), not
$\nhat_{HA}$. Both choices are deliberate; combined, they enforce
the two requirements above.

\subsection{Limit behaviours}

\begin{derivbox}[The PT vs.\ skeletal-mode discriminator (test\_sanity\_physics.py)]
\begin{itemize}
  \item \fn{test\_ha\_pure\_acceptor\_motion\_yields\_zero}:
        $\evec_N = \evec_H = 0$, $\evec_A \neq 0$ gives
        $\Delta r_{HA} = 0$. No PT character.
  \item \fn{test\_ha\_pure\_h\_motion\_toward\_acceptor\_yields\_positive}:
        $\evec_N = \evec_A = 0$, $\evec_H = +d\,\nhat_{NA}$ gives
        $\Delta r_{HA} = +d$. Genuine PT toward the acceptor.
  \item \fn{test\_ha\_pure\_n\_motion\_toward\_acceptor\_yields\_negative}:
        $\evec_H = \evec_A = 0$, $\evec_N = +d\,\nhat_{NA}$ gives
        $\Delta r_{HA} = -d$. \emph{Skeletal} mode (the N-H unit
        translates as a block toward A, the H-A distance shrinks
        but the proton has not moved relative to its donor).
\end{itemize}
The PT/skeletal-discriminator property is the principal reason for
choosing the reaction-coordinate definition over the naive
$(\evec_H - \evec_A) \cdot \nhat_{HA}$.
\end{derivbox}

\subsection{Choice of acceptor: nearest-neighbour heuristic}

\codref{reorganization.build\_channels\_from\_coord\_info}

For each protonated His-N-H, the acceptor $A$ is selected as the
nearest electronegative atom (N, O, S) within a configurable cutoff
\cfgkey{pcet\_max\_distance\_a} (default $3.5$~\AA{}) of the H atom,
subject to:
\begin{itemize}
  \item the candidate is not the donor itself,
  \item the candidate is not bonded (within $1.5$~\AA{}) to the
        donor,
  \item the candidate has a valid Gaussian centre and PDB residue
        assignment.
\end{itemize}
If no valid candidate is found, the H is recorded with
$A = $ \texttt{None}, $\Delta r_{HA} = 0$ for all modes, and a
\fn{UserWarning} is emitted on the first such occurrence in the run.

\subsection{Geometric weighting by H\texorpdfstring{$\cdots$}{...}A distance}

\codref{reorganization.build\_channels\_from\_coord\_info}

When a protein has multiple plausible H-bond acceptors, the code
sums their HA-contributions with a gaussian weight
$w_a = \exp(-(d_a - d_0)^2 / (2 \sigma_d^2))$,
$d_a = \lVert\rvec_H - \rvec_{A_a}\rVert$,
$d_0 = $ \cfgkey{pcet\_optimal\_distance\_a} (default $2.8$~\AA{}),
$\sigma_d = $ \cfgkey{pcet\_distance\_sigma\_a} (default $0.3$~\AA{}).
This makes the HA channel a weighted aggregate
\begin{equation}
  \Delta r_{HA}^\text{aggregated}
    \;=\;
  \sqrt{\sum_a w_a \cdot (\Delta r_{HA,a})^2}
  \label{eq:dr-ha-aggregated}
\end{equation}
which is the RSS aggregation rule of \cref{sec:rss-aggregation} applied
to the HA sub-channels.

\begin{workedexample}[The PT/skeletal discriminator on three motion patterns]
We compute $\Delta r_{HA}$ for three test eigenvectors on the
same His-N-H$\cdots$A geometry, to demonstrate the PT-vs-skeletal
discrimination property.

\textbf{Setup (geometry):}
\begin{center}
\begin{tabular}{lc}
\toprule
Atom & $\rvec$~[\AA{}] \\
\midrule
N (donor)     & $(0.00,\,0.00,\,0.00)$ \\
H (proton)    & $(0.10,\,0.00,\,1.00)$ \\
A (acceptor)  & $(0.30,\,0.20,\,2.80)$ \\
\bottomrule
\end{tabular}
\end{center}

Donor-to-acceptor vector:
$\rvec_A - \rvec_N = (0.30, 0.20, 2.80)$,
$\lVert\rvec_A - \rvec_N\rVert = 2.823$, so
$\nhat_{NA} = (0.1063, 0.0709, 0.9920)$
($\approx$ pointing along $+\zhat$, as expected for an in-line
H-bond).

\textbf{Test eigenvector 1 -- pure H motion toward acceptor:}
$\evec_N = (0, 0, 0)$, $\evec_H = (0.0106, 0.0071, 0.0992)$
(magnitude 0.10 along $\nhat_{NA}$), $\evec_A = (0, 0, 0)$.
\[
  \Delta r_{HA}
  = (\evec_H - \evec_N) \cdot \nhat_{NA}
  = (0.0106)(0.1063) + (0.0071)(0.0709) + (0.0992)(0.9920)
\]
\[
  = 0.00113 + 0.00050 + 0.09841
  = +0.1000.
\]
\textbf{Result: $+0.100$~\AA{} (full PT signal).} Sign positive
because H moves toward A.

\textbf{Test eigenvector 2 -- pure N motion toward A (skeletal):}
$\evec_N = (0.0106, 0.0071, 0.0992)$ (same magnitude 0.10 along
$\nhat_{NA}$), $\evec_H = (0, 0, 0)$, $\evec_A = (0, 0, 0)$.

Now the N-H bond \emph{translates as a rigid block} toward A.
The H-A distance shrinks (an L1-style ``HA modulation''
$|\evec_H - \evec_A| \cdot \nhat_{HA}$ would yield a large
positive number), but the proton has not actually moved
relative to its donor.
\[
  \Delta r_{HA}
  = (\evec_H - \evec_N) \cdot \nhat_{NA}
  = (-0.0106)(0.1063) + (-0.0071)(0.0709) + (-0.0992)(0.9920)
\]
\[
  = -0.1000.
\]
\textbf{Result: $-0.100$~\AA{} (skeletal motion correctly flagged).}
Sign \emph{negative}, signalling that this is the opposite of the
PT direction. A reviewer looking at the sign immediately knows
this mode is \emph{not} a PT-promoting mode -- the geometric mean
$C_\PCET(\omega) = \sqrt{M_\text{HA} \cdot M_\text{FeFe}}$ would
not classify this mode as PCET.

\textbf{Test eigenvector 3 -- pure A motion (acceptor wagging):}
$\evec_N = (0, 0, 0)$, $\evec_H = (0, 0, 0)$, $\evec_A = (0.05, 0, 0)$
(magnitude 0.05 along $\xhat$, perpendicular to the H-bond).
\[
  \Delta r_{HA}
  = (\evec_H - \evec_N) \cdot \nhat_{NA}
  = (0 - 0)(\nhat_{NA})_x + 0 + 0
  = 0.
\]
\textbf{Result: 0 (acceptor wagging correctly ignored).}
A skeletal mode of the acceptor that doesn't involve the H is
not a PT mode at all.

\textbf{Comparison with the naive formula:}
The discarded alternative
$(\evec_H - \evec_A) \cdot \nhat_{HA}$ would give:
\begin{center}
\begin{tabular}{lccc}
\toprule
Test                 & RC formula (correct) & Naive formula \\
\midrule
1. H toward A        & $+0.100$ \checkmark     & $+0.099$ (close) \\
2. N+H toward A      & $-0.100$ \checkmark     & $+0.099$ (\emph{wrong}: skeletal $\to$ PT) \\
3. A wags perp       & $0.000$ \checkmark      & $-0.049$ (\emph{wrong}: pure-A motion $\to$ PT) \\
\bottomrule
\end{tabular}
\end{center}
The reaction-coordinate formula passes all three; the naive
formula fails 2 and 3. This is the discriminator property cited
in \cref{sec:dr-ha} -- it is what makes the HA channel a
meaningful PCET indicator rather than a generic ``H-bond wiggle''
indicator.

\textbf{In the code:}
\fn{reorganization.compute\_mode\_modulations} computes $\Delta r_{HA}$
via the channels of type \texttt{HA}, using the donor-to-acceptor
axis stored as \fn{ChannelGeometry.r\_donor} (see
\cref{sec:dataclass-channelgeometry}).
\end{workedexample}

% -------------------------------------------------------------------------
\section{Fe-ligand bond decomposition}
\label{sec:fe-ligand-decomp}

\subsection{Three orthogonal degrees of freedom}

\codref{core.analyze\_fe\_ligand}

For an Fe-ligand bond (with the ligand donor atom $L \in \{S, N, O\}$),
the relative displacement is
\begin{equation}
  \Delta\evec \;\equiv\; \evec_L - \evec_{\text{Fe}}.
\end{equation}
This 3-vector contains all the information about how the bond
deforms in the given mode. We decompose it into three physically
distinct degrees of freedom:

\paragraph{Stretch along the bond axis.}
\begin{equation}
  \Delta r_\text{stretch} \;\equiv\; \Delta\evec \cdot \bhat,
  \qquad
  \bhat = \frac{\rvec_L - \rvec_{\text{Fe}}}{\lVert\rvec_L - \rvec_{\text{Fe}}\rVert}.
  \label{eq:fel-stretch}
\end{equation}
This is the signed bond-stretching component (positive = bond
lengthening).

\paragraph{Bend in the cluster plane.}
The bend is, by definition, the component of $\Delta\evec$
\emph{perpendicular} to $\bhat$:
\begin{equation}
  \Delta\evec_\perp
    \;\equiv\;
  \Delta\evec - \Delta r_\text{stretch}\,\bhat.
  \label{eq:fel-perp}
\end{equation}
This perpendicular vector lives in the 2D plane orthogonal to the
bond. We further decompose it relative to the cluster normal
$\nhat$:
\begin{align}
  \Delta r_\text{bend\_oop}
    &\;\equiv\; \Delta\evec_\perp \cdot \nhat,
      \qquad\text{(signed)} \\
  \Delta r_\text{bend\_inp}
    &\;\equiv\;
   \lVert\Delta\evec_\perp - \Delta r_\text{bend\_oop}\,\nhat\rVert.
   \label{eq:fel-bend-split}
\end{align}
The first projection is signed; the second is a length, so it is
non-negative. The Pythagorean identity
\begin{equation}
  \lVert\Delta\evec_\perp\rVert^2
    \;=\; (\Delta r_\text{bend\_oop})^2 + (\Delta r_\text{bend\_inp})^2.
  \label{eq:fel-pythagoras}
\end{equation}
follows immediately from the orthogonality of the cluster-normal to
the in-plane complement.

\paragraph{Why $\nhat$ as the bend-OOP axis.}
The cluster normal $\nhat$ from \cref{sec:cluster-normal} is the
natural OOP axis for a [2Fe-2S] cluster: it points perpendicular to
the rhombic Fe$_2$S$_2$ plane and is the same axis used by the
cluster-core OOP/INP classification (\cref{sec:oop-derivation}).
This choice makes the Fe-ligand bend\_oop / bend\_inp directly
comparable to the cluster-core OOP/INP fractions: a mode that is
cluster-OOP-dominant will tend to have Fe-ligand bend\_oop also
dominant.

\codref{core.analyze\_fe\_ligand}

\subsection{Percentages of the total bend}

When comparing modes with very different total amplitudes, the bend
components are converted to percentage shares of the total bend
power:
\begin{equation}
  \text{bend\_inp\_pct}
    \;=\;
  100 \cdot \frac{(\Delta r_\text{bend\_inp})^2}{
                (\Delta r_\text{bend\_inp})^2 + (\Delta r_\text{bend\_oop})^2},
\end{equation}
and analogously for \texttt{bend\_oop\_pct}. By
\cref{eq:fel-pythagoras} these sum to $100\%$ exactly.

For modes with very small total bend amplitude (below the
noise threshold $2\sigma_e^\text{therm}\sqrt{2}$;
\cref{sec:sigma-propagation}) the percentage is set to NaN, so
numerical noise does not get amplified into spurious
classifications.

\subsection{Sigma propagation}

Each of $\Delta r_\text{stretch}, \Delta r_\text{bend\_inp},
\Delta r_\text{bend\_oop}$ is a difference of two independent
Gaussian-noise components, so each carries the same
\begin{equation}
  \sigma \;=\; \sqrt{2}\,\sigma_e \cdot \urms(\omega_l)
  \label{eq:fel-sigma}
\end{equation}
where $\sigma_e = \cfgkey{sigma\_eigvec}$ is the per-component
eigenvector noise (default $5 \times 10^{-4}$). The $\sqrt{2}$
comes from variance addition for a difference; the $\urms$ factor
is the thermal-scaling convention (\cref{sec:urms-derivation}). These
$\sigma$ values are recorded in the columns \texttt{s\_stretch},
\texttt{s\_bend\_inp}, \texttt{s\_bend\_oop} of the
\file{Fe\_S\_Cys.xlsx} / \file{Fe\_N\_His.xlsx} workbooks.

\subsection{Physical interpretation for PCET studies}

The Fe-N(His) bend decomposition is particularly informative:

\begin{itemize}
  \item \emph{bend\_inp dominant} ($> 70\%$): the His ring tilts in
        the cluster plane. This modulates the Fe-Fe distance and
        the Fe-S core, which is the classical electron-transfer-
        active geometry.
  \item \emph{bend\_oop dominant} ($> 70\%$): the His ring tilts out
        of plane. This modulates the orientation of the N-H bond
        and therefore the H$\cdots$A donor-acceptor axis, which is
        the proton-transfer-coupling channel.
  \item \emph{mixed} (30--70\%): the mode contributes simultaneously
        to ET and PT pathways -- the canonical PCET signature.
\end{itemize}

The Fe-S(Cys) bend decomposition follows the same logic, with the
caveat that the S-C bond to the protein backbone is more flexible
than the N-C bond, so the bend\_oop component of Fe-S(Cys) often
carries longer-range allosteric coupling than the in-plane component.

% -------------------------------------------------------------------------
\section{His-NH stretching}
\label{sec:his-nh}

\subsection{Definition}

\codref{core.analyze\_his\_hn}

For each protonated histidine, the N-H bond modulation is computed
with the classical formula
\cref{eq:dr-classical-final}, with $a$ = N$_\text{His}$,
$b$ = H:
\begin{equation}
  \Delta r_\text{NH}
    \;=\;
  (\evec_N - \evec_H) \cdot \nhat_{NH},
  \qquad
  \nhat_{NH} \;=\; \frac{\rvec_N - \rvec_H}{\lVert\rvec_N - \rvec_H\rVert}.
\end{equation}

\codref{core.analyze\_his\_hn}

\subsection{Hydrogen inclusion}

The N-H stretching analysis requires the hydrogen eigenvector
components, which the default heavy-atom-only filter
(\cref{sec:conv-heavy}) discards. The analysis is therefore gated
by \cfgkey{include\_hn\_vibration}:
\begin{itemize}
  \item \cfgkey{include\_hn\_vibration} = true: a hydrogen-inclusive
        eigenvector loader is used; $\Delta r_\text{NH}$ is computed
        for every protonated His.
  \item \cfgkey{include\_hn\_vibration} = false (default): the
        His-NH stretch sheet is still produced but all
        $\Delta r_\text{NH}$ values are zero, and the sheet's
        existence serves as a placeholder for downstream tools.
\end{itemize}

\subsection{Sigma propagation -- the v1.0.4 fix}

\codref{core.analyze\_his\_hn}

Prior to v1.0.4, the sigma for $\Delta r_\text{NH}$ used the raw
$\sqrt{2}\,\sigma_e$ instead of the thermally scaled
$\sqrt{2}\,\sigma_e \cdot \urms$. For typical low-frequency modes
this made the recorded \texttt{s\_hn\_stretch} a factor of
$1/\urms \approx 1/0.05 = 20\times$ \emph{too large}, which forced
many genuine PCET signals below the visual significance threshold.
The v1.0.4 fix (FUND 6, \cref{sec:fund-6-8}) applies the thermal
scaling
$\sigma_e^\text{therm} = \sigma_e \cdot \urms(\omega_l)$, the
convention discussed in detail in
\cref{sec:sigma-propagation}; the recorded sigmas are now $\sim 20$
times smaller than in v1.0.2/v1.0.3, but the underlying values
$\Delta r_\text{NH}$ are unchanged.

% -------------------------------------------------------------------------
\section{Pair-wise reorganization $\lambda_X^{\text{pair}}$}
\label{sec:lambda-derivation}

\subsection{From harmonic oscillator energy to cm$^{-1}$}

\codref{reorganization.lambda\_pair\_cm1}

For a single bond with effective reduced mass $\mu$ and a single
normal mode of frequency $\omega$, the harmonic-oscillator energy
contribution of a bond-length displacement $\Delta r$ is
\begin{equation}
  E \;=\; \tfrac{1}{2}\,\mu\,\omega^2\,(\Delta r)^2 \quad [\text{Joule}].
  \label{eq:lambda-joule}
\end{equation}
In Marcus-Hush theory, this expression -- evaluated at the
displacement that takes the system from reactant to product
geometry -- is half the reorganization energy
$\lambda$ along that mode. For \emph{spectroscopic} bond
reorganization in a frozen-environment setting, we instead take
$\Delta r$ to be the thermally averaged amplitude of the mode
along the bond, $\Delta r = \urms(\omega) \cdot \alpha_X$ where
$\alpha_X$ is the dimensionless projection of the mode onto the
bond unit-vector. Then \cref{eq:lambda-joule} reads
\begin{equation}
  \lambda_X^{\text{pair}}
    \;=\;
  \tfrac{1}{2}\,\mu\,\omega^2\,(\Delta r_X)^2,
  \label{eq:lambda-joule-channel}
\end{equation}
which is the per-mode bond-reorganization energy along channel $X$.

To convert to spectroscopic cm$^{-1}$ units, divide by
$hc = 1.986\,445\,857 \times 10^{-23}$ J cm:
\begin{equation}
  \boxed{\;\lambda_X^{\text{pair}}[\text{cm}^{-1}]
    \;=\;
  \frac{1}{hc}\,\tfrac{1}{2}\,\mu\,\omega^2\,(\Delta r_X)^2\;}
  \codref{reorganization.lambda\_pair\_cm1}
  \label{eq:lambda-cm1}
\end{equation}

\subsection{Unit conversions used in the implementation}

The function \fn{reorganization.lambda\_pair\_cm1(dr\_a,
omega\_cm1, mu\_amu)} accepts SI-mixed inputs:
\begin{itemize}
  \item $\Delta r$ in \AA{}; internal conversion $\Delta r \cdot
        10^{-10}$ to metres,
  \item $\omega$ in cm$^{-1}$; internal conversion
        $\omega_\text{SI} = 2\pi c \cdot \tilde\nu$ with
        $c = 2.997\,924\,58 \times 10^{10}$ cm/s,
  \item $\mu$ in amu; internal conversion $\mu_\text{SI} =
        \mu_\text{amu} \cdot 1.660\,539\,066\,60 \times 10^{-27}$
        kg.
\end{itemize}
The output is $\lambda$ in cm$^{-1}$.

\subsection{Sanity checks}

\begin{derivbox}[Quadratic and scaling properties (test\_sanity\_physics.py)]
\begin{itemize}
  \item \fn{test\_lambda\_zero\_for\_zero\_displacement}: $\Delta r = 0$
        yields $\lambda = 0$ exactly.
  \item \fn{test\_lambda\_positive\_for\_any\_nonzero\_displacement}:
        $\lambda$ is non-negative and depends only on
        $|\Delta r|$ (sign-symmetric).
  \item \fn{test\_lambda\_quadratic\_in\_displacement}: doubling
        $\Delta r$ quadruples $\lambda$.
  \item \fn{test\_lambda\_quadratic\_in\_frequency}: doubling
        $\omega$ quadruples $\lambda$.
\end{itemize}
\end{derivbox}

\subsection{$T \to 0$ limit for pure stretching modes}

For a pure-stretching mode along bond $X$ with $\alpha_X = 1$, in
the $T \to 0$ limit:
\begin{equation}
  \Delta r_X
    \;=\;
  \urms \cdot 1
    \;\xrightarrow{T \to 0}\;
  \sqrt{\frac{\hbar}{2\mu\omega}},
\end{equation}
so
\begin{equation}
  \lambda_X^{\text{pair}}
    \;=\;
  \tfrac{1}{2}\mu\omega^2 \cdot \frac{\hbar}{2\mu\omega}
    \;=\;
  \frac{\hbar\omega}{4}.
  \label{eq:lambda-zero-T-derivation}
\end{equation}
That is, half the zero-point energy. In cm$^{-1}$ that is just
$\omega/4$ -- the formula is dimensionally exact.

\begin{derivbox}[Sanity check (test\_sanity\_physics.py)]
\fn{test\_lambda\_zero\_T\_pure\_stretching}: with $\omega = 320$
cm$^{-1}$, $\mu = 31.97$ amu, $T = 10^{-4}$ K, $\alpha_X = 1$, the
result matches $\omega/4 = 80.0$ cm$^{-1}$ to $< 10^{-3}$
relative error.
\end{derivbox}

\subsection{Relation to classical Marcus reorganization}

The spectroscopic bond reorganization $\lambda_X^\text{pair}$ in
\cref{eq:lambda-cm1} should not be conflated with the classical
Marcus solvent reorganization energy, even though they share the
$\tfrac{1}{2}\mu\omega^2(\Delta r)^2$ structure. The classical Marcus
quantity uses $\Delta r$ = static geometry shift between reactant
and product (a Born-Oppenheimer property), and the relevant degrees
of freedom are usually intermolecular solvent modes. The
spectroscopic quantity here uses $\Delta r$ = thermal/zero-point
displacement amplitude of an internal vibrational mode, and is a
\emph{frequency-resolved} measurement of which bond modes
``zappel'' the most. They become equal only in the special case
of a single classical promoting mode driving Born-Oppenheimer
hopping; in general the spectroscopic quantity is the more directly
NRVS-observable.

\begin{workedexample}[$\lambda^\text{pair}$ for the 320 cm$^{-1}$ Fe-S stretch]
We compute $\lambda_\text{FeS}^\text{pair}$ for the bond and mode
from the Fe-S stretch example above:
\begin{itemize}
  \item $\Delta r = 0.0217$~\AA{} = $2.17 \times 10^{-12}$~m,
  \item $\tilde\nu = 320$~cm$^{-1}$ $\to$
        $\omega = 6.026 \times 10^{13}$~rad/s (computed above),
  \item $\mu = 31.97$~amu = $5.309 \times 10^{-26}$~kg (using
        the reduced-mass-of-the-mode convention; in a more refined
        treatment one would use the bond-only reduced mass
        $\mu_\text{Fe-S} = m_\text{Fe} m_\text{S} / (m_\text{Fe} +
        m_\text{S}) = 20.34$~amu, but for the per-mode formula
        \cref{eq:lambda-cm1} the Gaussian-reported reduced mass is
        used).
\end{itemize}

\textbf{Step 1 -- $E$ in Joule:}
\begin{align*}
  E &= \tfrac{1}{2}\,\mu\,\omega^2\,(\Delta r)^2 \\
    &= \tfrac{1}{2} \cdot 5.309 \times 10^{-26}\,\text{kg}
       \cdot (6.026 \times 10^{13}\,\text{rad/s})^2
       \cdot (2.17 \times 10^{-12}\,\text{m})^2 \\
    &= \tfrac{1}{2} \cdot 5.309 \times 10^{-26}
       \cdot 3.631 \times 10^{27}
       \cdot 4.709 \times 10^{-24}\,\text{J} \\
    &= 4.539 \times 10^{-22}\,\text{J}.
\end{align*}

\textbf{Step 2 -- convert to cm$^{-1}$:}
\[
  \lambda^\text{pair} = E / (hc)
   = \frac{4.539 \times 10^{-22}}{1.986 \times 10^{-23}}
   = \boxed{22.85~\text{cm}^{-1}}.
\]

\textbf{Cross-check via the $T \to 0$ formula:}
For a pure stretch mode with $\alpha_X = 1$ (the displacement is
entirely along the bond), \cref{eq:lambda-zero-T-derivation} gives
$\lambda^\text{pair} = \omega/4 = 320/4 = 80$~cm$^{-1}$. Our 22.85
cm$^{-1}$ is well below this -- consistent with $\alpha_X
\approx 0.534$ (from the previous worked example), giving
$\lambda^\text{pair}_{\text{predicted}} \approx
80 \cdot 0.534^2 = 22.8$~cm$^{-1}$. The two routes agree
within rounding.

\textbf{In the code:}
\fn{reorganization.lambda\_pair\_cm1(0.0217, 320.0, 31.97)} returns
\texttt{22.85}.
\end{workedexample}

% -------------------------------------------------------------------------
\section{RSS aggregation across sub-channels}
\label{sec:rss-aggregation}

\subsection{Why RSS, not L1}

\codref{reorganization.aggregate\_by\_parent}

A parent channel $X \in \{\text{FeFe}, \text{FeN}, \text{FeS},
\text{NH}, \text{HA}\}$ generally has multiple \emph{sub-channels}
that aggregate to it:
\begin{itemize}
  \item FeN: two Fe-N(His) sub-bonds (one per His residue) for a
        Cys$_2$His$_2$ Rieske cluster; zero for a pure Cys$_4$
        cluster.
  \item FeS: four Fe-S sub-bonds for a Cys$_4$ cluster (two terminal
        Fe-S$_\text{Cys}$ per Fe), or two for a Cys$_2$His$_2$
        cluster, plus the two bridging Fe-S$_{\mu_2}$.
  \item NH: one N-H sub-bond per protonated His.
  \item HA: one H-A sub-bond per protonated His paired with a
        unique acceptor candidate (one His can have multiple).
\end{itemize}

For mode $i$, we have sub-channel modulations
$\{\Delta r_X^{(s)}(i)\}_{s \in \text{sub-channels}(X)}$ with
weights $\{w_s\}$. We need an aggregation rule
$\Delta r_X^\text{aggregated}(i) = f(\{\Delta r_X^{(s)}(i)\},
\{w_s\})$ that produces the parent-channel modulation. Two natural
candidates are:

\begin{enumerate}
  \item \emph{Signed L1 sum}:
        $f = \sum_s w_s \cdot \Delta r_X^{(s)}$. Problem: if two
        sub-bonds modulate in opposite directions (one stretches,
        the other contracts -- a typical antisymmetric mode), they
        cancel, and the parent appears stationary. This is the
        wrong physics: an antisymmetric stretch is still a
        modulation of FeS in the sense of total bond-energy
        change.
  \item \emph{Absolute L1 sum}:
        $f = \sum_s w_s \cdot |\Delta r_X^{(s)}|$. Problem: not
        energetically consistent. A pure antisymmetric stretch
        with $|\Delta r| = 0.05$ \AA{} on two equal bonds gives
        $f = 0.10$ \AA{}, but the energy contribution is
        $2 \cdot \tfrac{1}{2}\mu\omega^2(0.05)^2 \neq
        \tfrac{1}{2}\mu\omega^2(0.10)^2$.
\end{enumerate}

The energetically consistent choice is the \emph{root-sum-square}
(RSS):
\begin{equation}
  \boxed{\;\Delta r_X^\text{RSS}(i)
    \;\equiv\;
  \sqrt{\sum_{s \in \text{sub-channels}(X)}
        w_s\,(\Delta r_X^{(s)}(i))^2}\;}
  \codref{reorganization.aggregate\_by\_parent}
  \label{eq:dr-rss-derivation}
\end{equation}

\subsection{Energetic consistency}

\codref{reorganization.aggregate\_by\_parent}

The RSS aggregation has the property that the parent-channel
$\lambda$ -- defined as the weighted sum of sub-channel
$\lambda$'s -- can be written in the same harmonic-oscillator form
as a sub-channel $\lambda$:
\begin{align}
  \lambda_X^\text{parent}(i)
    &\;=\;
   \sum_s w_s \cdot \tfrac{1}{2}\mu_i\omega_i^2
       (\Delta r_X^{(s)}(i))^2 \\[2pt]
    &\;=\;
   \tfrac{1}{2}\mu_i\omega_i^2 \cdot
   \underbrace{\sum_s w_s (\Delta r_X^{(s)}(i))^2}_{=
     (\Delta r_X^\text{RSS})^2} \\[2pt]
    &\;=\;
   \tfrac{1}{2}\mu_i\omega_i^2 (\Delta r_X^\text{RSS}(i))^2.
  \label{eq:lambda-rss-consistency}
\end{align}
The choice of RSS makes \cref{eq:lambda-cm1} a valid
parent-level formula with $\Delta r$ replaced by
$\Delta r_X^\text{RSS}$. The L1 alternative breaks this
consistency.

\begin{derivbox}[Sanity check (test\_sanity\_physics.py)]
\fn{test\_rss\_aggregation\_two\_subchannels\_consistency\_with\_lambda}:
with two sub-channels at $\Delta r_1 = 0.008$, $\Delta r_2 = -0.012$
\AA{}, $w_1 = w_2 = 1$, $\omega = 350$ cm$^{-1}$, $\mu = 31.97$
amu, the two computations
$\lambda_\text{sum} = \lambda_\text{pair}(\Delta r_1) +
\lambda_\text{pair}(\Delta r_2)$ and
$\lambda_\text{via-RSS} = \lambda_\text{pair}(\Delta r^\text{RSS})$
match to $< 10^{-10}$ relative error.
\end{derivbox}

\subsection{Weights per channel}

\codref{reorganization.build\_channels\_from\_coord\_info}

\begin{itemize}
  \item FeFe, FeN, FeS, NH: weights $w_s = 1$ for all sub-channels
        (each Fe-N or Fe-S bond contributes equally to the parent
        channel's energetic budget).
  \item HA: weights are the gaussian H$\cdots$A-distance weights of
        \cref{sec:dr-ha},
        $w_a = \exp(-(d_a - d_0)^2/(2\sigma_d^2))$. The
        $\sqrt{w_a}$ factor in
        \cref{eq:dr-rss-derivation} downweights distant or weakly
        H-bonded acceptors.
\end{itemize}

\subsection{Single sub-channel limit}

For a parent channel with a single sub-channel ($|s| = 1$),
\cref{eq:dr-rss-derivation} reduces to
\begin{equation}
  \Delta r_X^\text{RSS}
    \;=\;
  \sqrt{w \cdot (\Delta r_X)^2}
    \;=\;
  \sqrt{w}\,|\Delta r_X|.
\end{equation}
For unit weight ($w = 1$, the FeFe and NH cases),
$\Delta r_X^\text{RSS} = |\Delta r_X|$ -- the absolute bond
modulation, lossless against the sign-cancellation of an L1 sum.

\begin{derivbox}[Sanity check (test\_sanity\_physics.py)]
\fn{test\_rss\_aggregation\_single\_subchannel}: with $|\Delta r| =
0.012$ \AA{}, $w = 1$, the RSS aggregate is exactly $0.012$~\AA{}.
\end{derivbox}

\begin{workedexample}[RSS aggregation: FeS parent over 4 sub-bonds]
A Cys$_4$ cluster has 4 Fe-S$_\text{Cys}$ terminal bonds plus 2 bridging
Fe-S$_{\mu_2}$ bonds = 6 sub-channels feeding the FeS parent (or
4 + 4 = 8 depending on whether one counts bridging-S as ``one bond''
or ``two''; the code counts each Fe-S as one bond). For this
worked example we take a typical antisymmetric Fe-S stretch mode
with the per-sub-channel modulations
\begin{center}
\begin{tabular}{lcc}
\toprule
Sub-channel       & $\Delta r_s$ [\AA{}] & $w_s$ \\
\midrule
FeS\_Fe1\_S1\_term  & $+0.0150$ & $1$ \\
FeS\_Fe1\_S2\_term  & $-0.0145$ & $1$ \\
FeS\_Fe2\_S3\_term  & $-0.0148$ & $1$ \\
FeS\_Fe2\_S4\_term  & $+0.0152$ & $1$ \\
\bottomrule
\end{tabular}
\end{center}

\textbf{Note the alternating signs} -- this is a typical
antisymmetric stretch: Fe$_1$-S$_1$ and Fe$_2$-S$_4$ stretch
while Fe$_1$-S$_2$ and Fe$_2$-S$_3$ contract.

\textbf{Step 1 -- Signed L1 sum (the discarded alternative):}
\[
  \Sigma \Delta r = 0.0150 - 0.0145 - 0.0148 + 0.0152
   = +0.0009~\text{\AA{}},
\]
which is \emph{essentially zero} -- the antisymmetric character
of the mode caused cancellation. This is exactly the failure
mode that disqualifies the signed L1 sum.

\textbf{Step 2 -- L1 sum of absolutes (the second discarded alternative):}
\[
  \Sigma |\Delta r| = 0.0150 + 0.0145 + 0.0148 + 0.0152
   = 0.0595~\text{\AA{}}.
\]
This survives the antisymmetry, but is energetically wrong because
$\tfrac{1}{2}\mu\omega^2 (0.0595)^2 = \tfrac{1}{2}\mu\omega^2 (3.540
\times 10^{-3}~\text{\AA}^2)$, whereas the true energy is the sum
of per-sub-channel energies $\sum_s \tfrac{1}{2}\mu\omega^2 \Delta
r_s^2$.

\textbf{Step 3 -- RSS aggregation (the correct formula):}
\[
  \Delta r_\text{FeS}^\text{RSS}
   = \sqrt{\sum_s w_s (\Delta r_s)^2}
   = \sqrt{0.0150^2 + 0.0145^2 + 0.0148^2 + 0.0152^2}
\]
\[
   = \sqrt{2.250 + 2.103 + 2.190 + 2.310} \cdot 10^{-4}
   = \sqrt{8.853 \times 10^{-4}}
   = \boxed{0.02975~\text{\AA{}}}.
\]

\textbf{Step 4 -- Verify energetic consistency:}
The per-mode parent reorganization should equal the sum of
per-sub-channel reorganizations:
\[
  \lambda_\text{FeS}(\text{via RSS})
   = \tfrac{1}{2}\mu\omega^2 (0.02975)^2 / (hc)
\]
should equal $\sum_s \lambda_\text{pair}^{(s)}$, where each
$\lambda_\text{pair}^{(s)} = \tfrac{1}{2}\mu\omega^2 (\Delta r_s)^2/(hc)$.

Using the conversion factor from the $\lambda^\text{pair}$ example
($\Delta r = 0.0217$~\AA{} gave $\lambda = 22.85$~cm$^{-1}$, so the
proportionality is $\lambda^\text{pair} / (\Delta r)^2 = 22.85 /
(0.0217)^2 = 4.852 \times 10^4$~cm$^{-1}$/\AA$^2$):
\[
  \lambda_\text{FeS}(\text{via RSS})
   = 4.852 \times 10^4 \cdot 0.02975^2
   = 42.96~\text{cm}^{-1}.
\]
And from the sum:
\[
  \sum_s \lambda_\text{pair}^{(s)}
   = 4.852 \times 10^4 \cdot (2.250 + 2.103 + 2.190 + 2.310) \cdot 10^{-4}
   = 4.852 \cdot 8.853
   = 42.96~\text{cm}^{-1}.
\]
The two values match to all digits. \emph{This is the
consistency property} that the RSS aggregation guarantees and
that the L1 alternatives violate.

\textbf{Sigma propagation:}
With $\sigma_e^\text{therm} = 5\!\times\!10^{-4} \cdot 0.0407 =
2.04 \times 10^{-5}$~\AA{}, and using \cref{eq:sigma-rss},
\[
  \sigma_\text{RSS}^\text{FeS}
  = \sqrt{2} \cdot 2.04 \times 10^{-5}
    \cdot \frac{\sqrt{(0.0150)^2 + (0.0145)^2 + (0.0148)^2 + (0.0152)^2}}
              {0.02975}
\]
\[
  = 2.88 \times 10^{-5} \cdot \frac{0.02975}{0.02975}
  = 2.88 \times 10^{-5}~\text{\AA{}}.
\]
The square-root factor reduces to 1 for equal-magnitude sub-channels.
Significance: $0.02975/2.88\!\times\!10^{-5} \approx 1030\sigma$ --
deep above the noise.

\textbf{In the code:}
\fn{reorganization.aggregate\_by\_parent} loops over sub-channel
\fn{ChannelResult}s and returns the parent \texttt{dr\_rss\_FeS}
= 0.02975.
\end{workedexample}

% -------------------------------------------------------------------------
\section{Per-mode and total reorganization}
\label{sec:lambda-total}

\subsection{Per-mode parent-channel $\lambda_X(i)$}

\codref{reorganization.compute\_mode\_modulations}

For mode $i$ and parent channel $X$, the per-mode reorganization
energy is the harmonic-oscillator energy of the RSS-aggregated
bond modulation:
\begin{equation}
  \lambda_X(i)
    \;\equiv\;
  \tfrac{1}{2}\,\mu_i\,\omega_i^2\,(\Delta r_X^\text{RSS}(i))^2
    \;=\;
  \sum_{s \in \text{sub-channels}(X)}
   w_s \cdot \lambda_X^{(s)}(i),
  \codref{reorganization.compute\_mode\_modulations}
  \label{eq:lambda-per-mode}
\end{equation}
where the equality with the weighted sub-channel sum follows from
\cref{eq:lambda-rss-consistency}.

\subsection{System-wide total $\Lambda_X^\text{total}$}

\codref{reorganization.compute\_total\_reorganization}

The total reorganization energy along channel $X$ is the sum of
the per-mode contributions over all modes in the frequency window:
\begin{equation}
  \boxed{\;\Lambda_X^\text{total}
    \;\equiv\;
  \sum_{i=1}^{N_\text{modes}} \lambda_X(i)\;}
  \codref{reorganization.compute\_total\_reorganization}
  \label{eq:lambda-total-derivation}
\end{equation}
This is the headline quantity used for system comparisons
(``does this $\mu$S-protonated state have a larger
$\Lambda_\text{FeS}$ than the corresponding deprotonated state?'',
or ``does the Cys$_2$His$_2$ Rieske-type cluster have larger
$\Lambda_\text{FeFe}$ than the Cys$_4$ ferredoxin-type?''). It is
computed once per cluster, per
frequency window, and recorded in the
\file{*\_analysis\_main.xlsx} workbook on the
\texttt{Reorganisation} sheet.

\subsection{Cumulative-Lambda curve $\Lambda_X(\omega_\text{cut})$}

\codref{reorganization.compute\_cumulative\_reorganization}

The cumulative curve is the partial sum truncated at
$\omega_\text{cut}$:
\begin{equation}
  \Lambda_X(\omega_\text{cut})
    \;\equiv\;
  \sum_{\substack{i \\ \omega_i \leq \omega_\text{cut}}}
   \lambda_X(i).
  \label{eq:lambda-cumulative}
\end{equation}
This is a monotonically non-decreasing step function that climbs
to $\Lambda_X^\text{total}$ as $\omega_\text{cut} \to \infty$. The
locations of large jumps in $\Lambda_X(\omega)$ identify the
spectral regions where bond channel $X$ accumulates most of its
reorganization energy -- directly comparable with NRVS peak
positions.

The cumulative curve is computed on a uniform frequency grid (default
$\Delta\omega = $ 0.05 cm$^{-1}$, set by \cfgkey{interp\_step}) and
written to the \texttt{Cumulative} sheet of the
\file{*\_analysis\_interp\#\#.xlsx} workbook.

% -------------------------------------------------------------------------
\section{Modulation spectra $M_X(\omega)$ and co-modulation}
\label{sec:modulation-spectra}

\subsection{Single-channel $M_X(\omega)$}

\codref{reorganization.compute\_modulation\_spectra}

The modulation spectrum along channel $X$ is a frequency-resolved
amplitude landscape obtained by gaussian-broadening every mode's
$\Delta r_X^\text{RSS}$ at its frequency:
\begin{equation}
  \boxed{\;M_X(\omega)
    \;\equiv\;
  \sum_{i=1}^{N_\text{modes}}
   \Delta r_X^\text{RSS}(i)\,
   \mathcal{G}(\omega - \omega_i,\,\sigma_M)\;}
  \codref{reorganization.compute\_modulation\_spectra}
  \label{eq:M-spectrum-derivation}
\end{equation}
with the normalised gaussian
\begin{equation}
  \mathcal{G}(x, \sigma) \;\equiv\;
  \frac{1}{\sigma\sqrt{2\pi}}\,\exp\!\left(-\frac{x^2}{2\sigma^2}\right).
\end{equation}
The broadening width $\sigma_M = \cfgkey{reorg\_spectrum\_sigma\_cm1}$
(default $5$ cm$^{-1}$) is chosen empirically to match the typical
DFT-vs-experiment frequency mismatch. $M_X(\omega)$ has units of
\AA{} (the broadening preserves the dimensional content of
$\Delta r_X^\text{RSS}$); it is recorded on the
\texttt{ModulationsSpektrum} sheet.

\subsection{Co-modulation as geometric mean}

\codref{reorganization.compute\_co\_modulation\_spectra}

For PCET analysis we want a frequency-resolved descriptor that
peaks where \emph{two} channels are simultaneously modulating. The
geometric mean of two $M$ spectra is the natural choice:
\begin{align}
  C_\PCET(\omega)
    &\;=\;
   \sqrt{M_{HA}(\omega)\, M_\text{FeFe}(\omega)}, \\
  C_{\PT\text{-FeN}}(\omega)
    &\;=\;
   \sqrt{M_{HA}(\omega)\, M_\text{FeN}(\omega)}, \\
  C_{\ET\text{-FeS}}(\omega)
    &\;=\;
   \sqrt{M_\text{FeFe}(\omega)\, M_\text{FeS}(\omega)}.
   \codref{reorganization.compute\_co\_modulation\_spectra}
  \label{eq:co-modulation}
\end{align}
Geometric-mean co-modulation has two attractive properties:

\begin{itemize}
  \item It is zero if either factor is zero (a mode active in only
        one channel does not contaminate the co-modulation).
  \item It is symmetric in the two channels (no preferred ordering).
  \item Its units are \AA{} (the same as the underlying $M$
        spectra).
\end{itemize}

The physical interpretation: $C_\PCET(\omega)$ peaks at frequencies
where modes simultaneously modulate the proton-transfer reaction
coordinate \emph{and} the Fe-Fe distance (i.e.\ the electron-
transfer reaction coordinate). These are the modes most likely to
serve as PCET promoting modes.

\subsection{Choice of broadening width}

The default $\sigma_M = 5$ cm$^{-1}$ is small enough to preserve
narrow features (peak resolution $\sim 2\sigma = 10$ cm$^{-1}$) and
large enough to bridge the typical DFT-vs-experimental frequency
mismatch (5--15 cm$^{-1}$ for harmonic DFT, 1--3 cm$^{-1}$ for
scaled DFT). Smaller $\sigma_M$ produces sharper peaks but is more
sensitive to vibrational anharmonicity that DFT misses; larger
$\sigma_M$ smooths over the band-by-band identification that is
the principal goal of $M_X(\omega)$. The user can override via
\cfgkey{reorg\_spectrum\_sigma\_cm1}.

\begin{workedexample}[$M_X(\omega)$ evaluated at three frequencies]
We compute $M_\text{FeS}(\omega)$ on a uniform 0.5-cm$^{-1}$ grid
using three mock modes:
\begin{center}
\begin{tabular}{lcc}
\toprule
Mode \# & $\omega_i$ [cm$^{-1}$] & $\Delta r^\text{RSS}_\text{FeS}(i)$ [\AA{}] \\
\midrule
1 & 280 & 0.012 \\
2 & 320 & 0.030 \\
3 & 365 & 0.022 \\
\bottomrule
\end{tabular}
\end{center}
with default broadening $\sigma_M = 5$~cm$^{-1}$.

\textbf{Step 1 -- Gaussian normalisation:}
$1/(\sigma_M \sqrt{2\pi}) = 1/(5 \cdot 2.507) = 0.0798$~cm.

\textbf{Step 2 -- evaluate at $\omega = 320$~cm$^{-1}$ (peak of mode 2):}
\begin{align*}
  \mathcal{G}(320-280, 5) &= 0.0798 \cdot \exp(-(40)^2/(2 \cdot 25))
   = 0.0798 \cdot e^{-32} \approx 1.13 \times 10^{-15}, \\
  \mathcal{G}(320-320, 5) &= 0.0798 \cdot \exp(0) = 0.0798, \\
  \mathcal{G}(320-365, 5) &= 0.0798 \cdot \exp(-(45)^2/50)
   = 0.0798 \cdot e^{-40.5} \approx 2.0 \times 10^{-19}.
\end{align*}
\[
  M_\text{FeS}(320) = 0.012 \cdot 1.13 \times 10^{-15}
                   + 0.030 \cdot 0.0798
                   + 0.022 \cdot 2.0 \times 10^{-19}
\]
\[
   \approx 0.030 \cdot 0.0798
   = \boxed{2.39 \times 10^{-3}~\text{\AA{}}\cdot\text{cm}}.
\]
Note units: the Gaussian has dimensions 1/cm$^{-1}$ = cm, so
$M_X$ inherits units of \AA{}$\cdot$cm. The Excel column is
labelled simply with the units \AA{} (the cm factor is the inverse
of the frequency resolution and is conventional).

\textbf{Step 3 -- evaluate at $\omega = 300$~cm$^{-1}$ (between modes 1 and 2):}
\begin{align*}
  \mathcal{G}(300-280, 5) &= 0.0798 \cdot e^{-(20)^2/50}
   = 0.0798 \cdot e^{-8} = 2.68 \times 10^{-5}, \\
  \mathcal{G}(300-320, 5) &= 0.0798 \cdot e^{-8}
   = 2.68 \times 10^{-5}.
\end{align*}
\[
  M_\text{FeS}(300) = 0.012 \cdot 2.68 \times 10^{-5}
                   + 0.030 \cdot 2.68 \times 10^{-5}
                   + 0 \approx 1.13 \times 10^{-6}~\text{\AA{}}\cdot\text{cm}.
\]
Three orders of magnitude smaller than the peak at $\omega = 320$
-- the Gaussian falls off rapidly outside $\pm 2\sigma$.

\textbf{Step 4 -- evaluate at $\omega = 280$~cm$^{-1}$ (peak of mode 1):}
\[
  M_\text{FeS}(280) = 0.012 \cdot 0.0798 = 9.6 \times 10^{-4}~\text{\AA{}}\cdot\text{cm}.
\]
Smaller than the peak at $\omega = 320$ because mode 1 has smaller
$\Delta r$.

\textbf{Cumulative $\Lambda_\text{FeS}(\omega)$ at the same three points:}
$\lambda(i) = 4.852 \times 10^4 \cdot (\Delta r_i)^2$ (using the
proportionality from the $\lambda^\text{pair}$ example):
$\lambda(1) = 6.99$, $\lambda(2) = 43.7$, $\lambda(3) = 23.5$
cm$^{-1}$.
\[
  \Lambda(280) = \lambda(1) = 6.99,
  \quad
  \Lambda(320) = 6.99 + 43.7 = 50.69,
  \quad
  \Lambda(365) = 50.69 + 23.5 = 74.19~\text{cm}^{-1}.
\]
The cumulative curve is a step function that increases at each
mode frequency. The total $\Lambda^\text{total}_\text{FeS} =
\Lambda(365) = 74.2$~cm$^{-1}$.

\textbf{In the code:}
\fn{reorganization.compute\_modulation\_spectra(per\_mode\_results,
parent="FeS", omega\_grid, sigma\_cm1=5)} returns the array
$M_\text{FeS}(\omega)$ on the grid.
\end{workedexample}

% -------------------------------------------------------------------------
\section{SCSD projection}
\label{sec:scsd-derivation}

\subsection{The $D_{2\mathrm{h}}$ symmetry group}

A planar four-atom rhombic Fe$_2$S$_2$ core with the convention
$\rvec_{\text{Fe}_1} = (+a, 0, 0)$, $\rvec_{\text{Fe}_2} = (-a, 0, 0)$,
$\rvec_{\text{S}_1} = (0, +b, 0)$, $\rvec_{\text{S}_2} = (0, -b, 0)$
has the symmetry group $D_{2\mathrm{h}}$, which has 8 irreducible
representations:
\begin{equation}
  \mathrm{D}_{2\mathrm{h}} \;=\;
  \{E, C_2(\zhat), C_2(\yhat), C_2(\xhat), i, \sigma(xy), \sigma(xz), \sigma(yz)\}
\end{equation}
with the irreps $A_g, B_{1g}, B_{2g}, B_{3g}, A_u, B_{1u}, B_{2u},
B_{3u}$, where the $g$/$u$ subscript records parity under inversion
$i$ and the $B_{kg/u}$ subscript records symmetry under the $C_2$
axes.

For a 4-atom cluster the displacement space is 12-dimensional, and
its decomposition into $D_{2\mathrm{h}}$ irreps is
$3 A_g + B_{1g} + 2 B_{2g} + 2 B_{3g} + A_u + 2 B_{1u} + B_{2u} + B_{3u}$,
totalling 13 dimensions -- one more than 12, which means one
linear combination is over-counted; this combination is exactly the
zero mode (uniform translation along $\zhat$) plus the rigid-rotation
pseudomodes that span 6 of the 12 dimensions; the remaining 6
dimensions are the genuine in-cluster vibrational modes.

Kingsbury \& Senge \cite{Kingsbury2024} have published a robust
projection procedure that handles the over-counting cleanly by
projecting onto the 8 irrep-pure subspaces directly and reporting
the projection magnitudes:
\begin{equation}
  \boxed{\;\mathbf{a}_l
    \;=\;
  \mathbf{P}_{\mathrm{D}_{2\mathrm{h}}}
   \bigl(\evec_l^\text{cluster} \cdot \urms\bigr),
   \qquad
   \mathbf{a}_l \in \R^8\;}
  \codref{core.compute\_scsd\_for\_mode\_full}
  \label{eq:scsd-projection-derivation}
\end{equation}
where $\mathbf{P}_{\mathrm{D}_{2\mathrm{h}}}$ is the $8 \times 12$
projection-matrix and $\evec_l^\text{cluster}$ denotes the
12-component cluster-core sub-vector of the mode-$l$ eigenvector,
listed in canonical order
$[\text{Fe}_1, \text{Fe}_2, \text{S}_1, \text{S}_2]$.

\subsection{Reference-vs-distorted difference method}

\codref{core.compute\_scsd\_for\_mode\_full}

The projection matrix is constructed for an idealised \emph{model
geometry} (\fn{core.\_SCSD\_MODEL\_COORDS}) with hard-coded
$\text{Fe-Fe} = 2.73$ \AA{} and $\text{Fe-S} = 2.20$ \AA{}. The
actual cluster geometry differs from this model. To avoid carrying
the geometry difference into the SCSD result, the code uses a
\emph{difference method}:
\begin{enumerate}
  \item Compute the SCSD projection of the model geometry,
        $\mathbf{a}_l^\text{ref}$.
  \item Compute the SCSD projection of the actual (Kabsch-aligned)
        cluster geometry, $\mathbf{a}_l^\text{dist}$.
  \item Report the difference,
        $\mathbf{a}_l^\text{output} =
         \mathbf{a}_l^\text{dist} - \mathbf{a}_l^\text{ref}$.
\end{enumerate}
The fixed-geometry component cancels by linearity of the
projection. The reported amplitudes are the \emph{distortion}
components, which are invariant under the choice of reference
(any other reference geometry with the same symmetry would give
the same answer).

\subsection{Sigma propagation -- the v1.0.4 fix}

\codref{core.compute\_scsd\_for\_mode\_full}

Before v1.0.4, the SCSD sigmas used hard-coded constants
($5 \times 10^{-7}$ for the reference, $5 \times 10^{-6}$ for the
distortion). The v1.0.4 fix (FUND 8, \cref{sec:fund-6-8}) replaces
these with formulas built from the user-configurable
$\sigma_e = \cfgkey{sigma\_eigvec}$ and
$\sigma_c = \cfgkey{sigma\_coord}$:
\begin{align}
  \sigma_\text{SCSD,ref}^{(\Gamma)}
    &\;=\;
   \sqrt{2}\,\sigma_c, \\[2pt]
  \sigma_\text{SCSD,dist}^{(\Gamma)}
    &\;=\;
   \sqrt{2}\,\sqrt{\sigma_c^2 + (\sigma_e \cdot \urms(\omega_l))^2}, \\[2pt]
  \sigma_\text{SCSD,diff}^{(\Gamma)}
    &\;=\;
   \sqrt{(\sigma_\text{SCSD,ref}^{(\Gamma)})^2
       + (\sigma_\text{SCSD,dist}^{(\Gamma)})^2}.
\end{align}
With the default $\sigma_c = 10^{-3}$ \AA{} and $\sigma_e = 5
\times 10^{-4}$, the new sigmas are about three orders of
magnitude larger than the hard-coded ones, which were physically
implausibly small. The underlying SCSD amplitudes are unchanged.

% -------------------------------------------------------------------------
\section{Kernel classification}
\label{sec:kernel-classification}

\subsection{Goal}

\codref{core.\_oop\_ring\_metrics, core.classify\_oop\_inp}

While SCSD gives a rigorous $D_{2\mathrm{h}}$-irrep decomposition,
many users want a more chemist-friendly summary in terms of the
classical cluster-core normal-mode names: ``Fe-S symmetric
stretch'', ``Fe-Fe antisymmetric stretch'', ``S-S breathing'',
``cluster umbrella'', etc. The kernel-classification heuristic
computes a score for each of 10 named mode types and reports
\texttt{kern\_primary} (the highest-scoring) and
\texttt{kern\_secondary} (the second-highest).

\subsection{The 10 named mode types}

\begin{longtable}{p{0.27\textwidth} p{0.13\textwidth} p{0.55\textwidth}}
\toprule
\textbf{Name} & \textbf{Symmetry} & \textbf{Description} \\
\midrule
\texttt{Fe\_Fe\_sym\_stretch}
  & $A_g$ & Both Fe atoms move symmetrically along $\xhat$
            (out from / into the cluster centre). \\
\texttt{Fe\_Fe\_antisym\_stretch}
  & $B_{1g}$ & Fe$_1$ along $+\xhat$, Fe$_2$ along $-\xhat$
              (out-of-phase). \\
\texttt{S\_S\_sym\_stretch}
  & $A_g$ & S atoms move along $\yhat$ symmetrically. \\
\texttt{S\_S\_antisym\_stretch}
  & $B_{2g}$ & S atoms move along $\yhat$ out-of-phase. \\
\texttt{Fe\_S\_breathing}
  & $A_g$ & All four atoms move radially outward (cluster expansion). \\
\texttt{Fe\_S\_anti\_breathing}
  & $B_{1g}$ & Antisymmetric breathing. \\
\texttt{cluster\_umbrella}
  & $A_u$ & Synchronous motion along $\zhat$ (out of plane). \\
\texttt{cluster\_seesaw}
  & $B_{2u}$ & Fe atoms move along $+\zhat$ while S along $-\zhat$. \\
\texttt{cluster\_rotation}
  & $B_{3u}$ & Rigid in-plane rotation. \\
\texttt{cluster\_translation}
  & --  & Uniform translation (should be zero in a converged
          frequency calculation, but is computed as a diagnostic). \\
\bottomrule
\end{longtable}

\subsection{Scoring}

\codref{core.classify\_oop\_inp}

Each mode is scored against each named type by computing the inner
product
\begin{equation}
  s_l^{(k)}
    \;=\;
  \left\lvert\frac{\evec_l^\text{cluster} \cdot
                   \mathbf{v}^{(k)}_\text{template}}
                  {\lVert\evec_l^\text{cluster}\rVert
                   \cdot \lVert\mathbf{v}^{(k)}_\text{template}\rVert}
       \right\rvert^2 \cdot 100\%
\end{equation}
where $\mathbf{v}^{(k)}_\text{template}$ is the unit-normalised
template vector for mode type $k$, and the inner product is over
the 12 cluster-core components. Scores are cosine-squared
similarities expressed as percentages. The two highest scores are
recorded as \texttt{kern\_primary} and \texttt{kern\_secondary},
together with the localisation score \texttt{kern\_loc} (the
fraction of the mode's total amplitude residing in the cluster
core).

\subsection{Caveats}

\begin{itemize}
  \item Mode types are not orthogonal in general -- the templates
        for $A_g$-symmetric stretches (Fe-Fe sym, S-S sym, Fe-S
        breathing) span a 3-dimensional $A_g$ subspace, so scores
        across these three sum to $\leq 100\%$ but a single mode
        can mix them.
  \item Pure cluster-translation should have zero score across all
        types except $\texttt{cluster\_translation}$. If the
        zero-frequency translation modes were not filtered out
        (\cref{sec:phase4-loop}), they would diagnostically
        highlight if the harmonic approximation has broken down.
  \item For the rigorous, mathematically orthogonal decomposition,
        the SCSD sheet (\cref{sec:scsd-derivation}) is the
        authoritative reference; the kernel sheet is a
        human-friendly summary.
\end{itemize}

% -------------------------------------------------------------------------
\section{$B$-factor (Debye-Waller) accumulation}
\label{sec:b-factor-derivation}

\subsection{From mean-square displacement to crystallographic $B$}

\codref{runner.\_run\_analysis\_single}

The crystallographic Debye-Waller factor for atom $i$ is
\begin{equation}
  B_i \;\equiv\; 8\pi^2\,\langle u_i^2 \rangle\,/3
\end{equation}
where $\langle u_i^2 \rangle = \langle\lVert\rvec_i(t) -
\bar\rvec_i\rVert^2\rangle$ is the time-averaged mean-square
displacement, and the factor $8\pi^2/3$ comes from the isotropic
gaussian convention in X-ray diffraction.

In a harmonic approximation the mean-square displacement
decomposes into the sum of per-mode contributions:
\begin{equation}
  \langle u_i^2 \rangle
    \;=\;
  \sum_{l=1}^{3N-6} \langle q_l^2 \rangle \cdot
       \lVert\evec_{l,i}\rVert^2
    \;=\;
  \sum_l \urms^2(\omega_l) \cdot \lVert\evec_{l,i}\rVert^2,
\end{equation}
where the second equality uses
$\langle q_l^2 \rangle = \urms^2$ from
\cref{sec:urms-derivation} and the orthogonal decomposition of
the cartesian displacement of atom $i$ over modes.

The pipeline accumulates the bracket sum mode by mode in Phase 4:
\begin{equation}
  \text{b\_accum}[i] \;\mathrel{+}=\;
   \sum_{\alpha = x, y, z} (\evec_{l,i}^{(\alpha)} \cdot \urms(\omega_l))^2
  \label{eq:b-factor-accum}
\end{equation}
and finalises after the loop:
\begin{equation}
  \boxed{\;B_i
    \;=\;
  \frac{8\pi^2}{3}\,\text{b\_accum}[i]\;}
  \codref{runner.\_run\_analysis\_single}
  \label{eq:b-factor-final}
\end{equation}
in units of \AA$^2$.

\subsection{Numerical details}

The accumulator is an \fn{np.ndarray} of length $N_\text{atoms}$,
indexed by the position in the heavy-atom list (\emph{not} by
Gaussian centre). The mapping from Gaussian centre to accumulator
index uses \fn{idx\_map[centre]} which is the 0-based position in
the atoms list. The accumulator is initialised to zero at the start
of the cluster run and the per-cluster reset in
\fn{reset\_warning\_state} also re-zeros \fn{b\_accum} implicitly
via the per-cluster pipeline state.

\subsection{Frequency window dependence}

Since the accumulation runs only over modes inside the user's
frequency window, the resulting $B$ factors are
\emph{window-dependent}: a $B$ factor reported for a 0--800 cm$^{-1}$
analysis will be smaller than the same atom's true full-spectrum
$B$ factor, because higher-frequency modes contribute zero-point
displacement amplitudes too. The window-restricted $B$ is
nonetheless useful diagnostically (the relative ordering of atoms
by $B$ is stable across windows; the absolute values are not).

\begin{workedexample}[$B$-factor for Fe$_1$ from three modes]
We accumulate the $B$-factor of Fe$_1$ across three modes, using
the test eigenvectors and frequencies from previous examples.
\begin{center}
\begin{tabular}{lccc}
\toprule
Mode & $\omega$~[cm$^{-1}$] & $\urms$~[\AA{}]
                            & $\evec_{l,\text{Fe}_1}$ (mass-unweighted) \\
\midrule
1 & 80    & 0.0727 & $(+0.10, +0.20, +0.40)$ (umbrella-like) \\
2 & 320   & 0.0407 & $(+0.42, 0,     0)$     (Fe-S stretch) \\
3 & 510   & 0.0244 & $(+0.05, +0.30, +0.15)$ (Fe-N bend) \\
\bottomrule
\end{tabular}
\end{center}

\textbf{Step 1 -- per-mode contributions to b\_accum[Fe$_1$]:}
\begin{align*}
  &\text{mode 1: } (0.10 \cdot 0.0727)^2 + (0.20 \cdot 0.0727)^2 + (0.40 \cdot 0.0727)^2 \\
  &\quad = (0.00727)^2 + (0.01454)^2 + (0.02908)^2
  = 5.3 \times 10^{-5} + 2.1 \times 10^{-4} + 8.5 \times 10^{-4} \\
  &\quad = 1.11 \times 10^{-3}~\text{\AA{}}^2, \\[4pt]
  &\text{mode 2: } (0.42 \cdot 0.0407)^2 + 0 + 0
  = (0.01709)^2 = 2.92 \times 10^{-4}~\text{\AA{}}^2, \\[4pt]
  &\text{mode 3: } (0.05 \cdot 0.0244)^2 + (0.30 \cdot 0.0244)^2 + (0.15 \cdot 0.0244)^2 \\
  &\quad = (0.00122)^2 + (0.00732)^2 + (0.00366)^2
  = 1.5 \times 10^{-6} + 5.4 \times 10^{-5} + 1.3 \times 10^{-5} \\
  &\quad = 6.86 \times 10^{-5}~\text{\AA{}}^2.
\end{align*}

\textbf{Step 2 -- sum:}
\[
  \text{b\_accum}[\text{Fe}_1] = 1.11 \times 10^{-3}
                              + 2.92 \times 10^{-4}
                              + 6.86 \times 10^{-5}
   = 1.47 \times 10^{-3}~\text{\AA{}}^2.
\]

\textbf{Step 3 -- multiply by $8\pi^2/3 = 26.32$:}
\[
  B_{\text{Fe}_1} = 26.32 \cdot 1.47 \times 10^{-3}
   = \boxed{0.0387~\text{\AA{}}^2}.
\]

\textbf{Sanity:} For a real protein at 80~K, $B$ factors at the
Fe site are typically $0.5$ to $2.0$~\AA{}$^2$. Our value
$0.039$~\AA{}$^2$ is much lower because we are accumulating over
only 3 mock modes; a full DFT spectrum would contribute hundreds
of modes summing into the realistic range.

\textbf{Observation:} The umbrella mode (mode 1, low frequency,
high $\urms$) dominates the $B$-factor with $76\%$ of the total,
even though its amplitude on Fe$_1$ is moderate. This is the
characteristic feature of low-frequency modes -- they have small
displacement amplitudes per mode but large $\urms$, and contribute
disproportionately to the time-averaged mean-square displacement.

\textbf{In the code:}
\fn{runner.\_run\_analysis\_single} loops over selected modes and
accumulates \fn{b\_accum} in-place; the final \fn{b\_factors}
array is exposed on the \texttt{B\_factors} sheet of the main
workbook.
\end{workedexample}

% =========================================================================
% Theory deep-dive sections (added v1.0.4 expansion)
% =========================================================================

\section{Connection to Marcus-Hush ET theory}
\label{sec:marcus-hush-deep}

The spectroscopic reorganization $\lambda_X^\text{pair}$ defined
in \cref{eq:lambda-cm1} shares the quadratic functional form
$\tfrac{1}{2}\mu\omega^2(\Delta r)^2$ with the classical Marcus
reorganization energy, but the two quantities are conceptually
different and should not be conflated. This section explains the
relationship in detail.

\subsection{Marcus classical reorganization energy}

In Marcus theory of electron transfer (ET), the reorganization
energy $\lambda_M$ is the work required to distort the equilibrium
nuclear configuration of the reactant state to the equilibrium
nuclear configuration of the product state, with the electron
held in place:
\begin{equation}
  \lambda_M \;=\; \sum_i \frac{1}{2}\,k_i\,(\Delta Q_i)^2
                  \;+\; \lambda_\text{solvent},
  \label{eq:marcus-classical}
\end{equation}
where $k_i = \mu_i \omega_i^2$ is the harmonic force constant of
internal mode $i$, $\Delta Q_i$ is the static (Born-Oppenheimer)
displacement of mode $i$'s equilibrium between the two electronic
states, and $\lambda_\text{solvent}$ is the dielectric
reorganization contribution. The ET rate constant in the Marcus
non-adiabatic regime is then
\begin{equation}
  k_\text{ET}
   = \frac{2\pi}{\hbar} |V_{DA}|^2
     \frac{1}{\sqrt{4\pi\lambda_M\kB T}}
     \exp\!\left(-\frac{(\Delta G^\circ + \lambda_M)^2}
                       {4\lambda_M\kB T}\right),
  \label{eq:marcus-rate}
\end{equation}
where $V_{DA}$ is the donor-acceptor electronic coupling and
$\Delta G^\circ$ is the standard free-energy of the reaction.

\subsection{Spectroscopic reorganization energy}

Our $\lambda_X^\text{pair}$, by contrast, uses $\Delta r$ =
\emph{thermal/zero-point displacement amplitude of a single
vibrational mode along a bond axis}, not the static geometry
shift between two electronic states. It is a frequency-resolved
measure of which bond modes ``zappel'' (oscillate) the most at
the chosen temperature.

The two reduce to each other only in a special limit: the
\emph{single-promoting-mode} regime of ET, where the entire
$\lambda_M$ is concentrated in one internal mode and that mode
also dominates the spectroscopic amplitude. In general the
spectroscopic $\lambda$ is a different observable -- it tells the
reader which modes \emph{currently} oscillate, while the Marcus
$\lambda$ tells which static geometric changes are induced by
\emph{electron motion}.

\subsection{The spin-Boson model bridge}

The spin-Boson model of ET treats the donor-acceptor system as a
two-level electronic system linearly coupled to a bath of harmonic
modes. The coupling is parameterised by the spectral density
$J(\omega) = \sum_i \tfrac{c_i^2}{2 m_i \omega_i} \delta(\omega -
\omega_i)$, where $c_i$ is the coupling strength of bath mode $i$
to the electronic system. The Marcus $\lambda_M$ is then
\begin{equation}
  \lambda_M = \frac{1}{\pi}\int_0^\infty
              \frac{J(\omega)}{\omega}\,d\omega.
\end{equation}
This is the integrated spectral density divided by frequency --
it weights low-frequency modes most heavily. Our
$\Lambda_X^\text{total}(\omega) = \sum_l \lambda_X^{(l)}$ is a
similar integral but with different weighting:
\begin{equation}
  \Lambda_X^\text{total}
   \;=\; \sum_l \tfrac{1}{2}\mu_l\omega_l^2 (\Delta r_l)^2 / (hc),
\end{equation}
where $\Delta r_l \propto \urms(\omega_l) = \sqrt{\hbar/(2\mu_l
\omega_l)\coth(\beta\hbar\omega_l/2)}$. Substituting:
\begin{equation}
  \Lambda_X^\text{total}
   \;=\; \frac{1}{hc} \sum_l
         \frac{\hbar \omega_l}{4}
         \coth(\beta\hbar\omega_l/2)\,\alpha_{X,l}^2,
\end{equation}
where $\alpha_{X,l}$ is the bond-axis projection coefficient
(\cref{eq:dr-classical-final} structure). At $T \to 0$,
$\coth \to 1$ and $\Lambda^\text{total} = \sum_l \omega_l
\alpha_{X,l}^2 / (4 hc)$, which is the spin-Boson reorganization
specific to channel $X$.

The mapping is:
\begin{itemize}
  \item Marcus $\lambda_M$: sum of $\tfrac{c_i^2}{2 k_i}$ over all
        modes coupled to the ET reaction coordinate.
  \item Our $\Lambda_X^\text{total}$: sum of $\tfrac{1}{4}\hbar
        \omega_l \alpha_{X,l}^2$ over all modes with displacement
        $\alpha_{X,l}$ along bond $X$ at temperature $T$.
\end{itemize}
The two coincide when bond $X$ \emph{is} the reaction coordinate
and at $T \to 0$.

\subsection{Huang-Rhys factors and the Franck-Condon distribution}

For a single mode of frequency $\omega_l$, the Huang-Rhys factor
$S_l$ is the dimensionless ratio of reorganization energy to
$\hbar\omega_l$:
\begin{equation}
  S_l = \frac{\lambda_l}{\hbar\omega_l}
      = \frac{1}{2}\frac{\mu_l\omega_l(\Delta Q_l)^2}{\hbar}
      = \frac{(\Delta Q_l)^2}{2\,\urms^2(\omega_l, T = 0)}.
\end{equation}
$S_l$ counts ``vibrational quanta worth'' of reorganization energy
in mode $l$. The Franck-Condon distribution of vibrational
overlap factors between displaced harmonic oscillators is then a
Poisson distribution with parameter $S_l$:
\begin{equation}
  P_l(n) = \frac{S_l^n}{n!}\,e^{-S_l},
\end{equation}
giving the relative intensity of the $0 \to n$ transition in
optical spectroscopy. For our typical Fe-S modes,
$S_l \sim 0.01$--$0.1$ (because the per-mode $\Delta r$ is small
compared to $\urms$), placing the modes in the weak-coupling
regime where the dominant Franck-Condon contribution is
$0 \to 0$ ($P = e^{-S} \approx 1$). This is consistent with the
narrow homogeneous linewidths observed in NRVS.

\subsection{Spectroscopic $\lambda$ as a Franck-Condon-weighted observable}

The NRVS cross-section at frequency $\omega$ (see
\cref{sec:nrvs-cross-section}) involves the Franck-Condon factor
of a 0-phonon transition, which for a single mode reduces to:
\begin{equation}
  \sigma_\text{NRVS}(\omega) \;\propto\; \sum_l e^{-2W_l}
                                          S_l\,\delta(\omega - \omega_l),
\end{equation}
where $e^{-2W_l}$ is the Lamb-Mössbauer factor and $S_l$ is the
Huang-Rhys factor. Crucially, the cross-section is proportional
to $S_l$ -- which is proportional to $(\Delta Q_l)^2 \omega_l$ --
which is exactly the integrand of our $\Lambda^\text{total}$.

\emph{Therefore: $\Lambda_X^\text{total}$, integrated over the
analysis window, is proportional to the integrated NRVS
cross-section attributed to bond modulations along axis $X$,
in the weak-coupling, low-$T$ limit.}

This is the link that makes $\Lambda$ a direct NRVS observable
rather than just a per-bond ``zappel-energy''. The proportionality
constant depends on $\langle Fe^{57} \rangle$ enrichment, the
Lamb-Mössbauer factor of Fe at the measurement temperature, and
the experimental resolution -- factors that absorb into the
overall NRVS intensity scale.

% -------------------------------------------------------------------------
\section{Lamb-Mössbauer factor and $B$-factors}
\label{sec:lamb-mossbauer}

\subsection{Definition}

The Lamb-Mössbauer factor $f_\text{LM}$ is the fraction of
recoil-free $\gamma$-ray transitions in a Mössbauer experiment:
\begin{equation}
  f_\text{LM} = e^{-\langle k^2 x^2 \rangle},
\end{equation}
where $k$ is the magnitude of the $\gamma$-ray's wave vector and
$\langle x^2 \rangle$ is the mean-square displacement of the
Mössbauer nucleus along the $\gamma$-ray direction. For Fe$^{57}$
at 14.4 keV, $k = 7.3 \times 10^{10}$~m$^{-1}$ = $7.3$~\AA{}$^{-1}$.

The mean-square displacement $\langle x^2 \rangle$ along a single
direction is one-third of the isotropic mean-square displacement
$\langle r^2 \rangle$:
\begin{equation}
  \langle x^2 \rangle = \frac{1}{3}\langle r^2 \rangle
                     = \frac{1}{3}\sum_l \urms^2(\omega_l)\,
                       \lVert\evec_{l,\text{Fe}}\rVert^2.
\end{equation}
Substituting into $f_\text{LM}$ and using $k^2 = 53.3$~\AA{}$^{-2}$:
\begin{equation}
  f_\text{LM}(T) = \exp\!\left(-\frac{k^2}{3}
                    \sum_l \urms^2(\omega_l)\,
                    \lVert\evec_{l,\text{Fe}}\rVert^2\right).
\end{equation}

\subsection{Connection to crystallographic $B$-factors}

The $B$-factor of a Mössbauer nucleus (\cref{sec:b-factor-derivation})
is exactly
\begin{equation}
  B_\text{Fe} = \frac{8\pi^2}{3}
                \sum_l \urms^2(\omega_l)\,
                \lVert\evec_{l,\text{Fe}}\rVert^2
              = 8\pi^2 \langle x^2 \rangle.
\end{equation}
Therefore:
\begin{equation}
  f_\text{LM}(T) = \exp\!\left(-\frac{k^2}{8\pi^2}\,B_\text{Fe}(T)\right)
                 = \exp(-0.0625 \cdot B_\text{Fe}~[\text{\AA}^2]).
\end{equation}
For a typical NRVS-quality system at 80~K with
$B_\text{Fe} = 0.7$~\AA{}$^2$:
\[
  f_\text{LM} = e^{-0.044} = 0.957.
\]
About 96\% of the absorption events are recoil-free at this
temperature -- explaining why cryogenic NRVS produces clean
phonon spectra. At 300~K, $B_\text{Fe} \approx 2.5$~\AA{}$^2$
gives $f_\text{LM} \approx 0.86$, still high but no longer
overwhelming.

\subsection{Implication for the pipeline}

The pipeline reports $B_\text{Fe}$ (and $B$ for every other atom)
on the \texttt{B\_factors} workbook sheet. A reviewer comparing
the output to NRVS measurements can independently verify the
Lamb-Mössbauer factor via the formula above. A pipeline run that
reports $B_\text{Fe} > 5$~\AA{}$^2$ at 80~K should be regarded
with suspicion -- the DFT-predicted modes are inconsistent with
the experimentally observed sharp NRVS lines.

% -------------------------------------------------------------------------
\section{NRVS cross-section vs.\ pDOS}
\label{sec:nrvs-cross-section}

A frequent source of confusion is the difference between the
\emph{NRVS cross-section} (what the instrument actually measures)
and the \emph{partial density of states} (pDOS, what theorists
often plot). The pipeline computes neither directly, but its
$M_X(\omega)$ modulation spectra are intermediate between them.

\subsection{NRVS cross-section (Sturhahn formula)}

The single-phonon NRVS absorption cross-section at energy
$E = \hbar\omega$ above the Mössbauer line is \cite{Sturhahn2004}:
\begin{equation}
  \sigma_\text{NRVS}(\omega)
   \;=\; \sigma_0\,f_\text{LM}\,\frac{\hbar\omega/(2 m_\text{Fe})}
                                 {\omega^2}
         (1 + n(\omega))\,D_\text{Fe}(\omega),
  \label{eq:sigma-nrvs}
\end{equation}
where $\sigma_0$ is the nuclear-resonance cross-section,
$f_\text{LM}$ is the Lamb-Mössbauer factor,
$n(\omega) = (e^{\beta\hbar\omega} - 1)^{-1}$ is the
Bose-Einstein occupation, and $D_\text{Fe}(\omega)$ is the
projected density of states on the Fe nucleus:
\begin{equation}
  D_\text{Fe}(\omega) = \sum_l e_{l,\text{Fe}}^2\,
                               \delta(\omega - \omega_l),
  \label{eq:pdos-fe}
\end{equation}
with $e_{l,\text{Fe}}$ the mass-weighted Fe component of mode $l$'s
eigenvector.

\subsection{Fe-pDOS = NRVS cross-section, modulo prefactors}

Equation \cref{eq:sigma-nrvs} is what the NRVS instrument
measures. The $D_\text{Fe}(\omega)$ factor is the Fe-projected
density of states. To extract the pDOS from a measured
cross-section, one divides out the prefactors $\sigma_0,
f_\text{LM}$, the kinematic factor $\hbar\omega/2 m_\text{Fe}/\omega^2$,
and the Bose-Einstein factor -- standard procedure in the NRVS
analysis literature.

The pipeline does \emph{not} compute the cross-section directly.
What it produces is $M_X(\omega)$, the bond-projected modulation
spectrum (\cref{sec:modulation-spectra}). The two are related
but distinct:

\begin{itemize}
  \item $D_\text{Fe}(\omega) \propto \sum_l \lVert\evec_{l,\text{Fe}}
        \rVert^2 \delta(\omega - \omega_l)$: projected onto a
        single \emph{atom}.
  \item $M_X(\omega) \propto \sum_l (\Delta r_{X,l})^2
        \mathcal{G}(\omega - \omega_l)$: projected onto a
        \emph{bond axis}.
\end{itemize}

The two coincide when the relevant atom moves only along one
bond axis (e.g.\ in a pure Fe-S stretch, $\evec_\text{Fe}$ is
parallel to $\bhat_\text{FeS}$). For mixed modes the two differ:
a mode that moves Fe perpendicular to all Fe-S bonds would
contribute to $D_\text{Fe}$ but not to $M_\text{FeS}$.

\subsection{Why bond-projection is more useful for PCET}

For PCET analysis, the bond-projected spectrum $M_X(\omega)$
identifies \emph{which bonds} modulate together with the ET
reaction coordinate -- the relevant question for the
``promoting modes'' that drive PCET. Pure pDOS would conflate
all Fe motion (including bond-irrelevant rocks and tilts) into
one curve. The bond-projection is the analysis-stage
specialisation that makes the pipeline useful for PCET.

% -------------------------------------------------------------------------
\section{Anharmonicity: what we do not correct for}
\label{sec:anharmonicity}

The entire pipeline assumes harmonic vibrations. Real molecular
vibrations are anharmonic; this section lists the corrections we
omit and the regime where the omission is justified.

\subsection{Diagonal anharmonicity}

A diagonal anharmonic mode has the potential
$V(Q) = \tfrac{1}{2}\mu\omega^2 Q^2 + \tfrac{1}{6}k_3 Q^3 +
\tfrac{1}{24}k_4 Q^4 + \dots$. The leading correction to $\urms$
is
\begin{equation}
  \urms^2 \approx (\urms^\text{harm})^2
   \left(1 - \frac{k_3^2}{\mu^2\omega^4}\cdot\frac{15\hbar}{2\mu\omega}
              + \frac{k_4}{2\mu^2\omega^4}\cdot\frac{3\hbar}{\mu\omega}\right),
\end{equation}
which for typical Fe-S systems amounts to a few percent at most.
We do not apply this correction. It would require the cubic
and quartic anharmonic constants from a higher-order DFT
calculation (\fn{Freq=Anharm}); the typical Gaussian frequency
job we read does not contain these.

\subsection{Off-diagonal anharmonicity (mode-mode coupling)}

The truly harmonic spectrum has $N_\text{modes}$ independent
oscillators; the truly anharmonic spectrum has $N_\text{modes}^2$
pair-couplings $V_{ij} = k_{ij} Q_i Q_j$ that mix modes. Strong
mode-mode coupling can shift peak frequencies, broaden lines, and
produce combination bands at $\omega_i \pm \omega_j$. We ignore
all of these. The pipeline's $M_X(\omega)$ peaks should be read
as ``the harmonic-DFT-predicted peaks before anharmonic
broadening''.

\subsection{Frequency scaling}

A common partial-anharmonic correction is to multiply all
DFT-predicted frequencies by a uniform scale factor (typically
$0.96$--$0.99$, functional-dependent) to bring them closer to
experimental peak positions. Our $\Lambda$ and $M_X$ outputs are
\emph{not} scaled; the user can apply a uniform scaling
post-hoc by multiplying all reported frequencies. The
$\Lambda$ values are quadratic in $\omega$ -- a 1\% frequency
scaling produces a 2\% $\Lambda$ change.

\subsection{Temperature-dependent harmonic frequencies}

The harmonic frequencies are computed at the equilibrium geometry
(DFT-optimised $T = 0$ structure). At finite temperature, the
average geometry shifts due to thermal expansion, and the
``effective'' harmonic frequencies shift correspondingly. This is
usually a small effect (few cm$^{-1}$), and the pipeline ignores
it. A user who needs precise comparison with high-$T$ experiment
should consider supplying a temperature-corrected DFT structure
as input.

\subsection{Practical regime of validity}

Anharmonic corrections become important when:
\begin{itemize}
  \item The mode is near a structural transition (the cluster
        being close to a bond-rearrangement, e.g.\ between
        rhombic and butterfly geometries).
  \item The mode is in a highly nonlinear coupling regime
        (multi-quantum overtones).
  \item Hydrogen-stretching modes ($> 2500$ cm$^{-1}$) -- our
        analysis focuses below 1000 cm$^{-1}$ where anharmonicity
        is mild.
\end{itemize}
For the [2Fe-2S] systems analysed by this pipeline at $T \leq 100$~K,
harmonic predictions are typically accurate to $\pm 5$\% in
frequency and $\pm 10$\% in NRVS-intensity-related observables.
The validation tolerances in \cref{sec:headline-numbers} are set
accordingly.

% =========================================================================
% PART III + IV + V + VI + VII + VIII + IX + X + XI -- STUBS
% =========================================================================
% =========================================================================
% PART III -- ALGORITHMS (complete pseudocodes)
% =========================================================================
\part{Algorithms (complete pseudocodes)}
\label{part:algorithms}

The pseudocodes in this part are intended as a reviewer's
high-fidelity reading of the code in each phase, condensed to the
control flow and data flow that matter physically. Each pseudocode
ends with a \codref{} marker; the actual implementation can be read
verbatim by following the marker. Pseudocode is given in a
language-agnostic syntax with the conventions of \cref{tab:pseudo-conv}.

\begin{table}[h]
\centering
\begin{tabular}{p{0.20\textwidth} p{0.72\textwidth}}
\toprule
\textbf{Convention} & \textbf{Meaning} \\
\midrule
\texttt{a <- b}            & Assignment. Multiple assignment in tuples permitted. \\
\texttt{a, b <- f(x)}      & Tuple-unpacking of the return value of $f$. \\
\texttt{for x in C}        & Iteration over collection $C$. \\
\texttt{foreach}           & As \texttt{for} but order is not significant. \\
\texttt{assert P}          & Halt with an error if predicate $P$ is false. \\
\texttt{require P}         & As \texttt{assert} but emits a clear user-facing message. \\
\texttt{yield x}           & Generator-style emission of $x$ to the caller. \\
\texttt{return}            & Return from the current procedure. \\
\texttt{...}               & Ellipsis: code irrelevant to the present logic. \\
\bottomrule
\end{tabular}
\caption{Pseudocode conventions used throughout Part~\ref{part:algorithms}.}
\label{tab:pseudo-conv}
\end{table}

% -------------------------------------------------------------------------
\section{Algorithm: \fn{scan\_log}}
\label{sec:algo-scan-log}

\codref{logio.scan\_log}

\fn{scan\_log} performs a streaming, byte-level scan of the Gaussian
\file{.log} file to find the byte offsets of three pattern types:
the Standard-Orientation headers, the normal-coordinate-block
headers, and the Frequencies headers. It must be \texttt{O}($N$)
in file size and constant in memory; the only state carried between
chunks is a small rolling tail buffer that catches patterns
spanning chunk boundaries.

\subsection{Byte-level pattern definitions}
\label{sec:scan-patterns}

The three regex patterns used by \fn{scan\_log} are:

\begin{bytebox}
pat\_so = re.compile(b"Standard orientation:")
pat\_nc = re.compile(b"and normal coordinates:")
pat\_fr = re.compile(b"\textbackslash{}s*Frequencies\textbackslash{}s*-\{2,\}",
                     re.MULTILINE)
\end{bytebox}

A few subtleties worth noting:
\begin{itemize}
  \item \texttt{pat\_so} and \texttt{pat\_nc} are \emph{plain
        substring matches} -- no regex meta-characters. The
        Gaussian writer is deterministic about how these lines are
        formatted, so a substring match is both faster and more
        robust than a regex.
  \item \texttt{pat\_fr} requires \emph{at least two} dashes after
        ``Frequencies'' (with optional whitespace). This matters
        because \texttt{hpmodes} blocks write three dashes
        (\texttt{Frequencies ---}) while standard blocks write two
        (\texttt{Frequencies --}). The \texttt{-\{2,\}} quantifier
        captures both. Less than two dashes would be a Gaussian
        format violation; we have not observed such logs in the
        wild.
  \item \texttt{pat\_fr} uses \texttt{re.MULTILINE} so the
        \texttt{\textbackslash{}s*} anchor at the start matches the
        per-line leading whitespace, not just the file start.
  \item All patterns are byte-strings (\texttt{rb"..."}), not
        unicode strings, so they apply directly to the raw bytes
        of the file -- no codec conversion needed during the scan.
\end{itemize}

\subsection{The 200-byte tail-overlap}
\label{sec:scan-tail}

Each chunk read from disk has size \texttt{cfg.scan\_chunk\_mb
{$\cdot$} 1024 {$\cdot$} 1024} bytes (default $16$ MiB). At chunk
boundaries, a pattern split across the boundary would be missed if
each chunk were processed independently. The remedy is to retain
the last 200 bytes of every chunk as a \emph{tail buffer}, prepended
to the next chunk:

\begin{lstlisting}[style=pseudo,
  caption={Streaming byte-offset scan of the Gaussian log with
           tail-overlap deduplication. Implementation:
           \protect\fn{logio.scan\_log}.}]
input:  filepath, chunk_size_bytes (16 MiB default), tail_overlap (200 B)
output: so_off, nc_off, fr_off    (sorted lists of byte offsets)

seen_so, seen_nc, seen_fr <- {}, {}, {}    # dedup sets
so_off, nc_off, fr_off <- [], [], []

tail   <- b""
offset <- 0          # bytes consumed so far (not counting tail)

with open(filepath, "rb") as fh:
    while True:
        block <- fh.read(chunk_size_bytes)
        if len(block) == 0: break
        data <- tail + block
        base <- offset - len(tail)
            # base = file-absolute byte index of data[0]

        foreach pattern, target_list, seen_set in
                [(pat_so, so_off, seen_so),
                 (pat_nc, nc_off, seen_nc),
                 (pat_fr, fr_off, seen_fr)]:
            foreach match in pattern.finditer(data):
                abs_off <- base + match.start()
                if abs_off not in seen_set:
                    seen_set.add(abs_off)
                    target_list.append(abs_off)

        tail   <- data[-200:]
        offset <- offset + len(block)

so_off.sort(); nc_off.sort(); fr_off.sort()
return so_off, nc_off, fr_off
\end{lstlisting}

\paragraph{Why 200 bytes?} The longest of the three patterns is
``\texttt{and normal coordinates:}'' (22 characters); the \texttt{pat\_fr}
regex with \texttt{\textbackslash{}s*Frequencies\textbackslash{}s*-\{2,\}}
matches at most a few dozen bytes. 200 bytes is generously above
the longest possible pattern, leaving headroom for future pattern
additions.

\paragraph{Why dedup with sets?} Without deduplication, every pattern
that falls inside the tail-overlap window of two adjacent chunks would
be reported twice (once at the end of chunk $N$, once at the start
of chunk $N+1$). The per-pattern \texttt{seen\_*} set discards
duplicates by absolute file offset.

\subsection{Worked example: byte layout of a Gaussian log}

For a typical [2Fe-2S] \fn{hpmodes} Gaussian log, the first 4 KiB
after the SCF-convergence point contain the equilibrium geometry
and its surrounding metadata. A representative excerpt from byte
offset $\approx$~280\,000 (immediately after the SCF block) is:

\begin{bytebox}
                          Standard orientation:                         \\
 ---------------------------------------------------------------------\\
 Center     Atomic      Atomic             Coordinates (Angstroms)\\
 Number     Number       Type             X           Y           Z\\
 ---------------------------------------------------------------------\\
      1          7           0       -4.679206    2.130175    2.693054\\
      2          6           0       -4.196084    1.896548    1.322426\\
      3          6           0       -5.362864    1.844689    0.374943\\
\end{bytebox}

The pattern \texttt{pat\_so} matches the literal text
\texttt{"Standard orientation:"} at the start of line 1 of this
excerpt. The byte offset reported by \fn{scan\_log} for this
match is the position of the \texttt{S} in
\texttt{"Standard"}, \emph{not} the start of the line (which has
26 leading spaces). \fn{read\_std\_orient} (\cref{sec:algo-read-meta})
relies on this offset and seeks to it before parsing.

For a frequency block, the layout is:

\begin{bytebox}
                     1                      2                      3  \\
                     A                      A                      A  \\
       Frequencies ---    11.3016   17.0158   19.3550   24.8910   25.6887\\
    Reduced masses ---     8.0483    8.4516    6.1589    8.1501    7.3141\\
   Force constants ---     0.0006    0.0014    0.0014    0.0030    0.0028\\
    IR Intensities ---     0.4340    0.0239    0.0809    0.7555    0.6085\\
 Coord Atom Element:\\
   1     1     7         -0.14255   0.17791  -0.01001   0.02676   0.09199\\
   2     1     7         -0.02809   0.05164  -0.03976   0.09096  -0.03367\\
\end{bytebox}

Several Reviewer-relevant observations:
\begin{itemize}
  \item The mode numbers (top: ``\texttt{1   2   3}'') appear on the
        line \emph{before} the symmetry labels (``\texttt{A A A A
        A}'') -- separated by the columns of 5 mode numbers (for
        hpmodes) or 3 (for standard). \fn{\_read\_block\_at\_freq\_line}
        seeks 300 bytes back from the \texttt{Frequencies ---} match
        and scans forward to recover the mode numbers from the
        all-digit line.
  \item The \texttt{Frequencies ---} pattern has \emph{three} dashes,
        not two; the \texttt{pat\_fr} regex matches both via
        \texttt{-\{2,\}}.
  \item The number of modes per block (5 in hpmodes, 3 in standard)
        is inferred from the count of frequencies parsed from the
        \texttt{Frequencies ---} line.
  \item The line \texttt{"Coord Atom Element:"} is the
        \texttt{hpmodes} signature: its presence sets \fn{is\_hp} = True.
        Standard blocks use the line \texttt{"Atom AN  X  Y  Z"} (or
        omit the column header entirely).
  \item Within the eigenvector body, columns 1-3 are
        \texttt{Coord} (1=$x$, 2=$y$, 3=$z$), \texttt{Atom number} and
        \texttt{Element atomic number}. The eigenvector components
        start in column 4.
\end{itemize}

\subsection{Edge cases handled by the implementation}
\label{sec:scan-edge-cases}

\paragraph{(E1) Duplicate matches from the tail-overlap region.}
Prevented by the \texttt{seen\_*} sets, as explained above.

\paragraph{(E2) Pattern at the very start of the file.}
The first chunk has \texttt{tail = b""}, so a match at byte 0 is
reported with offset 0 (no negative-offset bug).

\paragraph{(E3) Pattern in the very last 200 bytes of the file.}
The final chunk read is shorter than \texttt{chunk\_size\_bytes};
the loop reads as much as available and the final \texttt{tail}
update is harmless because the next \texttt{fh.read} returns
empty.

\paragraph{(E4) Empty file.}
\texttt{fh.read} returns \texttt{b""} immediately; the loop exits;
empty lists are returned. The caller (\cref{sec:algo-read-meta})
handles the all-empty case as a fatal error.

\paragraph{(E5) Missing \texttt{Frequencies} blocks entirely.}
If the user supplies a Gaussian log without a frequency calculation,
\texttt{fr\_off} comes back empty. The caller raises a clear error
``no normal modes found in the log -- did the frequency job
converge?''.

\paragraph{(E6) Truncated file (terminated mid-block).}
The scan completes (no error), but \fn{read\_all\_meta} may
produce \fn{BlockInfo}s with incomplete frequencies lists.
\fn{\_read\_block\_at\_freq\_line} returns \texttt{None} for any
block with zero frequencies, and the offending offsets are
dropped. The user is alerted in the run-log: ``$N$ frequency
offsets found but $M < N$ blocks parsable -- log may be truncated.''

\paragraph{(E7) Patterns inside arbitrary strings (false positives).}
The pattern \texttt{"Standard orientation:"} is unique to the
expected coordinate-block header in normal Gaussian output. A
user-supplied geometry comment line containing this substring
would falsely match, but Gaussian's writer never emits user
comments inside the SCF output region, and we have not
encountered a false positive in practice. The pattern
\texttt{pat\_fr} requires the leading-whitespace anchor plus the
two-dash terminator, which is a strong structural signature.

\paragraph{(E8) Multi-byte UTF-8 sequences in comments.}
The patterns are byte-strings, so they match byte-for-byte. A
UTF-8 encoded German umlaut in a user-supplied job title (e.g.\
``Mosbauer'' written as ``Mössbauer'' with \texttt{0xC3 0xB6})
takes two bytes -- but again, Gaussian never emits these in the
relevant blocks. The byte-pattern scan therefore never produces a
false negative.

\paragraph{(E9) Compressed input (\texttt{.log.gz}, \texttt{.log.xz}).}
\fn{scan\_log} does \emph{not} transparently handle compressed
logs. The user is expected to decompress first (this is documented
in the manual). A future extension could wrap the input with
\texttt{gzip.open}/\texttt{lzma.open}, but the resulting random-
access seek operations (needed by \fn{get\_eigvec}) would no longer
be O(1).

\subsection{Computational complexity and performance}

\paragraph{Time.}
$O(N_\text{bytes} \cdot N_\text{patterns})$, with
\fn{bytes.find}/\fn{re.finditer} both running at near memory-bandwidth
on modern hardware. The default 16 MiB chunk size empirically
gives the best throughput on NVMe SSDs ($\sim$~2--3 GiB/s). HDDs
benefit from larger chunks (32-64 MiB) to amortise seek overhead.

\paragraph{Memory.}
$O(\text{chunk\_size} + \text{tail\_overlap} + N_\text{matches})$.
For a 1 GiB log with $\sim$~10\,000 modes (i.e.\
$\sim$~10\,000 / 5 = 2000 hpmodes blocks, 1 Standard-Orientation,
1 normal-coordinates section), peak RAM is
$\sim$~$16 + 0.0002 + 0.024$~MiB ${} \approx 16$ MiB.

\paragraph{Live progress.}
The implementation prints a progress line every chunk:
\begin{bytebox}
Scan: 38.5\% 2614 MB/s ETA 7s
\end{bytebox}
The throughput is computed as \texttt{offset / elapsed}; the ETA
is \texttt{(total\_size - offset) / (offset / elapsed)}. The
progress is also logged to the \fn{RunLog} every 10\% for the
post-run audit trail.

\subsection{Cache integration}

The byte-offset triple \texttt{(so\_off, nc\_off, fr\_off)} is the
\emph{sole} cached artefact from the scan phase. The cache key is
the SHA-256 hash of:

\begin{lstlisting}[style=pseudo]
key = sha256(
    filepath +
    str(file_size_bytes) +
    str(file_mtime_ns) +
    f"py{sys.version_info.major}.{sys.version_info.minor}" +
    _CACHE_VERSION
)
\end{lstlisting}

The cache value is the pickled \texttt{(so\_off, nc\_off, fr\_off,
\_CACHE\_VERSION)} tuple, written to
\file{.scan\_cache\_\{key\}.pkl} in the current working directory.
On the next \fn{scan\_log} call:

\begin{itemize}
  \item If the cache file exists and the stored
        \texttt{\_CACHE\_VERSION} matches the current
        \texttt{\_CACHE\_VERSION}, the offsets are loaded directly.
  \item If the file mtime, file size, or Python version changes,
        the cache key changes -- the old cache file is simply
        ignored. (Disk-space cleanup is the user's responsibility;
        in practice the cache files are 50--500 KiB.)
  \item If \cfgkey{use\_cache} = False, the cache is bypassed
        entirely (neither read nor written).
\end{itemize}

This cache layer was the source of FUND 2
(\cref{sec:fund-2}): pre-v1.0.4 the file mtime was not included
in the key, so re-running an analysis after replacing the log
would silently use stale offsets.

% -------------------------------------------------------------------------
\section{Algorithm: \fn{read\_all\_meta}}
\label{sec:algo-read-meta}

\codref{logio.read\_all\_meta}

\fn{read\_all\_meta} iterates over the \fn{fr\_off} list from
\cref{sec:algo-scan-log} and parses the per-mode metadata of each
mode block. No eigenvector data are loaded here -- only header
information. The function also performs the assignment of mode
numbers to blocks (which is non-trivial because Gaussian's mode
numbering inside hpmodes blocks restarts from 1 in every section).

\subsection{Three-stage pipeline}

\begin{enumerate}
  \item \textbf{Per-block parsing.} For each
        \texttt{Frequencies ---}\ offset, call
        \fn{\_read\_block\_at\_freq\_line} to extract that block's
        frequencies, reduced masses, force constants, IR
        intensities, mode numbers, and the byte offset of the
        eigenvector data.
  \item \textbf{Mode-number resolution.} Recover correct global
        mode numbers using the section structure of the log (one
        section per ``\texttt{and normal coordinates:}'' marker).
  \item \textbf{HP/standard preference map.} Build the
        \texttt{best\_block} mapping that prefers hpmodes blocks
        over standard blocks for any mode that appears in both.
\end{enumerate}

\subsection{Per-block parsing (\fn{\_read\_block\_at\_freq\_line})}

The byte structure of one frequency block, as it appears in a
Gaussian \fn{hpmodes} log, is:

\begin{bytebox}
\hspace*{0.5em}offset $f$ \hfill (matched by \texttt{pat\_fr})\\
       Frequencies ---    11.3016   17.0158   ...\\
    Reduced masses ---     8.0483    8.4516   ...\\
   Force constants ---     0.0006    0.0014   ...\\
    IR Intensities ---     0.4340    0.0239   ...\\
 Coord Atom Element:  $\leftarrow$ HP-mode marker; \texttt{data\_offset} starts after this line\\
\hspace*{0.5em}offset $d$ \hfill (eigenvector body begins here)\\
   1     1     7         -0.14255   0.17791   ...\\
   2     1     7         -0.02809   0.05164   ...\\
\end{bytebox}

For a standard (non-\fn{hpmodes}) block, the equivalent layout is:

\begin{bytebox}
\hspace*{0.5em}offset $f$\\
                  Frequencies --    11.30   17.02   19.36\\
                  Red. masses --     8.05    8.45    6.16\\
                  Frc consts  --     0.00    0.00    0.00\\
                  IR Inten    --     0.43    0.02    0.08\\
 Atom AN      X      Y      Z       X      Y      Z       X      Y      Z\\
   1   7    -0.05  -0.05  -0.07    -0.06  -0.02   0.02    0.00  -0.01   0.03\\
\end{bytebox}

\fn{\_read\_block\_at\_freq\_line} parses this by the following procedure:

\begin{lstlisting}[style=pseudo,
  caption={Per-block metadata parser. Implementation:
           \protect\fn{logio.\_read\_block\_at\_freq\_line}.}]
input:  filepath, fr_offset
output: BlockInfo or None

bi <- BlockInfo(offset=fr_offset)

# (a) parse the Frequencies line itself
with open(filepath, "rb") as fh:
    fh.seek(fr_offset)
    line <- fh.readline().decode("utf-8", errors="replace")
    bi.is_hp <- bool(re.match(r"\s*Frequencies\s*---", line))
        # 3+ dashes = hpmodes; 2 dashes = standard
    bi.freqs <- _parse_dashes(line)
        # _parse_dashes strips "Frequencies ---" prefix and returns
        # a list of floats; missing or malformed values become NaN
    n <- len(bi.freqs)
    if n == 0:
        return None     # malformed block

    # (b) recover mode numbers from the line BEFORE fr_offset
    # Gaussian writes them as:    1     2     3     4     5
    # immediately preceding the symmetry labels (A A A ...)
    fh.seek(max(0, fr_offset - 300))
    prev_lines <- read all complete lines until fh.tell() > fr_offset
    foreach pline in reversed(prev_lines):
        parts <- pline.strip().split()
        if all parts are digit strings:
            bi.mode_nums <- [int(p) for p in parts]
            break
    if no all-digit line found or count mismatches n:
        bi.mode_nums <- [1, 2, ..., n]      # local fallback

    # (c) parse Red masses / Frc consts / IR ints + find data_offset
    fh.seek(fr_offset)
    fh.readline()    # skip Frequencies line
    for _ in range(60):           # 60-line window
        prev_pos <- fh.tell()     # B4-fix: record pos BEFORE readline
        line <- fh.readline()
        ll <- line.lower()
        if "red" in ll and "mass" in ll:
            bi.red_masses <- _parse_dashes(line)
        elif "frc" in ll or ("force" in ll and "const" in ll):
            bi.frc_consts <- _parse_dashes(line)
        elif "coord atom" in ll:
            bi.is_hp <- True
            bi.data_offset <- fh.tell()
            break
        elif "atom" in ll and (" an" in ll or "an " in ll):
            bi.is_hp <- False
            bi.data_offset <- fh.tell()
            break
        elif re.match(r"\s+\d+\s+\d+\s+[-\d.]", line):
            # first data row encountered without a column header
            bi.is_hp <- False
            bi.data_offset <- prev_pos    # B4-fix
            break

    # (d) pad missing optional fields to length n
    for lst in (bi.red_masses, bi.frc_consts):
        while len(lst) < n: lst.append(0.0)
    while len(bi.syms) < n: bi.syms.append("A")

return bi
\end{lstlisting}

\paragraph{The 60-line window.} Gaussian can include optional
extra header rows between the \texttt{Frequencies ---} line and the
eigenvector body, such as ``\texttt{Raman Activ --}'',
``\texttt{Depolar (P)}'', NMR shielding, etc. The 60-line window
is empirically large enough to cover all observed Gaussian-16
output variants. If 60 lines pass without finding either marker
or a first data row, the block is treated as malformed (\fn{None}
returned).

\paragraph{Bugfix B4: position before \texttt{readline}.} The
\texttt{prev\_pos} is recorded \emph{before} the \texttt{readline}
call. This is critical for the no-header standard-block case
(the \texttt{elif re.match} branch): when the first data row is
encountered without a preceding \texttt{Atom AN} header line, the
\texttt{data\_offset} must point at the start of that data row, not
at the position after reading it. The previous version of the
code used \texttt{fh.tell() - len(line.encode())}, which subtly
mis-counted by the encoded-vs-raw byte-length difference in the
presence of non-ASCII characters (rare but possible).

\paragraph{Bugfix B5: \texttt{data\_offset = -1} sentinel.} If
neither marker nor data row is found within the 60-line window,
\texttt{data\_offset} remains at its default $-1$. The downstream
\fn{get\_eigvec} explicitly checks this sentinel and raises a
clear error rather than seeking to byte $-1$ (which would silently
\texttt{seek\_end} on some operating systems).

\subsection{Mode-number resolution and the section structure}

Each \texttt{and normal coordinates:} marker in the log opens a new
\emph{section}. Inside a section, Gaussian numbers the modes
starting from 1 again. To recover global mode numbers across the
whole log, \fn{read\_all\_meta} performs this remapping:

\begin{lstlisting}[style=pseudo,
  caption={Mode-number resolution across sections. Implementation:
           \protect\fn{logio.read\_all\_meta}.}]
input:  raw_blocks (per-block parsed by _read_block_at_freq_line),
        nc_sorted  (sorted list of "and normal coordinates:" offsets)
output: blocks      (same blocks, mode_nums corrected)

section_count <- {}      # nc_offset -> number of blocks already seen
                         # in that section
foreach bi in raw_blocks:
    sec <- max(nc for nc in nc_sorted if nc <= bi.offset)
        # which section does this block belong to?
    section_count.setdefault(sec, 0)
    nm <- len(bi.mode_nums)

    # Heuristic: if the parsed mode numbers are [1, 2, ..., nm] AND
    # we've already seen blocks in this section, then the parser
    # found local mode numbers, not global. Renumber globally.
    fallback <- (bi.mode_nums == [1, 2, ..., nm]
                 and section_count[sec] > 0)

    # Same fallback if parser failed altogether (mode_nums empty
    # or starts with <= 0)
    if fallback or not bi.mode_nums or bi.mode_nums[0] <= 0:
        start <- section_count[sec] * nm + 1
        bi.mode_nums <- [start, start+1, ..., start+nm-1]

    section_count[sec] += 1
    blocks.append(bi)
\end{lstlisting}

\paragraph{Why the section-aware remapping.}
A Gaussian job that writes both an \fn{hpmodes} frequency analysis
and a separate \fn{Freq=Anharm} job will have two
\texttt{and normal coordinates:} markers, two restart-at-1
mode numberings, and possibly overlapping global mode numbers.
The remapping ensures that the global mode numbers are unique
across the whole log, which is what \fn{best\_block} and
\fn{cand\_map} require.

\subsection{HP/standard preference}

The final step builds the two output maps:

\begin{lstlisting}[style=pseudo]
best     <- {}    # mode_num -> preferred BlockInfo (HP wins)
cand_map <- {}    # mode_num -> list of all candidate BlockInfos
foreach bi in blocks:
    foreach mn in bi.mode_nums:
        if mn not in best:
            best[mn] <- bi
        elif bi.is_hp and not best[mn].is_hp:
            best[mn] <- bi           # promote HP over standard
        cand_map.setdefault(mn, []).append(bi)
\end{lstlisting}

\fn{best\_block} is used by the main mode-loop as the fast path;
\fn{cand\_map} is used by \fn{analyze\_mode\_with\_fallback}
when the preferred block's eigenvector data fail to parse (e.g.\
truncated file, corrupt block) -- the loop then tries the next-
best candidate from \fn{cand\_map}.

\subsection{Sentinels and missing data}

A missing reduced mass (Gaussian prints ``\texttt{********}'' when
the value overflows the fixed-width format -- happens for some very
soft modes) is parsed as $0.0$ by \fn{\_parse\_dashes} (which
returns the empty list on parse failure and is then padded with
zeros up to length $n$). Downstream code treats $\texttt{red\_mass}
= 0$ as fatal at the mode-selection stage -- the mode is skipped
with a warning, and the user is alerted in the run-log.

\paragraph{Why both \fn{best\_block} and \fn{cand\_map}:}
\fn{best\_block} is the fast path for the common case (one block per
mode). \fn{cand\_map} is the fallback list used by
\fn{analyze\_mode\_with\_fallback} (\cref{sec:algo-mode-loop}) if
the preferred block's eigenvector data are unreadable (truncated
file, etc.).

\subsection{Live progress reporting}

The function prints progress every 500 blocks and logs to
\fn{RunLog} at the same cadence. For a 200-MiB log with $\sim$~10\,000
frequency blocks, the metadata phase typically completes in
1--3 seconds on an NVMe SSD.

\subsection{Edge cases}

\paragraph{(M1) Block at byte 0.}
\fn{\_read\_block\_at\_freq\_line} clamps the back-scan to
\texttt{max(0, fr\_offset - 300)} so it does not seek before the
file start. If the block is at byte 0 there is no preceding mode-
number line; the fallback list \texttt{[1, 2, ..., n]} is used.

\paragraph{(M2) Frequency line with NaN values.}
A line containing ``\texttt{NaN}'' is parsed as float NaN
(Python's \texttt{float("nan")} accepts the string ``nan'').
NaN frequencies are filtered out at the mode-selection stage
(\cref{sec:algo-mode-loop}).

\paragraph{(M3) Standard block without column header.}
Some Gaussian versions (older) omit the \texttt{Atom AN} line and
go straight from the IR-Intensities line to the eigenvector body.
The fallback regex \texttt{\textbackslash{}s+\textbackslash{}d+
\textbackslash{}s+\textbackslash{}d+\textbackslash{}s+[-\textbackslash{}d.]}
detects a first data row by pattern (Coord, Atom, Element, then
numeric) and sets \texttt{is\_hp = False}, \texttt{data\_offset = prev\_pos}.

\paragraph{(M4) Concatenated logs.}
A user who concatenates two separate Gaussian frequency-job logs
(e.g.\ \texttt{cat job1.log job2.log > combined.log}) will have
two ``\texttt{and normal coordinates:}'' sections. The
section-aware remapping handles this correctly -- global mode
numbers are unique across the combined log.

% -------------------------------------------------------------------------
\section{Algorithm: cluster discovery}
\label{sec:algo-cluster-discovery}

\codref{geometry.find\_cluster, geometry.find\_all\_clusters}

The cluster-discovery algorithm finds every [2Fe-2S] cluster in a
molecular system using two cutoff parameters: a Fe-S bond cutoff
\cfgkey{fe\_s\_cutoff} (default $3.0$~\AA{}) and an Fe-Fe distance
cutoff \cfgkey{fe\_fe\_cutoff} (default $3.5$~\AA{}). The output is
a list of clusters; each cluster is a 4-tuple of Gaussian centres
$(\text{Fe}_1, \text{Fe}_2, \text{S}_1, \text{S}_2)$ in canonical
order.

\begin{lstlisting}[style=pseudo,
  caption={Multi-cluster discovery.
           Implementation: \protect\fn{geometry.find\_all\_clusters}.}]
input:  atoms (list of {symbol, x, y, z, center}),
        fe_s_cutoff, fe_fe_cutoff
output: clusters : list of (fe1_c, fe2_c, s1_c, s2_c)

# Step 1: collect all Fe and S indices
fe_indices <- [i for i, a in enumerate(atoms) if a.symbol == "Fe"]
s_indices  <- [i for i, a in enumerate(atoms) if a.symbol == "S"]

# Step 2: build the Fe-Fe adjacency by distance
fe_pairs <- []
foreach (i, j) in unordered_pairs(fe_indices):
    if euclidean(atoms[i], atoms[j]) <= fe_fe_cutoff:
        fe_pairs.append((i, j))

# Step 3: for each Fe-Fe pair, find the bridging sulfurs
clusters_raw <- []
foreach (i, j) in fe_pairs:
    bridging <- []
    foreach k in s_indices:
        d_ik <- euclidean(atoms[i], atoms[k])
        d_jk <- euclidean(atoms[j], atoms[k])
        if d_ik <= fe_s_cutoff and d_jk <= fe_s_cutoff:
            bridging.append((k, d_ik + d_jk))
    if len(bridging) >= 2:
        # Pick the two closest-bridging sulfurs
        bridging.sort(by=second-element)
        s1, s2 <- bridging[0][0], bridging[1][0]
        clusters_raw.append((i, j, s1, s2))

# Step 4: canonicalize ordering (smaller Gaussian centre first
# within Fe pair and within S pair)
clusters <- []
foreach (i, j, s1, s2) in clusters_raw:
    fe1_c, fe2_c <- canonical_order(atoms[i].center, atoms[j].center)
    s1_c,  s2_c  <- canonical_order(atoms[s1].center, atoms[s2].center)
    clusters.append((fe1_c, fe2_c, s1_c, s2_c))

# Step 5: deduplicate (the same cluster may have been discovered
# via multiple Fe-Fe orderings)
clusters <- unique(clusters, key=sorted-tuple-of-all-four-centres)

# Step 6: sort by smallest Gaussian centre across the cluster
# (deterministic cluster numbering across runs)
clusters.sort(by=min-centre-in-tuple)

return clusters
\end{lstlisting}

\paragraph{The \fn{strict\_cluster} option.}
With \cfgkey{strict\_cluster} = true, additional sanity checks are
applied:
\begin{itemize}
  \item The two bridging S atoms must form an Fe$_2$S$_2$ rhombus
        whose two diagonals (Fe-Fe and S-S) are within $\pm 30\%$
        of the standard rhombic-cluster geometry.
  \item The four atoms must be approximately coplanar
        (smallest SVD singular value $\leq 0.1$ \AA{}).
\end{itemize}
This guards against false-positive cluster detection in systems with
two adjacent [Fe(SCys)$_4$]-like centres that incidentally pass the
distance-only test but are not rhombic [2Fe-2S].

\paragraph{Cluster ordering for multi-cluster systems.}
For systems with $K > 1$ clusters (e.g.\ a Rieske + Mo-bisMGD enzyme
where two [2Fe-2S] clusters are present), the resulting list of
clusters is sorted by the smallest Gaussian centre of each tuple.
This makes the cluster index reproducible: across runs, ``cluster
1'' refers to the same physical cluster.

\subsection{Worked example: cluster discovery on a 2-cluster system}

We trace the algorithm on a synthetic 2-cluster geometry with 4 Fe
and 6 S atoms in the QM region:

\begin{center}
\begin{tabular}{ccccc}
\toprule
Centre & Sym & $x$ & $y$ & $z$ \\
\midrule
12 & Fe & $+1.35$  & $0$    & $0$    \\
13 & Fe & $-1.35$  & $0$    & $0$    \\
22 & Fe & $+1.35$  & $0$    & $8.00$ \\
23 & Fe & $-1.35$  & $0$    & $8.00$ \\
14 & S  & $0$      & $+1.08$ & $0$   \\
15 & S  & $0$      & $-1.08$ & $0$   \\
24 & S  & $0$      & $+1.08$ & $8.00$ \\
25 & S  & $0$      & $-1.08$ & $8.00$ \\
99 & S  & $+5.00$  & $0$    & $4.00$ \\
\bottomrule
\end{tabular}
\end{center}

\textbf{Step 2 -- Fe-Fe pairs:} All
$\binom{4}{2} = 6$ pairwise Fe distances:
\[
  d(12,13) = 2.70,\quad d(22,23) = 2.70,\quad
  d(12,22) = d(13,23) = 8.00, \dots
\]
With \cfgkey{fe\_fe\_cutoff} = 3.5~\AA{}, only the in-plane pairs
$(12, 13)$ and $(22, 23)$ pass. The cross-cluster pairs are filtered.

\textbf{Step 3 -- bridging S per Fe-Fe pair:}
For $(12, 13)$ check each S:
\begin{itemize}
  \item S 14: $d(12, 14) = \sqrt{1.35^2 + 1.08^2} = 1.73$,
        $d(13, 14) = 1.73$. Both $\leq 3.0$. \checkmark
  \item S 15: $d(12, 15) = 1.73$, $d(13, 15) = 1.73$. \checkmark
  \item S 24: $d(12, 24) = 8.14$. Rejected.
  \item S 99: $d(12, 99) = 5.66$. Rejected.
\end{itemize}
Bridging: $\{14, 15\}$ with $d_{ik} + d_{jk} = 3.46$ each. Two
closest: $14, 15$. Cluster raw: $(12, 13, 14, 15)$.

Same procedure for $(22, 23)$ yields $(22, 23, 24, 25)$.

\textbf{Step 4 -- canonicalize:}
For first cluster, Fe centres sorted: $(12, 13) \to (12, 13)$ (already
sorted); S centres sorted: $(14, 15) \to (14, 15)$.
Tuple: $(12, 13, 14, 15)$.

\textbf{Step 5 -- deduplicate:} No duplicates.

\textbf{Step 6 -- sort by min centre:} Cluster 1 has min-centre 12,
cluster 2 has min-centre 22. Final order:
\[
  \boxed{[(12, 13, 14, 15),\; (22, 23, 24, 25)]}.
\]

S atom 99 (a spectator S not bridging any Fe-Fe) is correctly
excluded.

\subsection{Edge cases handled}

\paragraph{(C1) Single isolated Fe atom.}
If only one Fe is present, the Fe-Fe-pair list is empty
($\binom{1}{2} = 0$). The function returns the empty list and the
caller raises ``no [2Fe-2S] clusters found in the geometry''.

\paragraph{(C2) Three Fe atoms in a triangle.}
A linear or near-linear arrangement of 3 Fe atoms (e.g.\ a partial
[3Fe-4S] cluster's projection) would yield 3 Fe-Fe pairs. Each
pair is independently checked for bridging S. The
\cfgkey{strict\_cluster} option filters out non-rhombic results.

\paragraph{(C3) Fe-Fe distance just above cutoff.}
A cluster with $d_\text{FeFe} = 3.51$~\AA{} would be missed at the
default cutoff. The cutoff is configurable; the user can raise it
to 4.0~\AA{} for unusually stretched clusters (e.g.\
super-reduced [2Fe-2S]$^0$ in some model compounds). The
\fn{strict\_cluster} mode then catches false positives at the
rhombicity stage.

\paragraph{(C4) Shared bridging sulfur between two Fe-Fe pairs.}
In a [4Fe-4S] cluster, the same S atom would bridge multiple
Fe-Fe pairs. The algorithm produces 4 candidate ``[2Fe-2S]''
clusters (one per Fe-Fe pair), each sharing 2 of its S atoms with
its neighbours. Step 5 (dedup by sorted-tuple) does \emph{not}
remove these because the tuples are not identical. The
\fn{strict\_cluster} mode's rhombicity check filters them out
(a [4Fe-4S] sub-rhombus is highly distorted from ideal). For
systems containing \emph{both} [4Fe-4S] and [2Fe-2S] clusters,
the user should run with \cfgkey{strict\_cluster} = True.

\paragraph{(C5) Hydrogenated $\mu$-S.}
A protonated bridging sulfur (a $\mu$-SH state, observed e.g.\ in
some Rieske-type and ferredoxin model studies of PCET) has the
same S atom, just with one more H bound. Cluster discovery is
unaware of H atoms and treats this S identically to its
non-protonated version. (Bridging-S protonation status is not
explicitly
classified by the pipeline; it manifests indirectly via the
PCET-channel modulations whenever an H atom sits within bonding
distance of $\mu$-S.)

\subsection{Computational complexity}

\paragraph{Time.}
$O(N_\text{Fe}^2 \cdot N_\text{S})$ for the cluster scan; for a
typical QM region with $N_\text{Fe} \leq 4$ and $N_\text{S} \leq 8$,
this is $\leq 128$ pairwise distance computations -- negligible.
For larger QM regions, the scan remains polynomial and well below
1\% of total runtime.

\paragraph{Memory.}
$O(K)$ where $K$ is the number of detected clusters; typically
$K \leq 2$.

% -------------------------------------------------------------------------
\section{Algorithm: PDB-Gaussian matching (Kabsch alignment)}
\label{sec:kabsch}

\codref{geometry.kabsch\_align}

When a PDB file is supplied, each PDB atom must be matched to the
Gaussian centre that represents it. The PDB atoms are usually a
subset of the Gaussian atoms (a typical QM/MM setup has the
[2Fe-2S]-binding pocket as QM atoms only). The matching is done in
two stages: first the rotation $R$ and translation $t$ that
optimally align the PDB heavy-atom positions onto the Gaussian
heavy-atom positions are computed via Kabsch \cite{Kabsch1976},
then each PDB atom is mapped to its nearest Gaussian neighbour
under the aligned coordinates.

\subsection{Stage 1: optimal alignment}

The Kabsch algorithm computes the rotation $R$ that minimises the
RMSD between two sets of $N$ points $\{\mathbf{p}_i\}$ and
$\{\mathbf{q}_i\}$:
\begin{equation}
  R^\star \;=\; \arg\min_{R \in SO(3)}
                 \frac{1}{N}\sum_{i=1}^{N}
                 \lVert R\mathbf{p}_i - \mathbf{q}_i\rVert^2.
\end{equation}
The closed-form solution uses an SVD of the cross-covariance matrix
$\mathbf{H} = \tilde{\mathbf{P}}^\top \tilde{\mathbf{Q}}$
(after centring both sets):
\begin{equation}
  \mathbf{H} = \Umat\,\Smat\,\Vmat^\top,
  \qquad
  d \;=\; \text{sign}(\det(\Vmat\,\Umat^\top)),
  \qquad
  R^\star = \Vmat\,\text{diag}(1, 1, d)\,\Umat^\top.
\end{equation}
The $d$ factor enforces $\det(R^\star) = +1$, ruling out
improper rotations (which would correspond to a chirality flip --
not physical for a protein backbone).

\begin{lstlisting}[style=pseudo,
  caption={Kabsch optimal-rotation algorithm. Implementation:
           \protect\fn{geometry.kabsch\_align}.}]
input:  pdb_points : (N, 3) array  -- PDB heavy-atom positions
        gauss_points : (N, 3) array -- corresponding Gaussian positions
output: R : (3, 3) rotation matrix
        t : (3,)   translation vector
        rmsd : float -- post-alignment RMSD

mean_p <- mean(pdb_points, axis=0)
mean_q <- mean(gauss_points, axis=0)
P_centered <- pdb_points - mean_p
Q_centered <- gauss_points - mean_q

H <- P_centered.T @ Q_centered    # (3, 3) cross-covariance

U, S, Vt <- np.linalg.svd(H)
d <- sign(det(Vt.T @ U.T))
correction <- diag([1, 1, d])
R <- Vt.T @ correction @ U.T

t <- mean_q - R @ mean_p

# diagnostics
aligned <- (R @ pdb_points.T).T + t
rmsd <- sqrt(mean(sum((aligned - gauss_points)**2, axis=1)))
return R, t, rmsd
\end{lstlisting}

\subsection{Stage 2: nearest-neighbour mapping}

After alignment, each PDB atom is mapped to its nearest Gaussian
neighbour, subject to the tolerance \cfgkey{coord\_match\_tol}
(default $1.0$~\AA{}). Matches beyond the tolerance are reported as
unmatched and excluded from the residue-label mapping.

\begin{lstlisting}[style=pseudo,
  caption={Nearest-neighbour mapping post-Kabsch.}]
input:  pdb_atoms (aligned), gauss_atoms, tol
output: pdb_to_center : dict[pdb_idx -> gaussian_center] or None

pdb_to_center <- {}
foreach pdb_idx, p in enumerate(pdb_atoms):
    best_g_c    <- None
    best_dist   <- +inf
    foreach g in gauss_atoms:
        d <- euclidean(p, g)
        if d < best_dist:
            best_dist <- d
            best_g_c  <- g.center
    if best_dist <= tol:
        pdb_to_center[pdb_idx] <- best_g_c
    else:
        # log a warning, keep mapping as None
        pdb_to_center[pdb_idx] <- None

return pdb_to_center
\end{lstlisting}

\paragraph{Multi-cluster behaviour.}
For multi-cluster systems, the Kabsch alignment is performed once
on the full heavy-atom set (typical proteins, where the QM region
covers all clusters), not per-cluster. This is important because
otherwise the alignment would over-fit to a single cluster and
mislabel atoms far from it.

\paragraph{Match-tolerance diagnostic.}
The post-Kabsch RMSD is recorded in the run-log. RMSD values larger
than $0.5$~\AA{} usually indicate either a wrong PDB file (different
chain or different protomer of the protein), a PDB with engineered
mutations, or a Gaussian geometry that diverged substantially from
the crystal structure during optimisation. The user is alerted in
the run-log.

\subsection{The chirality-flip determinant detail}

The factor $d = \text{sign}(\det(\Vmat \Umat^\top))$ is critical
and subtle. Without it, the SVD-based solution would sometimes
return an \emph{improper} rotation (a rotation composed with a
reflection), depending on the relative chirality of the two point
clouds. Such an improper rotation has $\det = -1$ and physically
corresponds to mirroring the molecule.

\paragraph{When does $d = -1$ occur?}
Mathematically: when the two point clouds are so different that
the optimal alignment is closer to a mirror image than to any
rotation. In practice this happens in two scenarios:

\begin{enumerate}
  \item Numerical: when the point cloud is near-coplanar and the
        smallest singular value of $\Hmat$ is comparable to
        machine epsilon. The sign of the determinant becomes
        random under FP arithmetic.
  \item Physical: the PDB and Gaussian structures genuinely have
        opposite chirality, e.g.\ if the user supplied the wrong
        enantiomer or if the PDB was processed through a
        chirality-flipping tool.
\end{enumerate}

In either case, the $d$ factor enforces $\det(R) = +1$ by
flipping the sign of the last column of $\Vmat$, choosing the
proper-rotation branch of the optimisation. The cost is a
slightly larger residual RMSD; the benefit is a physically valid
rotation matrix.

\subsection{Worked example: 4-atom alignment with two flips}

Two synthetic point clouds:

\begin{center}
\begin{tabular}{lccc|ccc}
\toprule
Atom & PDB $x$ & PDB $y$ & PDB $z$ & Gauss $x$ & Gauss $y$ & Gauss $z$ \\
\midrule
Fe$_1$ & $+1.0$ & $0$    & $0$    & $0$     & $+1.0$ & $0$    \\
Fe$_2$ & $-1.0$ & $0$    & $0$    & $0$     & $-1.0$ & $0$    \\
S$_1$  & $0$    & $+1.0$ & $0$    & $-1.0$  & $0$    & $0$    \\
S$_2$  & $0$    & $-1.0$ & $0$    & $+1.0$  & $0$    & $0$    \\
\bottomrule
\end{tabular}
\end{center}

The PDB is the Gaussian geometry rotated by $-90^\circ$ around the
$z$-axis. Both centroids are at the origin.

\textbf{Step 1 -- cross-covariance:}
\[
  \Hmat = \tilde\Pmat^\top \tilde\Qmat
   = \begin{pmatrix}
       (1)(0) + (-1)(0) + (0)(-1) + (0)(1) & \dots & \dots \\
       \dots & \dots & \dots
     \end{pmatrix}
\]
Computing all 9 entries:
\[
  \Hmat = \begin{pmatrix}
            0 & 0 & 0 \\
            -2 & 0 & 0 \\
            0 & 0 & 0
          \end{pmatrix}.
\]
Wait -- this can't be right. Let me recompute more carefully for
the $H_{12}$ entry (column 2 of $\tilde\Pmat^\top$ times column 1
of $\tilde\Qmat$):
\[
  H_{21} = \sum_i \tilde P_{i,2} \tilde Q_{i,1}
        = (0)(0) + (0)(0) + (1)(-1) + (-1)(1) = -2.
\]
\[
  H_{12} = \sum_i \tilde P_{i,1} \tilde Q_{i,2}
        = (1)(1) + (-1)(-1) + (0)(0) + (0)(0) = 2.
\]
So $\Hmat = \begin{pmatrix} 0 & 2 & 0 \\ -2 & 0 & 0 \\ 0 & 0 & 0 \end{pmatrix}$.

\textbf{Step 2 -- SVD:}
$\Hmat$ has singular values $\sigma_1 = \sigma_2 = 2$, $\sigma_3 = 0$.
$\Umat$, $\Vmat$ are rotation matrices that diagonalise it.
A valid decomposition is
\[
  \Umat = \begin{pmatrix} 1 & 0 & 0 \\ 0 & 1 & 0 \\ 0 & 0 & 1 \end{pmatrix},
  \quad
  \Smat = \text{diag}(2, 2, 0),
  \quad
  \Vmat^\top = \begin{pmatrix} 0 & 1 & 0 \\ -1 & 0 & 0 \\ 0 & 0 & 1 \end{pmatrix}.
\]

\textbf{Step 3 -- determinant correction:}
$\det(\Vmat \Umat^\top) = \det(\Vmat) = +1$ (in this case),
so $d = +1$ -- no correction needed.

\textbf{Step 4 -- rotation:}
$R = \Vmat\,\Umat^\top
   = \begin{pmatrix} 0 & -1 & 0 \\ 1 & 0 & 0 \\ 0 & 0 & 1 \end{pmatrix}$.

This is the rotation matrix for a $+90^\circ$ rotation about $\zhat$,
which is the inverse of the $-90^\circ$ that produced the PDB --
correct!

\textbf{Step 5 -- residual RMSD:}
After applying $R$ to each PDB point and comparing to Gaussian:
$R \cdot (1,0,0) = (0,1,0)$ \checkmark matches Gauss Fe$_1$. Similarly
for the others. RMSD = 0 exactly.

\paragraph{Variant with $d = -1$.}
If we had Gauss S$_1 = (+1, 0, 0)$ and S$_2 = (-1, 0, 0)$ instead
(swapped), the cross-covariance would yield $\det(\Vmat\Umat^\top)
= -1$. The correction would flip the sign of the last column of
the SVD product, giving a proper rotation with a residual RMSD of
$\sqrt{2}$ (the mirror cannot be fixed by rotation alone; the
algorithm picks the best-rotation-only fit). The user would see
``\texttt{coord\_match\_rmsd = 1.41 \AA{} -- likely wrong PDB
chirality}'' in the run-log.

\subsection{Edge cases}

\paragraph{(K1) Single-point or 2-point alignment.}
With $N < 3$ points, the alignment is geometrically under-
determined (a 2-point set has a free rotation around the
joining axis). The function refuses to align with $N < 3$ and
returns identity $R = \mathbf{I}, t = 0$, with a warning.

\paragraph{(K2) Co-linear points.}
With all points along a line, $\Hmat$ has rank 1 and the
alignment is determined only up to rotation around the line. The
SVD returns one large and two small singular values. The
algorithm proceeds (the resulting $R$ is one of the infinitely
many valid alignments) but the run-log warns
``\texttt{coord\_match singular values: 2.0, 1e-15, 1e-15}'',
alerting the user.

\paragraph{(K3) Co-planar points.}
With all points in a plane, $\sigma_3$ is small but $\Hmat$ has
rank 2; the alignment is well-determined. The chirality-flip
correction $d$ is the relevant safeguard.

\paragraph{(K4) Tolerance failure for one or two atoms.}
A single PDB atom that fails the tolerance check (e.g.\ because
of a renaming inconsistency or because the QM optimisation moved
it by $> 1$~\AA{}) is reported as unmatched. The function does not
abort -- it continues mapping the other atoms. The user is
alerted; downstream code treats the unmatched atom as having no
residue label.

\paragraph{(K5) Many-to-one matching.}
After Kabsch alignment, two PDB atoms can be nearest-neighbour-
matched to the same Gaussian centre if they are both within
\cfgkey{coord\_match\_tol} of it. This is the ``ambiguous match''
warning of FUND 9. The default behaviour is to log the
ambiguity and keep both mappings (the duplicate is post-filtered
in the residue-label resolution stage).

% -------------------------------------------------------------------------
\section{Algorithm: secondary-structure detection}
\label{sec:algo-sse}

\codref{geometry.detect\_sse\_dssp, geometry.detect\_sse\_phipsi}

Two parallel secondary-structure (SSE) classifiers are implemented.
DSSP is the gold standard but requires either the executable
installed locally or pre-computed records in the PDB file. As a
fallback, the lightweight $\phi/\psi$ dihedral-only classifier is
used.

\subsection{Algorithm: DSSP records from PDB}

\codref{geometry.detect\_sse\_dssp}

DSSP \cite{KabschSander1983} classifies each residue into one of
8 categories (H, B, E, G, I, T, S, blank), based on hydrogen-bond
patterns in the backbone. The implementation reads the SSE records
directly from the PDB header (most modern PDB files carry these in
the \texttt{HELIX} and \texttt{SHEET} records):

\begin{lstlisting}[style=pseudo,
  caption={DSSP-record extraction from PDB. Implementation:
           \protect\fn{geometry.detect\_sse\_dssp}.}]
input:  pdb_path
output: sse_per_residue : dict[(chain, resnum) -> sse_type]

sse_per_residue <- {}
with open(pdb_path) as f:
    foreach line in f:
        if line starts with "HELIX":
            chain     <- line[19]
            resnum_a  <- parse_int(line[21:25])
            resnum_b  <- parse_int(line[33:37])
            foreach r in range(resnum_a, resnum_b + 1):
                sse_per_residue[(chain, r)] <- "H"
        if line starts with "SHEET":
            chain     <- line[21]
            resnum_a  <- parse_int(line[22:26])
            resnum_b  <- parse_int(line[33:37])
            foreach r in range(resnum_a, resnum_b + 1):
                sse_per_residue[(chain, r)] <- "E"

return sse_per_residue
\end{lstlisting}

\subsection{Algorithm: $\phi/\psi$-only classifier}

\codref{geometry.detect\_sse\_phipsi}

If no SSE records are in the PDB and DSSP is not installed, the
fallback computes the $\phi$ and $\psi$ dihedral angles per residue
and applies the Ramachandran-region cutoffs:
\begin{itemize}
  \item $\alpha$-helix: $\phi \in [-100^\circ, -30^\circ]$,
        $\psi \in [-90^\circ, +20^\circ]$.
  \item $\beta$-strand: $\phi \in [-180^\circ, -60^\circ]$,
        $\psi \in [+90^\circ, +180^\circ]$ or $[-180^\circ, -120^\circ]$.
  \item All others: classified as ``loop''.
\end{itemize}
Additionally, a minimum-length filter ($\geq 3$ consecutive
$\alpha$ residues, $\geq 2$ consecutive $\beta$ residues) is
applied to avoid spurious single-residue classifications.

\begin{lstlisting}[style=pseudo,
  caption={$\phi/\psi$ secondary-structure classifier.
           Implementation: \protect\fn{geometry.detect\_sse\_phipsi}.}]
input:  pdb_atoms (with backbone N, CA, C identified)
output: sse_per_residue : dict[(chain, resnum) -> sse_type]

raw_sse <- {}
foreach residue (chain, resnum) in pdb_atoms:
    prev_C  <- atom(chain, resnum - 1, "C") or skip
    this_N  <- atom(chain, resnum, "N")
    this_CA <- atom(chain, resnum, "CA")
    this_C  <- atom(chain, resnum, "C")
    next_N  <- atom(chain, resnum + 1, "N") or skip
    phi <- dihedral(prev_C, this_N, this_CA, this_C)
    psi <- dihedral(this_N, this_CA, this_C, next_N)
    raw_sse[(chain, resnum)] <- classify(phi, psi)

# Apply min-length filter: drop isolated SSE-classified residues
sse_per_residue <- min_length_filter(raw_sse, alpha_min=3, beta_min=2)
return sse_per_residue
\end{lstlisting}

\paragraph{Comparison of the two classifiers.}
DSSP is more accurate (especially for distinguishing $\alpha$-helices
from 3$_{10}$- and $\pi$-helices), but the $\phi/\psi$ classifier
captures the majority of $\alpha$- and $\beta$-character correctly
for typical $> 30$-residue chains. Both classifiers feed the same
downstream SSE-amplitude analysis; the choice is logged.

\subsection{The fallback cascade}

The full SSE-detection pipeline tries three sources in order, stopping
at the first that produces a non-empty result:

\begin{enumerate}
  \item \textbf{Pre-computed records in the PDB header.} If the PDB
        contains \texttt{HELIX} and/or \texttt{SHEET} records, these
        are parsed directly. This is the fastest path (no
        computation, no external dependency). Most PDBs deposited
        in the RCSB Protein Data Bank since the 1990s carry these
        records.
  \item \textbf{External DSSP executable.} If the
        \cfgkey{dssp\_executable} path is set and the binary exists
        and is runnable, DSSP is invoked as a subprocess. The DSSP
        output (an 8-class assignment) is parsed; H, G, I are
        merged into ``$\alpha$''; E, B are merged into ``$\beta$'';
        the rest are ``loop''. Failure modes (non-zero exit code,
        empty stdout, missing residues) trigger the next fallback.
  \item \textbf{$\phi/\psi$ classifier on PDB coordinates.} The
        in-package fallback, requiring only the backbone N, CA, C
        atoms (which are always present in a standard PDB).
\end{enumerate}

The chosen source is recorded in the run-log as
``\texttt{sse\_source = pdb\_records}'', ``\texttt{sse\_source =
dssp\_executable}'', or ``\texttt{sse\_source = phipsi\_fallback}''.

\subsection{Edge cases}

\paragraph{(S1) PDB with no SSE records and no DSSP.}
The $\phi/\psi$ classifier is used. Some residues at the chain
termini (missing $\phi$ at residue 1, missing $\psi$ at the last
residue) get classified as ``loop'' regardless of their actual
backbone geometry. This is acceptable because terminal residues
are typically already disordered.

\paragraph{(S2) Cis-proline.}
A cis-Pro at $\omega \approx 0$ has $\psi$ shifted into the
``$\alpha$-like'' region; the classifier may mislabel it. This is
a known limitation; the DSSP-records source is more accurate.

\paragraph{(S3) Glycine.}
Gly residues lack a side chain and occupy a larger Ramachandran
region than other amino acids. The default cutoffs do not have a
Gly-specific branch; Gly residues in the ``left-handed
$\alpha$'' region (positive $\phi$) are classified as loop, which
is usually correct.

\paragraph{(S4) Missing backbone atoms.}
If the PDB is missing a backbone N, CA, or C atom for some
residue (e.g.\ a CA-only minimal PDB), the $\phi/\psi$ calculation
fails for that residue and adjacent residues. They are classified
as ``unknown'' (UNK in the output) and the user is warned.

\paragraph{(S5) Min-length filter removing real short helices.}
The 3-residue $\alpha$-min and 2-residue $\beta$-min filter
deliberately removes singletons. Real 3$_{10}$-helix turns and
$\beta$-bulges may be lost. The filter is necessary because the
$\phi/\psi$ classifier produces $\sim$~10--15\% spurious singletons.
The user can adjust the filter parameters via
\cfgkey{sse\_alpha\_min\_len} and \cfgkey{sse\_beta\_min\_len}.

\paragraph{(S6) Multiple chains.}
The classifier processes one chain at a time. Inter-chain
H-bonds (which DSSP would detect) are not visible to the
$\phi/\psi$ fallback, so cross-chain $\beta$-sheets are
classified residue-by-residue as if each chain were independent.
For most NRVS analyses this is acceptable (the cluster is in one
chain).

\subsection{Connection to the SSE-amplitude analysis}

The output \texttt{sse\_per\_residue} dict feeds two downstream
consumers:

\begin{enumerate}
  \item \texttt{export\_sse\_amplitude\_groups}: builds per-mode
        SSE-aggregated amplitude sheets. Residues classified
        ``$\alpha$'' are summed into the $\alpha$-amplitude column;
        ``$\beta$'' into the $\beta$-amplitude column; ``loop'' into
        the loop column. The sum-of-three equals the total
        protein amplitude (except for unclassified residues).
  \item \fn{geometry.build\_sse\_center\_map}: enriches the
        per-atom residue labels with their SSE class for the audit
        log.
\end{enumerate}

For both, an unclassified residue (UNK) is silently dropped from
the SSE-sum but kept in the total -- the resulting ``$\alpha + \beta
+ \text{loop} < \text{total}$'' inequality in the workbook signals
to the user that some residues lacked SSE classification.

% -------------------------------------------------------------------------
\section{Algorithm: main mode-analysis loop (Phase 4)}
\label{sec:algo-mode-loop}

\codref{runner.\_run\_analysis\_single}

The Phase 4 main loop is the most computationally expensive part of
the pipeline. For each mode that passes the frequency-window filter,
the loop:

\begin{enumerate}
  \item streams the mode's eigenvector from the log
        (via \fn{logio.get\_eigvec}),
  \item thermally scales it via $\urms(\omega_l)$,
  \item runs the full per-mode analysis chain (OOP/INP, SCSD,
        kernel-class, all bond-modulations, $B$-factor
        accumulation),
  \item collects the result-dict and appends it to the master list.
\end{enumerate}

\begin{lstlisting}[style=pseudo,
  caption={Main mode-analysis loop. Implementation:
           \protect\fn{runner.\_run\_analysis\_single},
           inner call: \protect\fn{core.analyze\_mode\_with\_fallback}.}]
input:  cfg, atoms, idx_map, all_blocks, best_block, cand_map,
        cluster_normal, fe_c, s_c, coord_info
output: results : list of per-mode dicts
        b_accum : (N_atoms,) array of B-factor contributions

# Step 1: build the selection list
selected <- []
foreach bi in all_blocks:
    foreach col, (mn, freq) in enumerate(zip(bi.mode_nums, bi.freqs)):
        # only count the preferred block for each mode number
        if best_block[mn] is not bi:
            continue
        # skip imaginary/zero modes
        if freq <= 0:
            log "skipping mode {mn} with freq={freq} cm^-1"
            continue
        # frequency-window filter
        if cfg.freq_min and freq < cfg.freq_min: continue
        if cfg.freq_max and freq > cfg.freq_max: continue
        if cfg.freq_windows is not None
                and not any(lo <= freq <= hi for (lo, hi) in cfg.freq_windows):
            continue
        selected.append((bi, col, mn, freq))

selected.sort(by=freq)
results  <- []
b_accum  <- zeros(N_atoms)

# Step 2: per-mode analysis
foreach (bi, col, mn, freq) in selected:
    # eigenvector loading with fallback to alternative blocks
    r, used_block <- analyze_mode_with_fallback(
        candidates  = cand_map[mn],
        col         = col,
        log_file    = cfg.log_file,
        atoms       = atoms,
        idx_map     = idx_map,
        normal      = cluster_normal,
        coord_info  = coord_info,
        fe_c        = fe_c,
        s_c         = s_c,
        cfg         = cfg,
        mode_num    = mn,
    )
    if r is None:
        # all blocks failed; log and skip
        log "failed to read eigenvector for mode {mn}, all blocks bad"
        continue

    # Accumulate B-factor: r["evg"] is the thermally scaled eigenvector
    foreach i in range(N_atoms):
        b_accum[i] += sum(r["evg"][i]**2)   # (e_x*u + e_y*u + e_z*u)^2

    results.append(r)

return results, b_accum
\end{lstlisting}

\paragraph{Inner helper: \fn{analyze\_mode\_with\_fallback}.}

\codref{core.analyze\_mode\_with\_fallback}

This helper wraps \fn{logio.get\_eigvec} with a fallback loop:

\begin{lstlisting}[style=pseudo,
  caption={Eigenvector loading with fallback. Implementation:
           \protect\fn{core.analyze\_mode\_with\_fallback}.}]
input:  candidates : list of BlockInfo, col : int, ...
output: result_dict, used_block

foreach bi in candidates_sorted_by_preference (HP first):
    try:
        evec_raw <- logio.get_eigvec(bi, col, ...)
            # raw cartesian-displacement eigenvector, shape (N_atoms, 3)
            # not yet thermally scaled
    except (CorruptDataError, ParseError):
        continue

    # Thermal scaling: each component of evec_raw becomes A x u_rms(omega_l)
    u_rms <- core.compute_thermal_amplitude(freq, red_mass, cfg.temp_k, cfg.amplitude)
    evg   <- evec_raw * u_rms        # thermally scaled, in Angstrom

    # Run all per-mode analyses
    result_dict <- {
        "number":    mode_num,
        "freq":      freq,
        "red_mass":  red_mass,
        "evec_raw":  evec_raw,
        "evg":       evg,
        ...,
    }
    result_dict.update(core.classify_oop_inp(evg, normal, fe_c, s_c, cfg))
    result_dict.update(core.analyze_fe_ligand(evg, normal, coord_info, cfg))
    result_dict.update(core.analyze_his_hn(evg, coord_info, cfg))
    result_dict.update(core.compute_scsd_for_mode_full(evec_raw, u_rms, atoms, fe_c, s_c, cfg))
    result_dict.update(reorganization.compute_mode_modulations(evec_raw, u_rms, freq, coord_info, cfg))

    return result_dict, bi

# All candidates failed
return None, None
\end{lstlisting}

\subsection{The candidate ordering}

\fn{cand\_map[mn]} contains all blocks whose mode-number list
includes \texttt{mn}. The fallback iterates them in this order:

\begin{enumerate}
  \item \textbf{Preferred block} (\fn{best\_block[mn]}): typically
        the HP block if any, otherwise the first standard block
        seen during scan.
  \item \textbf{HP blocks not chosen as preferred}: rare, but
        happens if two HP blocks have overlapping mode numbers
        (e.g.\ a corrupted log with duplicated sections).
  \item \textbf{Remaining standard blocks}: lowest precision,
        used only if the HP blocks all failed.
\end{enumerate}

This ordering is implemented by sorting \fn{cand\_map[mn]} with
\texttt{key=(not bi.is\_hp, bi.offset)} (\texttt{not bi.is\_hp}
is False for HP blocks, True for standard, putting HP first).

\subsection{What can cause a block to fail}

\fn{logio.get\_eigvec(bi, col, ...)} can raise three exception
types:

\begin{itemize}
  \item \texttt{CorruptDataError}: the block's eigenvector body
        ends before $3N$ rows are read (file truncated, or the
        block format violates expectations).
  \item \texttt{ParseError}: a row contains non-numeric data where
        floats are expected (e.g.\ ``\texttt{********}'' overflow
        markers in the unlikely case they appear in eigenvector
        columns).
  \item \texttt{IndexError}: the requested \texttt{col} is beyond
        the actual number of mode columns in the block (Gaussian
        sometimes writes fewer columns in the last block of a
        section if the total mode count is not a multiple of 5).
\end{itemize}

The fallback catches all three and tries the next candidate. If
all candidates fail, the mode is logged as ``\texttt{failed to read
eigenvector for mode \{mn\}, all blocks bad}'' and skipped. The
final results list excludes failed modes (FUND-13 synthetic-zero
modes are added separately).

\subsection{The B-factor accumulator}

The \texttt{b\_accum} array is filled in-place during the loop.
For each mode $l$ and each atom $i$, the per-mode contribution is
\[
  \Delta b_i^{(l)}
   = \sum_{c \in \{x,y,z\}} \left( \evec_{l,i,c} \cdot \urms(\omega_l) \right)^2
   = \urms^2 \cdot \lVert \evec_{l,i} \rVert^2.
\]
This is the per-atom contribution before the $8\pi^2/3$ scaling
(applied at the export stage,
\cref{sec:b-factor-derivation}).

Because the loop runs over modes that survived the
frequency-window filter, the resulting $B$-factors are
window-restricted (see \cref{sec:b-factor-derivation} for the
implications).

\subsection{Computational complexity}

\paragraph{Time per mode.} O($N_\text{atoms}$): one pass over the
eigenvector for OOP/INP, one for SCSD, one each for the
per-channel bond modulations. Total per-mode time scales linearly
with the QM region size.

\paragraph{Time per log.} O($N_\text{modes} \cdot N_\text{atoms}$).
For a typical [2Fe-2S] QM/MM system (here ${\sim}1260$ modes,
${\sim}422$ QM atoms): ${\sim}$~532k atom-operations per pass,
${\sim}$~3 minutes wall-clock on a 2024-era laptop CPU.

\paragraph{Memory.} O($N_\text{modes} \cdot N_\text{atoms}$): the
results list holds one dict per mode containing the eigenvector,
plus per-mode observables. For the same example system: ${\sim}$~50 MiB.

\subsection{Edge cases}

\paragraph{(L1) Imaginary frequencies.} Gaussian writes imaginary
frequencies as negative numbers (e.g.\ $-23.45$). The filter
\texttt{if freq <= 0: continue} drops these. The user is alerted
once with ``\texttt{N imaginary modes skipped}''.

\paragraph{(L2) Translation/rotation modes.} The 6 lowest modes of a
free molecule are translations and rotations with $\omega \approx 0$.
For QM/MM-optimised structures with the QM region embedded in the
MM region, these modes are partially constrained and have small
but non-zero frequencies. They are not filtered out by the
$\omega \leq 0$ check; if they fall in the user's frequency window,
they will be analysed (typically yielding mostly skeletal
displacement with little spectroscopic significance).

\paragraph{(L3) Multi-window scenarios.} If \cfgkey{freq\_windows}
is a list of (lo, hi) tuples instead of a single window, a mode
passes if it falls in \emph{any} window. The export stage produces
one workbook per window.

\paragraph{(L4) Mode-number collisions across HP and standard
blocks.} If both an HP and a standard block claim mode 100,
\fn{best\_block} stores the HP version; the loop processes the
HP version. The standard version is available via \fn{cand\_map}
as a fallback.

% -------------------------------------------------------------------------
\section{Algorithm: context-mode loading + FUND 12 + FUND 13}
\label{sec:algo-context}

\codref{runner.\_run\_analysis\_single, runner.\_make\_synthetic\_zero\_mode}

After the main loop has produced \fn{results}, the interpolation
workbook (\file{*\_analysis\_interp\#\#.xlsx}) requires that the
piecewise-linear interpolation of each per-mode observable has
well-defined anchors at the lower and upper edges of the frequency
window. Without anchors, the interpolation produces zero outside the
range of the modes inside the window -- which is generally not the
true spectral behaviour just above/below the window.

The remedy is to load a small number of \emph{context modes} just
outside the window. Three loading scenarios are distinguished:

\paragraph{Scenario A: at least one real mode within \cfgkey{interp\_context\_cm1} of the boundary.}
Load every such mode normally (\fn{analyze\_mode\_with\_fallback}).
This is the standard case for densely sampled DFT spectra.

\paragraph{Scenario B (FUND 12): no real mode within
\cfgkey{interp\_context\_cm1}, but real modes exist beyond.}
The nearest real mode beyond \cfgkey{interp\_context\_cm1} is loaded
as a single ``boundary anchor''. This guards against the failure mode
where the DFT spectrum happens to be sparse exactly at the window
boundary, which would leave the interpolation unanchored.

\paragraph{Scenario C (FUND 13): no real mode above the upper window
boundary at all.}
A synthetic null-mode is inserted at
$\omega_\text{anchor} = \cfgkey{freq\_max} + \cfgkey{interp\_context\_cm1}$
with all observables set to zero. This is the case when the
analysis window includes the highest frequency in the spectrum, and
the natural right-hand boundary of the spectrum is being asked to
serve as the interpolation's right boundary.

\begin{lstlisting}[style=pseudo,
  caption={Context-mode loading with FUND 12 + FUND 13 fallbacks.
           Implementation: \protect\fn{runner.\_run\_analysis\_single},
           helper \protect\fn{runner.\_make\_synthetic\_zero\_mode}.}]
input:  cfg, all_blocks, best_block, cand_map, ..., real_results
output: context_results_left, context_results_right

# RIGHT-EDGE (upper-frequency) CONTEXT
right_window_lo <- cfg.freq_max
right_window_hi <- cfg.freq_max + cfg.interp_context_cm1
context_right_modes <- [m for m in all_modes if right_window_lo < m.freq <= right_window_hi]

if len(context_right_modes) > 0:
    # Scenario A
    foreach m in context_right_modes:
        load_and_analyze(m); append to context_results_right
else:
    # Look for a single fallback above interp_context_cm1
    fallback <- min((m for m in all_modes if m.freq > right_window_hi), key=lambda m: m.freq, default=None)
    if fallback is not None:
        # Scenario B -- FUND 12
        load_and_analyze(fallback); append to context_results_right
        log "FUND-12 fallback: loaded mode {fallback.number} at {fallback.freq:.1f} cm^-1 as right boundary anchor (no mode within +{cfg.interp_context_cm1} cm^-1)"
    else:
        # Scenario C -- FUND 13
        synth <- _make_synthetic_zero_mode(omega_anchor=right_window_hi, atoms=atoms, ...)
        context_results_right.append(synth)
        log "FUND-13 fallback: inserted synthetic zero anchor at {right_window_hi:.1f} cm^-1 (no real mode above {cfg.freq_max})"

# LEFT-EDGE (lower-frequency) CONTEXT  -- only Scenarios A and B
left_window_lo <- cfg.freq_min - cfg.interp_context_cm1
left_window_hi <- cfg.freq_min
context_left_modes <- [m for m in all_modes if left_window_lo <= m.freq < left_window_hi]

if len(context_left_modes) > 0:
    # Scenario A
    foreach m in context_left_modes:
        load_and_analyze(m); append to context_results_left
else:
    fallback <- max((m for m in all_modes if m.freq < left_window_lo), key=lambda m: m.freq, default=None)
    if fallback is not None:
        # Scenario B -- FUND 12
        load_and_analyze(fallback); append to context_results_left
        log "FUND-12 fallback: loaded mode {fallback.number} at {fallback.freq:.1f} cm^-1 as left boundary anchor"
    # else: no left synthetic anchor -- interpolation's natural
    #       left=0 boundary is already physically correct

return context_results_left, context_results_right
\end{lstlisting}

\paragraph{Why no synthetic left anchor.}
The interpolation at the lower edge naturally extrapolates to
$\Delta r = 0$ (zero observables at zero frequency), because every
spectroscopic observable in the code vanishes at $\omega \to 0$
(the modulation is the displacement, which vanishes for a
translation/rotation pseudomode). There is therefore no
physically wrong behaviour at the lower edge that requires a synthetic
anchor; only the upper edge benefits from one.

\paragraph{Synthetic zero mode shape.}
The synthetic mode returned by \fn{\_make\_synthetic\_zero\_mode} has
the same dictionary shape as a real mode's result, with all
observable values set to zero and sentinel metadata
$\texttt{number} = -1$,
$\texttt{mode\_type} = \texttt{"synthetic\_zero"}$,
$\texttt{precision} = \texttt{"synthetic"}$. The eigenvector
matrix \fn{evg} and the SCSD coefficient list are deliberately
\emph{omitted} (set to \texttt{None}) so that the downstream
$B$-factor and SSE-amplitude loops detect the sentinel and skip the
mode -- it must not contribute to system-wide aggregates, only to
the interpolation anchor.

\subsection{The decision order in detail}

When the main loop finishes and the context-loading phase begins,
the implementation processes the right-edge first. The complete
decision tree:

\begin{enumerate}
  \item Build \texttt{candidates\_right = [m for m in all\_modes if
        cfg.freq\_max < m.freq <= cfg.freq\_max +
        cfg.interp\_context\_cm1]}. If non-empty, this is
        Scenario A; analyse each candidate and store as
        \texttt{context\_results\_right}. \emph{Done.}
  \item Otherwise build \texttt{fallback\_right = min(all\_modes
        with freq > cfg.freq\_max + cfg.interp\_context\_cm1, key
        = freq)}. If not None, this is Scenario B (FUND 12);
        load and analyse the single fallback mode; log the
        FUND-12 fallback message. \emph{Done.}
  \item Otherwise this is Scenario C (FUND 13); call
        \fn{\_make\_synthetic\_zero\_mode} at
        $\omega_\text{anchor} = \cfgkey{freq\_max} +
        \cfgkey{interp\_context\_cm1}$;
        log the FUND-13 fallback message. \emph{Done.}
\end{enumerate}

The left-edge processing is identical except for the missing
Scenario C (no left synthetic anchor).

The full FUND 12 + FUND 13 logic is exercised by three E2E
regression tests
(\texttt{test\_e2e\_terminal\_spectrum\_fund\_13},
\texttt{test\_synthetic\_zero\_mode\_has\_zero\_observables},
\texttt{test\_synthetic\_zero\_mode\_marked\_as\_synthetic};
\cref{sec:regression-tests}).

\subsection{Edge cases}

\paragraph{(X1) Window includes the full DFT spectrum.}
If \cfgkey{freq\_min} = 0 and \cfgkey{freq\_max} = $\omega_\text{max}^\text{DFT}$,
both edges are at the natural spectrum boundaries. Right-edge:
Scenario C (synthetic). Left-edge: empty (no synthetic; the
natural left boundary is correct).

\paragraph{(X2) Window of zero width.}
If \cfgkey{freq\_min} = \cfgkey{freq\_max} (a degenerate window),
the main loop produces no real modes. The context loader still
runs and may produce 1--2 anchors. The resulting interpolation is
constant zero. This is documented but not strictly disallowed.

\paragraph{(X3) Negative \cfgkey{interp\_context\_cm1}.}
The context-window construction
\texttt{[freq\_max, freq\_max + interp\_context\_cm1]} becomes
inverted (\texttt{freq\_max} > endpoint). The list of context
candidates is then empty. The fallback Scenario B/C is triggered.
A user who set \cfgkey{interp\_context\_cm1} = 0 will reliably
land in Scenario B (FUND 12) for any spectrum with modes above
\cfgkey{freq\_max}, and in Scenario C for a window that includes
the spectrum's maximum.

\paragraph{(X4) Multiple clusters, multiple windows.}
The context-loading is per-cluster, per-window. Each per-cluster,
per-window result has its own synthetic anchor if needed; the
sentinels do not bleed across clusters or windows.

% -------------------------------------------------------------------------
\section{Algorithm: multi-window loop}
\label{sec:algo-multi-window}

\codref{runner.run\_analysis}

When the user supplies \cfgkey{freq\_windows} (a list of disjoint
[lo, hi] intervals), Phase 4 is invoked once per window, producing
separate output workbooks (and separate \texttt{Reorganisation}
sheets). Each window's $\Lambda_X^\text{total}$ is independent.

\begin{lstlisting}[style=pseudo,
  caption={Multi-window dispatching. Implementation:
           \protect\fn{runner.run\_analysis}.}]
input:  cfg
output: dict[window -> workbook paths]

if cfg.freq_windows is None:
    # Single implicit window from (freq_min, freq_max)
    return _run_analysis_single(cfg)

results_per_window <- {}
foreach (lo, hi) in cfg.freq_windows:
    cfg_sub <- replace(cfg, freq_min=lo, freq_max=hi, freq_windows=None)
    cfg_sub.logname_suffix <- f"_{lo:.0f}_{hi:.0f}_cm"
    results_per_window[(lo, hi)] <- _run_analysis_single(cfg_sub)
return results_per_window
\end{lstlisting}

\paragraph{Window-boundary convention.}
The implementation uses \emph{half-open} intervals
$[\texttt{lo}, \texttt{hi})$: a mode at exactly $\omega = \texttt{hi}$
is \emph{not} included in the window, but \emph{is} included in the
window that starts at $\texttt{hi}$. This avoids double-counting
when windows are contiguous (e.g.\ $[0, 100), [100, 300)$). The
convention is documented in the manual and verified by
\fn{tests/test\_v104\_audit\_fixes.py::test\_freq\_window\_half\_open}.

\paragraph{Cache reuse across windows.}
Because the byte-offset scan and the \fn{read\_all\_meta} structures
are window-independent, the cache (\cref{sec:cache}) is shared
between window invocations: only the first window pays the
disk-scan cost.

\subsection{Edge cases for the multi-window loop}

\paragraph{(W1) Overlapping windows.}
The half-open convention only prevents double-counting at adjacent
window boundaries (\texttt{hi$_1$ = lo$_2$}). If the user supplies
overlapping windows (\texttt{lo$_2$ < hi$_1$}), modes in the overlap
appear in \emph{both} output workbooks. This is intentional --
sometimes a user wants two overlapping windows to compare
($\Lambda^\text{total}$ in $[0, 500]$ vs.\ $[0, 300]$ for example).
A warning ``\texttt{windows \{w1\} and \{w2\} overlap}'' is logged
once.

\paragraph{(W2) Window outside the spectrum.}
A window entirely above $\omega_\text{max}^\text{DFT}$ produces zero
real modes. The synthetic-zero anchor handles the right boundary;
the left boundary triggers Scenario B/C from the previous section.
The output workbook is well-defined but contains only the
anchor's zero observables.

\paragraph{(W3) Negative or zero-width windows.}
A window with \texttt{lo $\geq$ hi} is rejected at config-validation
time; the run aborts.

\paragraph{(W4) Five or more windows.}
The loop scales linearly with the number of windows. The byte-
offset scan is reused; only the per-mode analysis (Phase 4) is
repeated. For 5 windows on a 1260-mode spectrum, the total
runtime is ${\sim}$~15 minutes instead of ${\sim}$~3.

% -------------------------------------------------------------------------
\section{Algorithm: multi-cluster loop}
\label{sec:algo-multi-cluster}

\codref{runner.run\_analysis}

When \cfgkey{analyze\_all\_clusters} = true, the pipeline iterates
over every cluster found by \fn{find\_all\_clusters}. Each cluster
produces an independent output directory (suffix
\texttt{\_cluster\#\#}) with its own set of workbooks. Inter-cluster
aggregation is not currently performed.

\begin{lstlisting}[style=pseudo,
  caption={Multi-cluster dispatching. Implementation:
           \protect\fn{runner.run\_analysis}.}]
input:  cfg
output: dict[cluster_index -> window-result-dict]

clusters <- find_all_clusters(atoms, cfg.fe_s_cutoff, cfg.fe_fe_cutoff,
                              strict=cfg.strict_cluster)
if not cfg.analyze_all_clusters:
    cluster_subset <- [clusters[cfg.cluster_index]]
else:
    cluster_subset <- clusters

results_per_cluster <- {}
foreach (idx, cluster) in enumerate(cluster_subset):
    fe1_c, fe2_c, s1_c, s2_c <- cluster
    cfg_sub <- replace(cfg,
        output_dir = cfg.output_dir + f"_cluster{idx:02d}",
    )
    # Reset per-cluster warning state so each cluster gets a clean
    # set of WARNED_EMPTY_GROUP / WARNED_EMPTY_LIG / WARNED_HN_SKIP
    core.reset_warning_state()
    results_per_cluster[idx] <- _run_analysis_single(
        cfg_sub, fe_c=(fe1_c, fe2_c), s_c=(s1_c, s2_c), ...)
return results_per_cluster
\end{lstlisting}

\paragraph{Per-cluster state reset.}
The deduplication sets
$\fn{\_WARNED\_EMPTY\_GROUP}, \fn{\_WARNED\_EMPTY\_LIG}, \fn{\_WARNED\_HN\_SKIP}$
in \fn{core.py} are module-level and used to suppress duplicate
\fn{UserWarning} emissions during the inner per-mode analyses. They
must be reset at the start of each cluster run, otherwise a warning
that fired during cluster 1's analysis would silence the
corresponding warning in cluster 2 -- which would be misleading.
\fn{core.reset\_warning\_state} performs this reset; it is invoked
unconditionally at the start of every cluster, regardless of whether
the system actually has multiple clusters.

\subsection{Per-cluster output naming}

Each cluster's output is placed in a sibling directory of the
main output directory, with suffix \texttt{\_cluster\#\#} where \#\#
is the zero-padded cluster index (\texttt{\_cluster00},
\texttt{\_cluster01}, \dots). The cluster index follows the
canonical ordering of \fn{find\_all\_clusters} (sorted by smallest
Gaussian centre).

\subsection{Edge cases}

\paragraph{(M1) Single cluster with \cfgkey{analyze\_all\_clusters}.}
Setting \cfgkey{analyze\_all\_clusters} = True when only one
cluster exists is equivalent to setting it False with
\cfgkey{cluster\_index} = 0, except that the output directory
suffix is added. This is documented; the resulting outputs are
identical apart from the directory name.

\paragraph{(M2) No clusters found.}
\fn{find\_all\_clusters} returns an empty list. The runner raises
\texttt{NoClustersFoundError} with a clear message and exits
without producing any workbook.

\paragraph{(M3) Cluster index out of range.}
\cfgkey{cluster\_index} larger than \texttt{len(clusters) - 1}
raises \texttt{IndexError} at the cluster-subset construction
step. The error message includes the actual number of clusters
detected.

\paragraph{(M4) Cluster-shared atoms.}
A user may construct a system where one S atom is shared between
two ``adjacent'' [2Fe-2S] clusters (e.g.\ a [3Fe-4S]-like
arrangement projected through cluster discovery). The shared
atom appears in both per-cluster outputs; the per-cluster
$\Lambda^\text{total}$ values double-count it. The strict-cluster
mode of \fn{find\_all\_clusters} (\cref{sec:algo-cluster-discovery})
filters this case at the discovery stage.

% =========================================================================
% PART IV -- DATA STRUCTURES
% =========================================================================
\part{Data Structures}
\label{part:dataclasses}

This part enumerates every non-trivial data structure passed between
the phases of \cref{sec:dataflow}. Every field of every dataclass is
documented: its type, what it represents, and which downstream code
consumes it. A reviewer who has read this part should be able to
write a test that constructs a synthetic instance of any data
structure and feeds it into the consuming function.

% -------------------------------------------------------------------------
\section{Dataclass: \fn{BlockInfo}}
\label{sec:dataclass-blockinfo}

\codref{logio.scan\_log}

Defined in \file{logio.py}. Carries the per-frequency-block metadata
read by \fn{read\_all\_meta} from each \texttt{Frequencies --}
section in the Gaussian log.

\begin{longtable}{p{0.22\textwidth} p{0.30\textwidth} p{0.43\textwidth}}
\toprule
\textbf{Field} & \textbf{Type} & \textbf{Meaning} \\
\midrule
\texttt{offset} & \texttt{int}
  & Byte offset of the \texttt{Frequencies --} header. Used to
    \texttt{seek()} the file before reading the block. \\
\texttt{mode\_nums} & \texttt{List[int]}
  & Mode numbers contained in this block (3 modes per row for
    standard blocks, 5 modes per row for \texttt{hpmodes} blocks). \\
\texttt{freqs} & \texttt{List[float]}
  & Frequencies of the modes in cm$^{-1}$. Same length as
    \texttt{mode\_nums}. \\
\texttt{red\_masses} & \texttt{List[float]}
  & Reduced masses in amu. NaN for modes where Gaussian printed
    \texttt{********} (overflow). \\
\texttt{frc\_consts} & \texttt{List[float]}
  & Force constants in mDyne/\AA{}. Currently unused. \\
\texttt{syms} & \texttt{List[str]}
  & Irreducible-representation labels (\texttt{A}, \texttt{A'},
    $A_g$, etc.). Used by the kernel-classification scoring. \\
\texttt{is\_hp} & \texttt{bool}
  & True for \texttt{hpmodes} blocks (5-decimal precision), False
    for standard 3-decimal blocks. Used by \fn{best\_block} to
    prefer HP over standard for any mode that appears in both. \\
\texttt{data\_offset} & \texttt{int}
  & Byte offset of the first eigenvector data row in this block.
    Used by \fn{get\_eigvec} to stream the eigenvector matrix
    without re-parsing the header. \\
\bottomrule
\end{longtable}

% -------------------------------------------------------------------------
\section{Dataclass: \fn{LigandInfo}}
\label{sec:dataclass-ligandinfo}

\codref{geometry.find\_coordinating\_residues}

Defined in \file{geometry.py}. Carries per-ligand coordination
information for each Fe-ligand bond.

\begin{longtable}{p{0.22\textwidth} p{0.30\textwidth} p{0.43\textwidth}}
\toprule
\textbf{Field} & \textbf{Type} & \textbf{Meaning} \\
\midrule
\texttt{fe\_idx} & \texttt{int}
  & Index (0 or 1) of which Fe in the cluster this ligand binds. \\
\texttt{fe\_center} & \texttt{int}
  & Gaussian centre of that Fe atom. \\
\texttt{lig\_center} & \texttt{int}
  & Gaussian centre of the ligand donor atom. \\
\texttt{lig\_element} & \texttt{str}
  & Element symbol of the donor atom (``N'', ``S'', or ``O''). \\
\texttt{lig\_aname} & \texttt{str}
  & PDB atom name of the donor (\texttt{SG}, \texttt{ND1},
    \texttt{NE2}, etc.). \\
\texttt{res\_num} & \texttt{int}
  & PDB residue number of the donor. \\
\texttt{res\_name} & \texttt{str}
  & PDB residue name (\texttt{CYS}, \texttt{HIS}, \texttt{ASP},
    \texttt{GLU}, etc.). \\
\texttt{res\_label} & \texttt{str}
  & Human-readable label (e.g.\ \texttt{Cys11}, \texttt{His87}).
    Used as the column-key in the output spreadsheets. \\
\texttt{bond\_vec} & \texttt{np.ndarray}
  & 3-vector $\rvec_L - \rvec_{\text{Fe}}$. Used for the
    Fe-ligand stretch/bend axis. \\
\texttt{bond\_len} & \texttt{float}
  & $\lVert$\texttt{bond\_vec}$\rVert$ in \AA{}. \\
\texttt{h\_center} & \texttt{Optional[int]}
  & For histidines: Gaussian centre of the protonated-N-H proton, if
    any. \texttt{None} for non-histidines or unprotonated histidines. \\
\texttt{hn\_vec} & \texttt{Optional[np.ndarray]}
  & N-H bond vector $\rvec_H - \rvec_N$. \\
\texttt{hn\_len} & \texttt{Optional[float]}
  & N-H bond length in \AA{}. \\
\texttt{his\_protonated} & \texttt{bool}
  & True if this is a protonated histidine ligand. Drives the
    inclusion of this ligand in the His-NH and HA channels. \\
\texttt{his\_hn\_center} & \texttt{Optional[int]}
  & Synonym for \texttt{h\_center}; kept for backward compatibility
    with pre-v1.0.2 result dicts. \\
\texttt{coord\_n\_aname} & \texttt{Optional[str]}
  & For histidines: which side-chain N coordinates Fe
    (\texttt{ND1} or \texttt{NE2}). \\
\bottomrule
\end{longtable}

% -------------------------------------------------------------------------
\section{Dataclass: \fn{CoordInfo}}
\label{sec:dataclass-coordinfo}

\codref{geometry.find\_coordinating\_residues}

Defined in \file{geometry.py}. The master container for all
coordination information that the per-mode analysis needs. Built
once per cluster.

\begin{longtable}{p{0.22\textwidth} p{0.30\textwidth} p{0.43\textwidth}}
\toprule
\textbf{Field} & \textbf{Type} & \textbf{Meaning} \\
\midrule
\texttt{ligands} & \texttt{List[LigandInfo]}
  & Every Fe-coordinating ligand, one entry per Fe-L bond. Length
    is typically 4 (Cys$_4$ cluster), 4 (Cys$_2$His$_2$ Rieske, two
    Cys$_2$ + two His$_2$), or 5 (mixed-coordination systems). \\
\texttt{group\_map} & \texttt{Dict[str, List[int]]}
  & Pre-defined or auto-generated atom groups, key = group name,
    value = list of Gaussian centres in the group. Used by the
    per-group OOP/INP analysis. Default groups: \texttt{Cluster},
    \texttt{ProteinBackbone}, one \texttt{Res\_\#\#\#} per residue. \\
\texttt{pdb\_to\_center} & \texttt{Dict[int, int]}
  & Mapping from PDB atom index to Gaussian centre, populated by
    the Kabsch matching of \cref{sec:kabsch}. \\
\texttt{his\_ligand\_labels} & \texttt{List[str]}
  & Convenience list of \texttt{res\_label} strings for all
    protonated histidines. Drives the order of columns in the
    His-NH and HA output sheets. \\
\texttt{pcet\_info} & \texttt{Optional[object]}
  & Carries the PCET acceptor candidates per protonated histidine
    (built by \fn{reorganization.build\_channels\_from\_coord\_info}).
    Per-mode HA analysis reads this. \\
\texttt{et\_info} & \texttt{Optional[object]}
  & Carries the FeFe channel geometry (one per pair of Fe atoms in
    the cluster). \\
\texttt{atoms\_h} & \texttt{Optional[List[Dict]]}
  & Hydrogen-inclusive atom list. \texttt{None} when
    \cfgkey{include\_hn\_vibration} is False. \\
\texttt{idx\_map\_h} & \texttt{Optional[Dict[int, int]]}
  & Gaussian-centre $\rightarrow$ index in \texttt{atoms\_h}. \\
\bottomrule
\end{longtable}

% -------------------------------------------------------------------------
\section{Dataclass: \fn{ChannelGeometry}}
\label{sec:dataclass-channelgeometry}

\codref{reorganization.build\_channels\_from\_coord\_info}

Defined in \file{reorganization.py}. Carries the geometric input
to the per-mode reorganization calculation for one bond sub-channel.

\begin{longtable}{p{0.22\textwidth} p{0.30\textwidth} p{0.43\textwidth}}
\toprule
\textbf{Field} & \textbf{Type} & \textbf{Meaning} \\
\midrule
\texttt{name} & \texttt{str}
  & Unique sub-channel name (e.g.\ \texttt{FeS\_Fe1\_S1},
    \texttt{FeN\_His87}, \texttt{HA\_His87\_Asp42-OD1}). Used as
    column-key. \\
\texttt{idx\_a} & \texttt{int}
  & Position of atom $a$ in the heavy-atom (or all-atom) array. \\
\texttt{idx\_b} & \texttt{int}
  & Position of atom $b$. \\
\texttt{r\_a} & \texttt{np.ndarray}
  & Equilibrium position of atom $a$, 3-vector in \AA{}. \\
\texttt{r\_b} & \texttt{np.ndarray}
  & Equilibrium position of atom $b$. \\
\texttt{elem\_a} & \texttt{str}
  & Element symbol of atom $a$ (used for $\mu_\text{pair}$). \\
\texttt{elem\_b} & \texttt{str}
  & Element symbol of atom $b$. \\
\texttt{mu\_pair\_amu} & \texttt{float}
  & Reduced mass $\mu = m_a m_b / (m_a + m_b)$ in amu. Used
    directly in \fn{lambda\_pair\_cm1}. \\
\texttt{parent\_channel} & \texttt{str}
  & Which of the five parent channels this sub-channel feeds:
    \texttt{FeFe}, \texttt{FeN}, \texttt{FeS}, \texttt{NH}, or
    \texttt{HA}. \\
\texttt{weight} & \texttt{float}
  & Aggregation weight $w_s$ for the RSS aggregation
    (\cref{sec:rss-aggregation}). For HA channels this is the
    gaussian-distance weight; for all others $w_s = 1$. \\
\texttt{idx\_donor} & \texttt{Optional[int]}
  & For HA channels: index of the donor N$_\text{His}$ atom. Used
    by \cref{eq:dr-ha-final} as the displacement subtraction
    reference. \\
\texttt{r\_donor} & \texttt{Optional[np.ndarray]}
  & Equilibrium position of the donor; defines the $\nhat_{NA}$
    projection axis. \\
\bottomrule
\end{longtable}

% -------------------------------------------------------------------------
\section{Dataclass: \fn{ChannelResult}}
\label{sec:dataclass-channelresult}

\codref{reorganization.compute\_mode\_modulations}

Defined in \file{reorganization.py}. Carries the per-mode, per-sub-channel
result of the reorganization calculation, one instance per (mode,
sub-channel) pair.

\begin{longtable}{p{0.22\textwidth} p{0.30\textwidth} p{0.43\textwidth}}
\toprule
\textbf{Field} & \textbf{Type} & \textbf{Meaning} \\
\midrule
\texttt{name} & \texttt{str}
  & Sub-channel name, copied from the corresponding
    \texttt{ChannelGeometry}. \\
\texttt{parent\_channel} & \texttt{str}
  & One of FeFe, FeN, FeS, NH, HA. \\
\texttt{dr\_signed\_a} & \texttt{float}
  & Signed bond modulation $\Delta r_X^{(s)}$ in \AA{}, computed by
    \fn{signed\_dr\_along\_axis} (classical formula
    \cref{eq:dr-classical-final}) or \fn{compute\_mode\_modulations}
    (reaction-coordinate formula
    \cref{eq:dr-ha-final} for HA). \\
\texttt{lambda\_pair\_cm1} & \texttt{float}
  & Per-sub-channel reorganization energy in cm$^{-1}$, computed by
    \fn{lambda\_pair\_cm1} \cref{eq:lambda-cm1}. \\
\texttt{lambda\_mode\_cm1} & \texttt{float}
  & Per-mode aggregate $\lambda_X(\text{mode})$ via
    \cref{eq:lambda-per-mode}. Only set on the parent-aggregated
    instance, NaN on sub-channel instances. \\
\texttt{weight} & \texttt{float}
  & Aggregation weight $w_s$ for this sub-channel. \\
\bottomrule
\end{longtable}

% -------------------------------------------------------------------------
\section{Per-mode result dictionary}
\label{sec:dataclass-modedict}

\codref{core.analyze\_mode\_with\_fallback}

The per-mode result dictionary is the central data structure of the
pipeline -- one dict per mode in the frequency window, returned by
\fn{analyze\_mode\_with\_fallback} and accumulated by
\fn{\_run\_analysis\_single}. Every Excel column in every output
workbook traces back to one of these keys. This section enumerates
every key.

\subsection{Identification and frequency metadata}

\begin{longtable}{p{0.22\textwidth} p{0.20\textwidth} p{0.53\textwidth}}
\toprule
\textbf{Key} & \textbf{Type} & \textbf{Meaning} \\
\midrule
\texttt{number} & \texttt{int}
  & Gaussian mode number (1-based). Sentinel $-1$ for synthetic
    zero modes (FUND 13). \\
\texttt{freq} & \texttt{float}
  & Frequency in cm$^{-1}$. \\
\texttt{red\_mass} & \texttt{float}
  & Reduced mass in amu. \\
\texttt{frc\_const} & \texttt{float}
  & Force constant in mDyne/\AA{}. \\
\texttt{symmetry} & \texttt{str}
  & Irrep label as printed by Gaussian. \\
\texttt{precision} & \texttt{str}
  & One of \texttt{"hp"}, \texttt{"std"}, \texttt{"synthetic"}. \\
\texttt{u\_rms\_a} & \texttt{float}
  & Thermal RMS amplitude $\urms(\omega_l)$ in \AA{}, computed by
    \cref{eq:urms-final}. \\
\bottomrule
\end{longtable}

\subsection{Eigenvector and per-atom contributions}

\begin{longtable}{p{0.22\textwidth} p{0.20\textwidth} p{0.53\textwidth}}
\toprule
\textbf{Key} & \textbf{Type} & \textbf{Meaning} \\
\midrule
\texttt{evec\_raw} & \texttt{np.ndarray}
  & Raw (un-thermally-scaled) eigenvector, shape $(N_\text{atoms},
    3)$. \\
\texttt{evg} & \texttt{np.ndarray}
  & Thermally scaled eigenvector $\widetilde{\evec} = \evec \cdot
    \urms$, shape $(N_\text{atoms}, 3)$, in \AA{}. Used by all
    downstream geometric observables. \\
\bottomrule
\end{longtable}

\subsection{Cluster-core OOP/INP descriptors}

Every key uses the OOP/INP decomposition of
\cref{sec:oop-derivation}.

\begin{longtable}{p{0.27\textwidth} p{0.20\textwidth} p{0.48\textwidth}}
\toprule
\textbf{Key} & \textbf{Type} & \textbf{Meaning} \\
\midrule
\texttt{a\_oop\_pct} & \texttt{float}
  & Cluster-core $\alpha_{\OOP}$ in \%. \\
\texttt{a\_inp\_pct} & \texttt{float}
  & $\alpha_{\INP} = 100 - \alpha_{\OOP}$. \\
\texttt{mode\_type} & \texttt{str}
  & Binary classification (OOP / INP / mixed), per
    \cref{eq:oop-binary}. \\
\texttt{mode\_type\_detail} & \texttt{str}
  & Eight-bin classification (pure\_oop, strong\_oop, etc.). \\
\texttt{ca\_displacement\_a} & \texttt{float}
  & Cluster centre-of-mass displacement
    $\lVert \tfrac{1}{4}\sum_{i \in \text{cluster}} \widetilde{\evec}_i\rVert$
    in \AA{}. Diagnoses translational contamination. \\
\texttt{ca\_oop\_a} & \texttt{float}
  & Out-of-plane component of the cluster-COM displacement. \\
\texttt{cluster\_expansion\_a} & \texttt{float}
  & Mean radial cluster expansion (\AA{}). Positive = breathing
    outward; negative = contracting. \\
\texttt{cluster\_rotation\_deg} & \texttt{float}
  & In-plane rotation angle of the rhombic Fe$_2$S$_2$ unit, in
    degrees. \\
\bottomrule
\end{longtable}

\subsection{Fe-ligand bond descriptors}

For each entry in \fn{coord\_info.ligands} the following keys are
created (with the ligand's \texttt{res\_label} as the column suffix
in the output sheet -- e.g.\ \texttt{stretch\_Cys11}).

\begin{longtable}{p{0.27\textwidth} p{0.20\textwidth} p{0.48\textwidth}}
\toprule
\textbf{Key} & \textbf{Type} & \textbf{Meaning} \\
\midrule
\texttt{stretch} & \texttt{float}
  & Fe-L stretch component $\Delta r_\text{stretch}$
    (\cref{eq:fel-stretch}) in \AA{}. \\
\texttt{bend\_inp} & \texttt{float}
  & In-plane bend component, $\geq 0$. \\
\texttt{bend\_oop} & \texttt{float}
  & Signed out-of-plane bend component. \\
\texttt{bend\_inp\_pct} & \texttt{float}
  & Percentage of total bend power in plane. \\
\texttt{bend\_oop\_pct} & \texttt{float}
  & Percentage of total bend power out of plane. \\
\texttt{s\_stretch} & \texttt{float}
  & Propagated sigma for \texttt{stretch}
    ($\sqrt{2}\sigma_e \urms$, \cref{eq:fel-sigma}). \\
\texttt{s\_bend\_inp} & \texttt{float}
  & Propagated sigma for \texttt{bend\_inp}. \\
\texttt{s\_bend\_oop} & \texttt{float}
  & Propagated sigma for \texttt{bend\_oop}. \\
\bottomrule
\end{longtable}

\subsection{His-NH and HA descriptors}

\begin{longtable}{p{0.27\textwidth} p{0.20\textwidth} p{0.48\textwidth}}
\toprule
\textbf{Key} & \textbf{Type} & \textbf{Meaning} \\
\midrule
\texttt{hn\_stretch\_\textit{Hisn}} & \texttt{float}
  & N-H stretch modulation for each protonated histidine, in \AA{}. \\
\texttt{s\_hn\_stretch\_\textit{Hisn}} & \texttt{float}
  & Propagated sigma for \texttt{hn\_stretch}; v1.0.4 fix
    (\cref{sec:fund-6-8}) applies the $\urms$ factor. \\
\texttt{ha\_dr\_\textit{Hisn}\_to\_\textit{acceptor}} & \texttt{float}
  & Reaction-coordinate HA modulation
    (\cref{eq:dr-ha-final}), in \AA{}. \\
\texttt{s\_ha\_dr\_\textit{Hisn}\_to\_\textit{acceptor}} & \texttt{float}
  & Propagated sigma for the HA modulation. \\
\bottomrule
\end{longtable}

\subsection{Marcus-Hush reorganization per mode}

\begin{longtable}{p{0.27\textwidth} p{0.20\textwidth} p{0.48\textwidth}}
\toprule
\textbf{Key} & \textbf{Type} & \textbf{Meaning} \\
\midrule
\texttt{dr\_rss\_FeFe} & \texttt{float}
  & RSS-aggregated $\Delta r$ along the FeFe channel. \\
\texttt{dr\_rss\_FeN}, \texttt{dr\_rss\_FeS}, \texttt{dr\_rss\_NH},
\texttt{dr\_rss\_HA} & \texttt{float}
  & RSS-aggregated $\Delta r$ along each parent channel
    (\cref{eq:dr-rss-derivation}). \\
\texttt{lambda\_FeFe}, \texttt{lambda\_FeN}, \texttt{lambda\_FeS},
\texttt{lambda\_NH}, \texttt{lambda\_HA} & \texttt{float}
  & Per-mode reorganization energy in cm$^{-1}$ per parent channel
    (\cref{eq:lambda-per-mode}). \\
\texttt{s\_dr\_rss\_*} & \texttt{float}
  & Propagated sigma for each RSS modulation. \\
\texttt{s\_lambda\_*} & \texttt{float}
  & Propagated sigma for each per-mode $\lambda$. \\
\bottomrule
\end{longtable}

\subsection{SCSD amplitudes}

One key per $D_{2\mathrm{h}}$ irrep, for both the reference and the
distortion projections (\cref{sec:scsd-derivation}).

\begin{longtable}{p{0.27\textwidth} p{0.20\textwidth} p{0.48\textwidth}}
\toprule
\textbf{Key} & \textbf{Type} & \textbf{Meaning} \\
\midrule
\texttt{scsd\_Ag},   \texttt{scsd\_B1g}, \texttt{scsd\_B2g}, \texttt{scsd\_B3g},
\texttt{scsd\_Au},   \texttt{scsd\_B1u}, \texttt{scsd\_B2u}, \texttt{scsd\_B3u}
   & \texttt{float}
  & Distortion amplitude in each irrep, in \AA{}. \\
\texttt{s\_scsd\_*} & \texttt{float}
  & Propagated sigma for each irrep amplitude (v1.0.4 fix). \\
\bottomrule
\end{longtable}

\subsection{Kernel-classification scores}

\begin{longtable}{p{0.27\textwidth} p{0.20\textwidth} p{0.48\textwidth}}
\toprule
\textbf{Key} & \textbf{Type} & \textbf{Meaning} \\
\midrule
\texttt{kern\_primary} & \texttt{str}
  & Name of the highest-scoring kernel template. \\
\texttt{kern\_primary\_score} & \texttt{float}
  & Score (\%) of the primary kernel. \\
\texttt{kern\_secondary} & \texttt{str}
  & Name of the second-highest kernel. \\
\texttt{kern\_secondary\_score} & \texttt{float}
  & Score (\%) of the secondary. \\
\texttt{kern\_loc} & \texttt{float}
  & Localisation fraction (\%): fraction of the mode's total
    amplitude residing on the cluster-core atoms. \\
\bottomrule
\end{longtable}

\subsection{Group amplitudes and secondary-structure descriptors}

For each group key in \fn{coord\_info.group\_map}, the following are
appended (suffixed by group name -- e.g.\ \texttt{a\_oop\_pct\_Cys11}).

\begin{longtable}{p{0.32\textwidth} p{0.18\textwidth} p{0.45\textwidth}}
\toprule
\textbf{Key} & \textbf{Type} & \textbf{Meaning} \\
\midrule
\texttt{a\_oop\_pct\_\textit{group}} & \texttt{float}
  & Per-group $\alpha_{\OOP}$. \\
\texttt{a\_inp\_pct\_\textit{group}} & \texttt{float}
  & Per-group $\alpha_{\INP}$. \\
\texttt{amplitude\_a\_\textit{group}} & \texttt{float}
  & Total amplitude (RMS over atoms in the group). \\
\texttt{sse\_amplitude\_*} & \texttt{float}
  & Per-SSE-element amplitudes (only when SSE analysis is enabled). \\
\bottomrule
\end{longtable}

% -------------------------------------------------------------------------
\section{The \fn{Config} dataclass}
\label{sec:dataclass-config}

\codref{config.Config}

The \fn{Config} dataclass is the single source of truth for every
configurable parameter. All 48 fields are enumerated in
\cref{part:config}. The dataclass is built either from a TOML file
(via \fn{Config.from\_toml}, the classmethod) or constructed in
Python directly (useful for tests). Field defaults are listed
alongside the field definitions in \file{config.py} and reproduced
in \cref{part:config}.

\paragraph{Validation.}
After construction, \fn{Config.validate} runs a small set of
internal-consistency checks:
\begin{itemize}
  \item \cfgkey{freq\_max} > \cfgkey{freq\_min} (if both set).
  \item All cutoff distances positive.
  \item \cfgkey{temp\_k} either \texttt{None} or positive.
  \item Mutually exclusive: \cfgkey{analyze\_all\_clusters} and
        \cfgkey{cluster\_index} > 0 are not both meaningful.
\end{itemize}
Validation failures raise \fn{ConfigError} before any analysis
begins, ensuring the user sees the problem before the long
disk-scan phase.

% =========================================================================
% PART V -- CONFIGURATION REFERENCE
% =========================================================================
\part{Configuration Reference}
\label{part:config}

This part enumerates every \cfgkey{} field of the \fn{Config}
dataclass (\cref{sec:dataclass-config}) with its type, default,
effect on the analysis, valid range or set of accepted values, and
which functions consume it. The fields are grouped by topic, in the
order they appear in \file{src/modenanalyse\_2fe2s/config.py}; the
table at the start of each section gives the line-range in
\file{config.py} for cross-reference.

% -------------------------------------------------------------------------
\section{Input/output paths and frequency window}

\begin{longtable}{p{0.21\textwidth} p{0.13\textwidth} p{0.13\textwidth} p{0.45\textwidth}}
\toprule
\textbf{Field} & \textbf{Type} & \textbf{Default} & \textbf{Effect} \\
\midrule
\cfgkey{log\_file} & \texttt{str} & \texttt{""}
  & Path to the Gaussian \file{.log} or ORCA \file{.hess} input.
    Required. The choice between Gaussian and ORCA is automatic
    via filename extension. \\
\cfgkey{pdb\_file} & \texttt{str} & \texttt{""}
  & Path to a PDB file. Optional. Required for residue labels and
    SSE analysis. \\
\cfgkey{output\_dir} & \texttt{str} & \texttt{""}
  & Directory for all output workbooks and the run-log. Created
    if missing. \\
\cfgkey{freq\_min} & \texttt{Optional[float]} & \texttt{None}
  & Lower frequency bound of the analysis window, in cm$^{-1}$.
    Half-open: a mode at exactly $\omega = $ \cfgkey{freq\_min} is
    \emph{included}. \\
\cfgkey{freq\_max} & \texttt{Optional[float]} & \texttt{None}
  & Upper bound. Half-open: a mode at exactly $\omega = $
    \cfgkey{freq\_max} is \emph{excluded} (so contiguous windows
    can be defined without overlap). \\
\cfgkey{freq\_windows} & \texttt{Optional[List[Tuple]]} & \texttt{None}
  & Multi-window analysis: a list of (lo, hi) tuples. When set,
    \cfgkey{freq\_min}/\cfgkey{freq\_max} are ignored. Each window
    triggers a separate Phase-4 run. \\
\cfgkey{temp\_k} & \texttt{Optional[float]} & \texttt{None}
  & Temperature in K for $\urms$ computation. \texttt{None}
    triggers classical-amplitude mode (uses \cfgkey{amplitude}
    instead, see below). \\
\bottomrule
\end{longtable}

\subsection*{Notes}
\begin{itemize}
  \item \cfgkey{freq\_min} and \cfgkey{freq\_max} are
        \emph{half-open}; this convention is documented in the
        Manual and tested by
        \fn{test\_v104\_audit\_fixes::test\_freq\_window\_half\_open}.
  \item \cfgkey{freq\_windows} is mutually exclusive with the
        single-window pair: the runner picks
        \cfgkey{freq\_windows} if set, otherwise the
        \cfgkey{freq\_min}/\cfgkey{freq\_max} pair.
  \item For \cfgkey{temp\_k} = \texttt{None}, the
        \cfgkey{amplitude} value (default 1.0 \AA{}) is used in
        place of $\urms$ throughout the per-mode analysis. This is
        intended for benchmarking against tools that don't apply
        thermal scaling -- production runs should set
        \cfgkey{temp\_k} explicitly.
\end{itemize}

% -------------------------------------------------------------------------
\section{Coordination detection thresholds}

\begin{longtable}{p{0.21\textwidth} p{0.13\textwidth} p{0.13\textwidth} p{0.45\textwidth}}
\toprule
\textbf{Field} & \textbf{Type} & \textbf{Default} & \textbf{Effect} \\
\midrule
\cfgkey{pdb\_chain} & \texttt{str} & \texttt{"A"}
  & PDB chain ID for residue-label mapping. Affects
    \fn{geometry.find\_coordinating\_residues}. \\
\cfgkey{fe\_coord\_cutoff\_n} & \texttt{float} & \texttt{2.50}
  & Maximum Fe-N distance (\AA{}) for a residue's N to be
    considered an Fe-binding ligand. Used in
    \fn{geometry.find\_coordinating\_residues}. \\
\cfgkey{fe\_coord\_cutoff\_s} & \texttt{float} & \texttt{2.70}
  & Maximum Fe-S distance for Cys-S binding. \\
\cfgkey{fe\_coord\_cutoff\_o} & \texttt{float} & \texttt{2.50}
  & Maximum Fe-O distance for Asp/Glu-O binding. Rare; supports
    mixed-coordination clusters. \\
\cfgkey{include\_hn\_vibration} & \texttt{bool} & \texttt{True}
  & Whether to compute the His-NH stretch and HA channel. When
    False, the hydrogen-inclusive eigenvector loader is skipped
    and the NH/HA sheets are filled with zeros. \\
\bottomrule
\end{longtable}

% -------------------------------------------------------------------------
\section{Classification thresholds}

\begin{longtable}{p{0.27\textwidth} p{0.13\textwidth} p{0.15\textwidth} p{0.40\textwidth}}
\toprule
\textbf{Field} & \textbf{Type} & \textbf{Default} & \textbf{Effect} \\
\midrule
\cfgkey{mode\_type\_threshold} & \texttt{float} & \texttt{60.0}
  & $\alpha_{\OOP}$ (\%) threshold for the binary
    OOP/INP/mixed classification (\cref{eq:oop-binary}). \\
\cfgkey{mode\_type\_detail\_thresholds} & \texttt{Tuple[3]}
       & \texttt{(60, 75, 90)}
  & Three thresholds for the eight-bin detailed classification:
    boundary between mixed and majority, between majority and
    strong, between strong and pure. \\
\cfgkey{significance\_threshold\_low} & \texttt{float} & \texttt{1.0}
  & Value-to-sigma ratio threshold for the ``below noise''
    colouring in Excel output. \\
\cfgkey{significance\_threshold\_high} & \texttt{float} & \texttt{3.0}
  & Ratio threshold for ``clearly above noise''. \\
\bottomrule
\end{longtable}

% -------------------------------------------------------------------------
\section{Analysis switches}

\begin{longtable}{p{0.27\textwidth} p{0.13\textwidth} p{0.15\textwidth} p{0.40\textwidth}}
\toprule
\textbf{Field} & \textbf{Type} & \textbf{Default} & \textbf{Effect} \\
\midrule
\cfgkey{analysis\_compact} & \texttt{bool} & \texttt{True}
  & Produce the compact main workbook with the most-used sheets. \\
\cfgkey{analysis\_full} & \texttt{bool} & \texttt{False}
  & Produce the full workbook with every available descriptor.
    Doubles output size; off by default. \\
\cfgkey{analyze\_sse} & \texttt{bool} & \texttt{True}
  & Run secondary-structure analysis if a PDB is supplied. \\
\cfgkey{analyze\_scsd} & \texttt{bool} & \texttt{True}
  & Run SCSD projection per mode. Setting False saves
    \mbox{$\sim$}~10\% runtime on large logs. \\
\cfgkey{sse\_chain} & \texttt{str} & \texttt{""}
  & Restrict SSE analysis to a single chain. Empty = all chains. \\
\bottomrule
\end{longtable}

% -------------------------------------------------------------------------
\section{Cluster discovery}

\begin{longtable}{p{0.27\textwidth} p{0.13\textwidth} p{0.15\textwidth} p{0.40\textwidth}}
\toprule
\textbf{Field} & \textbf{Type} & \textbf{Default} & \textbf{Effect} \\
\midrule
\cfgkey{fe\_s\_cutoff} & \texttt{float} & \texttt{3.0}
  & Maximum Fe-S distance for cluster discovery. Used by
    \fn{geometry.find\_cluster}. \\
\cfgkey{fe\_fe\_cutoff} & \texttt{float} & \texttt{3.5}
  & Maximum Fe-Fe distance for cluster discovery. \\
\cfgkey{cluster\_index} & \texttt{int} & \texttt{0}
  & Which cluster to analyse (0 = first). Ignored when
    \cfgkey{analyze\_all\_clusters} = True. \\
\cfgkey{analyze\_all\_clusters} & \texttt{bool} & \texttt{False}
  & Iterate Phase 4 over every detected cluster, producing
    independent output directories with the
    \texttt{\_cluster\#\#} suffix. \\
\cfgkey{strict\_cluster} & \texttt{bool} & \texttt{False}
  & Apply additional rhombic-geometry sanity checks during cluster
    discovery (\cref{sec:algo-cluster-discovery}). \\
\bottomrule
\end{longtable}

% -------------------------------------------------------------------------
\section{PCET acceptor heuristic}

\begin{longtable}{p{0.27\textwidth} p{0.13\textwidth} p{0.15\textwidth} p{0.40\textwidth}}
\toprule
\textbf{Field} & \textbf{Type} & \textbf{Default} & \textbf{Effect} \\
\midrule
\cfgkey{pcet\_enabled} & \texttt{bool} & \texttt{True}
  & Whether to compute the HA channel at all. Setting False
    disables all PCET-related descriptors. \\
\cfgkey{pcet\_hbond\_cutoff\_a} & \texttt{float} & \texttt{4.0}
  & Maximum H$\cdots$A distance (\AA{}) for an acceptor candidate. \\
\cfgkey{pcet\_acceptor\_r0\_a} & \texttt{float} & \texttt{2.8}
  & Gaussian weight centre $d_0$ for H$\cdots$A distance weighting
    (\cref{sec:dr-ha}). \\
\cfgkey{pcet\_acceptor\_sigma\_a} & \texttt{float} & \texttt{0.4}
  & Gaussian weight width $\sigma_d$. \\
\bottomrule
\end{longtable}

% -------------------------------------------------------------------------
\section{Modulation-spectrum broadening}

\begin{longtable}{p{0.27\textwidth} p{0.13\textwidth} p{0.15\textwidth} p{0.40\textwidth}}
\toprule
\textbf{Field} & \textbf{Type} & \textbf{Default} & \textbf{Effect} \\
\midrule
\cfgkey{reorg\_spectrum\_sigma\_cm1} & \texttt{float} & \texttt{5.0}
  & Gaussian broadening width $\sigma_M$ for the $M_X(\omega)$
    spectra (\cref{eq:M-spectrum-derivation}). \\
\cfgkey{reorg\_spectrum\_step\_cm1} & \texttt{float} & \texttt{0.5}
  & Uniform $\omega$-grid step for the modulation spectra and
    cumulative $\Lambda_X(\omega)$ curves. \\
\bottomrule
\end{longtable}

% -------------------------------------------------------------------------
\section{Precision and noise model}

\begin{longtable}{p{0.27\textwidth} p{0.13\textwidth} p{0.15\textwidth} p{0.40\textwidth}}
\toprule
\textbf{Field} & \textbf{Type} & \textbf{Default} & \textbf{Effect} \\
\midrule
\cfgkey{include\_hydrogen} & \texttt{bool} & \texttt{True}
  & Whether the cluster-core analysis should include any hydrogen
    atoms. Reserved for special workflows; mostly equivalent to
    \cfgkey{include\_hn\_vibration} = True. \\
\cfgkey{sigma\_eigvec} & \texttt{float} & \texttt{5e-4}
  & Per-component eigenvector noise floor. Drives every Gauss-
    propagation calculation (\cref{part:uncertainty}). \\
\cfgkey{sigma\_coord} & \texttt{float} & \texttt{1e-3}
  & Per-component coordinate noise floor (e.g.\ from Gaussian's
    geometry-print precision). \\
\cfgkey{amplitude\_threshold} & \texttt{float} & \texttt{0.001}
  & Below this thermal amplitude (\AA{}), per-mode percentages
    are reported as NaN to prevent noise amplification. \\
\cfgkey{amplitude} & \texttt{float} & \texttt{1.0}
  & Classical-mode amplitude in \AA{} when \cfgkey{temp\_k} =
    \texttt{None}. Useful only for benchmarking. \\
\cfgkey{coord\_match\_tol} & \texttt{float} & \texttt{1.0}
  & Distance tolerance (\AA{}) for the post-Kabsch nearest-
    neighbour PDB-to-Gaussian mapping (\cref{sec:kabsch}). \\
\bottomrule
\end{longtable}

% -------------------------------------------------------------------------
\section{Performance}

\begin{longtable}{p{0.27\textwidth} p{0.13\textwidth} p{0.15\textwidth} p{0.40\textwidth}}
\toprule
\textbf{Field} & \textbf{Type} & \textbf{Default} & \textbf{Effect} \\
\midrule
\cfgkey{scan\_chunk\_mb} & \texttt{int} & \texttt{16}
  & Disk-scan chunk size in MiB. Larger = faster scan on SSD, more
    RAM use. \\
\cfgkey{show\_errors\_in\_excel} & \texttt{bool} & \texttt{True}
  & Whether to render mode-loading errors as red-coloured cells
    in the Excel output. When False, errors are only in the
    run-log. \\
\cfgkey{use\_cache} & \texttt{bool} & \texttt{True}
  & Read/write the scan cache (\cref{sec:cache}). Set False for
    one-off runs or to force re-scanning. \\
\bottomrule
\end{longtable}

% -------------------------------------------------------------------------
\section{Interpolation workbook}

\begin{longtable}{p{0.27\textwidth} p{0.13\textwidth} p{0.15\textwidth} p{0.40\textwidth}}
\toprule
\textbf{Field} & \textbf{Type} & \textbf{Default} & \textbf{Effect} \\
\midrule
\cfgkey{interp\_step} & \texttt{float} & \texttt{0.05}
  & Frequency grid step (cm$^{-1}$) for the interp workbook. \\
\cfgkey{interp\_boundary\_mode} & \texttt{str} & \texttt{"context"}
  & Boundary handling: \texttt{"context"} loads context modes
    (\cref{sec:algo-context}), \texttt{"zero"} forces zero outside
    the window, \texttt{"extend"} flatly extends the boundary
    value. \\
\cfgkey{interp\_edge\_extend} & \texttt{float} & \texttt{0.5}
  & Edge-extension distance (cm$^{-1}$) for
    \texttt{interp\_boundary\_mode} = \texttt{"extend"}. \\
\cfgkey{interp\_context\_cm1} & \texttt{float} & \texttt{5.0}
  & Width of the context-mode search window
    (\cref{sec:algo-context}). \\
\bottomrule
\end{longtable}

% -------------------------------------------------------------------------
\section{Embedding}

\begin{longtable}{p{0.27\textwidth} p{0.13\textwidth} p{0.15\textwidth} p{0.40\textwidth}}
\toprule
\textbf{Field} & \textbf{Type} & \textbf{Default} & \textbf{Effect} \\
\midrule
\cfgkey{umap\_n\_neighbors} & \texttt{Optional[int]} & \texttt{None}
  & UMAP neighborhood size. \texttt{None} = auto-adapt to
    $N_\text{modes}$. \\
\cfgkey{export\_embedding\_plots} & \texttt{bool} & \texttt{True}
  & Whether to render and save PNG plots of the UMAP-HDBSCAN
    clusterings. \\
\cfgkey{logname\_suffix} & \texttt{str} & \texttt{""}
  & Optional suffix appended to all output filenames (useful for
    differentiating windows or systems within one output dir). \\
\bottomrule
\end{longtable}

% =========================================================================
% PART VI -- OUTPUT REFERENCE
% =========================================================================
\part{Output Reference}
\label{part:output}

This part is the column-by-column documentation of every Excel
workbook produced by the pipeline. For each workbook the sheets are
enumerated; for each sheet, every column is given with its type,
units, source function, and any propagated-sigma column that
accompanies it. The implementation lives in \file{export.py}.

% -------------------------------------------------------------------------
\section{The main workbook (\file{*\_analysis\_main.xlsx})}
\label{sec:wb-main}

\codref{export.export\_all, export.export\_main\_excel}

The main workbook contains the per-mode descriptors and the
system-wide aggregates. Sheets are grouped thematically:

\begin{itemize}
  \item \emph{Overview} sheets: \texttt{Modes\_overview} (per mode:
        number, freq, red mass, mode type, primary kernel, classifier
        precision) -- the index sheet.
  \item \emph{Cluster-core} sheets: \texttt{Cluster\_OOP\_INP},
        \texttt{Cluster\_displacement},
        \texttt{Cluster\_expansion\_rotation}, \texttt{SCSD},
        \texttt{Kernel\_class}.
  \item \emph{Per-bond} sheets: \texttt{Fe\_S\_Cys},
        \texttt{Fe\_N\_His} (and \texttt{Fe\_O\_Asp} when present);
        each contains per-mode \texttt{stretch}, \texttt{bend\_inp},
        \texttt{bend\_oop} per ligand, plus their sigmas.
  \item \emph{Marcus-Hush} sheets: \texttt{Reorganisation\_per\_mode}
        (per-mode $\lambda_X$ and $\Delta r_X^\text{RSS}$),
        \texttt{Reorganisation\_total} (system-wide
        $\Lambda_X^\text{total}$ per channel),
        \texttt{Modulationsspektrum} (mode-resolved $M_X(\omega)$),
        \texttt{Cumulative\_Lambda}.
  \item \emph{Diagnostic} sheets: \texttt{Coordination} (the
        Gaussian-to-PDB atom mapping used throughout),
        \texttt{B\_factors} (the Phase-7 accumulated Debye-Waller
        factors per atom).
\end{itemize}

\subsection{Sheet: \texttt{Modes\_overview}}

\begin{longtable}{p{0.22\textwidth} p{0.13\textwidth} p{0.13\textwidth} p{0.45\textwidth}}
\toprule
\textbf{Column} & \textbf{Type} & \textbf{Units} & \textbf{Source} \\
\midrule
\texttt{number}            & int    & --            & \fn{logio.read\_all\_meta} \\
\texttt{freq}              & float  & cm$^{-1}$     & \fn{logio.read\_all\_meta} \\
\texttt{red\_mass}         & float  & amu           & \fn{logio.read\_all\_meta} \\
\texttt{u\_rms\_a}         & float  & \AA{}         & \fn{core.compute\_thermal\_amplitude} \\
\texttt{mode\_type}        & str    & --            & \fn{core.classify\_oop\_inp} \\
\texttt{a\_oop\_pct}       & float  & \%            & \fn{core.classify\_oop\_inp} \\
\texttt{kern\_primary}     & str    & --            & \fn{core.classify\_kernel\_mode\_from\_evg} \\
\texttt{kern\_primary\_score} & float & \%          & \fn{core.classify\_kernel\_mode\_from\_evg} \\
\texttt{precision}         & str    & --            & \fn{logio.get\_eigvec} \\
\bottomrule
\end{longtable}

\subsection{Sheet: \texttt{Fe\_S\_Cys}}

One row per mode. Columns repeated for each Cys ligand
(suffix = \texttt{res\_label}, e.g.\ \texttt{stretch\_Cys11},
\texttt{stretch\_Cys15}, etc.).

\begin{longtable}{p{0.30\textwidth} p{0.13\textwidth} p{0.13\textwidth} p{0.37\textwidth}}
\toprule
\textbf{Column} & \textbf{Type} & \textbf{Units} & \textbf{Source} \\
\midrule
\texttt{number}              & int    & --     & \\
\texttt{freq}                & float  & cm$^{-1}$ & \\
\texttt{stretch\_\textit{L}} & float  & \AA{}  & \fn{core.analyze\_fe\_ligand}, \cref{eq:fel-stretch} \\
\texttt{s\_stretch\_\textit{L}} & float & \AA{} & \cref{eq:fel-sigma} \\
\texttt{bend\_inp\_\textit{L}}  & float  & \AA{}  & \cref{eq:fel-bend-split} \\
\texttt{s\_bend\_inp\_\textit{L}}  & float  & \AA{}  & \cref{eq:fel-sigma} \\
\texttt{bend\_oop\_\textit{L}}  & float  & \AA{}  & \cref{eq:fel-bend-split} \\
\texttt{s\_bend\_oop\_\textit{L}}  & float  & \AA{}  & \cref{eq:fel-sigma} \\
\texttt{bend\_inp\_pct\_\textit{L}} & float  & \%   & \fn{core.analyze\_fe\_ligand} \\
\texttt{bend\_oop\_pct\_\textit{L}} & float  & \%   & \fn{core.analyze\_fe\_ligand} \\
\bottomrule
\end{longtable}

\subsection{Sheet: \texttt{Reorganisation\_per\_mode}}

One row per mode, one column-group per parent channel.

\begin{longtable}{p{0.27\textwidth} p{0.13\textwidth} p{0.18\textwidth} p{0.37\textwidth}}
\toprule
\textbf{Column} & \textbf{Type} & \textbf{Units} & \textbf{Source} \\
\midrule
\texttt{dr\_rss\_FeFe}     & float  & \AA{}     & \fn{reorganization.aggregate\_by\_parent}, \cref{eq:dr-rss-derivation} \\
\texttt{s\_dr\_rss\_FeFe}  & float  & \AA{}     & propagated, \cref{part:uncertainty} \\
\texttt{lambda\_FeFe}      & float  & cm$^{-1}$ & \fn{reorganization.lambda\_pair\_cm1}, \cref{eq:lambda-per-mode} \\
\texttt{s\_lambda\_FeFe}   & float  & cm$^{-1}$ & propagated \\
\texttt{dr\_rss\_FeN}      & float  & \AA{}     & same as FeFe row \\
\texttt{...\_FeS}, \texttt{...\_NH}, \texttt{...\_HA} & & & \\
\bottomrule
\end{longtable}

\subsection{Sheet: \texttt{Reorganisation\_total}}

A single row, one column per parent channel.

\begin{longtable}{p{0.27\textwidth} p{0.13\textwidth} p{0.18\textwidth} p{0.37\textwidth}}
\toprule
\textbf{Column} & \textbf{Type} & \textbf{Units} & \textbf{Source} \\
\midrule
\texttt{Lambda\_total\_FeFe} & float & cm$^{-1}$ & \cref{eq:lambda-total-derivation} \\
\texttt{Lambda\_total\_FeN}  & float & cm$^{-1}$ & \\
\texttt{Lambda\_total\_FeS}  & float & cm$^{-1}$ & \\
\texttt{Lambda\_total\_NH}   & float & cm$^{-1}$ & \\
\texttt{Lambda\_total\_HA}   & float & cm$^{-1}$ & \\
\bottomrule
\end{longtable}

\subsection{Sheet: \texttt{Modulationsspektrum}}

One row per $\omega$-grid point, one column per parent channel.

\begin{longtable}{p{0.27\textwidth} p{0.13\textwidth} p{0.18\textwidth} p{0.37\textwidth}}
\toprule
\textbf{Column} & \textbf{Type} & \textbf{Units} & \textbf{Source} \\
\midrule
\texttt{omega}     & float & cm$^{-1}$ & grid generated from
                                          \cfgkey{reorg\_spectrum\_step\_cm1} \\
\texttt{M\_FeFe}   & float & \AA{}     & \fn{reorganization.compute\_modulation\_spectra}, \cref{eq:M-spectrum-derivation} \\
\texttt{M\_FeN}, \texttt{M\_FeS}, \texttt{M\_NH}, \texttt{M\_HA} & & & \\
\texttt{C\_PCET}   & float & \AA{}     & \fn{reorganization.compute\_co\_modulation\_spectra}, \cref{eq:co-modulation} \\
\texttt{C\_PT\_FeN}, \texttt{C\_ET\_FeS} & & & \\
\texttt{Lambda\_FeFe} & float & cm$^{-1}$ & cumulative
                            $\Lambda_X(\omega)$ \cref{eq:lambda-cumulative} \\
\texttt{Lambda\_FeN}, \texttt{Lambda\_FeS}, \texttt{Lambda\_NH},
\texttt{Lambda\_HA} & & & \\
\bottomrule
\end{longtable}

\subsection{Sheet: \texttt{SCSD}}

One row per mode, one column-pair per $D_{2\mathrm{h}}$ irrep.

\begin{longtable}{p{0.27\textwidth} p{0.13\textwidth} p{0.18\textwidth} p{0.37\textwidth}}
\toprule
\textbf{Column} & \textbf{Type} & \textbf{Units} & \textbf{Source} \\
\midrule
\texttt{number}     & int    & --     & \\
\texttt{freq}       & float  & cm$^{-1}$ & \\
\texttt{scsd\_Ag}, \texttt{scsd\_B1g}, \texttt{scsd\_B2g},
\texttt{scsd\_B3g}, \texttt{scsd\_Au}, \texttt{scsd\_B1u},
\texttt{scsd\_B2u}, \texttt{scsd\_B3u}
   & float  & \AA{}     & \fn{core.compute\_scsd\_for\_mode\_full},
                          \cref{eq:scsd-projection-derivation} \\
\texttt{s\_scsd\_*} & float & \AA{} & propagated (\cref{sec:scsd-derivation}) \\
\bottomrule
\end{longtable}

\subsection{Sheet: \texttt{Coordination}}

The \texttt{Coordination} sheet documents the master ligand list
used by all per-mode analyses. One row per Fe-ligand bond.

\begin{longtable}{p{0.27\textwidth} p{0.13\textwidth} p{0.13\textwidth} p{0.42\textwidth}}
\toprule
\textbf{Column} & \textbf{Type} & \textbf{Units} & \textbf{Source} \\
\midrule
\texttt{Fe}             & str   & --     & e.g.\ \texttt{Fe1}, \texttt{Fe2}. \\
\texttt{Fe\_center}     & int   & --     & Gaussian centre of Fe. \\
\texttt{Lig\_center}    & int   & --     & Gaussian centre of ligand donor. \\
\texttt{Lig\_element}   & str   & --     & N, S, or O. \\
\texttt{Lig\_aname}     & str   & --     & PDB atom name. \\
\texttt{Residue}        & str   & --     & e.g.\ \texttt{Cys11}, \texttt{His87}. \\
\texttt{bond\_length\_a} & float & \AA{} & equilibrium Fe-L bond length. \\
\texttt{H\_center}      & int / NA & --  & if applicable (His-NH protonated). \\
\bottomrule
\end{longtable}

% -------------------------------------------------------------------------
\section{Interpolation workbook (\file{*\_analysis\_interp\#\#.xlsx})}

\codref{export.export\_interpolated\_excel}

Contains the per-mode descriptors re-sampled onto a uniform
$\omega$-grid of step \cfgkey{interp\_step}. The boundary-mode
behaviour is governed by \cfgkey{interp\_boundary\_mode} and
\cfgkey{interp\_context\_cm1}.

\subsection*{Sheets}

\begin{itemize}
  \item \texttt{Interpoliert\_per\_channel}: $\omega$ vs.\
        per-channel $\Delta r_X^\text{RSS}(\omega)$, $\lambda_X(\omega)$.
  \item \texttt{Cumulative}: $\omega$ vs.\ cumulative
        $\Lambda_X(\omega_\text{cut})$ per channel.
  \item \texttt{Interp\_bond\_metrics}: $\omega$ vs.\ per-bond
        \texttt{stretch}, \texttt{bend\_inp}, \texttt{bend\_oop}
        (one column per ligand).
\end{itemize}

The interp values are computed by linear interpolation between
adjacent modes (in $\omega$), anchored at the boundary by the
context modes loaded in \cref{sec:algo-context}.

% -------------------------------------------------------------------------
\section{SSE workbook (\file{*\_analysis\_SSE.xlsx})}

\codref{export.export\_sse\_excel}

Produced only when SSE analysis is enabled and a PDB with SSE records
(or DSSP records) is supplied. Sheets enumerate per-SSE-element
amplitudes per mode, plus axial-tilt diagnostics.

\subsection*{Sheets}

One sheet per descriptor, each a per-mode $\times$ per-SSE-element matrix.
The element's principal axis is taken from its C$_\alpha$ backbone;
\emph{axial} and \emph{lateral} denote the components parallel and
perpendicular to that axis. Displacements use the thermally scaled
eigenvector $\widetilde{\evec}$; amplitudes are in \AA{}, angles in degrees.

\begin{itemize}
  \item \texttt{SSE\_amplitude\_mean} / \texttt{SSE\_amplitude\_max}:
        mean and maximum, over the element's atoms, of the per-atom
        displacement magnitude.
  \item \texttt{SSE\_com\_amplitude}: magnitude of the \emph{mass-weighted}
        centre-of-mass displacement
        $\lVert \sum_i m_i \widetilde{\evec}_i / \sum_i m_i \rVert$
        (rigid-body translation of the element). \emph{(v1.0.5: now
        mass-weighted; previously the unweighted atom-mean.)}
  \item \texttt{SSE\_axial\_amplitude}: mean absolute displacement
        component along the principal axis.
  \item \texttt{SSE\_lateral\_amplitude}: mean displacement component
        perpendicular to the principal axis.
  \item \texttt{SSE\_lateral\_std}: standard deviation of the lateral
        component across the atoms (lateral inhomogeneity).
  \item \texttt{SSE\_stretching}: standard deviation of the axial
        component across the atoms (axial elongation/compression).
  \item \texttt{SSE\_tilting\_angle}: magnitude of the rigid-body rocking
        rotation perpendicular to the axis. It is the best-fit
        infinitesimal rotation $\bm{\omega}$ minimising
        $\sum_i m_i \lVert \mathbf{d}_i - \bm{\omega}\times\mathbf{r}_i \rVert^2$
        ($\mathbf{d}_i$ = displacement after removing the COM translation,
        $\mathbf{r}_i$ = position about the mass centre), reported as
        $\lVert\bm{\omega}_\perp\rVert$ in degrees. \emph{(v1.0.5: a genuine
        rigid-body tilt; previously the near-constant polar angle of the
        centroid translation.)}
  \item \texttt{SSE\_internal\_amplitude}: mean magnitude of the residual
        displacement after removing \emph{both} the rigid translation and
        the rigid rotation -- a TLS-style non-rigid internal deformation
        \emph{(new in v1.0.5)}.
  \item \texttt{SSE\_UMAP\_Cluster} (optional): present when the
        SSE-fingerprint UMAP--HDBSCAN clustering is computed
        (\fn{compute\_sse\_umap\_cluster}); HDBSCAN cluster ID per mode
        ($-1$ = noise). The clustering uses the decorrelated five-feature
        subset \texttt{com\_amplitude}, \texttt{tilting\_angle},
        \texttt{internal\_amplitude}, \texttt{axial\_amplitude},
        \texttt{lateral\_amplitude}.
\end{itemize}

Each amplitude sheet carries a matching \texttt{s\_}-sigma
(\cref{sec:sigma-propagation}): the mean-type amplitudes use the standard
error $\sigma_e^\text{therm}/\sqrt{n}$, the tilt angle uses
$\sigma_e^\text{therm}/R_g$ converted to degrees.

% -------------------------------------------------------------------------
\section{Embedding workbook (\file{*\_analysis\_embedding.xlsx})}

\codref{export.export\_embedding\_excel}

Three parallel UMAP-HDBSCAN clusterings, on three complementary
feature spaces:

\begin{itemize}
  \item \emph{Reorg-features}: per-mode
        $(\lambda_\text{FeFe}, \lambda_\text{FeN}, \lambda_\text{FeS},
          \lambda_\text{NH}, \lambda_\text{HA})$.
  \item \emph{SSE-fingerprint}: per-mode SSE-element amplitudes
        (only when SSE is enabled).
  \item \emph{C$_\alpha$-amplitude} fingerprint: per-mode mean
        $\alpha$-carbon amplitudes per residue group.
\end{itemize}

Each clustering produces 4 sheets: \texttt{Projektionen} (per-mode
UMAP-$x$, UMAP-$y$), \texttt{Ca\_Amplituden} (the underlying
feature matrix), \texttt{UMAP\_Cluster} (HDBSCAN cluster IDs per
mode; $-1$ = noise), \texttt{UMAP\_Cluster\_Profil} (mean feature
vector per cluster). Together these three workbooks form the
hierarchical-clustering signature of the system.

% =========================================================================
% PART VII -- UNCERTAINTY PROPAGATION COMPLETE
% =========================================================================
\part{Uncertainty Propagation Complete}
\label{part:uncertainty}
\label{sec:sigma-propagation}

This part derives every sigma field reported in the output workbooks
from the two user-controllable noise floors
$\sigma_e = \cfgkey{sigma\_eigvec}$ and
$\sigma_c = \cfgkey{sigma\_coord}$, applied to the relevant
quantities via standard linearised Gauss propagation.

% -------------------------------------------------------------------------
\section{Underlying assumptions}
\label{sec:uncertainty-overview}

Every linearised sigma in this section assumes:

\begin{itemize}
  \item Eigenvector components $\evec_{l,i}^{(\alpha)}$ are
        independent Gaussian random variables with zero mean and
        per-component standard deviation $\sigma_e$ (default
        $5 \times 10^{-4}$, motivated by the printed precision of
        Gaussian \texttt{hpmodes}).
  \item Atomic coordinates $r_i^{(\alpha)}$ are independent Gaussian
        random variables with per-component standard deviation
        $\sigma_c$ (default $10^{-3}$~\AA{}, motivated by the
        Standard-Orientation print format).
  \item Frequencies $\omega_l$ and reduced masses $m_l$ are
        treated as exact (their printed precision is sufficient
        for the linear-noise regime).
  \item Eigenvector and coordinate noises are uncorrelated.
\end{itemize}

The general rule for a smooth function
$y = f(\evec, \rvec)$ is
\begin{equation}
  \sigma_y^2 \;\approx\;
   \sum_i (\partial_{\evec_i} f)^2 \sigma_e^2
   + \sum_i (\partial_{\rvec_i} f)^2 \sigma_c^2.
\end{equation}
This expansion is applied below to each output.

% -------------------------------------------------------------------------
\section{Thermal scaling of $\sigma_e$}
\label{sec:sigma-therm}

\codref{core.compute\_thermal\_amplitude}

Every geometric observable consumed downstream uses the thermally
scaled eigenvector $\widetilde{\evec} = \evec \cdot \urms$. The
corresponding scaled per-component noise is
\begin{equation}
  \boxed{\;\sigma_e^\text{therm}(\omega_l)
    \;=\;
   \sigma_e \cdot \urms(\omega_l)\;}
  \label{eq:sigma-therm}
\end{equation}
This rescaling is the universally applied convention; failure to
apply it was the FUND 6/7/8 bug class in v1.0.4
(\cref{sec:fund-6-8}). For $\omega = 200$ cm$^{-1}$, $T = 80$~K,
$m_\text{red} = 30$~amu, $\urms \approx 0.05$~\AA{}, so
$\sigma_e^\text{therm} \approx 2.5 \times 10^{-5}$~\AA{}.

% -------------------------------------------------------------------------
\section{Bond-modulation sigmas: \texttt{s\_stretch}, \texttt{s\_bend\_inp}, \texttt{s\_bend\_oop}}

\codref{core.analyze\_fe\_ligand}

The bond stretch $\Delta r_\text{stretch}$
(\cref{eq:fel-stretch}) is the dot product of a 3-vector difference
with a unit vector:
\begin{equation}
  \Delta r_\text{stretch} \;=\; (\widetilde{\evec}_L - \widetilde{\evec}_\text{Fe}) \cdot \bhat.
\end{equation}
Treating $\bhat$ as constant (since coordinates have their own
sigma $\sigma_c$, treated separately) and the eigenvector
components as independent, the variance of the dot product is
\begin{align}
  \sigma_\text{stretch}^2
    &\;=\; \sum_\alpha b_\alpha^2 \cdot
              \text{Var}\!\left[(\widetilde{\evec}_L - \widetilde{\evec}_\text{Fe})_\alpha\right]
   \;=\; \sum_\alpha b_\alpha^2 \cdot 2 (\sigma_e^\text{therm})^2
   \;=\; 2 (\sigma_e^\text{therm})^2,
\end{align}
where the last equality uses $\sum_\alpha b_\alpha^2 = 1$ (unit
vector). Therefore
\begin{equation}
  \boxed{\;\sigma_\text{stretch}
    \;=\;
   \sqrt{2}\,\sigma_e \cdot \urms(\omega_l)\;}
  \label{eq:sigma-stretch}
\end{equation}
This is the formula in \cref{eq:fel-sigma}. The same argument
applies to $\Delta r_\text{bend\_inp}$ and $\Delta r_\text{bend\_oop}$:
they are obtained from the same difference vector, projected onto
different orthonormal axes; each carries the same
$\sqrt{2}\,\sigma_e^\text{therm}$.

% -------------------------------------------------------------------------
\section{N-H stretch sigma: \texttt{s\_hn\_stretch}}

\codref{core.analyze\_his\_hn}

Same form as \cref{eq:sigma-stretch}, with the v1.0.4 FUND 6 fix
(\cref{sec:fund-6-8}) applying $\urms$:
\begin{equation}
  \boxed{\;\sigma_\text{NH}
    \;=\;
   \sqrt{2}\,\sigma_e \cdot \urms(\omega_l)\;}
\end{equation}

% -------------------------------------------------------------------------
\section{HA reaction-coordinate sigma: \texttt{s\_ha\_dr}}

\codref{reorganization.compute\_mode\_modulations}

The HA reaction coordinate \cref{eq:dr-ha-final} is
\begin{equation}
  \Delta r_{HA} \;=\; (\widetilde{\evec}_H - \widetilde{\evec}_N) \cdot \nhat_{NA}.
\end{equation}
Again a unit-norm projection of a difference of two independent
eigenvectors:
\begin{equation}
  \boxed{\;\sigma_{HA}
    \;=\;
   \sqrt{2}\,\sigma_e \cdot \urms(\omega_l)\;}
\end{equation}

% -------------------------------------------------------------------------
\section{RSS-aggregate sigma: \texttt{s\_dr\_rss\_X}}

\codref{reorganization.aggregate\_by\_parent}

For the RSS-aggregated parent-channel modulation
$\Delta r_X^\text{RSS} = \sqrt{\sum_s w_s (\Delta r_s)^2}$, the
linearised propagation of independent $\Delta r_s$ sigmas (each
$= \sqrt{2}\sigma_e \urms$) gives
\begin{align}
  \sigma_\text{RSS}^2
   &\;=\; \sum_s \left(\frac{\partial \Delta r^\text{RSS}}{\partial \Delta r_s}\right)^2
                 \sigma_{\Delta r_s}^2 \\
   &\;=\; \sum_s \left(\frac{w_s \Delta r_s}{\Delta r^\text{RSS}}\right)^2
                  \cdot 2 (\sigma_e^\text{therm})^2 \\
   &\;=\; \frac{2 (\sigma_e^\text{therm})^2}{(\Delta r^\text{RSS})^2}
            \sum_s w_s^2 (\Delta r_s)^2.
\end{align}
So
\begin{equation}
  \boxed{\;\sigma_\text{RSS}^X
    \;=\;
   \sqrt{2}\,\sigma_e \cdot \urms(\omega_l) \cdot
   \frac{\sqrt{\sum_s w_s^2 (\Delta r_s)^2}}{\Delta r^\text{RSS}}\;}
  \label{eq:sigma-rss}
\end{equation}
For the single-sub-channel limit ($|s| = 1$, $w = 1$), the
square-root factor reduces to 1, and $\sigma_\text{RSS} =
\sqrt{2}\,\sigma_e^\text{therm}$ as expected.

For the special case of all-equal sub-channels ($w_s = w$,
$|\Delta r_s| = \Delta r_0$), the factor becomes $1/\sqrt{|s|}$ --
i.e.\ averaging over $|s|$ sub-channels reduces the relative
uncertainty by $\sqrt{|s|}$.

% -------------------------------------------------------------------------
\section{Per-mode $\lambda_X$ sigma: \texttt{s\_lambda\_X}}

\codref{reorganization.compute\_mode\_modulations}

$\lambda_X = \tfrac{1}{2}\mu\omega^2 (\Delta r^\text{RSS})^2 / (hc)$.
For exact $\mu, \omega, hc$, the variance is
\begin{equation}
  \sigma_{\lambda_X}^2
    \;=\;
   \left(\frac{\partial \lambda_X}{\partial \Delta r^\text{RSS}}\right)^2
     \cdot \sigma_\text{RSS}^2
    \;=\;
   \left(\frac{\mu\omega^2 \Delta r^\text{RSS}}{hc}\right)^2 \cdot \sigma_\text{RSS}^2.
\end{equation}
So
\begin{equation}
  \boxed{\;\sigma_{\lambda_X}
    \;=\;
   \frac{2\lambda_X}{\Delta r^\text{RSS}} \cdot \sigma_\text{RSS}\;}
  \label{eq:sigma-lambda}
\end{equation}
i.e.\ the relative uncertainty on $\lambda$ is twice the
relative uncertainty on the RSS modulation -- a direct
consequence of $\lambda \propto (\Delta r)^2$.

% -------------------------------------------------------------------------
\section{SCSD sigmas: \texttt{s\_scsd\_*}}

\codref{core.compute\_scsd\_for\_mode\_full}

The SCSD reference projection uses only equilibrium coordinates:
\begin{equation}
  \sigma_\text{SCSD,ref}^{(\Gamma)} \;=\; \sqrt{2}\,\sigma_c.
\end{equation}
The distorted projection uses both coordinates (for the centre-
of-mass term) and the thermally scaled eigenvector:
\begin{equation}
  \sigma_\text{SCSD,dist}^{(\Gamma)}
    \;=\;
   \sqrt{2}\,\sqrt{\sigma_c^2 + (\sigma_e \cdot \urms(\omega_l))^2}.
\end{equation}
The difference SCSD is the difference of the two; by variance
addition (independent),
\begin{equation}
  \boxed{\;\sigma_\text{SCSD,diff}^{(\Gamma)}
    \;=\;
   \sqrt{(\sigma_\text{SCSD,ref}^{(\Gamma)})^2 + (\sigma_\text{SCSD,dist}^{(\Gamma)})^2}\;}
\end{equation}
This is the v1.0.4 FUND 8 fix (\cref{sec:fund-6-8}): pre-v1.0.4,
hard-coded constants ($5 \times 10^{-7}, 5 \times 10^{-6}$) were
used instead.

% -------------------------------------------------------------------------
\section{Cluster expansion and rotation sigmas}

\codref{core.classify\_oop\_inp}

Both are derived from the cluster-COM displacement plus
post-projection cleanup. The COM is an average of 4 atomic
displacements, so its sigma is $\sigma_e^\text{therm}/2$. The
cluster-expansion observable subtracts this COM from each cluster
atom and reads the radial component:
\begin{equation}
  \sigma_\text{expansion}
    \;=\;
   \sqrt{(\sigma_e^\text{therm})^2 + (\sigma_e^\text{therm}/2)^2}
    \;=\;
   \tfrac{\sqrt{5}}{2}\,\sigma_e^\text{therm}.
\end{equation}
The rotation angle is a small-angle approximation of the in-plane
projected displacement, so its sigma is the in-plane sigma divided
by the mean cluster radius.

% -------------------------------------------------------------------------
\section{Aggregate sigmas in totals: \texttt{s\_Lambda\_total\_X}}

\codref{reorganization.compute\_total\_reorganization}

The system-wide $\Lambda_X^\text{total}$ is a sum of per-mode
$\lambda_X(i)$ values. Assuming the per-mode sigmas are
independent,
\begin{equation}
  \sigma_{\Lambda_X^\text{total}}
    \;=\;
   \sqrt{\sum_i \sigma_{\lambda_X(i)}^2}.
\end{equation}
For typical analyses with $\sim 100$ modes per window, the
relative sigma on the total is $\sim 10\%$ of the relative sigma
on a single mode.

% -------------------------------------------------------------------------
\section{The 24 \texttt{s\_*} fields enumerated}

The complete list of propagated sigmas reported in the output:

\begin{longtable}{p{0.32\textwidth} p{0.18\textwidth} p{0.45\textwidth}}
\toprule
\textbf{Sigma field} & \textbf{Formula} & \textbf{Comment} \\
\midrule
\texttt{s\_stretch\_\textit{L}}        & $\sqrt{2}\sigma_e\urms$ & per Fe-L bond \\
\texttt{s\_bend\_inp\_\textit{L}}      & $\sqrt{2}\sigma_e\urms$ & \\
\texttt{s\_bend\_oop\_\textit{L}}      & $\sqrt{2}\sigma_e\urms$ & \\
\texttt{s\_hn\_stretch\_\textit{His}}  & $\sqrt{2}\sigma_e\urms$ & FUND 6 fix \\
\texttt{s\_ha\_dr\_\textit{His}\_to\_\textit{A}} & $\sqrt{2}\sigma_e\urms$ & FUND 7 fix \\
\texttt{s\_dr\_rss\_FeFe}              & \cref{eq:sigma-rss} & \\
\texttt{s\_dr\_rss\_FeN}               & \cref{eq:sigma-rss} & \\
\texttt{s\_dr\_rss\_FeS}               & \cref{eq:sigma-rss} & \\
\texttt{s\_dr\_rss\_NH}                & \cref{eq:sigma-rss} & \\
\texttt{s\_dr\_rss\_HA}                & \cref{eq:sigma-rss} & \\
\texttt{s\_lambda\_FeFe}               & \cref{eq:sigma-lambda} & \\
\texttt{s\_lambda\_FeN}                & \cref{eq:sigma-lambda} & \\
\texttt{s\_lambda\_FeS}                & \cref{eq:sigma-lambda} & \\
\texttt{s\_lambda\_NH}                 & \cref{eq:sigma-lambda} & \\
\texttt{s\_lambda\_HA}                 & \cref{eq:sigma-lambda} & \\
\texttt{s\_scsd\_Ag, \ldots, B3u}      & \cref{sec:scsd-derivation} & 8 irrep amplitudes (FUND 8 fix) \\
\texttt{s\_expansion\_a}               & $\tfrac{\sqrt{5}}{2}\sigma_e\urms$ & \\
\texttt{s\_rotation\_deg}              & in-plane sigma / radius & \\
\texttt{s\_lambda\_total\_X} (5 channels) & RSS over modes & system-wide \\
\bottomrule
\end{longtable}

\paragraph{Significance colouring in Excel.}
The export layer compares each value to its sigma and assigns a
colour: below \cfgkey{significance\_threshold\_low}~$\sigma$ = grey
(below noise); between low and \cfgkey{significance\_threshold\_high}~$\sigma$
= yellow (marginal); above high~$\sigma$ = white (significant).
This makes the noise floor visually apparent in the workbook.

% =========================================================================
% PART VIII -- AUDIT TRAIL OF v1.0.4
% =========================================================================
\part{Audit Trail of v1.0.4}
\label{sec:audit-trail}

This part documents every code-level and documentation-level change
between v1.0.3 and v1.0.4. The two principal categories are
\emph{fundamental code findings} (\fn{FUND}~1 through \fn{FUND}~13)
and \emph{manual drift findings}. Each entry follows a uniform
template: \emph{symptom}, \emph{root cause}, \emph{fix},
\emph{regression test reference}, \emph{migration impact}.

% -------------------------------------------------------------------------
\section{FUND 1: \texttt{evec\_raw} mutation in place}
\label{sec:fund-1}

\paragraph{Symptom.}
After Phase 4, an attempt to re-use the original (un-scaled)
eigenvector for diagnostics revealed that \texttt{r["evec\_raw"]}
was numerically identical to \texttt{r["evg"]} (the scaled one).
The raw eigenvector had been clobbered.

\paragraph{Root cause.}
In \fn{core.analyze\_mode\_with\_fallback}, the line
\texttt{evg = evec\_raw * u\_rms} relied on numpy creating a new
array. But an earlier in-place line
\texttt{evec\_raw *= u\_rms} (left over from a refactor) was
modifying the source array in place. The result-dict's
\texttt{evec\_raw} key then pointed to the already-scaled array.

\paragraph{Fix.}
Replaced the in-place op with \texttt{evg = evec\_raw * u\_rms}
(explicit out-of-place multiplication). Added a defensive
\fn{.copy()} at the entry to \fn{analyze\_mode\_with\_fallback}.

\paragraph{Regression test.}
\fn{tests/test\_v104\_audit\_fixes.py::test\_evec\_raw\_preserved\_after\_scaling}.

\paragraph{Migration impact.}
None for numerical outputs. Affects only direct consumers of the
result-dict's \texttt{evec\_raw} field (which were always supposed
to read the unscaled eigenvector).

% -------------------------------------------------------------------------
\section{FUND 2: scan-cache key insensitive to file mtime}
\label{sec:fund-2}

\paragraph{Symptom.}
Re-running an analysis after updating the Gaussian log produced
output identical to the previous run -- the new log was not being
read.

\paragraph{Root cause.}
The cache key (\cref{sec:cache}) hashed only the file path and
size, not the modification time. A user who replaced their log
with a re-run of the same length got a false cache hit.

\paragraph{Fix.}
Added \texttt{os.stat(filepath).st\_mtime} to the cache-key hash.
Bumped \fn{logio.\_CACHE\_VERSION} to invalidate all pre-existing
caches.

\paragraph{Regression test.}
\fn{tests/test\_v104\_audit\_fixes.py::test\_cache\_key\_depends\_on\_mtime}.

\paragraph{Migration impact.}
First run after upgrading to v1.0.4 re-scans every log (no cache
hit). Subsequent runs are cached as before.

% -------------------------------------------------------------------------
\section{FUND 3: synthetic-zero \texttt{evg} causes B-factor crash}
\label{sec:fund-3}

\paragraph{Symptom.}
With \cfgkey{interp\_boundary\_mode} = \texttt{"context"} and a log
whose highest real mode is well within the analysis window, the
B-factor accumulation crashed on a NoneType.

\paragraph{Root cause.}
The FUND-13 synthetic zero mode was being created with
\texttt{evg = None}, but the Phase-7 B-factor loop
unconditionally summed \texttt{r["evg"]**2}.

\paragraph{Fix.}
Two-step guard: \fn{\_make\_synthetic\_zero\_mode} omits the
\texttt{evg} key entirely, and the B-factor loop checks
\texttt{r.get("evg") is not None} before summing. Same guard
applied to the SSE-amplitude and group-amplitude loops.

\paragraph{Regression test.}
\fn{test\_v104\_audit\_fixes.py::test\_synthetic\_zero\_skipped\_by\_b\_factor}.

% -------------------------------------------------------------------------
\section{FUND 4: cluster-normal sign flip between clusters}
\label{sec:fund-4}

\paragraph{Symptom.}
In a multi-cluster system, the same physical
$\Delta r_\text{bend\_oop}$ for an identical mode had opposite
signs in different clusters.

\paragraph{Root cause.}
The SVD-based cluster-normal was returning $\Vmat_{:,3}$ with
SVD's arbitrary sign. For some clusters this happened to point
$+\zhat$, for others $-\zhat$.

\paragraph{Fix.}
After SVD, check the sign of $\Vmat_{:,3} \cdot \zhat$ and flip
$\nhat$ if negative (\cref{sec:conv-normal-sign}). Applied
uniformly per cluster.

\paragraph{Regression test.}
\fn{test\_v104\_audit\_fixes.py::test\_cluster\_normal\_sign\_convention}.

\paragraph{Migration impact.}
For multi-cluster systems with clusters whose $\nhat$ previously
pointed $-\zhat$, the sign of \texttt{bend\_oop} flips. The
magnitudes are unchanged.

% -------------------------------------------------------------------------
\section{FUND 5: His-NH skip warning fires for unprotonated His}
\label{sec:fund-5}

\paragraph{Symptom.}
\fn{UserWarning} ``His-NH not found'' fired for every unprotonated
His ligand, polluting the run-log.

\paragraph{Root cause.}
The unprotonated check in \fn{core.analyze\_his\_hn} happened
\emph{after} the warning emission.

\paragraph{Fix.}
Re-ordered: protonation check first, warning only if the H atom
is expected (his.\fn{his\_protonated} = True) but not found.

\paragraph{Regression test.}
\fn{test\_v104\_audit\_fixes.py::test\_unprotonated\_his\_no\_warning}.

% -------------------------------------------------------------------------
\section{FUND 6, 7, 8: sigma-propagation drift}
\label{sec:fund-6-8}

These three findings share a single root cause: the sigma for
geometric observables had not been thermally rescaled by $\urms$,
despite the underlying observables having that scaling applied.

\paragraph{FUND 6 -- His-NH stretch sigma.}
Pre-v1.0.4: \texttt{s\_hn\_stretch} $= \sqrt{2}\,\sigma_e$ (raw,
in eigenvector-component units). Post-v1.0.4: $\sqrt{2}\,\sigma_e
\cdot \urms$. For typical $\urms \approx 0.05$~\AA{}, the old
sigma was $\sim 20\times$ too large, forcing many genuine PCET
signals below the significance threshold.

\paragraph{FUND 7 -- HA reaction-coordinate sigma.}
Same drift as FUND 6, applied to the \texttt{s\_ha\_dr} column.

\paragraph{FUND 8 -- SCSD sigmas.}
Pre-v1.0.4: hard-coded constants $5 \times 10^{-7}$ (reference)
and $5 \times 10^{-6}$ (distortion) regardless of mode. Post-
v1.0.4: built from $\sigma_e$, $\sigma_c$, and $\urms$ per the
formulas in \cref{sec:scsd-derivation}.

\paragraph{Root cause (all three).}
Sigma-propagation was added as an afterthought to the code; the
authors did not at the time appreciate that the underlying
observables were already $\urms$-scaled. The fix is to apply
$\urms$ uniformly to all sigmas of observables derived from
$\widetilde{\evec}$.

\paragraph{Regression tests.}
\fn{test\_v104\_audit\_fixes.py::test\_sigma\_hn\_stretch\_uses\_urms},\\
\fn{test\_v104\_audit\_fixes.py::test\_sigma\_ha\_dr\_uses\_urms},\\
\fn{test\_v104\_audit\_fixes.py::test\_sigma\_scsd\_uses\_eigvec\_coord\_noise}.

\paragraph{Migration impact.}
\texttt{s\_hn\_stretch}, \texttt{s\_ha\_dr}, \texttt{s\_scsd\_*}
columns are now numerically different (typically $\sim 20\times$
smaller for the first two; $\sim 1000\times$ larger for the third
relative to the hard-coded constants). The Excel significance-
colouring is now meaningful for these columns; previously it was
either uniformly above threshold (FUND 6, 7) or uniformly below
(FUND 8). The \emph{underlying} observables in the workbook are
unchanged.

% -------------------------------------------------------------------------
\section{FUND 9: \fn{coord\_match\_tol} not applied}
\label{sec:fund-9}

\paragraph{Symptom.}
PDB atoms far from any Gaussian atom (e.g.\ in regions of the
protein not in the QM zone) were being mapped to the nearest
Gaussian atom regardless of distance, producing nonsense
residue-label assignments.

\paragraph{Root cause.}
The nearest-neighbour loop in
\fn{geometry.\_build\_pdb\_to\_center\_map} computed \texttt{best\_dist}
correctly but never checked it against the tolerance.

\paragraph{Fix.}
Apply \texttt{best\_dist <= cfg.coord\_match\_tol}; mark
\texttt{pdb\_to\_center[idx] = None} for over-tolerance pairs.

\paragraph{Regression test.}
\fn{test\_v104\_audit\_fixes.py::test\_coord\_match\_tol\_filters\_distant\_atoms}.

% -------------------------------------------------------------------------
\section{FUND 10: \fn{ChannelGeometry.weight} ignored in
                 single-sub-channel parents}
\label{sec:fund-10}

\paragraph{Symptom.}
For a parent channel with a single sub-channel (e.g.\ NH for a
single-His system), the RSS aggregator returned $|\Delta r|$
instead of $\sqrt{w}|\Delta r|$.

\paragraph{Root cause.}
Fast-path for $|s| = 1$ in
\fn{reorganization.aggregate\_by\_parent} short-circuited the
weight multiplication.

\paragraph{Fix.}
Removed the fast-path; the general formula
$\sqrt{\sum w_s (\Delta r_s)^2}$ now applies for all
$|s| \geq 1$.

\paragraph{Regression test.}
\fn{test\_v104\_audit\_fixes.py::test\_weight\_applied\_single\_subchannel}.

\paragraph{Migration impact.}
Only affects channels with $w \neq 1$. In practice this is HA only
(gaussian distance weights). The numerical impact is small
(\mbox{$<$}~3\%) but the formula is now consistent across all
channel cardinalities.

% -------------------------------------------------------------------------
\section{FUND 11: Excel mode-overview omits \texttt{u\_rms\_a}}
\label{sec:fund-11}

\paragraph{Symptom.}
The \texttt{Modes\_overview} sheet did not show the per-mode
thermal amplitude $\urms$, making it hard to relate the sigma
values to the noise model.

\paragraph{Fix.}
Added \texttt{u\_rms\_a} column to \texttt{Modes\_overview} export.

\paragraph{Regression test.}
\fn{test\_v104\_audit\_fixes.py::test\_overview\_includes\_urms}.

% -------------------------------------------------------------------------
\section{FUND 12: interpolation boundary anchor}
\label{sec:fund-12}

\paragraph{Symptom.}
For an analysis window whose right edge falls in a sparsely
sampled region of the DFT spectrum, the interpolation workbook
showed an abrupt drop to zero at the right edge -- the
piecewise-linear curve had no anchor beyond the last mode in
the window.

\paragraph{Root cause.}
The context-mode loader looked only for modes within
$\cfgkey{interp\_context\_cm1}$ of the boundary. When the
nearest mode beyond the window was farther than that, the
fallback ``load a single anchor mode'' was missing.

\paragraph{Fix.}
Added the FUND-12 fallback (\cref{sec:algo-context} Scenario B):
when no real mode lies within \cfgkey{interp\_context\_cm1} of
the upper edge, load \emph{exactly one} real mode -- the closest
above the upper edge -- as the boundary anchor.

\paragraph{Regression test.}
\fn{test\_v104\_audit\_fixes.py::test\_fund\_12\_single\_anchor\_fallback}.

\paragraph{Migration impact.}
Interpolation workbook gains anchor points for affected windows.
The interpolated $\Lambda_X(\omega)$ curves no longer have the
spurious right-edge drop.

% -------------------------------------------------------------------------
\section{FUND 13: synthetic zero anchor}
\label{sec:fund-13}

\paragraph{Symptom.}
For analysis windows whose right edge coincides with the highest
frequency in the DFT log (no modes above), even the FUND-12
fallback can't help -- there is genuinely no real mode to load.
The right-edge of the interpolation remained unanchored.

\paragraph{Root cause.}
No fallback existed for the case of an absolutely terminal
spectrum.

\paragraph{Fix.}
Introduced \fn{runner.\_make\_synthetic\_zero\_mode}: when no
real mode exists above \cfgkey{freq\_max}, inject a synthetic
mode at $\cfgkey{freq\_max} + \cfgkey{interp\_context\_cm1}$ with
all observables set to zero (the physically expected behaviour
of a non-existent mode). Sentinel values
$\texttt{number} = -1$, $\texttt{mode\_type} =
\texttt{"synthetic\_zero"}$, $\texttt{precision} =
\texttt{"synthetic"}$ identify it. The synthetic mode contributes
to interpolation anchors but \emph{not} to system-wide
aggregates (B-factors, $\Lambda_X^\text{total}$).

\paragraph{Regression tests.}
\fn{test\_v104\_audit\_fixes.py::test\_fund\_13\_synthetic\_zero\_inserted\_at\_correct\_freq},\\
\fn{test\_v104\_audit\_fixes.py::test\_fund\_13\_synthetic\_zero\_skipped\_by\_aggregates},\\
\fn{test\_v104\_audit\_fixes.py::test\_fund\_13\_no\_synthetic\_at\_lower\_edge}.

\paragraph{Migration impact.}
For terminal-spectrum windows the interpolation tail is now
well-defined. Headline $\Lambda_X^\text{total}$ values are
unchanged (synthetic mode contributes zero).

% -------------------------------------------------------------------------
\section{FUND 14: doubled \texttt{His\_HN NICHT erkannt} warning}
\label{sec:fund-14}

\paragraph{Symptom.}
On a real-world [2Fe-2S]-Cys$_2$His$_2$ run (May 2026) where both
coordinating histidines were deprotonated, the
``\texttt{His\_HN NICHT erkannt: His <num>}'' warning appeared
\emph{twice} per ligand in the run log instead of once. The doubled
output suggested a duplicate analysis was running, which was
alarming on initial inspection. The same bug would trigger on any
system whose histidine-N coordinating ligand carries no detectable
H atom (Rieske-type Cys$_2$His$_2$ in its deprotonated state, or
mitoNEET-type Cys$_3$His after deprotonation).

\paragraph{Root cause.}
In \fn{geometry.\_add\_his\_hn\_info}, the
``\texttt{if not lig.his\_protonated:}'' warning emission was
placed \emph{inside} the outer
``\texttt{for use\_pdb\_only in (True, False):}'' loop. The outer
loop runs twice (PDB-only path, then Gaussian-fallback path); for
a deprotonated His, both paths fail to find an H atom, so the
unsuccessful warning fired in each iteration -- producing two
identical warning lines.

The dedup set \fn{\_WARNED\_HN\_SKIP} was not involved, because
the warning was a one-shot \fn{print}/\fn{runlog.warn} rather
than a deduplicated \fn{UserWarning}; the dedup set guards a
different code path.

\paragraph{Fix.}
Hoist the ``\texttt{if not lig.his\_protonated:}'' block out of
the outer loop, so it is executed exactly once per ligand after
both paths have been attempted.

\paragraph{Regression test.}
\fn{test\_v104\_audit\_fixes.py::test\_his\_hn\_warning\_emitted\_only\_once\_per\_ligand}
constructs a synthetic His ligand with no nearby H, runs
\fn{\_add\_his\_hn\_info}, and asserts that exactly one
\texttt{"His\_HN NICHT"}-warning is emitted, not two.

\paragraph{Migration impact.}
Numerical output unchanged. Run-log noise reduced by 50\% in the
specific case of deprotonated His ligands.

% -------------------------------------------------------------------------
\section{FUND 15: misleading PCET-disabled message for deprotonated His}
\label{sec:fund-15}

\paragraph{Symptom.}
On the deprotonated-His system from the FUND 14 run, the
PCET-decision line in the run log stated:
\begin{quote}
\texttt{PCET reorg: NOT active (no His ligands at cluster).}
\end{quote}
But the same run produced \texttt{Lambda\_FeN} $> 0$ and
\texttt{Fe\_N\_His} sheets populated for both histidines --
clearly His ligands were coordinating Fe. The PCET-decision
text was actively misleading. The same bug would mislead any
analysis of a [2Fe-2S]-Cys/His system where all histidines
happen to be deprotonated.

\paragraph{Root cause.}
In \fn{runner.\_run\_analysis\_single}, the PCET-disabled branch
selected its message based on
\fn{coord\_info.pcet\_info.n\_his == 0}, where \fn{n\_his} (from
\fn{pcet\_et.build\_pcet\_info}) counts only \emph{protonated} His
ligands (\fn{pcet\_et.py}, line 243-244, skips non-protonated
His). For a deprot system with His ligands but no detectable HN
bond, \fn{n\_his} is zero, so the ``no His ligands at cluster''
message fires -- even though His ligands are clearly present.

\paragraph{Fix.}
In the PCET-disabled branch, count protein-total His ligands
directly via \texttt{coord\_info.ligands} (independent of
protonation status) and use the count to differentiate the two
disjoint cases:
\begin{itemize}
  \item ``no His ligands at cluster'' (true Cys$_4$ system),
  \item ``$N$ His ligand(s) found but none protonated''
        (deprotonated state).
\end{itemize}

\paragraph{Regression test.}
\fn{test\_v104\_audit\_fixes.py::test\_pcet\_decision\_distinguishes\_no\_his\_from\_deprotonated\_his}
constructs two synthetic \fn{LigandInfo} lists (Cys$_4$ only,
Cys$_2$His$_2$) and verifies the counting logic.

\paragraph{Migration impact.}
Numerical output unchanged. Reviewer correctly informed about
the actual reason PCET is disabled.

% -------------------------------------------------------------------------
\section{FUND 16: spurious \texttt{no context modes} warning at full-spectrum runs}
\label{sec:fund-16}

\paragraph{Symptom.}
All systems in the same batch run emitted the warning
\begin{quote}
\texttt{interp\_boundary\_mode='context' but no context modes
available -- boundary values set to 0.}
\end{quote}
even though the runs had not set \cfgkey{freq\_min} or
\cfgkey{freq\_max} (full-spectrum analysis). For a full-spectrum
run there is no interpolation boundary that requires anchoring --
the natural spectrum end serves as the upper limit and
\fn{np.interp}'s right-extrapolation already returns the value
at the highest mode.

\paragraph{Root cause.}
In \fn{export.export\_interpolated\_excel}, the warning condition
was
\begin{verbatim}
if bmode == "context" and not (_ctx_l or _ctx_r):
    runlog.warn(...)
\end{verbatim}
without checking whether the user actually requested a window.
The context-loading code in \fn{runner.\_run\_analysis\_single}
silently skips its body when \cfgkey{freq\_max} is None (line
988), so $\_ctx\_r$ is empty for full-spectrum runs too -- and
the export-stage warning fires inappropriately.

\paragraph{Fix.}
Guard the warning with an additional check that the user has
explicitly set at least one window parameter:
\begin{verbatim}
_user_set_window = (cfg.freq_min is not None
                    or cfg.freq_max is not None
                    or cfg.freq_windows is not None)
if bmode == "context" and not (_ctx_l or _ctx_r)
        and _user_set_window:
    runlog.warn(...)
\end{verbatim}

\paragraph{Regression test.}
\fn{test\_v104\_audit\_fixes.py::test\_interp\_no\_context\_warning\_suppressed\_when\_no\_window}
constructs a default Config, verifies that the
\texttt{\_user\_set\_window} flag is False, and then verifies it
flips to True after setting \cfgkey{freq\_max}.

\paragraph{Migration impact.}
Numerical output unchanged. Spurious warning silenced for
full-spectrum runs. The warning still fires correctly when the
user actually sets a window with no context modes available.

% -------------------------------------------------------------------------
\section{FUND 17: REPORT.txt missing \cfgkey{freq\_windows}}
\label{sec:fund-17}

\paragraph{Symptom.}
The batch run used \cfgkey{freq\_windows} =
\texttt{[(0,100), (100,300), (300,500)]} but the resulting
REPORT.txt showed
\begin{quote}
\texttt{freq\_filter: - - - cm-1}
\end{quote}
making the active frequency filter invisible in the run audit
trail.

\paragraph{Root cause.}
In \texttt{RunLog.write\_befund} (the report-writer method in
\file{logio.py}), the report writer only formatted
\cfgkey{freq\_min} and \cfgkey{freq\_max} (both None
for a windows-only run); \cfgkey{freq\_windows} was never
emitted. The user had to read the original Config to learn what
windows were active.

\paragraph{Fix.}
After the \texttt{freq\_filter} line, emit an additional
\texttt{freq\_windows: [...]} line whenever
\fn{cfg.freq\_windows} is non-empty.

\paragraph{Regression test.}
\fn{test\_v104\_audit\_fixes.py::test\_report\_includes\_freq\_windows}
verifies the string-format produced by the report writer.

\paragraph{Migration impact.}
Numerical output unchanged. REPORT.txt audit trail now fully
records which windows were active.

% -------------------------------------------------------------------------
\section{Manual drift findings}
\label{sec:manual-drift}

In addition to the code-level FUND findings, an audit of the
v1.0.3 manual identified 9 drift items where the manual's
description had diverged from the code:

\begin{enumerate}
  \item ``OOP fraction is kinetic-energy fraction'' -- the manual
        claimed mass-weighted aggregation. Reality: geometric-
        cartesian (\cref{sec:oop-cartesian-cavet}). Fixed.
  \item ``HA reaction coordinate is $(\evec_H - \evec_A) \cdot
        \nhat_{HA}$'' -- wrong. Reality:
        $(\evec_H - \evec_N) \cdot \nhat_{NA}$
        (\cref{eq:dr-ha-final}). Fixed with derivation.
  \item ``$\sigma$ scaling is constant'' -- did not flag the
        FUND 6/7/8 fix. Fixed.
  \item ``window boundaries are closed'' -- conflicts with code's
        half-open convention. Fixed.
  \item ``cluster normal points away from the protein bulk'' --
        not enforced anywhere. Reality: $+\zhat$ half-space of the
        Gaussian frame (\cref{sec:cluster-normal-sign}). Fixed.
  \item ``L1 sum aggregation'' -- wrong. Reality: RSS
        (\cref{sec:rss-aggregation}). Fixed with consistency
        proof.
  \item ``\fn{interp\_boundary\_mode} = \texttt{"context"} works
        for all windows'' -- did not document the FUND 12/13
        fallbacks. Fixed.
  \item ``$\alpha_{\OOP}$ is always in $[0\%, 100\%]$'' -- the
        per-atom decomposition is, but aggregates over atom groups
        can exceed $100\%$ if displacements differ in sign. Fixed.
  \item Missing \texttt{Coordination} sheet documentation. Fixed.
\end{enumerate}

% -------------------------------------------------------------------------
\section{Migration notes for v1.0.4}
\label{sec:migration}

\paragraph{Columns with numerical differences (output-affecting):}
\begin{itemize}
  \item \texttt{s\_hn\_stretch\_*}: $\sim 20\times$ smaller
        (FUND 6).
  \item \texttt{s\_ha\_dr\_*}: $\sim 20\times$ smaller (FUND 7).
  \item \texttt{s\_scsd\_*}: $\sim 1000\times$ larger
        (FUND 8) -- previously implausibly small.
  \item \texttt{bend\_oop} in some multi-cluster systems: sign
        flipped (FUND 4). Magnitude unchanged.
  \item Interpolation workbook tails: now anchored
        (FUND 12, 13). Headline totals unchanged.
\end{itemize}

\paragraph{Columns with no numerical difference but improved
documentation:}
\begin{itemize}
  \item All RSS-aggregated $\Delta r$ columns -- now derived
        rigorously from energetic-consistency arguments
        (\cref{sec:rss-aggregation}).
  \item All \texttt{a\_oop\_pct\_*} columns -- now documented as
        geometric-cartesian fraction (\cref{sec:oop-cartesian-cavet}).
  \item HA channel -- now documented as reaction-coordinate
        formula (\cref{sec:dr-ha}).
\end{itemize}

\paragraph{No breaking config changes.}
The TOML schema is backward-compatible. All existing v1.0.3
configs run unchanged under v1.0.4.

% =========================================================================
% PART IX -- TEST SUITE
% =========================================================================
\part{Test Suite}
\label{sec:test-suite}

This part documents the test suite of v1.0.4: 145 unit tests + 4
end-to-end (E2E) tests across four categories. The total count
makes the package's regression coverage roughly proportional to
its line count (one test per 87 lines of source on average).

% -------------------------------------------------------------------------
\section{Overview: 149 tests across four categories}

\begin{longtable}{p{0.32\textwidth} p{0.10\textwidth} p{0.50\textwidth}}
\toprule
\textbf{Suite} & \textbf{Tests} & \textbf{Purpose} \\
\midrule
\fn{test\_v104\_audit\_fixes.py} & 19
  & Regression guards for the 13 FUND findings + 3 FUND-13 synthetic-zero
    sub-tests + manual-drift sanity. \\
\fn{test\_sanity\_physics.py} & 19
  & Physical property tests on $\urms$, OOP/INP, $\Delta r$,
    $\lambda$, RSS -- run against the implementation without
    requiring a Gaussian log. \\
\fn{test\_manual\_code\_refs.py} & 5
  & Cross-validation that every \codref{} in Manual.tex and
    Supplement.tex points to a real function in the source. \\
\fn{test\_geometry.py}, \fn{test\_logio.py}, etc.
  & 102 & The legacy unit suites for individual modules. \\
\fn{e2e/} & 4 & End-to-end smoke tests on the example
                Cys$_2$His$_2$ (Apd1-WT) and Cys$_4$
                (mitoNEET-H87C) benchmark systems. \\
\midrule
\textbf{Total} & \textbf{154} & \\
\bottomrule
\end{longtable}

% -------------------------------------------------------------------------
\section{Regression suite: \fn{test\_v104\_audit\_fixes.py}}
\label{sec:regression-tests}

19 tests, one per FUND finding plus three FUND-13 sub-tests. Each
test is named after the finding it guards
(\texttt{test\_fund\_\#\#\_\textit{description}}). The tests
construct minimal synthetic inputs (no Gaussian log required) and
assert the post-fix behaviour. A regression in any FUND is
detected by the corresponding test.

\subsection{Test patterns used in this suite}

\begin{itemize}
  \item \emph{Round-trip}: construct a fake mode, run the full
        \fn{analyze\_mode\_with\_fallback}, assert the result-dict
        keys and values match the expected post-fix output.
  \item \emph{Property-based}: assert that the post-fix formula
        produces the expected limit value (e.g.\ FUND 13 synthetic
        zero must have \texttt{number} = $-1$).
  \item \emph{State-isolation}: assert that the per-cluster reset
        cleans up between cluster runs (FUND 5 dedup state).
\end{itemize}

\subsection{Per-test specifications}

The following table specifies, for each test, what it constructs,
what it asserts, and which FUND finding it guards. The
construction column shows the key synthetic inputs; the assertion
column shows the threshold or invariant; the FUND column refers to
the audit-trail section.

\begin{longtable}{p{0.30\textwidth} p{0.30\textwidth} p{0.30\textwidth}}
\toprule
\textbf{Test} & \textbf{Construction \& assertion} & \textbf{Protects} \\
\midrule
\texttt{test\_his\_hn\_sigma\_scales\_with\_u\_rms}
 & 1-His-ligand stub, $\evec_H = (0.1, 0, 0)$, run
   \fn{analyze\_his\_hn} once with $\urms = 1.0$ and once with
   $\urms = 0.05$. Assert
   $\sigma_1 = \sqrt{2}\cdot 5\!\times\!10^{-4}$ and
   $\sigma_2 = \sqrt{2}\cdot 5\!\times\!10^{-4}\cdot 0.05$;
   ratio $\sigma_2/\sigma_1 = 0.05$ exact to $10^{-10}$.
 & FUND 6 (\cref{sec:fund-6-8}) -- $\sigma$ should scale with
   $\urms$. Pre-fix the ratio was $1.0$ instead of $0.05$. \\
\midrule
\texttt{test\_his\_hn\_value\_and\_sigma\_have\_same\_scaling}
 & Same stub with $\evec$ unchanged across two $\urms$ values
   ($1.0$ and $0.05$). Assert
   $|\Delta r_\text{NH}^{(1)} / \Delta r_\text{NH}^{(2)}|
    = \sigma_1 / \sigma_2$. Verifies value and sigma scale
   \emph{together}.
 & FUND 6 -- consistency of value-and-sigma propagation. \\
\midrule
\texttt{test\_sse\_element\_sigma\_responds\_to\_sigma\_eigvec}
 & Construct a small SSE-element with 3 atoms, run the export-
   amplitude function with \cfgkey{sigma\_eigvec} = $10^{-4}$ and
   then $5\!\times\!10^{-4}$. Assert sigma ratio = 5.
 & FUND 7 -- SSE-element sigmas should be user-configurable. \\
\midrule
\texttt{test\_scsd\_sigmas\_use\_cfg}
 & Set \cfgkey{sigma\_eigvec} = 1e-3, \cfgkey{sigma\_coord} = 2e-3;
   compute SCSD for a synthetic 4-atom cluster; assert
   $\sigma_\text{SCSD,ref} = \sqrt{2}\cdot 2\!\times\!10^{-3}$ and
   $\sigma_\text{SCSD,dist}$ matches \cref{sec:scsd-derivation}.
 & FUND 8 -- SCSD sigmas built from Config (not hard-coded). \\
\midrule
\texttt{test\_analyze\_fe\_ligand\_warns\_on\_missing\_fe\_center}
 & Build a LigandInfo whose \texttt{fe\_center} doesn't appear in
   the c2l mapping. Run \fn{analyze\_fe\_ligand} and capture the
   warning. Assert exactly one \fn{UserWarning} with the substring
   ``Fe center'' is emitted.
 & FUND 5 -- warning is correctly raised on missing center. \\
\midrule
\texttt{test\_analyze\_fe\_ligand\_warning\_deduplicated}
 & Same stub, call the function twice. Assert that only
   \emph{one} warning is emitted across both calls (the second
   call should be silent thanks to the
   \fn{\_WARNED\_EMPTY\_LIG} dedup set).
 & FUND 5 -- warning deduplication via module-level set. \\
\midrule
\texttt{test\_analyze\_his\_hn\_warns\_on\_protonated\_lookup\_failure}
 & Protonated His ligand whose \texttt{his\_hn\_center} is not in
   c2l. Assert exactly one warning emitted with substring ``His-NH''.
 & FUND 5 -- warning fires only for protonated case. \\
\midrule
\texttt{test\_analyze\_his\_hn\_does\_not\_warn\_for\_deprotonated}
 & Same stub but \texttt{his\_protonated = False}. Assert
   \emph{no} warning emitted (the function early-returns before
   the lookup).
 & FUND 5 -- warning is gated on protonation status. \\
\midrule
\texttt{test\_ws\_fe\_bindung\_warns\_on\_all\_zero\_ligand}
 & Construct a group whose \texttt{c2l} maps map all atoms to
   indices that lie outside the eigvec array (degenerate group).
   Assert warning emitted.
 & FUND 5 -- empty-group warning. \\
\midrule
\texttt{test\_ws\_fe\_bindung\_no\_warning\_when\_data\_present}
 & Valid group with 2 atoms inside the eigvec array. Assert no
   warning.
 & FUND 5 -- false-positive guard for empty-group warning. \\
\midrule
\texttt{test\_ws\_his\_hn\_warns\_only\_for\_protonated\_zero}
 & Two His ligands: one protonated, one not. Both have
   $\evec_H = 0$. Assert exactly one warning, fired for the
   protonated case only.
 & FUND 5 -- selective warning based on protonation. \\
\midrule
\texttt{test\_torsion\_loop\_robust\_to\_c2l\_gaps}
 & Build a torsion definition $(i, j, k, l)$ where atom $k$ is
   missing from c2l. Run the analysis loop; assert it does not
   raise, returns NaN for the affected torsion, and emits one
   warning.
 & General robustness against malformed inputs. \\
\midrule
\texttt{test\_build\_pcet\_info\_warns\_on\_missing\_h\_index}
 & PCET donor with $\texttt{h\_center}$ not in c2l. Assert warning
   substring ``H index'' present.
 & FUND 5 + PCET path. \\
\midrule
\texttt{test\_pdb\_matching\_warns\_on\_high\_ambiguity}
 & Two PDB atoms within $\cfgkey{coord\_match\_tol} = 1.0$~\AA{} of
   the same Gaussian atom. Assert warning fired (``ambiguous
   match'').
 & FUND 9 -- the tolerance check. \\
\midrule
\texttt{test\_pdb\_matching\_no\_warning\_for\_low\_ambiguity}
 & PDB atoms each at distinct Gaussian atoms. Assert no warning.
 & FUND 9 -- false-positive guard. \\
\midrule
\texttt{test\_window\_boundaries\_are\_half\_open}
 & Build synthetic blocks with modes at exactly $\omega = 100$
   and $\omega = 200$ cm$^{-1}$. Run with $\cfgkey{freq\_min} = 100$,
   $\cfgkey{freq\_max} = 200$. Assert mode-100 is included,
   mode-200 is excluded.
 & General half-open-interval convention. \\
\midrule
\texttt{test\_synthetic\_zero\_mode\_has\_zero\_observables}
 & Call \fn{runner.\_make\_synthetic\_zero\_mode(omega=410.0,
   n\_atoms=10)}. Assert every observable
   (\texttt{a\_oop\_pct}, \texttt{a\_inp\_pct}, all
   \texttt{dr\_rss\_*}, all \texttt{lambda\_*}) is exactly $0.0$.
 & FUND 13 (\cref{sec:fund-13}) -- synthetic mode is observable-free. \\
\midrule
\texttt{test\_synthetic\_zero\_mode\_marked\_as\_synthetic}
 & Call same function. Assert \texttt{number} = $-1$,
   \texttt{mode\_type} = \texttt{"synthetic\_zero"},
   \texttt{precision} = \texttt{"synthetic"}.
 & FUND 13 -- sentinel values. \\
\midrule
\texttt{test\_synthetic\_zero\_mode\_safe\_in\_interp}
 & Run the full interpolation export with a mix of real modes
   and one synthetic. Assert the interp curves are continuous and
   the synthetic mode anchors the right edge at zero.
 & FUND 13 -- end-to-end safety. \\
\bottomrule
\end{longtable}

\subsection{Patterns of synthetic input construction}

A typical regression-test setup constructs three things:

\begin{enumerate}
  \item A \emph{minimal coord\_info stub}: a Python class with only
        the attributes needed by the function under test
        (e.g.\ \texttt{ligands}, \texttt{group\_map},
        \texttt{atoms\_h}). The stub does not implement
        \fn{CoordInfo} fully -- it duck-types only enough to satisfy
        the function's accesses.
  \item A \emph{c2l mapping}: a Python \texttt{dict} from Gaussian
        centre (1-based) to position in the eigenvector array.
  \item A \emph{synthetic eigenvector matrix}: a small NumPy array
        of shape $(N_\text{atoms}, 3)$ that encodes the
        property under test (e.g.\ pure stretch, pure translation,
        zero everywhere except one atom).
\end{enumerate}

This pattern decouples the test from the full pipeline: no log
file is read, no PDB matching is performed, no I/O happens. The
test runs in milliseconds and exercises only the function under
test.

\subsection{Pytest collection}

All 19 tests in this file are picked up by the default
\fn{pytest} test discovery (function names starting with
\texttt{test\_}). The file is in the \fn{tests/} directory and is
included by \fn{pyproject.toml}'s \fn{tool.pytest.ini\_options}
testpath. The full file is run by:

\begin{lstlisting}[style=pycode]
pytest tests/test_v104_audit_fixes.py -v
\end{lstlisting}

Runtime: $\sim$~1.5 seconds for all 19 tests.

% -------------------------------------------------------------------------
\section{Sanity-physics suite: \fn{test\_sanity\_physics.py}}

\codref{core.compute\_thermal\_amplitude, core.classify\_oop\_inp,
        reorganization.signed\_dr\_along\_axis,
        reorganization.lambda\_pair\_cm1,
        reorganization.aggregate\_by\_parent}

19 tests on the physical properties of the core formulas. These
tests run against the implementation, without using any Gaussian
log, by constructing synthetic eigenvectors and geometries that
exhibit the property under test. Unlike the regression suite
above (which protects against bugs that already happened), these
tests protect against \emph{future} formula refactors that would
silently break a physical property.

\subsection{Per-test specifications: OOP/INP physics}

\begin{longtable}{p{0.34\textwidth} p{0.30\textwidth} p{0.27\textwidth}}
\toprule
\textbf{Test} & \textbf{Construction} & \textbf{Assertion} \\
\midrule
\texttt{test\_pure\_oop\_mode\_yields\_alpha\_oop\_100}
 & Cluster $\nhat = (0,0,1)$, all 4 atoms with $\evec_i = (0, 0, +1)$.
 & $\alpha_{\OOP} = 100.00\%$ exact. \\
\midrule
\texttt{test\_pure\_inp\_mode\_yields\_alpha\_oop\_0}
 & Cluster $\nhat = (0,0,1)$, all 4 atoms with
   $\evec_i = (1, 0, 0)$ (in-plane breathing).
 & $\alpha_{\OOP} = 0.00\%$ exact. \\
\midrule
\texttt{test\_mixed\_oop\_inp\_satisfies\_pythagoras}
 & Random eigenvector (seed=42), $\nhat = (0,0,1)$.
 & $\alpha_{\OOP} + \alpha_{\INP} = 1.0$ to $< 10^{-12}$
   absolute error. \\
\bottomrule
\end{longtable}

\subsection{Per-test specifications: bond modulation physics}

\begin{longtable}{p{0.40\textwidth} p{0.27\textwidth} p{0.24\textwidth}}
\toprule
\textbf{Test} & \textbf{Construction} & \textbf{Assertion} \\
\midrule
\texttt{test\_pure\_translation\_mode\_yields\_zero\_bond\_modulation}
 & Atoms $a, b$ at arbitrary positions; $\evec_a = \evec_b
   = (1.0, 0.5, 0)$ (pure translation).
 & $\Delta r_{ab} = 0$ exact. \\
\midrule
\texttt{test\_pure\_rotation\_mode\_yields\_zero\_bond\_stretch}
 & Atoms aligned along $\xhat$; $\evec_a = (0, +\delta, 0)$,
   $\evec_b = (0, -\delta, 0)$ (rigid in-plane rotation).
 & $\Delta r_{ab} = 0$ to first order. \\
\midrule
\texttt{test\_pure\_bond\_stretching\_mode\_yields\_full\_modulation}
 & Atoms aligned along $\xhat$, separation $2$~\AA{};
   $\evec_a = (+\Delta, 0, 0)$, $\evec_b = (-\Delta, 0, 0)$
   with $\Delta = 0.05$.
 & $\Delta r_{ab} = +0.10$~\AA{} exact. \\
\bottomrule
\end{longtable}

\subsection{Per-test specifications: $\urms$ physics}

\begin{longtable}{p{0.30\textwidth} p{0.32\textwidth} p{0.30\textwidth}}
\toprule
\textbf{Test} & \textbf{Construction} & \textbf{Assertion} \\
\midrule
\texttt{test\_urms\_zero\_T\_limit}
 & $\omega = 500$~cm$^{-1}$, $m = 30$~amu, $T = 10^{-3}$~K.
 & $\urms$ matches $\sqrt{\hbar/(2 m\omega)} = 0.0327$~\AA{}
   to $< 10^{-6}$ relative error. \\
\midrule
\texttt{test\_urms\_classical\_limit}
 & $\omega = 50$~cm$^{-1}$, $m = 30$~amu, $T = 2000$~K
   ($\hbar\omega/\kB T = 0.036 \ll 1$).
 & $\urms$ matches $\sqrt{\kB T/(m\omega^2)}$ to $< 1\%$ relative
   error. \\
\midrule
\texttt{test\_urms\_monotonic\_in\_T}
 & $\omega = 100$~cm$^{-1}$, $m = 30$~amu. Compute $\urms(T)$ at
   $T = 1, 10, 100, 300, 1000, 3000$~K.
 & Strictly increasing in $T$ (assert
   $\urms(T_i) < \urms(T_{i+1})$ for all $i$). \\
\bottomrule
\end{longtable}

\subsection{Per-test specifications: HA reaction-coordinate physics}

These three tests jointly verify the PT/skeletal discriminator
property (\cref{sec:dr-ha}).

\begin{longtable}{p{0.38\textwidth} p{0.30\textwidth} p{0.27\textwidth}}
\toprule
\textbf{Test} & \textbf{Construction} & \textbf{Assertion} \\
\midrule
\texttt{test\_ha\_pure\_acceptor\_motion\_yields\_zero}
 & N-H aligned along $\zhat$ at $(0,0,0)$ and $(0,0,1)$,
   acceptor A at $(0,0,3)$; $\evec_N = \evec_H = 0$,
   $\evec_A = (0.05, 0, 0)$.
 & $\Delta r_{HA} = 0$ exact (no PT character). \\
\midrule
\texttt{test\_ha\_pure\_h\_motion\_toward\_acceptor\_yields\_positive}
 & Same geometry; $\evec_N = \evec_A = 0$,
   $\evec_H = (0, 0, +0.1)$ (H moves toward A).
 & $\Delta r_{HA} = +0.1$~\AA{} (positive sign = PT). \\
\midrule
\texttt{test\_ha\_pure\_n\_motion\_toward\_acceptor\_yields\_negative}
 & Same geometry; $\evec_H = \evec_A = 0$,
   $\evec_N = (0, 0, +0.1)$ (N moves toward A, skeletal).
 & $\Delta r_{HA} = -0.1$~\AA{} (negative sign = skeletal).
   The naive formula $(\evec_H - \evec_A)\cdot\nhat_{HA}$ would
   yield $-0.1$~\AA{} \emph{with the same sign} as the genuine PT
   case -- this test would fail for that formula. \\
\bottomrule
\end{longtable}

\subsection{Per-test specifications: $\lambda$ and RSS physics}

\begin{longtable}{p{0.38\textwidth} p{0.30\textwidth} p{0.27\textwidth}}
\toprule
\textbf{Test} & \textbf{Construction} & \textbf{Assertion} \\
\midrule
\texttt{test\_lambda\_zero\_T\_pure\_stretching}
 & Cluster pair with $\alpha_X = 1$ (pure stretch along bond),
   $\omega = 320$~cm$^{-1}$, $\mu = 31.97$~amu, $T = 10^{-4}$~K.
 & $\lambda^\text{pair} = 80.0$~cm$^{-1}$ (= $\omega/4$) to
   $< 10^{-3}$ relative error
   (\cref{eq:lambda-zero-T-derivation}). \\
\midrule
\texttt{test\_lambda\_zero\_for\_zero\_displacement}
 & $\Delta r = 0$, any $\omega$, $\mu$.
 & $\lambda^\text{pair} = 0$ exact. \\
\midrule
\texttt{test\_lambda\_positive\_for\_any\_nonzero\_displacement}
 & $\Delta r = \pm 0.01$~\AA{}, $\omega = 200$~cm$^{-1}$,
   $\mu = 20$~amu.
 & $\lambda^\text{pair}$ identical for $+0.01$ and $-0.01$;
   both positive. \\
\midrule
\texttt{test\_lambda\_quadratic\_in\_displacement}
 & Two $\Delta r$ values $0.01, 0.02$~\AA{}, same $\omega, \mu$.
 & $\lambda(0.02) / \lambda(0.01) = 4.0$ to $< 10^{-10}$. \\
\midrule
\texttt{test\_lambda\_quadratic\_in\_frequency}
 & Two $\omega$ values $100, 200$~cm$^{-1}$, same $\Delta r, \mu$.
 & $\lambda(200) / \lambda(100) = 4.0$. \\
\midrule
\texttt{test\_rss\_aggregation\_single\_subchannel}
 & One sub-channel with $|\Delta r| = 0.012$~\AA{}, $w = 1$.
 & RSS aggregate = $0.012$~\AA{} exact. \\
\midrule
\texttt{test\_rss\_aggregation\_two\_subchannels\_consistency\_with\_lambda}
 & Sub-channels $\Delta r_1 = 0.008$, $\Delta r_2 = -0.012$~\AA{},
   $w_1 = w_2 = 1$, $\omega = 350$~cm$^{-1}$, $\mu = 31.97$~amu.
   Compute $\lambda_\text{sum} = \lambda(\Delta r_1) +
   \lambda(\Delta r_2)$ and $\lambda_\text{RSS}$.
 & $\lambda_\text{sum} = \lambda_\text{RSS}$ to $< 10^{-10}$
   relative. (Energetic consistency,
   \cref{eq:lambda-rss-consistency}.) \\
\bottomrule
\end{longtable}

\subsection{The role of sanity-physics tests}

These tests are the longest-lived part of the test suite. They do
not protect against any specific past bug; they assert
\emph{theorems} of the underlying physics. A future developer who
refactors \fn{compute\_thermal\_amplitude} into a more elegant
form, or migrates \fn{aggregate\_by\_parent} to a NumPy-vectorised
implementation, will only know they got it right if these tests
keep passing. They are the \emph{contract} between the code and
the physics; the regression suite above is the contract between
the code and yesterday.

% -------------------------------------------------------------------------
\section{Manual-code-ref suite: \fn{test\_manual\_code\_refs.py}}

5 tests verifying that every \texttt{\textbackslash codref\{module.function\}}
entry in \fn{Manual.tex} and \fn{Supplement.tex} points to a real
function in the source tree, and that the module name is a real
module under \file{src/modenanalyse\_2fe2s/}.

\subsection{Per-test specifications}

\begin{longtable}{p{0.40\textwidth} p{0.55\textwidth}}
\toprule
\textbf{Test} & \textbf{What it asserts} \\
\midrule
\texttt{test\_manual\_en\_codrefs\_exist}
 & Every \texttt{\textbackslash codref} entry in
   \file{Manual.tex} resolves to a function or class definition
   somewhere in \file{src/modenanalyse\_2fe2s/<module>.py}. Test
   fails with a list of unresolved entries. \\
\midrule
\texttt{test\_manual\_en\_codrefs\_nonempty}
 & Manual contains at least 8
   \texttt{\textbackslash codref} entries
   (current count: 14). Test fails if a future refactor of the
   Manual accidentally removes the formula-to-code annotations. \\
\midrule
\texttt{test\_manual\_en\_codrefs\_use\_known\_modules}
 & All Manual \texttt{\textbackslash codref} module names appear in
   \file{src/modenanalyse\_2fe2s/} as \texttt{.py} files. Catches
   typos like \texttt{cores.foo} (instead of \texttt{core.foo}). \\
\midrule
\texttt{test\_supplement\_codrefs\_exist}
 & Same as Manual test above, applied to \file{Supplement.tex}
   (current count: 93). Skipped if \file{Supplement.tex} is
   absent. \\
\midrule
\texttt{test\_supplement\_codrefs\_use\_known\_modules}
 & Same as Manual module test, applied to
   \file{Supplement.tex}. \\
\bottomrule
\end{longtable}

\subsection{Why this matters}

Without these tests, a refactor that renames a function (e.g.\
\texttt{logio.get\_eigvec\_orca} $\rightarrow$
\texttt{orca\_io.get\_eigvec\_orca}, which actually happened during
the v1.0.4 audit) would silently leave the manual referring to a
non-existent function. The reader of the manual would chase a
phantom reference.

The tests also act as a forcing function: every time a function is
renamed or moved, the developer is reminded -- via test failure --
that the documentation needs updating in lock-step. This is the
same principle as a type-checker: a structural error becomes
mechanically detectable rather than relying on human discipline.

\subsection{Implementation: \fn{\_extract\_codrefs} and \fn{\_function\_exists}}

The tests share two private helpers. \fn{\_extract\_codrefs(tex\_path)}
runs a regex over the TeX source to find all
\texttt{\textbackslash codref\{...\}} occurrences, splits comma-
separated entries, unescapes LaTeX \texttt{\textbackslash\_} back to
\texttt{\_}, and returns a list of \texttt{(module, function)}
pairs. LaTeX comment lines (starting with \texttt{\%}) are skipped
so the macro's own definition (which contains an example
\texttt{\textbackslash codref\{module.function\}}) is not counted.

\fn{\_function\_exists(module, name)} reads
\file{src/modenanalyse\_2fe2s/\{module\}.py} and searches with four
patterns: \texttt{def \{name\}(}, \texttt{async def \{name\}(},
\texttt{class \{name\}}, and the private-prefix variant
\texttt{def \_\{name\}(} (for cases where the manual references the
public name but the implementation uses a private prefix).

% -------------------------------------------------------------------------
\section{Legacy unit suites}

The legacy unit suites cover individual modules and predate the
v1.0.4 audit. They are kept running to maintain coverage of the
older code paths.

\begin{longtable}{p{0.32\textwidth} p{0.10\textwidth} p{0.50\textwidth}}
\toprule
\textbf{File} & \textbf{Tests} & \textbf{Coverage} \\
\midrule
\fn{test\_config.py} & $\sim$12 & TOML round-trip,
   default values, field validation. \\
\fn{test\_core\_physics.py} & $\sim$15 & Pre-v1.0.4
   physics tests; partly subsumed by
   \fn{test\_sanity\_physics.py}. \\
\fn{test\_reorganization.py} & $\sim$18 & Per-mode
   $\lambda_X$ computations on hand-crafted geometries. \\
\fn{test\_pcet\_et.py} & $\sim$8 & PCET acceptor-search
   heuristic edge cases. \\
\fn{test\_multi\_cluster.py} & $\sim$10 & Cluster discovery
   and isolation between cluster runs. \\
\fn{test\_group\_map\_indexing.py} & $\sim$6 & Group
   atom-index resolution and gaps. \\
\fn{test\_helpers\_v35.py} & $\sim$8 & Helper functions
   from the v0.3.5 era. \\
\fn{test\_ca\_umap\_and\_exports.py} & $\sim$10 & UMAP
   embedding and Excel export shape tests. \\
\fn{test\_smoke.py} & 4 & End-to-end smoke runs on the
   bundled small test log. \\
\bottomrule
\end{longtable}

The legacy suites do not have per-test documentation in this
supplement; they are tested as a group and serve as additional
regression coverage. Each test is self-documenting via its
docstring -- a reviewer interested in any specific test can read
the source directly.

% -------------------------------------------------------------------------
\section{End-to-end suite}

4 E2E tests run the full pipeline on small synthetic systems and
assert observable invariants:

\begin{longtable}{p{0.32\textwidth} p{0.62\textwidth}}
\toprule
\textbf{Test} & \textbf{What it does} \\
\midrule
\texttt{test\_e2e\_4fe\_dummy\_cys4}
 & Run pipeline on a 4-Fe + 4-Cys-S synthetic log. Verifies that
   the pipeline produces the expected workbook structure (sheet
   names, columns), executes Phase 4 over all modes, and writes
   non-empty outputs. \\
\texttt{test\_e2e\_4fe\_rieske\_cys2his2}
 & 2-Cys + 2-His ligation. Verifies that the NH and HA sheets
   populate only for the His ligands; verifies that the
   coord\_info correctly identifies both protonated histidines. \\
\texttt{test\_e2e\_two\_clusters}
 & 2 independent [2Fe-2S] clusters in one Gaussian log. Verifies
   multi-cluster dispatch, per-cluster output directories, and
   correct isolation of the per-cluster warning-dedup state. \\
\texttt{test\_e2e\_terminal\_spectrum\_fund\_13}
 & Log whose highest mode falls inside the analysis window.
   Verifies that FUND 13 synthetic-zero anchor is correctly
   inserted and that the interpolation workbook is well-defined
   at the right edge. \\
\bottomrule
\end{longtable}

E2E tests are marked with the \fn{pytest.mark.e2e} marker and are
\emph{not} run by the default \fn{pytest} invocation (because they
take several seconds each). They are run explicitly via:

\begin{lstlisting}[style=pycode]
pytest tests/ -m e2e -v
\end{lstlisting}

% =========================================================================
% PART X -- VALIDATION STRATEGY
% =========================================================================
\part{Validation Strategy}
\label{part:validation}

This part documents the validation strategy of the pipeline -- the
external (experimental and literature) data against which the
pipeline's headline numbers can be compared, the DFT setup used to
produce the benchmark inputs, and the explicit acceptance
tolerances per benchmark quantity.

\paragraph{Scope of applicability.}
The pipeline targets any [2Fe-2S] cluster with Cys/His
coordination, covering the major biological subclasses:
\begin{itemize}
  \item Ferredoxin-type Cys$_4$ coordination
        (plant-type and bacterial-type ferredoxins; the H87C
        variant of mitoNEET is a synthetic example).
  \item Rieske-type Cys$_2$His$_2$ coordination
        (Rieske of cytochrome bc$_1$ and bc$_6$f, ARH1, BphA1,
        IscR, Apd1).
  \item mitoNEET-type Cys$_3$His coordination
        (mitoNEET, NAF-1, MiNT, IscA, BolA-like).
\end{itemize}
For \emph{all} subclasses, the analysis observables (OOP/INP,
$\Delta r_X$ per channel, $\lambda_X^\text{pair}$, $M_X(\omega)$,
SCSD irreps, $B$-factors, embedding) are computed by the same
code path; only the ligand-specific channels populate or remain
empty depending on which N/S coordination is present (NH and
HA are skipped for Cys$_4$, FeN is skipped if no His-N
coordination exists, etc.).

\paragraph{Choice of benchmark systems below.}
The DFT setups and acceptance tolerances in
\cref{sec:validation-dft,sec:headline-numbers} cover three
example systems that span the three coordination classes above
(Apd1 in two protonation states, and the synthetic mitoNEET-H87C
Cys$_4$ variant). These are \emph{example validation systems},
not the only systems for which the pipeline is intended. Any
other [2Fe-2S]-Cys/His system can be analysed without
modification; for such systems, the headline numbers below
provide order-of-magnitude expectations rather than strict
acceptance bounds.

% -------------------------------------------------------------------------
\section{DFT calculation setup for benchmark logs}
\label{sec:validation-dft}

The benchmark Gaussian logs distributed with the package were
produced with the following setup. A reviewer who wishes to
regenerate them can use these settings; for other [2Fe-2S]
systems, the same setup parameters (functional, basis, charge
adjusted to ligation, broken-symmetry spin state) are a
reasonable default starting point.

\subsection{Example 1: Apd1 wild-type (Cys$_2$His$_2$, reduced [2Fe-2S]$^{1+}$, $S = 1/2$)}

\begin{longtable}{p{0.30\textwidth} p{0.65\textwidth}}
\toprule
\textbf{Setting} & \textbf{Value} \\
\midrule
Software             & Gaussian 16 Revision C.02 \\
Functional           & B3LYP-D3 (with Becke-Johnson damping) \\
Basis set, Fe        & TZVP (Ahlrichs) \\
Basis set, S         & TZVP with diffuse functions \\
Basis set, others    & 6-31G(d,p) \\
Charge               & $-2$ ([Fe$_2$S$_2$(SCys)$_4$]$^{2-}$ reduced
                       form) \\
Spin                 & broken-symmetry singlet ($S_z = 0$ with
                       antiferromagnetic Fe-Fe coupling) \\
SCF convergence      & $10^{-8}$ Ha density \\
Geometry             & QM/MM-optimised; QM region = cluster +
                       coordinating Cys side-chains + 6~\AA{} of
                       backbone; MM region = remainder of protein
                       with TIP3P water \\
Frequency option     & \texttt{Freq=hpmodes} \\
Temperature          & Thermochemistry at $T = 80$~K (cryogenic
                       NRVS condition) \\
Print orientation    & Standard orientation (default) \\
\bottomrule
\end{longtable}

The frequency output contains $3N - 6 = 1260$ modes for $N = 422$
QM atoms. The file size is approximately 200 MiB; the pipeline
processes it in $\sim$~3 minutes on a 2024-era laptop CPU.

\subsection{Example 2: Apd1 protonated ($\mu$S-protonated Cys$_2$His$_2$, $S = 0$)}

Same setup as above except:
\begin{itemize}
  \item Charge $= -1$ (one of the bridging $\mu$-S atoms carries
        an extra H$^+$).
  \item Spin: closed-shell singlet (the protonation removes the
        antiferromagnetic coupling).
  \item Geometry: as a separate QM/MM-optimised structure, since
        protonation re-organises the cluster.
\end{itemize}

\subsection{Example 3: mitoNEET H87C variant (Cys$_4$, oxidised [2Fe-2S]$^{2+}$)}

\begin{longtable}{p{0.30\textwidth} p{0.65\textwidth}}
\toprule
\textbf{Setting} & \textbf{Value} \\
\midrule
Software             & ORCA 5.0.4 (yields \texttt{.hess} file) \\
Functional           & TPSSh \\
Basis set            & def2-TZVP everywhere \\
Charge               & $-2$ ([Fe$_2$S$_2$(SCys)$_4$]$^{2-}$
                       oxidised) \\
Spin                 & broken-symmetry singlet \\
Reference            & The mutated mitoNEET structure from
                       \cite{Paddock2007}, with His87 mutated to
                       Cys in silico. \\
\bottomrule
\end{longtable}

This system uses the ORCA \fn{orca\_io.py} parser (\cref{sec:algo-read-meta}
counterpart). The pipeline path is identical from
Phase 4 onward.

% -------------------------------------------------------------------------
\section{Headline numbers per benchmark system with acceptance tolerances}
\label{sec:headline-numbers}

The following tables specify expected values \emph{and} acceptance
tolerances for our three example validation systems. A pipeline
release ``passes validation'' on a benchmark system if every entry
in its table is matched within the stated tolerance. The
tolerances are chosen to be 2-3$\sigma$ of the DFT-vs-experiment
systematic uncertainty.

\paragraph{Convention.}
All headline $\Lambda^\text{total}$ values are reported in the
\emph{pair-$\mu$ convention} (\fn{Lambda\_total\_pair\_cm1} column
in the \texttt{Reorg\_total} sheet), which uses the reduced mass of
the underlying bond ($\mu_{ab} = m_a m_b / (m_a + m_b)$, e.g.\
$\mu_\text{Fe-S} \approx 20.4$~amu) for each per-mode contribution.
This is the convention closest to the spectroscopic reorganization
energies reported in the NRVS+DFT literature
\cite{Gee2021,Mitra2017}. The pipeline additionally reports a
\emph{mode-$\mu$ convention} (\fn{Lambda\_total\_mode\_cm1} column)
using the Gaussian-printed reduced mass of the normal mode itself;
this gives values typically a factor 5--10 smaller and is useful
for cross-checks against vibrational-energy budgets but is not the
quantity meant by ``the reorganization energy of bond $X$''.

\paragraph{Use of these numbers for other systems.}
The numerical expectations below are calibrated for our three
specific benchmark inputs (analysis window $0\text{--}500$
cm$^{-1}$, $T = 5$~K, B3LYP-D3/TZVP-Ahlrichs/QM-MM, broken-symmetry
singlet) and should \emph{not} be treated as universal acceptance
criteria for arbitrary [2Fe-2S]-Cys/His systems. A user analysing
a different protein (e.g.\ a plant-type Cys$_4$ ferredoxin, the
native mitoNEET Cys$_3$His ligation, or a non-standard $\mu$-SH
state) should expect \emph{order-of-magnitude} agreement with the
appropriate coordination class, not exact match. The structural
patterns that \emph{must} hold across all systems (cluster
planarity, $\Lambda_\text{NH} = 0$ for unprotonated states,
$\alpha_{\OOP}$ near $100\%$ for true umbrella modes, etc.) are
noted explicitly in the tables.

\paragraph{Sensitivity to window and temperature.}
The $\Lambda^\text{total}$ values are window-dependent (they
accumulate only modes inside the analysis window) and weakly
temperature-dependent (via the $\coth$ factor in $\urms$). For
the $0$--$500$ cm$^{-1}$ window at $T = 5$~K the values are
zero-point-dominated and within $1\%$ of their $T = 0$ limit. A
full-spectrum analysis (window $0$ -- $\omega_\text{max}^\text{DFT}$)
would give substantially \emph{larger} values for the protonated
case (the N-H stretch at $\sim 3300$~cm$^{-1}$ contributes
strongly to $\Lambda_\text{NH}$); for the deprot/Cys$_4$ cases
the increase is smaller because no high-frequency channel-active
mode is missed.

\paragraph{$\Lambda$ in cm$^{-1}$ is an \emph{energy}, not a frequency bound.}
A common point of confusion: $\Lambda_\text{FeS}^\text{total} =
555$~cm$^{-1}$ for an analysis run with
$\cfgkey{freq\_max} = 500$~cm$^{-1}$ may look paradoxical (``how
can $\Lambda$ exceed the cut-off?''). The answer is that
$\Lambda^\text{total}$ is not a frequency at all: it is the
\emph{cumulative reorganization energy}, expressed in
spectroscopic cm$^{-1}$ units (where $1$~cm$^{-1} =
1.986 \times 10^{-23}$~J $\equiv 0.124$~meV). It is built up as
\[
  \Lambda_X^\text{total}
  \;=\; \sum_{i=1}^{N_\text{modes in window}}
        \tfrac{1}{2}\,\mu_X\,\omega_i^2\,(\Delta r_i)^2 ,
\]
i.e.\ a sum over all modes in the window. Even though no single
mode contributes more than $\sim 20$~cm$^{-1}$, the sum over
$\sim 5900$ modes easily reaches several hundred cm$^{-1}$. The
analogue is total kinetic energy in classical mechanics: if every
particle has at most $k_B T$ of kinetic energy, the system total
is $\tfrac{3}{2} N k_B T$, which can vastly exceed $k_B T$.
Concretely for Apd1\_deprot: the strongest single-mode FeS
contribution is $18.3$~cm$^{-1}$ at $433$~cm$^{-1}$, the average
is $0.09$~cm$^{-1}$, summing over $5931$ contributing modes gives
$\Lambda_\text{FeS}^\text{total} \approx 555$~cm$^{-1}$.

\subsection{Example 1: Apd1 wild-type (Cys$_2$His$_2$, deprotonated)}

Calibrated against an actual production run (Apd1\_deprot,
$3N - 6 \approx 18900$ modes, analysis window $0\text{--}500$
cm$^{-1}$, $T = 5$~K, B3LYP-D3/TZVP/QM-MM, $\sim$5958 modes inside
the window).

\begin{longtable}{p{0.32\textwidth} p{0.18\textwidth} p{0.12\textwidth} p{0.31\textwidth}}
\toprule
\textbf{Quantity (pair-$\mu$ convention)} & \textbf{Expected} & \textbf{Tol.} & \textbf{Source / rationale} \\
\midrule
$\Lambda_\text{FeS}^\text{total}$
                          & $555$~cm$^{-1}$  & $\pm 50$~cm$^{-1}$
                          & Production-run observation, summed over
                            6 Fe-S sub-bonds (2 Cys terminal + 4
                            bridging). Literature comparable values
                            $\sim 320\text{--}360$~cm$^{-1}$ from
                            \cite{Gee2021} (different DFT setup, smaller
                            cluster model). \\
$\Lambda_\text{FeFe}^\text{total}$
                          & $70$~cm$^{-1}$   & $\pm 15$~cm$^{-1}$
                          & Single Fe-Fe bond contribution. The
                            cumulative $\Lambda(\omega)$ should
                            saturate by $\omega \approx 380$~cm$^{-1}$. \\
$\Lambda_\text{FeN}^\text{total}$
                          & $48$~cm$^{-1}$   & $\pm 15$~cm$^{-1}$
                          & 2 Fe-N(His) sub-bonds, comparable
                            magnitudes ($\sim 24$ each). \\
$\Lambda_\text{NH}^\text{total}$
                          & 0                & exact
                          & Deprotonated WT has no N-H bonds. \\
$\Lambda_\text{HA}^\text{total}$
                          & 0                & exact
                          & No protonated donor $\Rightarrow$ HA
                            channel inactive. \\
$\alpha_{\OOP}$ for the cluster-umbrella mode ($\sim$~80 cm$^{-1}$)
                          & $\geq 92\%$      & --
                          & Property of the $A_u$ irrep. \\
Cluster planarity $\sigma_3/\sigma_1$
                          & $\leq 0.02$      & --
                          & Rhombic Fe$_2$S$_2$ is genuinely planar. \\
\bottomrule
\end{longtable}

\subsection{Example 2: Apd1 protonated ($\mu$S-protonated Cys$_2$His$_2$)}

Calibrated against an actual production run (Apd1\_prot,
analysis window $0\text{--}500$ cm$^{-1}$, $T = 5$~K,
B3LYP-D3/TZVP/QM-MM, $\sim$5941 modes inside the window).

\begin{longtable}{p{0.32\textwidth} p{0.18\textwidth} p{0.12\textwidth} p{0.31\textwidth}}
\toprule
\textbf{Quantity (pair-$\mu$ convention)} & \textbf{Expected} & \textbf{Tol.} & \textbf{Source / rationale} \\
\midrule
$\Lambda_\text{FeS}^\text{total}$
                          & $570$~cm$^{-1}$  & $\pm 50$~cm$^{-1}$
                          & Production-run observation, sum over
                            6 Fe-S sub-bonds. Slightly higher than
                            the deprot state (555 cm$^{-1}$) due to
                            protonation-induced bond stiffening. \\
$\Lambda_\text{FeFe}^\text{total}$
                          & $64$~cm$^{-1}$   & $\pm 15$~cm$^{-1}$
                          & Lower than deprot WT (70 cm$^{-1}$)
                            because protonation breaks the
                            antiferromagnetic Fe-Fe coupling. \\
$\Lambda_\text{FeN}^\text{total}$
                          & $58$~cm$^{-1}$   & $\pm 15$~cm$^{-1}$
                          & Higher than deprot WT (48 cm$^{-1}$),
                            consistent with stiffer Fe-N(His)
                            geometry in the protonated state. \\
$\Lambda_\text{NH}^\text{total}$ (in 0--500 cm$^{-1}$ window)
                          & $\lesssim 0.05$~cm$^{-1}$ & --
                          & The full N-H stretch sits at $\sim$3300
                            cm$^{-1}$, far outside the analysis
                            window. Only low-frequency librations
                            of the N-H vector contribute; the
                            value is therefore numerically very
                            small. \\
$\Lambda_\text{HA}^\text{total}$ (in 0--500 cm$^{-1}$ window)
                          & $10$--$25$~cm$^{-1}$ & $\pm 5$~cm$^{-1}$
                          & PCET signature: the H$\cdots$A reaction
                            coordinate modulates with frequencies
                            below 500 cm$^{-1}$, dominated by one
                            His (here His259 with
                            $\Lambda_\text{HA}\approx 11$~cm$^{-1}$
                            pair-$\mu$ or 22~cm$^{-1}$ mode-$\mu$);
                            the second His contributes
                            $< 1$~cm$^{-1}$. \\
\bottomrule
\end{longtable}

\paragraph{PCET diagnostic.}
The asymmetric distribution between the two histidines
(His259 carries $> 90\%$ of $\Lambda_\text{HA}$, His255 carries
$< 10\%$) is the expected PCET fingerprint: only one of the two
H-bond acceptor environments is geometrically favourable for the
proton transfer. The cumulative $\Lambda_\text{HA}(\omega)$ curve
should show a step in the $\sim 200\text{--}320$ cm$^{-1}$ range
where the $\mu$-S-H wagging modes are active.

\subsection{Example 3: Cys$_4$ system (e.g.\ mitoNEET H87C or Apd1 H255C/H259C)}

Calibrated against a Cys$_4$ production run (Apd1\_HH255\_259CC\_deprot,
4 Cys ligands, no His, analysis window $0\text{--}500$ cm$^{-1}$,
$T = 5$~K, B3LYP-D3/TZVP/QM-MM, $\sim$5931 modes inside the window).

\begin{longtable}{p{0.32\textwidth} p{0.18\textwidth} p{0.12\textwidth} p{0.31\textwidth}}
\toprule
\textbf{Quantity (pair-$\mu$ convention)} & \textbf{Expected} & \textbf{Tol.} & \textbf{Source / rationale} \\
\midrule
$\Lambda_\text{FeN}^\text{total}$
                          & 0                & exact
                          & Cys$_4$ has no His ligand; FeN channel
                            is structurally absent. \\
$\Lambda_\text{FeS}^\text{total}$
                          & $715$~cm$^{-1}$  & $\pm 60$~cm$^{-1}$
                          & \emph{Higher} than the Cys$_2$His$_2$
                            cases ($\sim 555\text{--}570$~cm$^{-1}$)
                            because replacing 2 Fe-N coordinations
                            with 2 additional Fe-S coordinations
                            redistributes the cluster reorganization
                            energy into the FeS channel. \\
$\Lambda_\text{FeFe}^\text{total}$
                          & $73$~cm$^{-1}$   & $\pm 15$~cm$^{-1}$
                          & Comparable to the Cys$_2$His$_2$ cases
                            (single Fe-Fe bond, ligand-environment-
                            independent to first order). \\
$\Lambda_\text{NH}^\text{total}$
                          & 0                & exact
                          & No protonated His; NH channel absent. \\
$\Lambda_\text{HA}^\text{total}$
                          & 0                & exact
                          & No PCET donor in Cys$_4$. \\
Number of detected clusters
                          & 1                & exact
                          & Single [2Fe-2S]. \\
Cluster planarity $\sigma_3/\sigma_1$
                          & $\leq 0.03$      & --
                          & Cys$_4$ systems are typically slightly
                            \emph{less} planar than Cys$_2$His$_2$
                            (folding residual $\sim 0.025$~\AA{}
                            in the production run vs $0.009$~\AA{}
                            for the WT). \\
\bottomrule
\end{longtable}

\subsection{The headline-number test (acceptance criterion)}

A pipeline release is considered to ``pass validation'' on a
given benchmark system if all numerical entries above are met
within their tolerances. The procedure is identical for each of
the three example systems; below we illustrate it for the
Cys$_2$His$_2$ wild-type case:

\begin{lstlisting}[style=pycode]
# Run pipeline on the canonical Cys2His2 deprotonated input
import modenanalyse_2fe2s as mna
result = mna.run_analysis(cfg)

# Read the headline values from the Reorg_total sheet
# (columns: 'Channel', 'Lambda_total_pair_cm1', 'Lambda_total_mode_cm1',
#  'n_modes_contributing'); each row corresponds to one channel.
import openpyxl
wb = openpyxl.load_workbook(result.main_xlsx_path)
ws = wb["Reorg_total"]
lambdas = {row[0]: row[1] for row in ws.iter_rows(min_row=2, values_only=True)}

# Check tolerances against the Cys2His2-WT (deprotonated) reference values
assert abs(lambdas["FeS"]  - 555) <= 50,  f"FeS  off: {lambdas['FeS']}"
assert abs(lambdas["FeFe"] -  70) <= 15,  f"FeFe off: {lambdas['FeFe']}"
assert abs(lambdas["FeN"]  -  48) <= 15,  f"FeN  off: {lambdas['FeN']}"
assert lambdas["NH"] == 0.0   # deprotonated WT has no N-H
assert lambdas["HA"] == 0.0   # no PCET donor
\end{lstlisting}

This procedure is run as part of the E2E test suite
(\cref{sec:test-suite}) before every release.

% -------------------------------------------------------------------------
\section{Reproducibility}

\subsection{Determinism guarantees}

With \cfgkey{n\_jobs} = 1 (default), every run produces bit-
identical output for the same Gaussian log, PDB, and Config. This
is enforced by avoiding all sources of non-determinism in the
analysis pipeline:

\begin{itemize}
  \item Python $\geq$ 3.7 preserves dict insertion order, so
        iteration over the per-mode result dict is deterministic.
  \item No hash-based set iteration is used in the analysis
        pipeline -- where a set is needed for deduplication
        (e.g.\ \fn{\_WARNED\_EMPTY\_LIG}), the set is only used
        for \fn{in}-checks, never iterated.
  \item NumPy operations are deterministic by default on a fixed
        BLAS configuration. We do not use any \fn{np.random}
        operations in the pipeline (only in the UMAP embedding,
        which has a fixed seed; see below).
  \item Floating-point summations are done in a fixed order (sum
        over modes is in ascending frequency order; sum over
        atoms is in Gaussian-centre order).
\end{itemize}

\subsection{UMAP and HDBSCAN determinism}

The UMAP embedding uses a fixed \texttt{random\_state = 42} for
reproducibility. With this seed, the umap-learn library produces
deterministic output across runs (provided the input feature
matrix is identical and the umap-learn version is identical).

\paragraph{Version sensitivity.}
UMAP's internal algorithm has changed between umap-learn versions:
0.4.x, 0.5.x, and (when released) 0.6.x will produce different
embeddings even with the same seed and inputs. We pin
\texttt{umap-learn == 0.5.5} in \file{pyproject.toml} to avoid
this. A reviewer can verify the version with:

\begin{lstlisting}[style=pycode]
import umap; print(umap.__version__)   # should be 0.5.5
\end{lstlisting}

The HDBSCAN clustering is deterministic given a fixed UMAP output,
so once UMAP is pinned, the cluster assignments are reproducible.

\subsection{Cache integrity}

Cache invalidation by mtime + file-size + cache-format version
(\cref{sec:cache}) ensures that updating the source Gaussian log
or upgrading the package version triggers a re-scan. The cache
key is robust against:

\begin{itemize}
  \item File replacement (mtime changes).
  \item File extension (the cache key is keyed by absolute path).
  \item Python major-version upgrades (cache key includes Python
        version).
  \item Cache-format changes within the package (cache key
        includes \fn{logio.\_CACHE\_VERSION}, bumped manually on
        format changes).
\end{itemize}

\subsection{Version pinning}

The exact dependency versions used to produce the example
benchmark results above are pinned in \file{pyproject.toml}:

\begin{longtable}{p{0.27\textwidth} p{0.18\textwidth} p{0.50\textwidth}}
\toprule
\textbf{Package} & \textbf{Version} & \textbf{Determinism note} \\
\midrule
\texttt{numpy}        & $\geq$~1.24  & BLAS configuration affects bit-
                                       level reproducibility across
                                       machines (Intel MKL vs OpenBLAS
                                       can differ in the last bit). \\
\texttt{scipy}        & $\geq$~1.10  & Used for SVD; same BLAS
                                       considerations apply. \\
\texttt{openpyxl}     & $\geq$~3.1   & Output format only. \\
\texttt{scikit-learn} & $\geq$~1.3   & Not used directly; pulled in
                                       by umap-learn. \\
\texttt{umap-learn}   & $=$~0.5.5    & Hard-pinned for embedding
                                       reproducibility. \\
\texttt{hdbscan}      & $=$~0.8.40   & Hard-pinned. \\
\bottomrule
\end{longtable}

\subsection{Floating-point reproducibility caveats}

Bit-identical output across machines requires the same BLAS
configuration. In practice we observe:

\begin{itemize}
  \item Linux x86-64 (Intel MKL, Anaconda): bit-identical across
        Intel and AMD CPUs of the same generation.
  \item Linux x86-64 (OpenBLAS): differs from MKL in the last bit
        of summations involving long sequences ($> 100$ terms).
  \item macOS arm64 (Apple Accelerate): differs from x86-64 in
        accumulated SVD operations beyond the last 2-3 bits.
\end{itemize}

For the validation acceptance criteria above, the tolerances are
chosen to be far larger than any BLAS-induced differences (typical
floating-point differences are at the $10^{-12}$ level; tolerances
are at the cm$^{-1}$ level, i.e.\ relative $10^{-2}$).

\subsection{Logging and audit trail}

Every run produces a \file{run-log.txt} that records:

\begin{itemize}
  \item Package version and dependency versions at run time.
  \item Path and mtime of the Gaussian log.
  \item Cache hit or miss.
  \item Number of modes scanned, selected, and analysed.
  \item All warnings emitted by the analysis (deduplication
        notwithstanding).
  \item FUND-12 and FUND-13 fallbacks if triggered.
  \item Final headline numbers (a copy of the
        \texttt{Reorganisation\_total} row).
\end{itemize}

A reviewer auditing a pipeline output should consult the run-log
first -- it contains all the information needed to reconstruct the
pipeline's decisions on a given input.

% =========================================================================
% Replace stub of Part XI -- already begun above; just add glossary
% =========================================================================

\part{References and Glossary}

\section*{Bibliography}
\addcontentsline{toc}{section}{Bibliography}

The references are grouped thematically: (i)~vibrational and
quantum-mechanical foundations, (ii)~NRVS theory and experimental
work on iron-sulfur proteins, (iii)~Marcus-Hush electron-transfer
theory, (iv)~proton-coupled electron transfer (PCET) theory and
biology, (v)~symmetry-coordinate decomposition and group theory,
(vi)~geometry algorithms and structural-biology data structures,
(vii)~dimension-reduction and clustering methods.

\begin{thebibliography}{99}

% ---------- (i) Vibrational and QM foundations ---------------------
\bibitem{WilsonDeciusCross1955}
E.~B. Wilson, J.~C. Decius, P.~C. Cross,
\emph{Molecular Vibrations: The Theory of Infrared and Raman
Vibrational Spectra}, McGraw-Hill, New York, 1955.
The canonical reference for the harmonic-oscillator basis of
vibrational spectroscopy and for the position-variance derivation
underlying \cref{eq:urms-final}.

\bibitem{LandauLifshitzQM1977}
L.~D. Landau, E.~M. Lifshitz,
\emph{Quantum Mechanics: Non-Relativistic Theory} (3rd ed.),
Pergamon Press, Oxford, 1977.
\S~23 derives $\langle q^2\rangle_T$ for the canonical-ensemble
quantum harmonic oscillator.

\bibitem{McQuarrie2000}
D.~A. McQuarrie,
\emph{Statistical Mechanics}, University Science Books, 2000.
Chapter 6 derives the harmonic-oscillator partition function and
thermal averages used in \cref{sec:urms-derivation}.

\bibitem{Wilson1980}
E.~B. Wilson, J.~C. Decius, P.~C. Cross,
\emph{Decius rules for normal-mode analysis under symmetry},
in \emph{Annual Review of Physical Chemistry} \textbf{31}, 1980.
Foundational discussion of how irrep decomposition organises
normal-mode amplitudes.

% ---------- (ii) NRVS theory and experimental work ----------------
\bibitem{Achterhold2002}
K.~Achterhold, C.~Keppler, A.~Ostermann, U.~van~B\"urck,
W.~Sturhahn, E.~E.~Alp, F.~G. Parak,
\emph{Vibrational dynamics of myoglobin determined by the phonon-
assisted M\"o\ss bauer effect},
\textbf{Phys. Rev. E} \textbf{65}, 051916 (2002).
Foundational application of NRVS to an Fe-protein system.

\bibitem{Sturhahn2004}
W.~Sturhahn,
\emph{Nuclear resonant spectroscopy},
\textbf{J. Phys.: Condens. Matter} \textbf{16}, S497--S530 (2004).
Theory of NRVS and connection to the Fe-projected partial density
of states (pDOS).

\bibitem{Kingsbury2024NRVS}
C.~J. Kingsbury, F.~Neese, M.~O. Senge,
\emph{The Fe-projected vibrational density of states from
nuclear-resonance vibrational spectroscopy and quantum chemistry},
\textbf{J. Chem. Theory Comput.} \textbf{20}, 7212--7228 (2024).
Recent benchmark study connecting computed Fe-projected pDOS to
NRVS spectra of [2Fe-2S] proteins.

\bibitem{Gee2021}
L.~B. Gee, C.~Y. Lin, F.~E. Jenney~Jr., M.~W.~W. Adams, Y.~Yoda,
R.~Masuda, T.~Mitsui, M.~Hu, J.~Zhao, J.~T. Sage, S.~P. Cramer,
\emph{Comparison of Iron Hydrogenase Vibrational Spectra to those
of [2Fe-2S] Model Compounds via NRVS and DFT},
\textbf{Inorg. Chem.} \textbf{60}, 9714--9727 (2021).
Cluster-vibration assignments in the 200--420 cm$^{-1}$ window
used as benchmark for the Cys$_2$His$_2$ wild-type validation
(\cref{part:validation}).

\bibitem{Mitra2017}
D.~Mitra, V.~Pelmenschikov, Y.~Guo, D.~A. Case, H.~Wang, W.~Dong,
M.-L. Tan, T.~Ichiye, F.~E. Jenney~Jr., M.~W.~W. Adams, Y.~Yoda,
J.~Zhao, S.~P. Cramer,
\emph{Dynamics of the [4Fe-4S] cluster in {\em Pyrococcus furiosus}
ferredoxin via NRVS and DFT},
\textbf{Biochemistry} \textbf{50}, 5220--5235 (2017).
Comprehensive NRVS/DFT comparison for iron-sulfur cluster modes.

\bibitem{Tinberg2008}
C.~E. Tinberg, R.~Y. Tonzetich, H.~Wang, L.~H. Do, Y.~Yoda,
S.~P. Cramer, S.~J. Lippard,
\emph{Characterization of iron dinitrosyl species formed in the
reaction of nitric oxide with a biological Rieske center},
\textbf{J. Am. Chem. Soc.} \textbf{130}, 14550--14551 (2008).
Rieske [2Fe-2S] benchmark.

\bibitem{Iwata1996}
S.~Iwata, M.~Saynovits, T.~A. Link, H.~Michel,
\emph{Structure of a water soluble fragment of the {`Rieske'}
iron-sulfur protein of the bovine heart mitochondrial cytochrome
bc$_1$ complex determined by {MAD} phasing at 1.5 \AA{} resolution},
\textbf{Structure} \textbf{4}, 567--579 (1996).
The structural reference for the Cys$_2$His$_2$ Rieske coordination
mode.

\bibitem{Paddock2007}
M.~L. Paddock, S.~E. Wiley, A.~T. Axelrod, A.~K. Cohen, M.~Roy,
E.~C. Abresch, D.~Capraro, A.~N. Murphy, R.~Nechushtai,
J.~E. Dixon, P.~A. Jennings,
\emph{MitoNEET is a uniquely folded 2Fe-2S outer mitochondrial
membrane protein stabilized by pioglitazone},
\textbf{Proc. Natl. Acad. Sci. USA} \textbf{104}, 14342--14347 (2007).
The Cys$_3$His Mn coordination of mitoNEET used as a benchmark
system.

% ---------- (iii) Marcus-Hush electron-transfer theory ------------
\bibitem{Marcus1956}
R.~A. Marcus,
\emph{On the theory of oxidation-reduction reactions involving
electron transfer. I.},
\textbf{J. Chem. Phys.} \textbf{24}, 966--978 (1956).
The foundational Marcus paper. Introduces the reorganization-
energy concept that \cref{eq:lambda-cm1} specialises to a single
vibrational mode.

\bibitem{Marcus1964}
R.~A. Marcus,
\emph{Chemical and electrochemical electron-transfer theory},
\textbf{Annu. Rev. Phys. Chem.} \textbf{15}, 155--196 (1964).

\bibitem{Marcus1985}
R.~A. Marcus, N.~Sutin,
\emph{Electron transfers in chemistry and biology},
\textbf{Biochim. Biophys. Acta} \textbf{811}, 265--322 (1985).
Definitive review covering the biological context of intramolecular
ET. The inner-sphere reorganization $\lambda_i$ is the closest
analogue to the spectroscopic
$\Lambda_X^\text{total}$ computed by this code.

\bibitem{Hush1961}
N.~S. Hush,
\emph{Adiabatic rate processes at electrodes},
\textbf{J. Chem. Phys.} \textbf{28}, 962--972 (1961).
Hush's analogous derivation; the Marcus-Hush nomenclature.

\bibitem{HopfieldET1974}
J.~J. Hopfield,
\emph{Electron transfer between biological molecules by thermally
activated tunneling},
\textbf{Proc. Natl. Acad. Sci. USA} \textbf{71}, 3640--3644 (1974).
Quantum vibrational version of the Marcus rate including
Huang-Rhys factors -- the link from this code's $\Lambda_X$ to
biological ET rates.

\bibitem{JortnerET1976}
J.~Jortner,
\emph{Temperature dependent activation energy for electron
transfer between biological molecules},
\textbf{J. Chem. Phys.} \textbf{64}, 4860--4867 (1976).
Theoretical extension covering the quantum regime $\hbar\omega \gg
\kB T$ relevant for low-temperature ET in iron-sulfur proteins.

% ---------- (iv) PCET theory and biology --------------------------
\bibitem{HammesSchifferPCET2010}
S.~Hammes-Schiffer, A.~A. Stuchebrukhov,
\emph{Theory of coupled electron and proton transfer reactions},
\textbf{Chem. Rev.} \textbf{110}, 6939--6960 (2010).
Comprehensive PCET-theory review. The HA reaction-coordinate
formula (\cref{eq:dr-ha-final}) is derived in the spirit of
\S~III.B.

\bibitem{CukierNocera1998}
D.~R. Cukier, D.~G. Nocera,
\emph{Proton-coupled electron transfer},
\textbf{Annu. Rev. Phys. Chem.} \textbf{49}, 337--369 (1998).
The earlier conceptual framework that separates PT and ET reaction
coordinates -- the basis for our HA-vs-FeFe co-modulation
descriptor.

\bibitem{ReeceNocera2009}
S.~Y. Reece, D.~G. Nocera,
\emph{Proton-coupled electron transfer in biology: results from
synergistic studies in natural and model systems},
\textbf{Annu. Rev. Biochem.} \textbf{78}, 673--699 (2009).
PCET in iron-sulfur and other metalloproteins.

\bibitem{HirstReview2010}
J.~Hirst,
\emph{Mitochondrial complex I},
\textbf{Annu. Rev. Biochem.} \textbf{82}, 551--575 (2013).
PCET-driven ET via [2Fe-2S] and [4Fe-4S] clusters in complex I.

% ---------- (v) Symmetry-coordinate decomposition -----------------
\bibitem{Kingsbury2024}
C.~J. Kingsbury, M.~O. Senge,
\emph{Quantifying near-symmetric molecular distortion using
symmetry-coordinate structural decomposition},
\textbf{Chem.\ Sci.} \textbf{15}, 13638--13649 (2024).
DOI: 10.1039/D4SC01670J. Web implementation:
\url{https://www.kingsbury.id.au/scsd}.
The reference for \cref{eq:scsd-projection-derivation}.

\bibitem{KingsburyPreprint2024}
C.~J. Kingsbury, M.~O. Senge,
\emph{Quantifying near-symmetric molecular distortion using
symmetry-coordinate structural decomposition},
ChemRxiv preprint, doi:10.26434/chemrxiv-2024-6b25s (2024).

\bibitem{BradleyCracknell1972}
C.~J. Bradley, A.~P. Cracknell,
\emph{The Mathematical Theory of Symmetry in Solids},
Clarendon Press, Oxford, 1972.
Group-theoretic reference for the $D_{2\mathrm{h}}$ character
table and irrep notation used in \cref{sec:scsd-derivation}.

\bibitem{Cotton1990}
F.~A. Cotton,
\emph{Chemical Applications of Group Theory} (3rd ed.),
Wiley, New York, 1990.
Practitioner-friendly reference for normal-mode irreps in
molecular spectroscopy.

% ---------- (vi) Geometry algorithms and structural biology -------
\bibitem{Kabsch1976}
W.~Kabsch,
\emph{A solution for the best rotation to relate two sets of
vectors},
\textbf{Acta Cryst.} \textbf{A32}, 922--923 (1976).
The SVD-based optimal-rotation algorithm used in
\cref{sec:kabsch}.

\bibitem{Kabsch1978}
W.~Kabsch,
\emph{A discussion of the solution for the best rotation to relate
two sets of vectors},
\textbf{Acta Cryst.} \textbf{A34}, 827--828 (1978).
Follow-up addressing the chirality sign-correction.

\bibitem{KabschSander1983}
W.~Kabsch, C.~Sander,
\emph{Dictionary of protein secondary structure: pattern
recognition of hydrogen-bonded and geometrical features},
\textbf{Biopolymers} \textbf{22}, 2577--2637 (1983).
The DSSP algorithm referenced in \cref{sec:algo-sse}.

\bibitem{Ramachandran1963}
G.~N. Ramachandran, C.~Ramakrishnan, V.~Sasisekharan,
\emph{Stereochemistry of polypeptide chain configurations},
\textbf{J. Mol. Biol.} \textbf{7}, 95--99 (1963).
The $\phi/\psi$ classification used in the SSE-detection fallback.

\bibitem{GoldenSection}
G.~H. Golub, C.~F. Van~Loan,
\emph{Matrix Computations} (4th ed.),
Johns Hopkins University Press, 2013.
Reference for the SVD algorithms used throughout.

% ---------- (vii) Dimension reduction and clustering --------------
\bibitem{Pearson1901}
K.~Pearson,
\emph{On lines and planes of closest fit to systems of points in
space},
\textbf{Philos. Mag.} \textbf{2}, 559--572 (1901).
PCA, used as a baseline reduction technique in the embedding
workbook.

\bibitem{vanDerMaaten2008}
L.~van~der~Maaten, G.~Hinton,
\emph{Visualizing data using t-SNE},
\textbf{J. Mach. Learn. Res.} \textbf{9}, 2579--2605 (2008).

\bibitem{McInnes2018}
L.~McInnes, J.~Healy, J.~Melville,
\emph{UMAP: Uniform Manifold Approximation and Projection for
Dimension Reduction}, arXiv:1802.03426 (2018).

\bibitem{Becht2019}
E.~Becht, L.~McInnes, J.~Healy, C.-A. Dutertre, I.~W.~H. Kwok,
L.~G. Ng, F.~Ginhoux, E.~W. Newell,
\emph{Dimensionality reduction for visualizing single-cell data
using UMAP},
\textbf{Nat. Biotechnol.} \textbf{37}, 38--44 (2019).

\bibitem{Campello2013}
R.~J.~G.~B. Campello, D.~Moulavi, J.~Sander,
\emph{Density-based clustering based on hierarchical density
estimates},
in \emph{Advances in Knowledge Discovery and Data Mining},
LNCS \textbf{7819}, 160--172 (2013).
The HDBSCAN algorithm.

% ---------- Software dependencies ---------------------------------
\bibitem{NumPy}
C.~R. Harris et al.,
\emph{Array programming with NumPy},
\textbf{Nature} \textbf{585}, 357--362 (2020).

\bibitem{SciPy}
P.~Virtanen et al.,
\emph{SciPy 1.0: fundamental algorithms for scientific computing
in Python},
\textbf{Nat. Methods} \textbf{17}, 261--272 (2020).

\bibitem{Matplotlib}
J.~D. Hunter,
\emph{Matplotlib: A 2D graphics environment},
\textbf{Comput. Sci. Eng.} \textbf{9}, 90--95 (2007).

\bibitem{openpyxl}
\emph{openpyxl -- A Python library to read/write Excel 2010 xlsx
files}, \url{https://openpyxl.readthedocs.io}.
The output backend used by \fn{export.py}.

\bibitem{Gaussian16}
M.~J. Frisch et al.,
\emph{Gaussian 16}, Revisions A.03--C.02,
Gaussian, Inc., Wallingford CT, 2016.
The primary frequency-job source. The \texttt{Freq=hpmodes} option
produces the 5-decimal eigenvectors that the parser expects.

\bibitem{ORCA50}
F.~Neese,
\emph{Software update: the ORCA program system, version 5.0},
\textbf{WIREs Comput. Mol. Sci.} \textbf{12}, e1606 (2022).
The alternative input format (\file{.hess}) handled by
\file{orca\_io.py}.

\end{thebibliography}

% =========================================================================
\section*{Symbol glossary}
\addcontentsline{toc}{section}{Symbol glossary}

\subsection*{Physical and mathematical symbols}

\begin{longtable}{p{0.16\textwidth} p{0.18\textwidth} p{0.62\textwidth}}
\toprule
\textbf{Symbol} & \textbf{Units} & \textbf{Meaning} \\
\midrule
$\hbar$ & J\,s & Reduced Planck constant. \\
$\kB$ & J/K & Boltzmann constant. \\
$hc$ & J\,cm & $1.986\,446 \times 10^{-23}$ J\,cm; used for
                  cm$^{-1}$-to-Joule conversion. \\
$T$ & K & Temperature. Defaults to none (classical mode); for
              production analysis, set via \cfgkey{temp\_k}. \\
$\omega_l, \tilde\nu_l$ & rad/s, cm$^{-1}$ & Angular and wavenumber
                          frequency of mode $l$. \\
$m_l, m_\text{red}, \mu$ & amu, kg & Reduced mass: \fn{red\_mass}
                                       in the code. \\
$\evec_{l,i}$ & dimensionless & Cartesian-displacement eigenvector
                              of mode $l$ on atom $i$. \\
$\widetilde{\evec}_{l,i}$ & \AA{} & Thermally scaled eigenvector,
                                  $\evec_{l,i} \cdot \urms(\omega_l)$. \\
$\urms(\omega, m, T)$ & \AA{} & QHO thermal RMS amplitude,
                                \cref{eq:urms-final}. \\
$\nhat$ & dimensionless & Unit cluster-normal,
                            \cref{eq:cluster-normal-svd}. \\
$\bhat$ & dimensionless & Unit bond vector. \\
$\alpha_{\OOP}, \alpha_{\INP}$ & \% & OOP/INP fractions,
                                       \cref{eq:oop-derivation}. \\
$\Delta r_X$ & \AA{} & Bond modulation along channel $X$,
                       \cref{eq:dr-classical-final}
                       or \cref{eq:dr-ha-final}. \\
$\Delta r_X^\text{RSS}$ & \AA{} & RSS aggregate over sub-channels,
                                    \cref{eq:dr-rss-derivation}. \\
$\lambda_X(\text{mode})$ & cm$^{-1}$ & Per-mode Marcus-Hush
                                        reorganization,
                                        \cref{eq:lambda-per-mode}. \\
$\Lambda_X^\text{total}$ & cm$^{-1}$ & System-wide total,
                                        \cref{eq:lambda-total-derivation}. \\
$\Lambda_X(\omega)$ & cm$^{-1}$ & Cumulative-Lambda curve,
                                   \cref{eq:lambda-cumulative}. \\
$M_X(\omega)$ & \AA{} & Modulation spectrum,
                          \cref{eq:M-spectrum-derivation}. \\
$C_X(\omega)$ & \AA{} & Co-modulation spectrum,
                          \cref{eq:co-modulation}. \\
$B_i$ & \AA$^2$ & Debye-Waller / B-factor for atom $i$,
                    \cref{eq:b-factor-final}. \\
$\sigma_e$ & dimensionless & Per-component eigenvector noise,
                                \cfgkey{sigma\_eigvec}. \\
$\sigma_c$ & \AA{} & Per-component coordinate noise,
                       \cfgkey{sigma\_coord}. \\
$\sigma_e^\text{therm}$ & \AA{} & $\sigma_e \cdot \urms$,
                                   \cref{eq:sigma-therm}. \\
\bottomrule
\end{longtable}

\subsection*{Subscripts and superscripts}

\begin{longtable}{p{0.16\textwidth} p{0.80\textwidth}}
\toprule
\textbf{Notation} & \textbf{Meaning} \\
\midrule
$X \in \{\text{FeFe, FeN, FeS, NH, HA}\}$ & The five Marcus-Hush
   channels (electron-transfer reaction coordinate, two metal-
   ligand variants, proton-donor coordinate, proton-acceptor
   coordinate). \\
$l$ & Mode index. \\
$i$ & Atom index. \\
$\alpha \in \{x, y, z\}$ & Cartesian-component index. \\
$s$ & Sub-channel index within a parent channel $X$. \\
$\OOP, \INP$ & Out-of-plane / in-plane component, relative to
                $\nhat$. \\
$G$ & Atom group (e.g.\ cluster, residue, SSE-element). \\
\bottomrule
\end{longtable}

\subsection*{Code identifiers}

\begin{longtable}{p{0.30\textwidth} p{0.66\textwidth}}
\toprule
\textbf{Identifier} & \textbf{Meaning} \\
\midrule
\fn{evec\_raw}        & Raw eigenvector, not thermally scaled.
                          Carried unchanged through the result-dict
                          (FUND 1 fix). \\
\fn{evg}              & Thermally scaled eigenvector,
                          \fn{evec\_raw}~$\cdot~\urms$, in \AA{}. \\
\fn{fe\_c, s\_c}      & Tuples of Gaussian centres for the two
                          Fe atoms and two S atoms of a cluster, in
                          canonical order. \\
\fn{c2l, c2l\_h}      & Cluster-to-ligand atom centre lists
                          (heavy-only vs hydrogen-inclusive). \\
\fn{coord\_info}      & The \fn{CoordInfo} dataclass instance,
                          \cref{sec:dataclass-coordinfo}. \\
\fn{normal}           & The cluster normal $\nhat$,
                          \cref{eq:cluster-normal-svd}. \\
\fn{u\_rms}           & The $\urms$ value for the current mode. \\
\fn{HC\_JCM}          & The constant $hc = 1.986446 \times 10^{-23}$
                          J\,cm, used for cm$^{-1}$ conversion. \\
\fn{b\_accum}         & B-factor accumulator array,
                          shape $(N_\text{atoms},)$. \\
\bottomrule
\end{longtable}

\subsection*{Glossary of terms}

\begin{longtable}{p{0.24\textwidth} p{0.72\textwidth}}
\toprule
\textbf{Term} & \textbf{Definition} \\
\midrule
Cluster-core mode
 & A normal mode whose dominant displacement amplitude resides on
   the four \mbox{[2Fe-2S]} cluster atoms.
   \fn{kern\_loc} quantifies this. \\
Context mode
 & A mode just outside the analysis window, loaded for the sole
   purpose of anchoring the interpolation workbook
   (\cref{sec:algo-context}). \\
Synthetic zero mode
 & A FUND 13 artefact mode at the upper window edge, all observables
   = 0, sentinel \texttt{number} = $-1$. Used only for
   interpolation anchoring; never contributes to aggregates. \\
Half-open frequency window
 & $[\texttt{lo}, \texttt{hi})$. A mode at $\omega = $ hi is
   excluded; this prevents double-counting in contiguous-window
   analyses. \\
RSS aggregation
 & Root-sum-square,
   $\sqrt{\sum_s w_s (\Delta r_s)^2}$, the energetically consistent
   aggregation rule for sub-channels
   (\cref{sec:rss-aggregation}). \\
PCET
 & Proton-coupled electron transfer. Detected in the pipeline by
   simultaneous activity of the HA and FeFe (or FeS) channels at
   the same frequency. \\
SCSD
 & Symmetry-coordinate structural decomposition
   (Kingsbury \& Senge 2024); projection onto $D_{2\mathrm{h}}$
   irreps for a planar 4-atom cluster. \\
Kernel classification
 & The 10-template heuristic that names cluster-core modes by their
   classical normal-mode character
   (``Fe-Fe sym stretch'', ``cluster umbrella'', etc.). \\
Significance colouring
 & Per-cell colour code in the Excel output that flags values
   below \cfgkey{significance\_threshold\_low}~$\sigma$ (grey),
   between low and high~$\sigma$ (yellow), or above
   high~$\sigma$ (white). \\
\bottomrule
\end{longtable}

% =========================================================================
\section*{Code-function index}
\addcontentsline{toc}{section}{Code-function index}

Alphabetical index of every function referenced via \codref{} in
this document, with the file and line range as of v1.0.4. The
companion test \fn{tests/test\_manual\_code\_refs.py} verifies
that each entry resolves to a real function definition; this
table is generated semi-automatically and refreshed for every
release. Subsequent releases may shift the line numbers; the
function names themselves are stable.

\begin{longtable}{p{0.42\textwidth} p{0.18\textwidth} p{0.13\textwidth}}
\toprule
\textbf{Reference} & \textbf{File} & \textbf{Lines} \\
\midrule
\texttt{config.Config}                                & \texttt{config.py}         & 34--699   \\
\texttt{core.\_oop}                                   & \texttt{core.py}           & 193--197  \\
\texttt{core.\_oop\_ring\_metrics}                    & \texttt{core.py}           & 318--368  \\
\texttt{core.analyze\_fe\_ligand}                     & \texttt{core.py}           & 467--634  \\
\texttt{core.analyze\_his\_hn}                        & \texttt{core.py}           & 635--738  \\
\texttt{core.analyze\_mode\_with\_fallback}           & \texttt{core.py}           & 1701--1841 \\
\texttt{core.classify\_oop\_inp}                      & \texttt{core.py}           & 924--1026 \\
\texttt{core.compute\_scsd\_for\_mode\_full}          & \texttt{core.py}           & 1995--2122 \\
\texttt{core.compute\_thermal\_amplitude}             & \texttt{core.py}           & 131--192  \\
\texttt{core.reset\_warning\_state}                   & \texttt{core.py}           & 88--130   \\
\texttt{export.export\_all}                           & \texttt{export.py}         & 118--328  \\
\texttt{export.export\_embedding\_excel}              & \texttt{export.py}         & 645--708  \\
\texttt{export.export\_interpolated\_excel}           & \texttt{export.py}         & 709--883  \\
\texttt{export.export\_main\_excel}                   & \texttt{export.py}         & 329--592  \\
\texttt{export.export\_sse\_excel}                     & \texttt{export.py}         & 593--644  \\
\texttt{geometry.cluster\_normal}                     & \texttt{geometry.py}       & 336--371  \\
\texttt{geometry.detect\_sse\_dssp}                    & \texttt{geometry.py}       & 1240--1422 \\
\texttt{geometry.detect\_sse\_phipsi}                  & \texttt{geometry.py}       & 1423--1566 \\
\texttt{geometry.find\_all\_clusters}                 & \texttt{geometry.py}       & 80--192   \\
\texttt{geometry.find\_cluster}                       & \texttt{geometry.py}       & 193--279  \\
\texttt{geometry.find\_coordinating\_residues}        & \texttt{geometry.py}       & 752--917  \\
\texttt{geometry.kabsch\_align}                       & \texttt{geometry.py}       & 423--573  \\
\texttt{logio.get\_eigvec}                            & \texttt{logio.py}          & 1209--1255 \\
\texttt{logio.load\_scan\_cache}                      & \texttt{logio.py}          & 549--600  \\
\texttt{logio.read\_all\_meta}                        & \texttt{logio.py}          & 878--971  \\
\texttt{logio.read\_std\_orient}                      & \texttt{logio.py}          & 658--766  \\
\texttt{logio.save\_scan\_cache}                      & \texttt{logio.py}          & 601--657  \\
\texttt{logio.scan\_log}                              & \texttt{logio.py}          & 421--520  \\
\texttt{orca\_io.get\_eigvec\_orca}                   & \texttt{orca\_io.py}       & 298--320  \\
\texttt{orca\_io.parseresult\_to\_atoms}              & \texttt{orca\_io.py}       & 249--274  \\
\texttt{reorganization.aggregate\_by\_parent}         & \texttt{reorganization.py} & 491--540  \\
\texttt{reorganization.build\_channels\_from\_coord\_info} & \texttt{reorganization.py} & 971--1096 \\
\texttt{reorganization.compute\_co\_modulation\_spectra} & \texttt{reorganization.py} & 854--891 \\
\texttt{reorganization.compute\_cumulative\_reorganization} & \texttt{reorganization.py} & 892--970 \\
\texttt{reorganization.compute\_mode\_modulations}    & \texttt{reorganization.py} & 356--440  \\
\texttt{reorganization.compute\_modulation\_spectra}  & \texttt{reorganization.py} & 799--853  \\
\texttt{reorganization.compute\_total\_reorganization} & \texttt{reorganization.py} & 744--798 \\
\texttt{reorganization.lambda\_pair\_cm1}             & \texttt{reorganization.py} & 246--291  \\
\texttt{reorganization.signed\_dr\_along\_axis}       & \texttt{reorganization.py} & 177--211  \\
\texttt{runner.\_make\_synthetic\_zero\_mode}         & \texttt{runner.py}         & 132--228  \\
\texttt{runner.\_run\_analysis\_single}               & \texttt{runner.py}         & 229--1577 \\
\texttt{runner.run\_analysis}                         & \texttt{runner.py}         & 1578--1712 \\
\bottomrule
\end{longtable}

\paragraph{How this index is maintained.}
The function names are stable across patch releases (v1.0.x). Line
ranges are refreshed at every release by re-running
\fn{tools/build\_codref\_index.py} against the source tree;
the test \fn{tests/test\_manual\_code\_refs.py::test\_supplement\_-
codrefs\_exist} ensures that every reference resolves to a real
function regardless of line numbers, providing a stable
hyperlink even when the underlying line numbers shift.

\end{document}
